\documentclass[a4paper]{report}
\usepackage[germanb]{babel}
\usepackage[final]{graphicx}

\selectlanguage{austrian}

% \addtolength{\textwidth}{4cm}
% \addtolength{\oddsidemargin}{-1.7cm}
% \addtolength{\evensidemargin}{-1.7cm}
% \addtolength{\topmargin}{-80pt}
% \addtolength{\textheight}{150pt}
% \addtolength{\marginparsep}{-5pt}

\makeindex
\sloppy
\frenchspacing
% \pagestyle{headings}

\setcounter{secnumdepth}{3}
\setcounter{tocdepth}{3}

\documentclass[a4paper]{report}
\usepackage[germanb]{babel}
\usepackage[final]{graphicx}

\selectlanguage{austrian}

% \addtolength{\textwidth}{4cm}
% \addtolength{\oddsidemargin}{-1.7cm}
% \addtolength{\evensidemargin}{-1.7cm}
% \addtolength{\topmargin}{-80pt}
% \addtolength{\textheight}{150pt}
% \addtolength{\marginparsep}{-5pt}

\makeindex
\sloppy
\frenchspacing
% \pagestyle{headings}

\setcounter{secnumdepth}{3}
\setcounter{tocdepth}{3}

\documentclass[a4paper]{report}
\usepackage[germanb]{babel}
\usepackage[final]{graphicx}

\selectlanguage{austrian}

% \addtolength{\textwidth}{4cm}
% \addtolength{\oddsidemargin}{-1.7cm}
% \addtolength{\evensidemargin}{-1.7cm}
% \addtolength{\topmargin}{-80pt}
% \addtolength{\textheight}{150pt}
% \addtolength{\marginparsep}{-5pt}

\makeindex
\sloppy
\frenchspacing
% \pagestyle{headings}

\setcounter{secnumdepth}{3}
\setcounter{tocdepth}{3}

\documentclass[a4paper]{report}
\usepackage[germanb]{babel}
\usepackage[final]{graphicx}

\selectlanguage{austrian}

% \addtolength{\textwidth}{4cm}
% \addtolength{\oddsidemargin}{-1.7cm}
% \addtolength{\evensidemargin}{-1.7cm}
% \addtolength{\topmargin}{-80pt}
% \addtolength{\textheight}{150pt}
% \addtolength{\marginparsep}{-5pt}

\makeindex
\sloppy
\frenchspacing
% \pagestyle{headings}

\setcounter{secnumdepth}{3}
\setcounter{tocdepth}{3}

\input{facharb.inc}

\begin{document}

\title{Design einer Programmiersprache\\
       und theoretische Besprechung\\
       eines Interpreters anhand der\\
       Programmiersprache SEL}
\author{Mark Probst}

\maketitle

\tableofcontents

\chapter{Die Programmiersprache}

Dieses Kapitel hat das Ziel, das Design der Programmiersprache
SEL, wie in \cite{selkr} beschrieben, zu erkl"aren bzw. zu
begr�nden, und zwar in Hinblick auf den Bau eines Interpreters
und eines Compilers. Aus diesem Grund setze ich eine gewisse
Kenntnis sowohl einer strukturierten Programmiersprache,
insbesondere nat"urlich der Sprache SEL, als auch der
grundlegenden Funktionsweise eines Interpreters bzw. eines
Compilers beim Leser voraus. Der Leser sollte also zumindest
\cite{selkr} gelesen haben, da"s, obwohl es als Kurzreferenz
konzipiert ist, durchaus auch geeignet ist, einen "Uberblick
"uber die Programmiersprache SEL zu vermitteln.

\section{Der Begriff der Programmiersprache}

Grunds"atzlich kann eine Programmiersprache definiert werden als

\begin{quotation}
Sprache zur Formulierung von Rechenvorschriften, d. h. von
Datenstrukturen und Algorithmen, die von einem Computer
ausgef"uhrt werden k"onnen.\footnote{siehe \cite{duden},
Seite~544}
\end{quotation}

Die in dieser Arbeit behandelte Programmiersprache SEL geh"ort
zur Gruppe der imperativen Programmiersprachen, deren wichtigste
Untergruppe wohl die der strukturierten Programmiersprachen ist.
Imperative Programmiersprachen werden im Duden Informatik
wie folgt definiert:

\begin{quotation}
Bei diesen Sprachen besteht ein Programm aus einer Folge von
Befehlen an den Computer, wie z.~B. "`Schreibe in die Variable a
den Wert 3"', "`Springe an die Stelle x im
Programm"'.\footnote{siehe \cite{duden}, Seite~545}
\end{quotation}

Strukturierte Programmiersprachen stellen f"ur den Programmierer
zus"atzlich Hilfen in Form von Kontrollanweisungen, Schleifen und
benutzerdefinierten Funktionen zur Verf"ugung. Diese machen
sowohl das Schreiben der Programme einfacher und schneller wie
auch die Programmtexte leichter lesbar. Zur Zeit sind sie die am
meisten verwendete Art der Programmiersprachen.

\section{Ziele f\protect{"ur} das Design von SEL}

Bevor man an die Arbeit an einem so komplexen Projekt wie das Design
einer Programmiersprache herangeht, sollte man sich im Klaren
dar"uber sein, welche Ziele man eigentlich verfolgt. So habe ich
auch an meine zuk"unftige Programmiersprache einige Anforderungen
gestellt, von denen ich die meisten in der gegenw"artigen Version
der Sprache verwirklicht sehe.

Die wichtigste Entscheidung bei einem Projekt ist die nach dem
Zweck desselben. Im Fall von SEL wollte ich weniger eine
Programmiersprache f"ur das Schreiben von kommerzieller Software
schaffen, als vielmehr eine kleine, feine Sprache, die ich in
vielen Situationen benutzen k"onnte. Haupts"achlich schwebte mir
dabei die Verwendung als Makrosprache in diversen Programmen vor.
Eine solche Makrosprache sollte vor allem leicht erlernbar sein
und keine gro"se Umstellung von Seiten des Benutzers verlangen,
sollte aber andererseits nicht so primitiv aufgebaut sein, da"s
sie eine vern"unftige Verwendung nicht mehr m"oglich macht, was
die meisten Makrosprachen leider sind.

Um eine geringe Umstellungszeit m"oglich zu machen, mu"s die
Programmiersprache an eine bereits bestehende, weit verbreitete
Sprache angelehnt sein. Im Fall von SEL habe ich mich
haupts"achlich an C und \pascal\ gehalten, weil mir diese am
gel"aufigsten und vermutlich die am weitest verbreiteten
Sprachen sind.

Da ich ein gro"ser Freund von C und ihrer pr"agnanten und kurzen
Ausdrucksweise bin, habe ich mir zum Ziel gesetzt, auch SEL zu
einer effizienten, kurzen Ausdrucksweise zu verhelfen. Eine
Konsequenz aus dieser Entscheidung ist es, anstatt von
Schl"usselw"ortern, so wie in \pascal , Symbole wie in C zu
verwenden, soferne dies nicht die "Ubersichtlichkeit und
Konsistenz innerhalb der Sprache zu sehr gef"ahrdet. Die zweite
Konsequenz ist es, redundante Sprachteile zu entfernen, solange
auch das nicht die Lesbarkeit der Programmtexte verletzt.

Ein weiteres wichtiges Kriterium bei einer Makrosprache ist die
Erweiterungsf"ahigkeit, da Makrosprachen in so verschiedenen
Anwendungen wie z.~B. Textverarbeitungen und Zeichenprogrammen
eingesetzt werden, da"s es schlichtweg unm"oglich ist, eine
Sprache zu konstruieren, die f"ur alle diese Anwendungsgebiete
passend ist. Diese Erweiterungsf"ahigkeit sollte einerseits durch
die Erweiterung mit Hilfe von Funktionen wie z.~B.
f"ur das Zeichnen von Linien oder Kreisen und andererseits durch
die M"oglichkeit der "Anderung in der Grammatik, wie etwa
f"ur das flexiblere Handling von Vektoren und Matrizen, gegeben
sein. Einige Vorschl"age f"ur "Anderungen der letzteren Art sind
in Kapitel \ref{langext} gegeben.

Zuletzt soll die Programmiersprache SEL so aufgebaut sein, da"s
es nicht nur m"oglich ist, einen Interpreter, sondern auch mit
vertretbarem Aufwand einen Compiler zu implementieren. Nicht nur,
weil das bei einer Programmiersprache immer eine gute Idee ist,
sondern auch in Hinblick darauf, da"s ich selbst einmal einen
Compiler schreiben m"ochte.

Als Abschlu"s dieses Unterkapitels seien nochmals alle Forderungen
kurz zusammengefa"st, um in sp"ateren Abschnitten darauf
zur"uckgreifen zu k"onnen.

\begin{itemize}
\item Geringer Sprachumfang, jedoch auch f"ur mittelgro"se
Programme geeignet
\item Anlehnung an C und \pascal\
\item Verwendung von Symbolen anstatt von Schl"usselw"ortern
(wenn m"oglich)
\item Vermeidung von Redundanzen
\item Konsistenz innerhalb der Sprache
\item Erweiterungsf"ahigkeit
\item Sowohl f"ur Interpreter als auch f"ur Compiler geeignet
\end{itemize}

\section{Das Programm}

Das gr"o"ste in sich abgeschlossene Objekt, da"s ein Compiler bzw.
Interpreter zu behandeln hat, ist das Programm. Uns stellt sich nun die
Frage, wie dieses Programm aufgebaut sein soll.

Zun"achst ist zu
entscheiden, ob das Programm in Module eingeteilt werden soll oder
nicht. Eine solche Einteilung hat den Vorteil, da"s der Programmierer
Implementationsdetails eines Programmteils oder einer Funktionssammlung
vom Rest des Programms abschirmen kann und den Anwender dieser
Funktionssammlung gar nicht erst in Versuchung bringt, in die internen
Abl"aufe einzugreifen. Nat"urlich ist das nicht der einzige Vorteil
dieser Technik. So kann es z.~B. bei einer Modularisierung des Programms
auch kaum zu Namens"uberschneidungen bei Funktionen und Variablen
kommen. Last but not least wird auch das Programm leichter
"uberschaubar, weil man sich jeweils nur auf ein Modul konzentrieren
mu"s. Ich habe mich f"ur eine Modularisierung entschieden. Erstens, weil
es die M"oglichkeiten des Programmierers um einen nicht zu
untersch"atzenden Faktor erweitert und zweitens, weil es nicht schwer
ist, ein solches Modulkonzept in einen Interpreter oder einen Compiler
einzubauen---jedenfalls nicht das hier verwendete.

In SEL wird ein Modul in mehrere Bl"ocke eingeteilt, die in beliebiger
Reihenfolge auftreten d"urfen, wobei es drei Arten
von Bl"ocken zu unterscheiden gibt:

\begin{itemize}
\item Die \rw{global}-Bl"ocke werden verwendet, um programmglobale Variablen,
Funktionen oder Datentypen
zu deklarieren, also solche, die im ganzen Programm bekannt
sind.
\item Die \rw{local}-Bl"ocke werden analog zu \rw{global}-Bl"ocken
verwendet, um
modulglobale Variablen, Funktionen oder Datentypen zu deklarieren.
\item Funktionen werden "ahnlich wie in C und \pascal\ definiert.
\end{itemize}

Der Hauptgrund, warum ich mich f"ur diese Art der Modulstruktur
entschieden habe ist der, da"s es f"ur einen Interpreter weit
praktikabler ist, zur Initialisierung der Variablen vor dem Start eines
Programms einfach nach den Schl"usselw"ortern \rw{global} und
\rw{local} zu suchen, als zwischen den Funktionsdefinitionen. Dieses
Problem h"atte nat"urlich auch eine Programmstruktur wie die in \pascal ,
wo Variablen- und Typendeklarationen ihren festen Platz haben, gel"ost,
doch erscheint mir dieser Ansatz als zu starr, vermutlich, weil ich die
Freiheiten von C gewohnt bin.

Noch ein kleines Wort zur allgemeinen Struktur eines Moduls: Wie
bereits oben erw"ahnt, besteht ein Modul aus einer beliebigen
Abfolge von Bl"ocken, wobei es durchaus m"oglich ist, da"s
mehrere Bl"ocke ein und desselben Typs vorkommen. F"ur das
gesamte Modul gilt jedoch eine Regel: Ein Objekt ist erst ab der
Zeile bekannt, in der es deklariert wurde, das hei"st, da"s es
z.~B. nicht m"oglich ist, eine Variable in Zeile 154 zu
verwenden, wenn sie erst in Zeile 186 deklariert wird.
Diese Regelung macht es einem Compiler m"oglich, ein Programm in
einem Durchgang zu "ubersetzen ({\em One-Pass Compiler}). Eine
Ausnahme bilden die Strukturen (siehe Kapitel \ref{structs}).

\subsection{global-Bl\protect{"ocke}}

Ein \rw{global}-Block\footnote{siehe Anhang \ref{grammar} oder
\cite{selkr}, Seite 9} wird durch das Schl"usselwort \rw{global}
eingeleitet und durch \rw{end} beendet. Dazwischen befindet sich
eine beliebige Abfolge von Variablen-, Funktions- und
Strukturdeklarationen. Alle in diesem Abschnitt deklarierten
Objekte sind dazu bestimmt, auch in anderen Modulen des Programms
verwendet zu werden, was jedoch nicht unbedingt bedeutet, da"s
sie in diesen auch ohne weiteres verf"ugbar sind. Soll ein
programmglobales Objekt auch in einem anderen Modul verwendet
werden, so ist es n"otig, es auch dort zu deklarieren und---im
Falle von Variablen---als \rw{extern} zu bezeichnen. Der Grund
f"ur diese Vorgangsweise ist, da"s sich ein Compiler, der die
Deklarationen aller anderen Module verwenden kann, nicht
konstruieren l"a"st, ohne nicht ein eigenes Objektformat zu
benutzen, wie das z.~B. bei Turbo \pascal\ der Fall ist. Das
Standard-Objektformat von MS-DOS und Windows NT ist dazu nicht in
der Lage. Au"serdem w"are es erheblich schwerer, Module
zu schreiben, die gegenseitig voneinander abh"angig sind, soda"s
ich von dieser L"osung Abstand genommen habe.

Auffallend an den drei Arten der Deklarationen, die innerhalb
eines \rw{global}-Blocks vorkommen d"urfen, ist, da"s die
Variablendeklaration als einzige mit einem Semikolon
abgeschlossen wird, und das, obwohl es ohne Gef"ahrdung der
Eindeutigkeit m"oglich w"are, dieses wegzulassen. Der Grund
hierf"ur ist einerseits die Einhaltung der Konsistenz der
Sprache\footnote{siehe auch Kapitel \ref{funcs}} und andererseits
die Offenhaltung f"ur Erweiterungen der Sprache. So ist es im
Moment zwar lediglich m"oglich, eine Variable mit einer \tchar -,
\tint -, \tfloat - oder Stringkonstante zu initialisieren, jedoch
ist es durchaus denkbar, dies in zuk"unftigen Versionen auf
Arrays, Strukturen und Zeiger auszuweiten, wo die Eindeutigkeit
nicht mehr unbedingt gegeben sein mu"s.

Der Grund, warum Variablen, die in einem anderen Modul definiert
worden sind und in einem bestimmten Modul verwendet werden, als
\rw{extern} zu deklarieren sind und dies bei Strukturen und
Funktionen nicht n"otig ist, ist einfach: Bei einer
Strukturdeklaration wird dem Compiler bzw. Interpreter lediglich
"`mitgeteilt"', wie diese aufgebaut ist. Es wird in diesem Fall
kein vom Programm verwendetes Objekt im Speicher erzeugt. Bei
einer Variablendefinition hingegen wird dem Compiler sowohl der
Typ der Variable mitgeteilt als auch der Befehl, sie im Speicher
zu erzeugen. Ohne Unterscheidung zwischen Variablendefinition und
\rw{extern}-Deklaration w"u"ste der Compiler nicht, ob er nun
eine Variable im Speicher zu erzeugen h"atte oder sich lediglich
vormerken sollte, da"s diese Variable in einem anderen Modul
existiert und in diesem verwendet wird. Bei Funktionen ist es
kaum anders: Eine Deklaration innerhalb eines \rw{local}- oder
\rw{global}-Blocks sagt dem Compiler lediglich, da"s es
"`irgendwo anders"' eine solche Funktion gibt und da"s sie hier
verwendet wird. Eine Funktionsdeklaration innerhalb eines
\rw{global}-Blocks dient weiterhin dazu, zu kennzeichnen, da"s
diese Funktion, falls sie in diesem Modul implementiert ist, f"ur
die Verwendung in anderen Modulen bestimmt ist. Eine zweite
M"oglichkeit, dies zu erzielen, w"are, dem Funktionskopf bei der
Implementation noch ein spezielles Schl"usselwort hinzuzuf"ugen,
doch erscheint mir die in SEL verwendete L"osung als die bessere
und auf jeden Fall als die "ubersichtlichere.

Variablen d"urfen in \rw{local} bzw. \rw{global}-Bl"ocken
lediglich mit Konstanten, in
Funktionen hingegen auch mit Ausdr"ucken initialisiert werden.
Dies ist einerseits der Einfachheit halber so, andererseits, weil
es immer eine riskante Sache ist, Werte zu berechnen, noch bevor
die einzelnen Teile des Programms initialisiert sind. Dar"uber hinaus w"are
es notwendig, eine Regelung zu erstellen, die angibt, welche
Variablen als erste initialisiert werden und welches Modul einem
anderen vorangeht. Dieser Aufwand erscheint mir allerdings etwas
"ubertrieben, nicht zuletzt deshalb, weil eine solche
Initialisierung leicht in der \fu{main}-Funktion\footnote{siehe
Kapitel \ref{main}} nachgeholt werden kann.

\subsection{local-Bl\protect{"ocke}}

Ein \rw{local}-Block\footnote{siehe Anhang \ref{grammar} oder
\cite{selkr}, Seite 9} unterscheidet sich von einem
\rw{global}-Block, abgesehen davon, da"s er durch ein anderes
Schl"usselwort eingeleitet wird, dadurch, da"s eine
\rw{extern}-Deklaration nicht erlaubt ist. Dies ist auch ganz
verst"andlich, denn es ist ein Widerspruch in sich, eine
modulglobale Variable, die in einem anderen Modul definiert
ist, zu verwenden.

Die Tatsache, da"s Funktionen, die nicht explizit in einem
\rw{global}-Block aufscheinen, automatisch modulglobal sind,
verleitet zum Schlu"s, da"s Funktionsdeklarationen in
\rw{local}-Bl"ocken redundant sind. Es k"onnte jedoch ohne
weiteres die Notwendigkeit entstehen, da"s zwei Funktionen sich
gegenseitig aufrufen. In diesem Fall w"are jedoch ein
Aufruf illegal, da Funktionen erst ab dem Punkt ihrer
Deklaration im Modul bekannt sind. Dieser Konflikt ist sehr
einfach zu l"osen, indem man eine der Funktionen (oder der
"Ubersichtlichkeit halber beide) in einem \rw{local}-Block
deklariert und diese erst nach der anderen implementiert. Ein
Beispiel: Die Funktion \fu{A} und \fu{B} rufen sich gegenseitig
auf. Nun wird vor der Implementation die Funktion \fu{A} als
\rw{local} deklariert, soda"s sie im Rest des Moduls bekannt ist.
Nun wird die Funktion \fu{B} implementiert. Der Aufruf von \fu{A}
stellt kein Problem dar, der oben geschilderten Deklaration
wegen. Die Implementation von \fu{A} ist nun auch kein
Problem mehr, da an diesem Punkt schon beide Funktionen bekannt
sind.

\subsection{Funktionen}
\label{funcs}

Im Gegensatz zu \pascal\ wird in SEL, so wie in C, nicht zwischen
Funktionen und Prozeduren unterschieden. Dies hat---anders als es
zu erwarten w"are---fast nur Vorteile. Auf der einen Seite ist es
m"oglich, Prozeduren durch Funktionen mit dem R"uckgabetyp
\tvoid \footnote{siehe Kapitel \ref{elemtypes}} nachzubilden,
andererseits vereinfacht es den Bau eines
Interpreters bzw. eines Compilers. Der einzige Nachteil ist ein
geringer Verlust an Lesbarkeit.

Eine Funktion\footnote{siehe Kapitel \ref{exprfuncs} und
Anhang \ref{grammar} oder
\cite{selkr}, Seiten 7--9 und 17} gliedert sich augenscheinlich in
drei gro"se Teile:
Den Funktionskopf, die Definitionen der direkten
Funktionsvariablen\footnote{siehe Kapitel \ref{variables}}
und den Funktionsk"orper.

Der Funktionskopf besteht aus dem Schl"usselwort \rw{function},
dem Namen der Funktion, einer in Klammern eingeschlossenen Liste
der Parameter inklusive deren Typen und dem R"uckgabetyp. Der
Doppelpunkt vor dem Typ soll andeuten, da"s eine Funktion in
gewisser Weise einer Variablen eines bestimmten Typs "ahnelt.
Es sei angemerkt, da"s dieser Doppelpunkt ohne weiteres
ausgelassen werden k"onnte, ohne die Eindeutigkeit der Sprache zu
gef"ahrden.\footnote{Tats"achlich k"onnte sogar die schlie"sende
Klammer ausgelassen werden, doch diente das sicherlich nicht der
Lesbarkeit.}

Bei genauerer Betrachtung des Funktionskopfes und der
Funktionsdeklaration in einem \rw{global} bzw. \rw{local}-Block
f"allt auf, da"s in einem Funktionskopf die Namen der
Parameter aufscheinen, in den beiden letztgenannten Bl"ocken dies
jedoch nicht der Fall ist. Betrachten wir zur Begr"undung dieser
Entscheidung einen Interpreter und einen Compiler. Der
Interpreter kann alle Module des Programms auf einmal betrachten,
er kennt also die Namen der Parameter auch dann, wenn sie in
einem anderen Modul erw"ahnt sind als in dem, von dem aus die
Funktion aufgerufen wird. Ein Compiler jedoch besitzt diesen
Vorteil nicht: Er beschr"ankt sich jeweils auf ein einzelnes
Modul. Bei einem Funktionsaufruf ist aber das Einzige, was der
Compiler zu tun hat, zu "uberpr"ufen, ob die Parameter die
richtigen Typen haben, um sie dann in der richtigen Reihenfolge
auf dem Stack abzulegen. Wirkliche Bedeutung haben die Namen der
Parameter nur f"ur die eigentliche Funktion und insofern stellt
die Beschr"ankung auf die Typen keine Einschr"ankung dar, es ist
jedoch durchaus denkbar, da"s es der "Ubersichtlichkeit wegen in
einer zuk"unftigen Version von SEL m"oglich sein wird, auch die
Namen anzugeben, wobei diese Informationen aber vom Interpreter
bzw. Compiler ignoriert werden.

Dem Funktionskopf folgen die Definitionen der direkten
Funktionsvariablen.
Diese werden bei jedem Funktionsaufruf neu auf dem Stack angelegt
und belegen daher bei einer rekursiven Funktion in den
verschiedenen Instanzen nicht denselben Speicherbereich---erst
diese Tatsache macht Rekursionen so effektiv. Es ist m"oglich,
die Variablen mit Werten zu initialisieren, die sich aus Ausdr"ucken ergeben.
Dies stellt weder einen Compiler noch einen Interpreter vor
Probleme, soda"s ich hier nicht n"aher darauf eingehen m"ochte.
Das Semikolon am Ende jeder Variablendefinition w"are nicht
unbedingt notwendig, ergibt jedoch sowohl in Hinsicht auf die
"Ubersichtlichkeit als auch auf die Erweiterungsf"ahigkeit der
Sprache durchaus Sinn.

Der Funktionsk"orper wird durch das Schl"usselwort \rw{begin}
eingeleitet und durch \rw{end} beendet. Dazwischen befindet sich
eine beliebige Anzahl von Anweisungen, die bei Aufruf der
Funktion ausgef"uhrt werden.

\subsubsection{Die main-Funktion}
\label{main}

Jedes Programm mu"s einen Punkt haben, an dem es beginnt, soda"s
nicht unkontrolliert irgendwelche Funktionen ausgef"uhrt werden.
Hierf"ur gibt es grunds"atzlich zwei M"oglichkeiten. Entweder,
man f"uhrt bestimmte Konstrukte ein, um diesen Hauptteil des
Programms zu kennzeichnen, so wie es z.~B. in \pascal\
praktiziert wird, oder man verwendet eine bestimmte Funktion
dazu, wie in C. Ich habe mich f"ur letzteren
L"osungsansatz entschieden, weil er es erlaubt, das gesamte
Programm als eine Art Funktion zu sehen, die vom Anwender
aufgerufen werden kann. Es wird auch m"oglich, die
Kommandozeilenparameter als einfache Funktionsparameter zu
"ubergeben---wenn dies auch in der aktuellen Version von SEL
noch nicht eingef"uhrt ist. Weiters kann durch dieses Prinzip dem
Betriebssystem ein Wert zur"uckgegeben werden---der sogenannte
{\em Exit-Code}. In \pascal\ sind hierf"ur bestimmte Funktionen
bzw. Prozeduren notwendig.

Der Einfachheit halber und wegen der Verwandschaft mit C wird die
Hauptfunktion des Programms \fu{main} genannt. Sie nimmt keine
Parameter an und hat entweder den R"uckgabetyp \tvoid\ oder
\tint . Ist ersteres der Fall, wird immer der Exit-Code
\fu{0} zur"uckgegeben, andernfalls wird dieser durch die
\rw{return}-Anweisung bestimmt.

\section{Anweisungen}
\label{stmts}

Anweisungen kontrollieren das Geschehen innerhalb eines
Programms. Sie befehlen\footnote{daher die Bezeichnung
"`imperative Programmiersprache"'} dem Computer, was er tun soll.
In SEL wird grunds"atzlich zwischen drei Arten von Anweisungen
unterschieden: Kontrollanweisungen, Schleifen und Ausdr"ucken.
W"ahrend die ersten beiden in diesem Unterkapitel besprochen
werden, widmet sich das gesamte n"achste Unterkapitel den
Ausdr"ucken. Der Korrektheit wegen sei hier gesagt, da"s es
eigentlich noch eine vierte Art von Anweisungen gibt: Die {\em
Nullanweisung}, die jedoch, da sie keinen bestimmten Zweck
erf"ullt, au"ser vielleicht die Lesbarkeit des Programms zu
erh"ohen, f"ur unsere Betrachtungen nicht weiter wichtig ist.

Es ist auffallend, da"s SEL keine Anweisungen zur Verf"ugung
stellt, die in irgendeiner Weise mit der Hardware des Computers
kommunizieren k"onnen, wie z.~B. Anweisungen zum Einlesen von der
Tastatur, oder zur Ausgabe von Zeichen auf den Bildschirm oder in
eine Datei. Die Gr"unde hierf"ur sind vielf"altig. Erstens haben
wir es uns zum Ziel gesetzt, den Sprachumfang von SEL m"oglichst
gering zu halten und Ein/Ausgabeanweisungen alleine
w"urden---sollten sie einigerma"sen vern"unftig sein---die
Grammatik von SEL betr"achtlich aufbl"ahen. Weiters habe ich SEL
mit dem Ziel vor Augen entworfen, sie einmal als Makrosprache
einzusetzen. Makrosprachen allerdings m"ussen in vielen
verschiedenen Umgebungen sinnvoll einzusetzen sein, wie z.~B. in
einer Textverarbeitung oder in einem Zeichenprogramm, das unter
einer graphischen Benutzeroberfl"ache l"auft. Nun sind aber die
M"oglichkeiten zur Textausgabe in einer graphischen
Benutzeroberfl"ache grundverschieden zu denen von z.~B. DOS.
Ausgabeanweisungen zu entwerfen, die f"ur alle diese
M"oglichkeiten halbwegs geeignet sind, erscheint mir fast
unm"oglich. Deshalb ist es unumg"anglich, jede Implementation von
SEL mit verschiedenen Grundfunktionen auszustatten, die auf die
jeweilige Umgebung zugeschnitten sind. Eine Standardisierung
dieser Grundfunktionen ist immer eine gute Idee, doch sicher sind
diese Standardfunktionen nicht "uberall anwendbar.\footnote{Ein
Beispiel hierf"ur w"are der ANSI-Standard von C, dessen
Ein/Ausgabefunktionen auf eine textorientierte Oberfl"ache
zugeschnitten sind. Die \fu{printf}-Funktion z.~B. ist unter der
graphischen Benutzeroberfl"ache Windows nutzlos.}

\bigskip

An dieser Stelle scheint es mir wichtig, ein Konzept aus anderen
strukturierten Programmiersprachen zu erw"ahnen. Das
Blockkonzept\footnote{nicht zu verwechseln mit den Bl"ocken in
SEL} erm"oglicht es, eine Anweisung entweder als eben nur eine
Anweisung oder als eine von zwei Schl"usselw"ortern oder Symbolen
begrenzte, beliebige Abfolge von Anweisungen zu sehen. So verlangt
eine \rw{if}-Anweisung in C oder \pascal\ nur eine Anweisung als
Zweig, es ist jedoch kein Problem, mit Hilfe der Blockstruktur
auch mehrere Anweisungen ausf"uhren zu lassen. Dieses
Blockkonzept ist in manchen Situationen sehr praktisch, aber es
bl"aht die Grammatik einer Programmiersprache auch etwas auf und
macht es so nat"urlich schwerer, einen Interpreter oder Compiler
zu schreiben. Aus diesem Grund habe ich mich entschieden, es in
SEL nicht einzubauen. Die Quelltexte sind so zwar etwas l"anger,
doch es entsteht keine wirkliche Beeintr"achtigung.

\subsection{Kontrollanweisungen}

Kontrollanweisungen haben die Aufgabe, zu kontrollieren, ob und
was ausgef"uhrt wird. In SEL gibt es drei Anweisungen dieser Art.
Die \rw{if}-Anweisung "uberpr"uft eine Aussage und f"uhrt, wenn
sie einen positiven Wahrheitswert ergibt, bestimmte Anweisungen
aus. Auf Wunsch werden auch Anweisungen ausgef"uhrt, wenn das
Ergebnis der "Uberpr"ufung negativ ist. Die \rw{switch}-Anweisung
f"uhrt f"ur bestimmte Werte eines Ausdrucks verschiedene
Anweisungen aus. Die \rw{return}-Anweisung springt zur
aufrufenden Funktion zur"uck.

\subsubsection{Die if-Anweisung}

Die \rw{if}-Anweisung\footnote{siehe Anhang \ref{grammar} oder
\cite{selkr}, Seite 10} wird durch das Schl"usselwort \rw{if}
eingeleitet, auf welches ein in Klammern eingeschlossener Ausdruck
folgt. Die Alternative zu dieser L"osung w"are gewesen, den
Ausdruck sofort auf das \rw{if} folgen zu lassen und durch ein
bestimmtes Schl"usselwort (z.~B. \rw{then}) oder ein Symbol zu
beenden. Der Doppelpunkt jedoch w"urde Mehrdeutigkeiten
hervorrufen und ein Semikolon erscheint mir als unpassend, und da
wir uns entschienden haben, Symbole zu verwenden, wo nicht
unbedingt Schl"usselw"orter erforderlich sind, f"allt auch das
\rw{then} weg. Die einfachste und unauff"alligste L"osung in
diesem Fall ist die Klammerung.

Auf die schlie"sende Klammer folgt eine beliebige Anzahl von
Anweisungen, beendet entweder durch \rw{end} oder \rw{else}. Da
auch \rw{if} eine Anweisung ist, ist es m"oglich, \rw{if}s zu
verschachteln. Im zweiten Kapitel wird sich zeigen, da"s es f"ur
einen Interpreter bzw. Compiler kein Problem ist, mit einer
solchen Schachtelung fertigzuwerden. Folgt auf die Anweisungen ein
\rw{end}, so ist die \rw{if}-Anweisung beendet. Es w"are
m"oglich, auch das \rw{end} durch ein Symbol zu ersetzen,
doch alle sinnvollen Symbole, bis auf die Klammern, sind bereits
f"ur andere Zwecke vergeben. Nat"urlich k"onnte man auch eine
schlie"sende Klammer benutzen, doch um ein solches Konstrukt
sinnvoll zu machen, ben"otigt man auch eine "offnende Klammer.
Eine runde Klammer w"urde den falschen Eindruck erwecken, eckige
Klammern sind bereits vergeben, es bleiben die geschwungenen
Klammern, die auch in C zu diesem Zweck verwendet werden. Da ich
aber auch in allen anderen F"allen einem \rw{end} bzw. einem
\rw{begin}-\rw{end} Paar den Vorrang gegeben habe, f"uhrte eine
Verwendung hier h"ochstens zu einer Verwirrung des Benutzers.

Folgt auf die Anweisungen ein \rw{else}, dann folgen auf
dieses weitere Anweisungen, durch ein \rw{end} begrenzt.
Nat"urlich k"onnte man auch anstelle des \rw{else} ein Symbol
verwenden, doch das w"urde den Zweck der Konstruktion zu sehr
verschleiern:

\begin{verbatim}
if x == 1
{
  do_what_you_want_to();
::
  do_something_else();
}
\end{verbatim}

Ich glaube, bei Betrachtung dieses Beispiels zeigt sich, warum
ich die Schl"usselw"orter den Symbolen vorgezogen habe.

Manche Programmiersprachen, wie z.~B. \modula , stellen
zus"atzlich zum normalen \rw{else} ein sogenannte \rw{elseif}
oder auch \rw{elsif} bereit. Es erlaubt, sollte die
\rw{if}-Bedingung negativ sein, noch eine zweite Bedingung zu
"uberpr"ufen und, sollte auch diese logisch falsch sein, eine
dritte u. s. w. Ich habe mich entschieden, ein solches
\rw{elseif} nicht in SEL einzubauen, aus zwei Gr"unden.
Erstens l"a"st sich diese \rw{if}-\rw{elseif} Kombination in
vielen F"allen durch ein besser passendes
\rw{switch}\footnote{siehe Kapitel \ref{switch}} ersetzen und
zweitens ist es m"oglich, den \rw{elseif}-Effekt durch ein
verschachteltes \rw{if} zu erzeugen---\rw{elseif} selbst ist
nichts anderes als eine Vereinfachung dieser Konstruktion.

\subsubsection{Die switch-Anweisung}
\label{switch}

Die \rw{switch}-Anweisung\footnote{siehe Anhang \ref{grammar} oder
\cite{selkr}, Seite 11} erlaubt es, abh"angig vom Wert eines
Ausdrucks verschiedene Anweisungen auszuf"uhren. Dazu wird
zun"achst der Ausdruck evaluiert und ein \rw{case} gesucht,
dessen Wert mit dem des Ausdrucks "aquivalent ist. Ist kein
passendes \rw{case} vorhanden, wird zu einem evtl. vorhandenen
\rw{default} gesprungen. Dort wird mit der Ausf"uhrung
fortgefahren und erst bei einem \rw{break} oder dem zum
\rw{switch} geh"orenden \rw{end} beendet, also nach dem \rw{end}
fortgesetzt. Ist erst einmal ein bestimmter \rw{case}- oder
\rw{default}-Zweig eingeschlagen, werden alle weiteren
\rw{case}'s und \rw{default}s ignoriert. Das erm"oglicht es, f"ur
mehrere Werte ein und dieselben Anweisungen auff"uhren zu lassen,
indem man die einzelnen \rw{case}s hintereinander anreiht.

Die \rw{case}- und \rw{default}-Schl"usselw"orter werden in SEL,
ebenso wie in C, als Pr"afixe f"ur Anweisungen behandelt: bei
normaler Ausf"uhrung werden sie einfach ignoriert. Sie sind also
im Grunde nichts anderes als eine Art Sprunglabel, so wie sie in
vielen Programmiersprachen f"ur die
\rw{goto}-Anweisung\footnote{siehe Kapitel \ref{goto}}
verwendet werden. Eine andere---auf den ersten Blick
nat"urlichere---L"osung w"are, auf die \rw{switch}-Anweisung eine
Liste von \rw{case}'s oder \rw{default}s folgen zu lassen, die
ihrerseits jeweils von Anweisungen gefolgt sind. Um es m"oglich
zu machen, da"s eine Anweisungsfolge f"ur mehrere verschiedene
Werte ausgef"uhrt wird, m"u"ste man statt einem \rw{case} jeweils
mehrere \rw{case}'s erwarten. Welche der beiden L"osungen man
bevorzugt, ist Geschmackssache. Ich habe mich f"ur die
Pr"afix-L"osung entschieden, weil ich sie erstens von C gewohnt
bin und sie zweitens ein wenig leistungsf"ahiger ist.

Das Schl"usselwort \rw{case} vor dem jeweiligen Wert ist
notwendig, um Zweideutigkeiten in der Sprache zu vermeiden. Um
dieses Schl"usselwort auslassen zu k"onnen, m"u"sten wir die oben
besprochene zweite L"osung w"ahlen, jedoch m"u"sten die einzelnen
Anweisungsfolgen wiederum durch ein anderes Schl"usselwort, z.~B.
\rw{end} beendet werden, soda"s sich keine wirkliche Ersparnis
erg"abe.

Auch daf"ur, da"s in \rw{case}-Zweigen nur Konstanten erlaubt
sind, gibt es Gr"unde. Erstens h"atten wir durch den Doppelpunkt
Zweideutigkeiten---er w"urde als Typecast\footnote{siehe Kapitel
\ref{casts}} interpretiert werden---doch dieses Problem ist durch
Wahl eines anderen Symbols ohne weiteres zu "ubergehen. Der weit
gewichtigere Grund ist, da"s die f"ur die \rw{switch}-Anweisung
typische Effizienz verloren ginge: Bis der Interpreter zum
richtigen Zweig gelangt, m"u"ste er eine Vielzahl von Ausdr"ucken
auswerten, was sehr viel mehr Zeit in Anspruch nimmt, als der
simple Vergleich mit einer Konstante. Au"serdem w"urde diese
Entscheidung das Schreiben eines Compilers etwas verkomplizieren.
In F"allen, in denen eine solche \rw{switch}-Anweisung n"otig
w"are, sollte man auf \rw{if}-\rw{else}-Verschachtelungen umsteigen.

\subsubsection{Die return-Anweisung}

Die \rw{return}-Anweisung hat die Aufgabe, zur aufrufenden
Funktion zur"uckzukehren und evtl. einen Wert zu "ubermitteln,
den sog. {\em Return-Code}. Eine andere M"oglichkeit w"are, den
return-code durch eine Zuweisung an eine Variable, die den Namen
der Funktion hat, zu bestimmen und mittels einer anderen
Anweisung bei Bedarf zur aufrufenden Funktion zur"uckzuspringen.
Da man im Normalfall beides an der gleichen Stelle tut, habe ich
mich f"ur die \rw{return}-L"osung entschieden.

Die einfachste Methode w"are jedoch, "uberhaupt keine Anweisung
bereitzustellen, mit der in die aufrufende Funktion
zur"uckgekehrt werden kann und den R"uckgabewert am Ende der
Funktion festzulegen. F"ur einen rekursiven Interpreter ist dies
die beste L"osung, doch wie wir im n"achsten Kapitel sehen
werden, stellt uns auch der \rw{return}-Ansatz vor keine gro"sen
Probleme.

\subsection{Schleifen}

Schleifen kontrollieren nicht, was, sondern wie oft etwas
ausgef"uhrt wird. Die drei Arten der Schleifen, die in SEL
erlaubt sind, entsprechen den in den meisten strukturierten
Programmiersprachen vorhandenen Schleifen.

\subsubsection{Die repeat-Schleife}

Die \rw{repeat}-Schleife in SEL entspricht der
\rw{repeat}-Schleife in \pascal , bis auf das Detail, da"s die
\rw{until}-Bedingung in Klammern zu setzen ist und kein Semikolon
danach erforderlich ist. Diese Konzeption deckt sich mit den
restlichen Anweisungen von SEL und wurde daher von mir bevorzugt.

In manchen Programmiersprachen, wie z.~B. C, existiert die
\rw{repeat}-Schleife in einer abgewandelten Form, n"amlich der
der \rw{do-while}-Schleife. Der einzige Unterschied
besteht darin, da"s die \rw{repeat}-Schleife eine
Abbruchbedingung erwartet, die \rw{do-while}-Schleife jedoch eine
Fortf"uhrungsbedingung, das hei"st, eine Schleife kann durch die
andere ersetzt werden, wenn die Bedingung negiert wird. Welche
Schleife die bessere ist, kommt auf die Situation an, wirkliche
Vor- oder Nachteile gibt es nicht, die Auswahl ist also reine
Geschmackssache.

\subsubsection{Die while-Schleife}

Die \rw{while}-Schleife ist eng verwandt mit der
\rw{repeat}-Schleife. Der einzige Unterschied besteht darin, da"s
die \rw{repeat}-Schleife mindestens einmal durchlaufen wird, die
\rw{while}-Schleife jedoch nicht unbedingt. Syntaktisch gesehen
ist die \rw{while}-Schleife sehr "ahnlich der von C, die
Unterschiede sind in der in Kapitel \ref{stmts} angesprochenen
Blockstruktur von C begr"undet. So folgt auf die schlie"sende
Klammer der Schleifenbedingung in C in der Regel eine "offnende
geschwungene Klammer, in SEL jedoch ist diese ausgespart, da
nicht notwendig. Die schlie"sende geschwungene Klammer, die in C
die Schleife beendet, ist mit dem Schl"usselwort \rw{end} in SEL
zu vergleichen. Die Gr"unde, warum ich die Konstruktion mittels
Schl"usselw"ortern vorgezogen habe, ist bereits in vorhergehenden
Kapiteln er"ortert, ich gehe deshalb hier nicht n"aher darauf
ein.

\subsubsection{Die for-Schleife}

Die \rw{for}-Schleife in SEL erlaubt es, Anweisungen in einer
Anzahl, die schon im voraus bekannt ist, zu wiederholen.
Damit unterscheidet sie sich signifikant von den \rw{while} und
der \rw{repeat}-Schleifen, bei denen die Anzahl der Wiederholungen
nicht unbedingt bekannt sind. Das verhilft der
\rw{for}-Schleife---zumindest auf einem Interpreter---zu einer
sehr schnellen Ausf"uhrung, da die Bedingungen f"ur die Schleife
nur einmal berechnet werden m"ussen, w"ahrenddessen dies bei den
beiden anderen Schleifenarten bei jedem Durchlauf geschehen mu"s.

Dieses eine Merkmal der \rw{for}-Schleife ist vielen
Programmiersprachen gemeinsam, doch eine bekannte Sprache macht
eine Ausnahme: C. In C werden nicht eine Z"ahlervariable und
deren untere und obere Grenzen und evtl. noch die Schrittweite
angegeben, sondern es werden drei Ausdr"ucke verwendet: Der erste
wird lediglich beim Start der Schleife evaluiert, der dritte nach
jedem Durchlauf. Der zweite Ausdruck schlie"slich gibt die
Bedingung an, bei der die Schleife nochmals ausgef"uhrt werden
soll---ist sie logisch falsch, wird die Schleife abgebrochen. In
C ist also eine \rw{for}-Schleife nichts anderes als eine
erweiterte \rw{while}-Schleife und hat den Vorteil, da"s mit
ihrer Hilfe auch verkettete Listen und "ahnliches leicht
durchlaufen werden k"onnen. Auch ein "`klassische"'
\rw{for}-Schleife l"a"st sich so leicht nachbilden, doch diese
Art der Schleifenkonstruktion hat f"ur einen Interpreter den
gro"sen Nachteil, da"s bei jedem Schleifendurchlauf gleich zwei
Ausdr"ucke evaluiert werden m"ussen.

Ungl"ucklicherweise habe ich diesen Nachteil nicht sofort bemerkt und
f"uhrte auch in SEL die von C gewohnte \rw{for}-Konstruktion ein. Erst
sp"ater merkte ich, um wie viel effizienter diese Schleife f"ur einen
Interpreter w"are, w"urde ich sie "ahnlich wie in \pascal\ aufbauen. Die
"Anderungen am Interpreter und an der Dokumentation lie"sen sich zum
Gl"uck in weniger als einer Stunde durchf"uhren, soda"s Verz"ogerungen
f"ur den Zeitplan dieser Arbeit ausblieben und ich zur Erkenntnis
gelangte, da"s man aus Fehlern doch am besten lernt.

Noch ein kleines Wort zur Syntax: Die Bedingungen der Schleife sind in
Klammern gesetzt, weil das in SEL auch bei anderen Anweisungen der Fall
ist und dadurch die Konsistenz gewahrt wird, und die Grammatik ohne
Begrenzung nicht eindeutig w"are. Die Z"ahlervariable wird als
\fu{primary} angegeben, weil es auch m"oglich sein sollte, Array- und
Strukturelemente als Z"ahler zu benutzen, und nicht als \fu{expr}, weil
dann der folgende Zuweisungsoperator als ein Teil dieses Ausdrucks
interpretiert w"urde. Sollte es wirklich notwendig sein, einen \fu{expr}
zu benutzen, so kann dieser in Klammer gesetzt werden, und entspricht
damit einem \fu{primary}.\footnote{siehe Anhang \ref{grammar}} Anstatt
der Schl"usselw"orter \rw{to} bzw. \rw{downto} Symbole zu verwenden,
w"are erstens un"ubersichtlich und w"urde zweitens---au"ser bei
Verwendung neuer Symbole---die Eindeutigkeit zunichte machen, weshalb
ich die oben angesprochenen Schl"usselw"orter verwendete. Gleiches gilt
f"ur \rw{step}.

\subsection{Die goto-Anweisung}
\label{goto}

Die "Uberschrift dieses Absatzes mag darauf deuten, da"s es in SEL eine
\rw{goto}-Anweisung gibt, doch der Schein tr"ugt: Aus Gr"unden, die in
Kapitel \ref{parser} aufgef"uhrt sind und aus dem Grund, da"s diese
Anweisung in dem meisten F"allen ohnedies nicht notwendig ist und nur zu
einem schlechten Programmierstil verf"uhrt, habe ich mich gegen eine
Einbindung in SEL entschieden.    

\section{Ausdr\protect{"ucke}}

Ausdr"ucke sind ein Konzept, das in fast allen strukturierten
Programmiersprachen verwendet wird und das nicht ohne Grund: Es
vereinfacht die Grammatik der Sprache und macht damit das Schreiben
eines Interpreters bzw. eines Compilers einfacher. Au"serdem erleichtert
es dem Anwender der Sprache das Schreiben von Programmen sehr und
unterst"utzt die Lesbarkeit. So k"onnen Ausdr"ucke in vielen
verschiedenen Situationen verwendet werden: Als Bedingungen in
Schleifen, als Funktionsaufrufe oder als die rechten Seiten von
Zuweisungen.\footnote{Tats"achlich ist in SEL sogar die Zuweisung ein
Ausdruck, im Gegensatz zu \pascal\ oder \modula , doch dazu sp"ater
mehr.}

Ein Ausdruck in SEL kann "ahnlich einem Ausdruck im
mathematischen Sinn aufgefa"st werden, er ist eine Verkn"upfung
von Operanden durch Operatoren. Jeder Operand f"ur sich ist
selbst ein Ausdruck\footnote{bis auf die Operanden der Operatoren
\fu{new} und \fu{sizeof}, die Typen sind.} und das Ergebnis der
Verkn"upfung mit einem Operator ist wieder ein Ausdruck. So ist
es in SEL m"oglich, Ausdr"ucke beliebig tief zu schachteln. Diese
Eigenschaften der Ausdr"ucke legen es nahe, einen Ausdruck als
Baum\footnote{siehe \cite{sedge}, Seiten 57--72} darzustellen.
Dabei stellt jeder innere Knoten sowie die Wurzel einen Operator
dar und jedes Blatt einen Wert. Betrachtet man nun jeden
einzelnen Knoten (und dessen Nachfolger), so gelangt man zum
Ergebnis, da"s dieser auch wieder ein Ausdruck f"ur sich ist. Ein
Beispiel f"ur diese Darstellung gibt Sedgewick in \cite{sedge},
Seite 63, Abbildung 4.4.

Jeder Ausdruck in SEL hat nicht nur einen Wert, sondern auch
einen Typ\footnote{siehe auch Kapitel \ref{types}}, mit einer Ausnahme:
Ein Ausdruck des Typs \tvoid\ hat keinen Wert, was auch gar keinen Sinn
h"atte.

\subsection{Operatoren}

Grunds"atzlich gibt es vier Arten, in SEL Operatoren zu
klassifizieren\footnote{genauer gesagt vier "`vern"unftige"' Arten}:
nach der Anzahl ihrer Argumente, nach ihrer funktionellen
Zugeh"origkeit, nach ihrer Priorit"at und nach ihrem Ergebnistyp.

\paragraph{Klassifikation nach der Anzahl der Argumente}

Auf den ersten Blick gesehen gibt es in SEL nach dem Gesichtspunkt der
Zahl der Argumente zwei Arten von Operatoren: un"are und bin"are. Zu den
un"aren z"ahlen z.~B. die Operatoren \fu{+} und \fu{<<}, zu den bin"aren
z"ahlen \fu{!} und \fu{-}. Der letztgenannte Operator existiert auch in
einer bin"aren Version, man k"onnte also sagen, da"s es Operatoren gibt,
die sowohl un"ar als auch bin"ar sind. Ich pers"onlich betrachte diese
beiden "`Versionen"' als zwei verschiedene Operatoren, da auch ihre
Funktionen verschieden sind. Betrachtet man Ausdr"ucke in SEL genauer,
wird man bemerken, da"s ein
Funktionsaufruf durchaus mehrere Argumente erwarten kann, jedoch stellt
sich hier die Frage, ob man einen Funktionsaufruf (besser: eine Funktion) als
Operator auffassen sollte. Ich bin der Meinung, eine Funktion ist
weniger ein Operator als ein Wert, und habe ihr deshalb ein eigenes
Unterkapitel gewidmet.\footnote{siehe Kapitel \ref{exprfuncs}}

\paragraph{Klassifikation nach der funktionellen Zugeh"origkeit}

Auf diese Weise kann man die Operatoren von SEL in mehrere
Gruppen einteilen, wobei die Grenzen eher flie"send sind. Mir
erscheint folgende Einteilung am vern"unftigsten: arithmetische
Operatoren wie \fu{+} oder \fu{<<}, logische Operatoren wie
\fu{\&\&} oder \fu{!}, Vergleichsoperatoren wie \fu{>=} oder
\fu{==}, Operatoren zur Speicherverwaltung wie \fu{new} oder
\fu{[]} und als zwei eigene Gruppen der Zuweisungsoperator \fu{=}
und der Typecasting-Operator \fu{:}. Wie gesagt sind die
"Uberg"ange zwischen diesen Gruppen sehr flie"send. So k"onnte
man beispielsweise meinen, der Operator \fu{[]} geh"ore zur
Gruppe der arithmetischen Operatoren, da er ja in gewisser Weise
Vektorarithmetik m"oglich macht. Auch w"are es m"oglich, den
Typecasting-Operator zur Speicherverwaltung z"ahlen, da er auch
den Speicher modifiziert.

\paragraph{Klassifikation nach der Priorit"at}

Operatoren in SEL besitzen verschiedene
Priorit"atsstufen\footnote{siehe \cite{selkr}, Seite 21}, um eine
"uberm"a"sige Klammerung zu vermeiden. Die Priorit"aten in SEL sind
stark an die von C angelehnt\footnote{siehe \cite{kandr}, Seite 54},
weichen jedoch teilweise von diesen ab. Meistens im selben Atemzug mit
der Priorit"at der Operatoren wird die Richtung der Auswertung genannt.
Diese gibt an, ob der Operator links- oder rechts-assoziativ ist. Ein
Beispiel f"ur einen rechts-assoziativen Operator ist \fu{=}: Der
Ausdruck \fu{a=b=0} ist "aquivalent zu \fu{a=(b=0)}. So ist es ohne
Klammerung m"oglich, mehreren Variablen in einer Anweisung denselben
Wert zuzuweisen. Ein
links-assoziativer Operator ist das bin"are \fu{-}. W"are er das nicht,
w"urde der Ausdruck \fu{5-3-1} nicht zu \fu{(5-3)-1}, sondern zu
\fu{5-(3-1)} "aquivalent sein und so das falsche Ergebnis, n"amlich
\fu{3} anstatt \fu{1}, liefern.

\paragraph{Klassifikation nach dem Ergebnistyp}

Diese Einteilung ist wohl die komplizierteste und es w"are wohl besser,
man w"urde als Kriterium nicht den Ergebnistyp selbst, sondern das, 
was den Ergebnistyp bestimmt, heranziehen. 
So sind dies etwa beim Operator \fu{+} die
Typen seiner Operanden, beim Operator \rw{new} hingegen ist es der Typ,
mit dem er verbunden wird. Weiters gibt es Operatoren, die immer denselben
Ergebnistyp haben, wie z.~B. die logischen Operatoren oder der Operator
\rw{sizeof}. Es gibt sogar einen Operator, dessen Ergebnistyp immer
\tvoid\ ist, n"amlich \rw{delete}.

\bigskip

In der folgenden Aufstellung der Operatoren verwende ich die
Klassifikation nach der funktionellen Zugeh"origkeit, weil sie
f"ur diesen Zweck am geeignetsten ist.

\subsubsection{Arithmetische Operatoren}

SEL bietet neben den numerischen Operatoren, zu denen die vier
Grundrechnungsoperatoren f"ur Addition, Subtraktion, Multiplikation und
Division z"ahlen, zus"atzlich der Modulo-Operator, der den Rest bei einer
ganzzahligen Division ergibt, und der Negationsoperator, der
einer Umkehr des Vorzeichens gleichkommt, auch
Operatoren zur Bitmanipulation, n"amlich bitweise Links- und
Rechtsverschiebung und bitweise Und-, Oder- und
XOR-Verkn"upfungen. Die bin"aren numerischen Operatoren, also
\fu{+}, \fu{-}, \fu{*}, \fu{/} und \fu{\%}, besitzen eine
niedrigere Priorit"at als alle un"aren Operatoren, jedoch eine
h"ohere Priorit"at als alle anderen bin"aren Operatoren.
Untereinander sind die Additions- und Subtraktionsoperatoren auf
derselben Priorit"atsstufe, die um einen Grad niedriger steht,
als die der anderen numerischen bin"aren Operatoren, in Anlehnung
an die in der Mathematik "ublichen Notation. Die un"aren
Operatoren haben deshalb eine h"ohere Priorit"at, weil das
erstens unseren Erwartungen entspricht---man betrachtet intuitiv einen
un"aren Operator einem Wert als "`n"aherstehender"', als einen
bin"aren---und es zweitens in fast allen F"allen praktischer ist.
So ist z.~B. das Ergebnis des Ausdrucks \fu{-x*y} dem von
\fu{-(x*y)} "aquivalent. Bes"a"sen aber die bin"aren Operatoren
eine h"ohere Priorit"at, m"u"ste man anstatt \fu{-x-y} den
Ausdruck \fu{(-x)-y} schreiben, was etwas unnat"urlich aussieht.
Hingegen k"onnte man den Ausdruck \fu{-(x+y)} zwar bei anderer
Operatorenrangfolge ohne Klammern schreiben, doch selbst ohne die
Rangfolge zu ver"andern, k"onnte man durch "Anderung des \fu{+} in
ein \fu{-} die Klammerung vermeiden: \fu{-x-y}. Die Rangfolge
erweist sich also als nat"urlich und praktisch. Die
Vergleichsoperatoren haben deshalb niedrigere Priorit"at als die
numerischen Operatoren, weil es kaum vorkommt, da"s man mit
Vergleichsergebnissen rechnet. F"ur die logischen Operatoren gilt
dieselbe Begr"undung. Die Festsetzung der Priorit"at der
Bitmanipulations-Operatoren relativ zu den numerischen Operatoren
ist eher willk"urlich. Es w"urde kaum ein Vor- oder Nachteil
entstehen, h"atten die Operatoren zur Bitmanipulation h"ohere
Priorit"at. Der Grund, warum ich diese Operatoren nicht auf
dieselbe Priorit"atsstufe gesetzt habe ist der, da"s sie zu
verschieden sind, als da"s dies dienlich w"are.

Die numerischen Operatoren haben eine besondere Eigenschaft: Sie
stellen sich auf die Typen ihrer Operanden ein und ver"andern
diese sogar, wenn n"otig. So stellt sich zum Beispiel die Frage,
welchen Typ eine Addition eines \tint\ und eines \tfloat\ haben
sollte. Vern"unftigerweise wandelt der Additionsoperator in
diesem Fall den \tint\ in einen \tfloat\ um, und f"uhrt erst
anschlie"send die Addition aus. Auf diese Weise gehen keine
Nachkommastellen des \tfloat\ verloren. Gleiches passiert mit
einem \tchar\ und einem \tint : Der \tchar\ wird in den Typ
\tint\ umgewandelt, da nur so gew"ahrleistet ist, da"s mit dem
richtigen Wert des \tint\ gearbeitet wird. F"ur die Typen bei
einer numerischen Operation gilt in SEL also die Regel \hbox{$
\tchar < \tint < \tfloat $}.\footnote{siehe \cite{selkr}, Seite
18} Zus"atzlich zu diesen Typen k"onnen die Operatoren \fu{+} und
\fu{-} auch mit Zeigern operieren, doch dazu mehr in Kapitel
\ref{memmgrops}.

Die Operatoren zur Bitmanipulation arbeiten in SEL lediglich mit
dem Typ \tint\, und das aus guten Gr"unden: Erstens ist es f"ur
einen Computer zwar leicht m"oglich, mit den Bits eines \tfloat\
zu operieren, doch es ist f"ur einen Menschen ungleich
schwieriger, diese Bits auch richtig zu interpretieren---auf den
Wert der Zahl bezogen. Zweitens sind die "ublichen 8 Bit eines
\tchar\ im Regelfall zu wenig f"ur die Anwendungen der
Bitmanipulation und drittens vereinfacht es nicht nur den Bau
eines Interpreters bzw. Compilers, es macht diese auch schneller.

\subsubsection{Logische Operatoren}

Neben den Operatoren zur logischen Und- und Oder-Verkn"upfung
bietet SEL auch einen Operator zur logischen XOR-Verkn"upfung an.
Ich habe mich dazu entschieden, weil es leicht in einen
Interpreter bzw. Compiler einzubauen ist und andererseits der
Operator in manchen Situationen sehr praktisch ist und logische
Ausdr"ucke verk"urzen kann. Die logischen Operatoren sind auf der
zweitniedrigsten Priorit"atsstufe angesiedelt---unter ihnen
befinden sich lediglich die Operatoren \fu{=} und \fu{delete}.
Diese Stufe habe ich gew"ahlt, weil logische Operatoren
haupts"achlich in Schleifen- und \rw{if}-Bedingungen eingesetzt
werden und es dort am praktischsten ist, wenn sie die niedrigste
Priorit"at haben. Nur selten ist es notwendig, da"s einer
Variablen ein logischer Wert zugewiesen wird, doch auch dieser
Fall ist ohne Klammerung zu bew"altigen, da ja der
Zuweisungsoperator eine noch niedrigere Priorit"at hat. Der
einzige Fall, in dem sich diese Rangfolge als unpraktisch
herausstellen w"urde ist, da"s in einer Schleifen- oder
\rw{if}-Bedingung ein logischer Wert einer Variablen zugewiesen
wird und dieser Wahrheitswert mit einem anderen mittels eines
logischen Operators verkn"upft werden mu"s. Ein Beispiel hierf"ur
w"are \fu{if (x \&\& (y = z))}. Tats"achlich tritt dieser Fall so
selten auf---in meiner bislang 4-j"ahrigen C-Laufbahn bin ich
noch nie in eine solche Situation gelangt---da"s man ihn getrost
au"ser Acht lassen kann.

Der logische Negationsoperator \fu{!} liegt auf derselben
Priorit"atsstufe wie die Operatoren zur bitweisen und
arithmetischen Negation und die beiden Operatoren \fu{*} und
\fu{\&}, die im Kapitel \ref{memmgrops} besprochen werden, welche
um eine Stufe tiefer liegt als die der Operatoren \fu{->}, \fu{:}
und \fu{[]}. Um das zu begr"unden, betrachten wir den Ausdruck
\fu{!x->y}. Offenbar wird hier versucht, aus der Struktur \fu{x}
auf das Element \fu{y} zuzugreifen und dieses zu negieren. H"atte
der Operator \fu{!} eine h"ohere Priorit"at als \fu{->}, so
w"urde sich die Negation auf die Struktur beziehen, und es ist
nicht m"oglich, eine Struktur zu negieren.

Logische Operatoren arbeiten in SEL nur mit dem Typ \tint .
Zwar w"urde es auch gen"ugen, den Typ \tchar\ zu diesem Zweck
heranzuziehen, doch da in kaum einer Anwendung Wahrheitswerte
gespeichert werden und \tint\ normalerweise der Gr"o"se des
Maschinenwortes des jeweiligen Prozessors entspricht, und daher
schneller zu bearbeiten ist, als ein \tchar , ist er die bessere
Wahl. Sollte es jedoch n"otig sein, Wahrheitswerte abzuspeichern,
kann das ohne weiteres mit Hilfe von \tchar s geschehen, die durch
ein Typecasting entstanden sind, oder aber durch Bitfelder, was
noch speicherplatzsparender ist.

\subsubsection{Vergleichsoperatoren}

SEL bietet insgesamt sechs Vergleichsoperatoren an. Diese sind um
eine Priorit"atsstufe h"oher als die logischen Operatoren
angesiedelt, da sie meistens in Schleifen- und
\rw{if}-Bedingungen eingesetzt werden und dort diese Rangfolge
von Vorteil ist. So wird es wahrscheinlich oft vorkommen, da"s
"uberpr"uft werden soll, ob zwei Werte zwei anderen Werten
entsprechen, aber eher selten wird man "uberpr"ufen wollen, ob
die Oder-Verkn"upfungen jeweils zweier Werte dieselben Ergebnisse
haben.

Vergleichsoperatoren k"onnen in SEL lediglich Werte der Typen
\tchar , \tint\ und \tfloat\ vergleichen, wobei das im vorletzten
Unterkapitel vorgestellte "`Gr"o"senverh"altnis"' gilt. Es ist
hingegen nicht m"oglich, z.~B. Strings miteinander zu
vergleichen, weil das, um es vern"unftig zu machen, eine
Ver"anderung im Typkonzept von SEL mit sich z"oge und au"serdem
sehr gut von einer Funktion erledigt werden kann---SEL stellt
alle n"otigen Konstrukte hierf"ur bereit.

\subsubsection{Speicherverwaltungsoperatoren}
\label{memmgrops}

Unter den Begriff "`Speicherverwaltungsoperatoren"' fallen all
jene Operatoren, die mit Zeigern, Arrays und Strukturen arbeiten.
Dazu z"ahlen nicht nur die Operatoren zur Ermittlung der
Speicheradresse einer Variablen und zur Dereferenzierung eines
Operators, \fu{\&} und \fu{*}, zum Zugriff auf ein Array- bzw.
Strukturelement, \fu{[]} und \fu{->}, sondern auch die Operatoren
zum Allokieren und Freigeben von Speicher, \fu{new} und
\fu{delete}, und der Operator zum Ermitteln der Gr"o"se eines
Typen, \fu{sizeof}. Au"serdem fallen unter diesen Begriff die
beiden bin"aren Operatoren \fu{+} und \fu{-}, mit denen
zus"atzlich zur numerischen Arithmetik auch Zeigerarithmetik
betrieben werden kann.

Die h"ochste Priorit"atsstufe aller Operatoren haben \fu{new} und
\fu{sizeof}, auf der n"achsten Stufe liegen \fu{->} und \fu{[]}.
Es h"atte keinen Sinn, diese Rangfolge zu "andern, denn die
Operanden von \fu{new} und \fu{sizeof} sind keine Ausdr"ucke,
sondern Typen, und es ist nicht m"oglich, den Operator \fu{->}
oder \fu{[]} auf einen Typ anzuwenden, soda"s dies schon von Haus
aus unterbunden werden sollte. Die beiden un"aren Operatoren
\fu{*} und \fu{\&} liegen um eine Priorit"atsstufe tiefer,
gemeinsam mit allen anderen un"aren Operatoren. Eine andere Stufe
f"ur diese Operatoren w"are nicht angebracht, aus Gr"unden, die
ich schon genannt habe.

Ein in gewisser Hinsicht au"sergew"ohnlicher Operator---gemessen
an seinen Gegenst"ucken aus C und \pascal ---ist der Operator
\fu{->}. Er wird benutzt, um auf Elemente einer Struktur
zuzugreifen. Derselbe Operator wird in C und \pascal\ durch den
Punkt (\fu{.}) symbolisiert. In C jedoch gibt es zus"atzlich den
Operator~\fu{->}. Er wird benutzt, um auf Elemente aus einem
Zeiger auf eine Struktur zuzugreifen---in \pascal\ mu"s hierf"ur
zuerst der Zeiger dereferenziert werden. Der SEL-Operator ist nun
insofern au"senstehend, als er die Funktionalit"at des
Punktoperators und des Pfeiloperators aus C vereinigt. So
erm"oglicht er es, angewendet auf eine Struktur, auf ein Element
daraus zuzugreifen. Angewendet auf einen Zeiger auf eine Struktur
dereferenziert er ersteren und greift dann erst auf das jeweilige
Element zu. Diese Vereinigung macht insofern Sinn, als da"s
keine Funktionalit"at verlorengeht, denn der Punktoperator
hat, angewendet auf einen Zeiger, keine Funktion. Auch das Design
eines Interpreters bzw. eines Compilers wird dadurch nicht
erschwert, denn diese wissen, wenn sie auf den Pfeiloperator
treffen, mit welchem Typ gearbeitet wird. Ein Compiler m"u"ste
also nur, tr"afe er auf einen Pfeiloperator, dessen linker Operand
ein Zeiger ist, den Code zur Dereferenzierung einbauen.

\subsubsection{Der Zuweisungsoperator}

Der Zuweisungsoperator von SEL ist ein Sonderfall unter den
Operatoren, aus mehreren Gr"unden. Der erste Grund ist, da"s es
in manchen Programmiersprachen, darunter auch \pascal , gar
keinen Zuweisungsoperator in diesem Sinne gibt. In \pascal\ ist
eine Zuweisung eine Anweisung und kein Ausdruck. Das macht es in
\pascal\ auf der einen Seite schwerer, den Zuweisungsoperator
vom Compiler bzw. Interpreter unentdeckt mit dem
Vergleichsoperator zu verwechseln, denn der Zuweisungsoperator
ist in Ausdr"ucken nicht erlaubt und Ausdr"ucke sind in \pascal\
keine g"ultigen Anweisungen. Andererseits jedoch ist es nicht
m"oglich, mehreren Variablen in einer Anweisung denselben Wert
zuzuweisen, wie man es von C oder auch aus der Mathematik gewohnt
ist (\fu{x=y=0}). Au"serdem ist es nicht m"oglich, z.~B. in einer
\rw{if}-Anweisung gleichzeitig eine Zuweisung vorzunehmen und den
Wert der zuzuweisenden Variablen zu "uberpr"ufen (z.~B.
\hbox{\fu{if ((x=y) == 3)}}). Manche legen das sogar als Vorteil
aus, weil es eine unsch"one Programmierung erschwert, und
insofern gebe ich ihnen Recht, doch ich pers"onlich finde solche
Konstrukte in manchen Situationen doch sehr praktisch, und da der
Aufwand f"ur einen Zuweisungsoperator etwa dem einer
Zuweisungsanweisung gleichkommt, oder sogar noch geringer ist,
habe ich ihn in SEL eingebaut.

Der zweite Grund f"ur die Besonderheit des \fu{=}-Operators ist,
da"s er mit allen Typen zusammenarbeitet. Anders als der
Vergleichsoperator, der auf die Grundtypen beschr"ankt ist, kann
der Zuweisunsoperator eine Variable jeden beliebigen Typs einer
anderen zuweisen---Typgleichheit vorausgesetzt. Ich habe diese
"`Allm"achtigkeit"' des Operators deswegen in SEL erm"oglicht, weil
es leicht ist, das in einen Interpreter einzubauen---im Fall der
Konzeption, die wir im zweiten Kapitel besprechen werden, sogar
noch leichter, als sich auf die Grundtypen zu beschr"anken. Auch
f"ur einen Compiler ist das nicht weiter schwer: Den Code f"ur die
Grundtypen kann der Compiler direkt einf"ugen, f"ur Strukturen
oder Arrays k"onnte eine Funktion aus der Laufzeitumgebung
aufgerufen werden. Hier dr"angt sich nat"urlich die Frage auf,
warum ich dieses Konzept nicht auf den Vergleichsoperator
\fu{==} ausgeweitet habe. Die Antwort ist einfach: H"atte ich das
getan, w"are es unausweichlich gewesen, das auch f"ur die
Operatoren \fu{<}, \fu{>}, \fu{<=}, \fu{>=} und \fu{!=} zu tun.
Mit letzterem h"atte ich keine Probleme gehabt, doch wann ist
eine Struktur, die die Adresse einer Person enth"alt, gr"o"ser
als eine andere desselben Typs? Byte-f"ur-Byte-Vergleich w"are
hier sicher nicht angebracht.\footnote{In C++ kann dieses Problem
gel"ost werden, weil es dort m"oglich ist, Operatoren zu
"uberladen, d. h., da"s der Benutzer imstande ist, eine Funktion
zu schreiben, die zwei Strukturen (beispielsweise des Typs
\fu{Adresse}) vergleicht, die dann anstatt des
Vergleichsoperators aufgerufen wird.}

\subsubsection{Der Typecasting-Operator}

Der Typecasting-Operator \fu{:} wird benutzt, um einen Wert
explizit in einen anderen Typ umzuwandeln. Er wird genauer in
Kapitel \ref{casts} beschrieben. Seine Priorit"at ist mit der der
beiden Operatoren \fu{->} und \fu{[]} "aquivalent, erstens, weil
er ein un"arer Operator ist und zweitens, weil es eher selten
vorkommt, da"s gr"o"sere Ausdr"ucke gecastet werden m"ussen. Wenn
n"otig, wird dies sowieso von SEL "ubernommen.

\subsection{Konstanten}
\label{const}

Grunds"atzlich gibt es in SEL vier verschiedene Arten von
Konstanten: \tchar -Konstanten, \tint -Konstanten,
\tfloat -Konstanten und String-Konstanten. Eine \tchar -Konstante
ist ein unter einfache Anf"uhrungsstriche gesetztes Zeichen bzw.
Steuerzeichen und hat---es ist nicht anders zu erwarten---den Typ
\tchar . Eine \tint -Konstante ist eine simple Abfolge von
Ziffern. Zahlenkonstanten im lexikalischen Sinn sind in SEL immer
positiv---negative Zahlen macht erst der un"are Operator \fu{-}
m"oglich, und insofern ergibt sich daraus kein Nachteil. Der
Grund, warum ich Konstanten nicht auch auf negative Zahlen
ausgeweitet habe ist der, da"s sich die lexikalische
Analyse\footnote{siehe Kapitel \ref{scanner}}
dadurch sehr erschweren w"urde. Nehmen wir als Beispiel den
Ausdruck \fu{7-4}. Der Scanner erkennt als erstes die
\tint -Konstante \fu{7}. Danach m"u"ste ein Operator folgen oder
der Ausdruck zu Ende sein, was aber der Scanner nicht wissen
kann, da er vom Parser in dieser Hinsicht abgekapselt ist. Aus
diesem Grund w"urde er anstatt den Operator \fu{-} zu erkennen
eine weitere \tint -Konstante, n"amlich \fu{-4} lesen, was
nat"urlich falsch ist. Das Gesagte gilt auch f"ur
\tfloat -Konstanten.

String-Konstanten sind in gewisser Hinsicht etwas Besonderes,
n"amlich in Anbetracht der Tatsache, da"s es einen Stringtyp in
SEL gar nicht gibt. Normalerweise werden Strings als Arrays von
\tchar s behandelt, die durch ein Byte mit dem Wert \fu{0}
abgeschlossen sind, es w"are allerdings nicht schwer, Funktionen
zu programmieren, die \pascal -Strings verarbeiten. Da jedoch die
erstere Stringart die verbreitetere ist, wollen wir bei dieser
bleiben. Wie wir wissen, kann ein Array in SEL in vielen
Situationen wie ein Zeiger behandelt werden. Was also w"are
falsch daran, einen String als Pointer-Konstante zu sehen?
Tats"achlich wird das nicht nur in SEL, sondern auch in C so
praktiziert, von wo ich dieses Konzept "`importiert"' habe. Dieser
Pointer zeigt auf einen Speicherbereich, in dem der eigentliche
String gespeichert ist. Nat"urlich mu"s dieser Speicherbereich
w"ahrend der gesamten Programmausf"uhrung an derselben
Speicheradresse bleiben, und nicht irgendeiner L"oschaktion zum
Opfer fallen, soda"s die Pointerkonstante im ganzen
Programm zu jeder Zeit "`aktuell"' ist.\footnote{In Wirklichkeit
mu"s der String nicht w"ahrend der gesamten Programmausf"uhrung
im Speicher sein, sondern erst nach seinem ersten Auftreten im
Code, doch der Einfachheit halber wird er die ganze Zeit im
Speicher gehalten.}

\subsection{Variablen}
\label{variables}

In SEL k"onnen zwei Arten von Variablen unterschieden werden, die
ihrerseits in jeweils zwei Untertypen unterteilt werden k"onnen.
Lokale Variablen sind nur innerhalb einer bestimmten Funktion
bekannt und existieren f"ur jede Instanz dieser Funktion
getrennt.\footnote{Mehrere Instanzen einer Funktion sind dann
m"oglich, wenn sich diese direkt oder indirekt selbst aufruft.}
Die Gruppe der lokalen Variablen werden weiter unterteilt in
Parameter und direkte Funktionsvariablen.\footnote{Eine bessere
Bezeichnung als "`direkte Funktionsvariablen"' habe ich nicht
gefunden.} Die Werte der Parameter werden von der aufrufenden
Funktion bestimmt, wohingegen die direkten Funktionsvariablen der
Funktion alleine unterstehen. Lokale Variablen werden auf dem
Stack gespeichert. Wird eine Funktion aufgerufen, werden die
Parameter von der aufrufenden Funktion auf den Stack gelegt und
anschlie"send zum Code der Funktion
gesprungen.\footnote{Nat"urlich wird auch die R"ucksprungadresse
auf dem Stack gespeichert} Die Funktion ihrerseits legt die
direkten Funktionsvariablen auf den Stack und f"uhrt dann ihre
Anweisungen aus. Wird sie beendet, entfernt sie sowohl
direkte Funktionsvariablen als auch Parameter vom Stack
und springt dann zur aufrufenden Funktion zur"uck. Das Konzept,
den Stack als Speicher f"ur lokale Variablen zu benutzen macht
es m"oglich, da"s bei einer rekursiven Funktion jede Instanz ihre
eigenen Variablen zur Verf"ugung hat.

Globale Variablen werden in SEL unterteilt in modulglobale und
programmglobale Variablen. Erstere sind nur in der Quelldatei
(Modul) bekannt, in der sie definiert sind, wohingegen letztere
auch in anderen Modulen benutzt werden k"onnen. Sowohl modul- als
auch programmglobale Variablen werden im statischen Speicher
aufbewahrt. Daf"ur, da"s modulglobale Variablen nur in ihrem
jeweiligen Modul referenziert werden, haben der Interpreter bzw.
der Compiler oder der Linker zu sorgen.

Ich glaube, "uber die Verwendung von Variablen in Ausdr"ucken
mu"s ich hier nicht viele Worte verlieren: das geschieht ganz
einfach---wie in fast allen Programmiersprachen "ublich---durch
die Benutzung des Variablennamens.

\subsection{Funktionen}
\label{exprfuncs}

Die Aufrufkonvention f"ur Funktionen ist einfach: Der
Funktionsname gefolgt von in Klammern gesetzten Parametern,
jeweils durch einen Beistrich getrennt, bilden einen Funktionsaufruf.
Hat die Funktion keine Parameter, sind trotzdem "offnende und
schlie"sende Klammern anzugeben, aus mehreren Gr"unden. Erstens
kann der Interpreter bzw. Compiler daran ohne nachzusehen
erkennen, da"s es sich um einen Funktionsaufruf handelt und nicht
um die Verwendung einer Variablen. Zweitens ist dies in Hinblick
auf die evtl. Erweiterung der Sprache durch
Funktionszeiger\footnote{siehe Kapitel \ref{langext}}, also
Zeigervariablen, die die Adressen von Funktionen aufnehmen,
n"otig. In diesem Fall signalisiert das Klammerpaar, da"s wir
vorhaben, diese Funktion aufzurufen und nicht, sie als Zeiger zu
verwenden.

Die Anzahl und die Typen der Funktionsparameter sind nat"urlich
einzuhalten. Die Regelung bei letzteren ist nicht so streng:
Stimmen die Typen nicht "uberein, versucht SEL mittels impliziten
Typecastings\footnote{siehe Kapitel \ref{casts}}, diese in die
richtigen umzuwandeln, was jedoch nicht unbedingt gelingen mu"s,
da sich der Interpreter bzw. der Compiler an bestimmte Regeln zur
Typumwandlung halten mu"s.

\section{Typen}
\label{types}

In diesem Abschnitt m"ochte ich das Typkonzept von SEL erkl"aren.
Zun"achst stellt sich beim Design einer Programmiersprache die
Frage, welches Typkonzept "uberhaupt zur Anwendung kommen soll.
So gibt es z.~B. in \basic\ einige wenige Typen, aus denen der
Benutzer ausw"ahlen kann, es ist dort jedoch nicht m"oglich,
selber Typen zu definieren. Zwar k"onnen Arrays definiert werden,
doch ist ein Array in \basic\ nichts anderes als eine Menge von
Variablen, deren Namen sich jeweils durch eine Nummer
unterscheiden---ein Array als Ganzes zu behandeln, ist in \basic\
nicht m"oglich.

Dann gibt es Programmiersprachen---\dbase\ ist
ein Beispiel hierf"ur---in denen Variablen zwar einen bestimmten
Typ haben, dieser sich aber jederzeit "andern kann. Dort ist es
z.~B. m"oglich, einer Variablen einen numerischen Wert zuzuweisen
und mit ihr Berechnungen anzustellen. Sp"ater kann man ihr bei
Bedarf einen String zuweisen und diesen am Bildschirm ausgeben.
Zwar bietet dieses Konzept auf den ersten Blick f"ur den Benutzer
Vorteile\footnote{Auch diese Vorteile sind nur scheinbar, denn
sie verf"uhren zu einem sehr schmutzigen Programmierstil} und es
ist f"ur eine reine Interpretersprache eine gute Wahl, doch einem
Compiler macht ein solches Konzept gro"se Probleme: In
Programmiersprachen wie C oder \pascal\ haben Variablen einen
vorgegebenen, konstanten Typ. Vom Compiler werden diese Variablen
an eine bestimmte Speicheradresse gelegt, die sich im Laufe des
Programms nicht ver"andert, was ohne weiteres m"oglich ist, da
sich ja die Speicheranforderungen der Variablen auch nicht
ver"andern. K"onnen sich aber die Typen der Variablen "andern,
"andern sich auch die Gr"o"sen der Variablen, d. h., da"s es
nicht m"oglich ist, sie an eine konstante Stelle im Speicher zu
setzen. W"ahrend ein C-Compiler, wenn er die Variable zu
referenzieren hat, ganz einfach ihre Speicheradresse einf"ugen
kann, mu"s ein compiliertes \dbase -Programm eine Tabelle
verwalten, in der alle Variablen und deren momentane Adressen und
Typen verzeichnet sind. Wird nun der Ausdruck \fu{x+1}
abgearbeitet, mu"s zun"achst in der Tabelle nachgesehen werden,
an welcher Adresse sich die Variable \fu{x} befindet. Zus"atzlich
mu"s auch gepr"uft werden, welchen Typ diese hat, da z.~B. eine
ganzzahlige Addition einen anderen Code verlangt, als eine
Addition mit Gleitkommazahlen. Erst jetzt kann die Addition
ausgef"uhrt werden. Es ist also nur verst"andlich, da"s Compiler
f"ur eine solche Sprache einen Code erzeugen, der zehnmal
langsamer\footnote{Die Zahl 10 mag wie eine "Ubertreibung
klingen, tats"achlich ist es nicht selten, da"s solche Codes noch
um einiges langsamer sind. So ben"otigt z.~B. ein Programm, da"s
einen QuickSort-Algorithmus ausf"uhrt in Clipper (einem
\dbase -Compiler) sage und schreibe 120 mal l"anger als ein
vergleichbares Programm in C} ist als der eines C- oder
\pascal -Compilers.

Eine weitere wichtige Kategorie ist die der typenlosen
Programmiersprachen---AWK f"allt teilweise darunter. Dort ist es
m"oglich, einer Variablen einen String zuzuweisen, der eine Zahl
enth"alt und in der n"achsten Zeile mit dieser Variable zu
rechnen als w"are sie numerischen Typs. Dieses Konzept ist
insofern unpraktisch, als man f"ur verschiedene Typen (obwohl
es diese eigentlich nicht gibt) auch verschiedene Operanden
braucht. Bei einem Ausdruck \fu{x+y} w"are also nicht klar, ob
die beiden Strings zu konkatenieren sind oder ob deren numerische
Werte addiert werden sollen. Auch in Bezug auf den Bau eines
Compilers ist dieses Konzept von Nachteil, aus "ahnlichen
wie im vorigen Absatz dargelegten Gr"unden.

Manche alte Programmiersprachen stellen als einzigen Typ das
Maschinenwort zur Verf"ugung, was einerseits die Gefahr birgt,
da"s in einer solchen Sprache geschriebene Programme sehr vom
Prozessor abh"angig sind und andererseits den Nachteil hat, da"s
viele Aufgaben, die der Interpreter bzw. Compiler ohne 
Effizienzverluste "ubernehmen k"onnte, vom Programmierer 
bestritten werden m"ussen.

Wie sich aus den vorigen Abs"atzen leicht herauslesen l"a"st,
habe ich mich f"ur das Typkonzept entschieden, das auch in C
oder \pascal\ verwendet wird. Dort mu"s jede Variable vor ihrer
Verwendung definiert werden, wobei gleichzeitig ihr Typ
festgelegt wird, der sich im Laufe des Programms nicht "andert.
Au"serdem ist es dem Benutzer m"oglich, aus den Basistypen der
Programmiersprache eigene Typen herzuleiten, wenn auch nicht
unbeschr"ankt.

\subsection{Elementare Typen}
\label{elemtypes}

Wie bereits gesagt, ben"otigt man zur Konstruktion von
komplizierteren Typen eine Menge an Basistypen, mit denen man
mittels der Konstrukte, die SEL bereitstellt, arbeiten kann.
Programmiersprachen wie C oder \pascal\ bieten eine relativ
gro"se Zahl solcher Basistypen an. So gibt es in C zwei
verschiedene Typen, die ein einzelnes Zeichen repr"asentieren,
drei Grundtypen f"ur ganze Zahlen, wobei von jedem dieser Typen
ein vorzeichenloser sowie ein vorzeichenbehafteter Typ
exisitiert, und mehrere Gleitkomma-Typen. Die Vielfalt in
\pascal\ ist "ahnlich gro"s, wobei zus"atzlich noch ein Typ f"ur
logische Werte exisitiert (\type{boolean}). Ich habe mich
entschieden, in SEL nur vier Elementartypen aufzunehmen:
\tchar , \tint , \tfloat\ und \tvoid .

Der Typ \tchar\ repr"asentiert ein einzelnes Zeichen und nimmt
dementsprechend nur ein Byte ein. Arrays aus \tchar s k"onnen
z.~B. als Strings interpretiert werden.

Ein \tint\ repr"asentiert eine ganze Zahl mit Vorzeichen und hat
in der Regel die Gr"o"se eines Maschinenwortes.\footnote{Eine
Ausnahme bildet MS-DOS, wo ein Maschinenwort nur 16 Bits lang
ist, d.~h. der Wertbereich w"urde von etwa -32000 bis 32000
reichen, was zwar f"ur viele Anwendungen reicht, aber ungen"ugend
ist, da es ja nur einen Typ f"ur ganze Zahlen gibt. Auf MS-DOS
Maschinen ist deshalb die Gr"o"se eines \tint\ zwei
Maschinenworte, also 32 Bits.} Zus"atzlich hat \tint\ die
Aufgabe, logische Werte zu repr"asentieren, wobei der Wert \fu{0}
einem "`logisch falsch"' und ein Wert ungleich \fu{0} "`logisch
wahr"' entspricht. Unter anderem hat diese Konzeption den
Vorteil, da"s ein Ausdruck des Typs \tint\ nicht nur den Wert des
\tint\ wiederspiegelt, sondern auch die Aussage, ob dieser
ungleich \fu{0} ist oder nicht. So ist es z.~B. m"oglich, mit
einer sehr kurzen \rw{if}-Anweisung zu pr"ufen, ob ein Zeiger ein
Nullzeiger ist. W"ahrend hierf"ur in \pascal\ die Anweisung
\fu{if p <> nil then} benutzt wird, kann in SEL einfach
\fu{if (p)} geschrieben werden. Zieht man es aus Gr"unden der
"Ubersichtlichkeit vor, trotzdem einen boolee'schen Typ zu haben,
kann man diesen in SEL problemlos mit den Direktiven

\begin{verbatim}
#define boolean        int
#define false          0
#define true           1
\end{verbatim}

\noindent simulieren.\footnote{Die Betonung liegt auf "`simulieren"',
da wir uns mit diesen Direktiven keinen eigenen Typ schaffen sondern
lediglich \fu{boolean} zu einem Pseudonym f"ur \rw{int} machen.}

Der Typ \tfloat\ repr"asentiert eine vorzeichenbehaftete
Gleitkommazahl. Ihre Gr"o"se im Speicher h"angt vom jeweiligen
Prozessor ab.

Einen Spezialfall unter den Typen stellt \tvoid\ dar: Er
repr"asentiert gar nichts! F"ur sein Dasein sprechen zwei
Gr"unde. Erstens macht er es m"oglich, Funktionen ohne
R"uckgabewert zu konstruieren, also Konstrukte, die mit den
Prozeduren in \pascal\ oder \modula\ vergleichbar sind. Au"serdem
ist er der Grundtyp f"ur typenlose Zeiger, also Zeiger, die auf
etwas nicht Bestimmtes zeigen. Verwendung finden solche Zeiger
z.~B. in Funktionen zur Speicherallokation. Diese geben einen
Zeiger auf einen Speicherbereich bestimmter Gr"o"se zur"uck,
wobei aber nicht bekannt ist, f"ur welchen Typ dieser Bereich
reserviert wird.

\subsection{Arrays}

Ein Array ist ein Zusammenschlu"s mehrerer Variablen desselben
Typs, die im Speicher sinnvollerweise aufeinander folgen. Ich
habe diesen Typ in SEL aufgenommen, weil er erstens nicht schwer
zu implementieren und zweitens in sehr vielen Situationen
unverzichtbar ist. Auffallend an der Schreibweise des Arraytyps
ist, da"s die Dimension des Arrays dem Grundtyp nicht nach- sondern
vorgestellt sind. Gem"a"s unserem Grundsatz, Symbole anstatt
Schl"usselw"ortern zu verwenden, wo es nicht st"ort, habe ich hier
auf Ausschm"uckungen wie das in \pascal\ verwendete \fu{array
... of ...} verzichtet und statt dessen eine in eckige Klammern
geschlossene Array-Dimension benutzt. Auf ein Angeben der
Grenzen, wie in \pascal\ "ublich, habe ich verzichtet, das erste
Element besitzt den Index \fu{0}. Diese Vereinfachung hat zwei
Gr"unde: Einerseits ist es nicht schwer, auf speziell indizierte
Arrays zu verzichten, andererseits w"urde sich das auf die
Konsistenz in Bezug auf die Verwendung von Zeigern anstatt Arrays
auswirken: Angenommen, wir h"atten ein Array \fu{a : [5..10]int}.
In SEL ist es nun m"oglich, dieses Array als Zeigervariable zu
betrachten und so zu behandeln, d. h., wir k"onnten mit ihm
Zeigerarithmetik betreiben. Per definitionem ist ein Zeiger
\fu{p} gleich \fu{p + 0}, was bedeutet, da"s wir mit \fu{*(a + 0)} auf das
erste Element des Arrays zugreifen k"onnen. Wir k"onnen auf
dieses Element auch direkt "uber das Array zugreifen, also mit
\fu{a[5]}. Ich denke, hier zeigt sich schon, worauf ich
hinauswill: Die Zahlen, die zur Indizierung verwendet werden,
stimmen nicht "uberein. Um das zu umgehen, m"u"ste man entweder
auf die Gleichsetzung von Array und Zeiger verzichten, was ich
nicht wollte, oder auch in die Zeiger eine komplizierte
Nummerierung einbauen, was mir und sicher auch dem Benutzer der
Programmiersprache zu umst"andlich ist.

Zur"uck zu der Dimension des Arrays, die, wie oben erw"ahnt, dem
Grundtyp vorangestellt ist. Der Grund hierf"ur ist in erster
Linie beim Typecasting\footnote{siehe Kapitel \ref{casts}}
von SEL zu suchen. Nehmen wir als Beispiel den Ausdruck
\fu{a : *int[10]}. Hier w"are es unklar, ob die eckige Klammer noch
zum Typ geh"ort, oder ob versucht wird, auf das zehnte Element
des Zeigers (der als Array behandelt wird) zuzugreifen.
Nat"urlich kann auch dieses Problem durch Klammerung umgangen
werden, doch warum etwas komplizierter machen, als es ist? Die
vorangestellte Arraydimension in SEL hat n"amlich noch einen
Vorteil: Es ist ohne eine Rangordnung oder Klammerung m"oglich,
einen Typ eindeutig zu bestimmen. Als Beispiel soll der C-Typ
\fu{*int[10]} dienen. Ohne den beiden Operatoren eine
unterschiedliche Priorit"at zuzuweisen, ist es unm"oglich zu
sagen, ob der Typ nun \fu{*(int[10])} oder \fu{(*int)[10]} ist.
In SEL ist die Sache eindeutig: \fu{[10]*int} kann nur als Array
aus zehn Zeigern auf \tint\ gedeutet werden.

\subsection{Zeiger}

Zeiger stellen ein wichtiges Konstrukt dar, mit dessen Hilfe es
m"oglich ist, sehr komplexe, verkettete und verzweigte
Datenstrukturen zu bilden, die ohne sie nur sehr schwer zu
verwirklichen w"aren. Der Bau eines Interpreters oder gar eines
Compilers z.~B. w"are ohne die Verwendung von Zeigern kaum zu
bewerkstelligen. Aus diesem Grund habe ich Zeiger in SEL
eingebaut, was f"ur eine Interpretersprache eigentlich sehr
ungew"ohnlich ist.

Sehr eng in Verbindung mit Zeigern steht in SEL der Operator
\fu{*}, der einerseits in der Typdefinition eines Zeigers
verwendet wird und andererseits benutzt wird, einen Zeiger zu
dereferenzieren. Letzteres geschieht, indem dem Zeiger der
Sternoperator vorangestellt wird, anders als in \pascal , wo der
Dereferenzierungsoperator als Suffix verwendet wird. Der Grund,
warum in SEL der Operator vorangestellt wird, ist die Tatsache,
da"s er dasselbe Symbol hat wie der Multiplikationsoperator, und
der Parser\footnote{siehe Kapitel \ref{parser}} nicht w"u"ste,
wie etwa der Ausdruck \fu{p * - x} gemeint ist. Einerseits k"onnte es
die Differenz zwischen der Variablen, auf die \fu{p} zeigt und
\fu{x} sein, andererseits aber auch \fu{p} multipliziert mit
\fu{-x}. Nat"urlich h"atte ich auch ein anderes Symbol f"ur die
Dereferenzierung verwenden k"onnen, doch das wollte ich nicht,
weil ich den Operator \fu{*} von C gewohnt bin, au"serdem fand
ich es schon immer ein wenig verwirrend, da"s bei der
Typdeklaration der Operator vorangestellt, bei der
Dereferenzierung aber nachgestellt sein mu"s.

Der zweite wichtige Operator in Verbindung mit Zeigern ist
\fu{\&}. Wird er einer Variablen vorangestellt, dann ergibt er
einen Zeiger auf diese Variable. Mit seiner Hilfe ist es also
m"oglich, einen Zeiger auf eine bestimmte Variable zeigen zu
lassen. Verst"andlicherweise ist es nicht m"oglich, \fu{\&} in
Verbindung mit einer Konstante zu benutzen, denn eine Konstante
hat keinen bestimmten Platz im Speicher, sie wird vom Interpreter
direkt in einen Wert umgewandelt und von einem Compiler direkt in
den Code eingef"ugt. Zu verwechseln ist dies jedoch nicht mit den
String-Konstanten\footnote{siehe Kapitel \ref{const}}, die
eigentlich nicht konstante Strings sind, sondern konstante Zeiger
auf einen String, der an einer bestimmten Stelle im Speicher
liegt. Doch zur"uck zum Operator \fu{\&}: Sein zweiter
Verwendungszweck ist die Kennzeichnung eines Parameters einer
Funktion, der {\em Call-By-Reference} "ubergeben werden soll. In
\pascal\ z.~B. "ubernimmt diese Aufgabe das Schl"usselwort
\rw{var}, in C gibt es etwas Derartiges nicht, erst in C++ wurden
Referenzparameter eingef"uhrt. Auch dort werden sie mit \fu{\&}
gekennzeichnet. In einer Programmiersprache, in der Zeiger zur
Verf"ugung stehen, sind Referenzparameter nicht zwingend
notwendig: Es k"onnten Zeiger auf die betreffenden Variablen
"ubergeben werden, wie es in C auch praktiziert wird. Ich habe
mich aus zwei Gr"unden entschlossen, Referenzparameter doch in
SEL aufzunehmen. Erstens k"onnen sie in manchen Situationen
Programme um sehr viel "ubersichtlicher und leichter wartbar
machen. Zweitens, und das war f"ur mich wohl der gewichtigere
Grund, arbeitet der im n"achsten Kapitel theoretisch beschriebene
Interpreter intern ohnedies mit Referenzen, soda"s es keine
Schwierigkeit war, Referenzparameter einzubauen.

\subsubsection{Dynamische Variablen}

Eng verbunden mit den Zeigern ist das Konzept der dynamischen
Variablen. Dynamisch deshalb, weil sie erzeugt werden k"onnen,
wenn man sie braucht und aus dem Speicher entfernt, wenn sie
"uberfl"ussig werden. Statische Variablen belegen
Speicherbereiche, die vom Compiler bzw. Interpreter im voraus
bestimmt werden, wohingegen der Speicherbereich f"ur eine
dynamische Variable erst w"ahrend der Programmausf"uhrung
bestimmt werden kann. Das impliziert nat"urlich, da"s dem
Benutzer geeignete Mittel in die Hand gegeben werden m"ussen, um
dynamische Variablen zu erzeugen. Grunds"atzlich gibt es hierf"ur
zwei M"oglichkeiten. Entweder, man baut das Konzept der
dynamischen Speicherverwaltung nicht fest in die
Programmiersprache ein, sondern "uberl"a"st es dem Benutzer bzw.
dem Designer der Laufzeitbibliothek, geeignete Funktionen zu
implementieren, wie es in C der Fall ist, oder man
ber"ucksichtigt dynamische Speicherverwaltung schon beim Design
der Programmiersprache und stellt Konstrukte zur Verf"ugung, um
dynamische Variablen zu erzeugen, wie z.~B. in C++.

Ich habe mich f"ur den zweiten L"osungsansatz entschieden, weil
ich finde, da"s das Konzept der dynamischen Speicherverwaltung zu
wichtig ist, um es beim Design einer Programmiersprache nicht
miteinzubeziehen. Ein daraus resultierender Vorteil ist, da"s
Programme leichter lesbar werden und es leicher ist, sauberer zu
programmieren, als z.~B. in C.\footnote{Was nicht hei"sen will,
da"s es in C nicht m"oglich ist, sauber zu programmieren, es ist
nur nicht so leicht wie etwa in C++} Angelehnt an das Vorbild C++
werden in SEL zur dynamischen Speicherverwaltung die beiden
un"aren Operatoren \rw{new} und \rw{delete}\footnote{siehe
Kapitel \ref{memmgrops}} verwendet. Der Operand von \rw{new} ist ein
Typ, also im Gegensatz zu denen der meisten anderen Operatoren
kein Ausdruck. Der Ergebnistyp ist ein Zeiger auf den
"ubergebenen Typ, der auf den neu allokierten Speicherbereich
zeigt. Dieser Bereich kann nun so lange wie n"otig benutzt
werden. Wird er nicht mehr verwendet, sollte er mit dem Operator
\rw{delete} freigegeben werden, um sp"ater anderen dynamischen
Variablen zur Verf"ugung zu stehen. Der Grund, warum ich aus
\rw{new} und \rw{delete} keine Anweisungen gemacht
habe---besonders bei \rw{delete} w"urde sich das anbieten---ist,
da"s alle anderen Anweisungen in SEL zur Flu"skontrolle im
Programm verwendet werden, soda"s dies gewisserma"sen ein Bruch
in der Konsistenz der Sprache w"are.

\subsection{Strukturen}
\label{structs}

Eine Struktur\footnote{Der Name "`Struktur"' stammt von C, in
\pascal\ wird f"ur dasselbe der Name "`Record"' verwendet} ist
der Zusammenschlu"s mehrerer Variablen, auch verschiedenen Typs,
unter einem Namen. Sie bilden zusammen mit Zeigern und
dynamischen Variablen die Basis f"ur weit komplexere, verzweigte
und verkettete Datenstrukturen. "Uber ihre Definition m"ochte ich
nicht viele Worte verlieren, da sie einfach und zweckm"a"sig ist
und es au"serdem nicht viele andere M"oglichkeiten gibt,
Strukturen zu definieren.  Was jedoch viel interessanter ist, ist
die Benutzung von Zeigern auf eine Struktur in der Struktur
selbst. Ein Beispiel:

\begin{verbatim}
struct List
  i    : int;
  Next : *List;
end
\end{verbatim}

Diese Struktur ist eine m"ogliche Definition einer einfach
verketteten Liste. Neben einem oder mehreren Werten (in diesem
Fall ein \tint ) enth"alt sie einen Zeiger auf eine Struktur
ihres eigenen Typs. Einen Interpreter oder Compiler, der mit nur
einem Pass arbeitet und diese Struktur akzeptiert, zu schreiben,
ist noch leicht: Man k"onnte sich schon beim Lesen des auf das
Schl"usselwort \rw{struct} folgenden Bezeichners \fu{List}
notieren, da"s es diesen Typ von Struktur gibt. Beim Zeiger
\fu{Next} kann diese Notiz benutzt werden, um seinen Typ
festzulegen. Beim Lesen des Schl"usselwortes \rw{end} wird die
Notitz anschlie"send zu einer kompletten Definition erg"anzt. Was
tun wir jedoch, wenn es zwei Strukturen gibt, die jeweils einen
Zeiger auf die andere Struktur enthalten, wie es nicht selten
vorkommt? Eine der beiden Strukturen mu"s zuerst gelesen werden,
und dort treffen wir auf einen unbekannten Typen. In C wird dieses
Problem dadurch gel"ost, da"s man dem Compiler vor der Definition
der beiden Strukturen einen Hinweis darauf gibt, da"s diese
Strukturen irgendwo einmal definiert werden, soda"s er schon vor
der Definition seine Notiz machen kann. In SEL jedoch k"onnen wir
dieses Problem weit einfacher l"osen: Da es nicht m"oglich ist,
einem Typ in SEL einen anderen Namen zu geben, steht, wenn ein
unbekannter Typ gelesen wird, zweifelsfrei fest, da"s es sich
dabei um eine Struktur handeln mu"s, soda"s unser Compiler bzw.
Interpreter bereits hier die Notiz vornehmen kann.

\subsection{Typecasting}
\label{casts}

Bei einer Gegen"uberstellung der Typkonzepte von C und \pascal\
f"allt auf, da"s \pascal\ in der Typ"uberpr"ufung weit strenger
ist als C. So ist es in C ohne weiteres m"oglich, einem \tchar\
einen numerischen Wert zuzuweisen, wohingegen in \pascal\ zu
diesem Zweck ein eigener Typ existiert. In der Sprache \pascal\
ist eine Umwandlung eines Typs in einen anderen gar nicht
vorgesehen,\footnote{au"ser in F"allen, wo die beiden Typen zur
gleichen Gruppe geh"oren, wie z.~B. \tint\ und \type{longint}} in
C wird sie in vielen F"allen sogar vom Compiler "ubernommen, an
anderen Stellen kann sie der Benutzer mittels Typecasting
durchf"uhren. Der Grund, warum ich in SEL ein flexibles
Typkonzept eingebaut und nicht die strenge "Uberpr"ufung von
\pascal\ "ubernommen habe, ist, da"s es bei vier Basistypen sehr
schwer ist, ohne Typumwandlung auszukommen. Au"serdem k"onnen
Programme durch implizites Typecasting um einiges k"urzer und
"ubersichtlicher werden, man denke nur an die Umwandlung zwischen
\tchar\ und \tint , wof"ur in \pascal\ eigene Funktionen
notwendig sind.

Um aber nicht in einem vollst"andigen Chaos zu enden, darf nicht
wahllos jeder Typ implizit in jeden anderen Typen umgewandelt
werden, nur damit im Programm ohne Fehlermeldung fortgefahren
werden kann. Die Regeln, die in SEL f"ur implizites sowie
explizites Typecasting gelten, sind in \cite{selkr}, Seite 24
zusammengefa"st.

Der Grund, warum ich nicht das in C "ubliche Typecasting, einem
Ausdruck den neuen Typ in Klammern voranzustellen, "ubernommen
habe, ist der, da"s ich einfach etwas Neues ausprobieren wollte,
und das Typecasting mittels des \fu{:}-Operators exzellent zu den
"ubrigen Konstrukten in SEL (Variablen- und Strukturdefinition,
Parameterdefinition in Funktionen u.~s.~w.) pa"st. Es zeigt sich,
da"s Programme bei richtiger Verwendung des Typecasting-Operators
weit "ubersichtlicher sind als in C, meine Wahl hat sich also in
keiner Weise als Fehlgriff erwiesen.

\subsection{Mengen, Listen, B\protect{"aume} etc.}

Komplexere Typen, wie z.~B. Mengen, Listen oder B"aume sind in
SEL nicht vorhanden, aus guten Gr"unden. Mengen einerseits
werden relativ selten verwendet und lassen sich leicht
vermeiden. Au"serdem sind sie sehr leicht mit den Mitteln von
SEL zu implementieren, soda"s das Einbeziehen in SEL nur
"au"serst wenig Sinn h"atte. Listen und B"aume hingegen werden
zwar oft verwendet, sind allerdings so komplex und existieren in
einer derartig gro"sen Anzahl von Variationen, da"s es mit
einfachen Mitteln kaum m"oglich ist, sie vern"unftig zu
unterst"utzen.

\section{Bezeichner und Schl\protect{"usselw\protect{"orter}}}

Wie in vielen Programmiersprachen "ublich, sind auch in SEL
Schl"usselw"orter reserviert, d.~h., sie k"onnen nicht als
Bezeichner verwendet werden, aus zwei Gr"unden. Erstens wird
dadurch die Programmierung eines Interpreters bzw. Compilers
erheblich erschwert, und zweitens w"aren dadurch sehr
verwirrende Konstruktionen wie z.~B.

\begin{verbatim}
for (for = to to step step end)
  ...
end
\end{verbatim}

\noindent m"oglich, was nicht im Sinn des Designers sein kann.

Schl"usselw"orter und Bezeichner sind in SEL case-sensitiv,
d.~h., da"s zwei W"orter, die sich nur dadurch unterscheiden,
da"s ein oder mehrere Buchstaben statt gro"s klein geschrieben
sind oder umgekehrt, als unterschiedlich betrachtet werden. Das
hat drei Vorteile: Erstens wird dadurch die Entwicklung der
lexikalischen Analyse\footnote{siehe Kapitel \ref{scanner}}
ein wenig einfacher, zweitens k"onnen dadurch verschiedene Arten
von Bezeichnern (Bezeichner f"ur Funktionen, Variablen, Makros
u.~s.~w.) durch verschiedene Schreibweise gekennzeichnet werden.
Drittens wird durch eine vorgegebene Schreibweise bei
unerfahrenen oder schlampigen Programmierern die
"Ubersichtlichkeit der Quelltexte gesteigert.

Die Syntax der Bezeichner ist relativ simpel: Sie bestehen aus
einem Buchstaben oder Unterstrich, gefolgt von einer beliebigen
Anzahl von Buchstaben, Ziffern oder Unterstrichen. Der Grund,
warum Bezeichner nicht mit Ziffern beginnen d"urfen ist der, da"s
dadurch eine Eindeutigkeit der Sprache nicht mehr gew"ahrleistet
w"are. So k"onnte z.~B. \fu{12to} einerseits als Zahl \fu{12}
gefolgt vom Schl"usselwort \rw{to}, andererseits aber auch als
Bezeichner gedeutet werden. Dieser Konflikt lie"se sich zwar
durch Verwendung von Leerzeichen vermeiden, w"urde aber die
lexikalische Analyse ungleich schwerer machen. Warum ich nicht
auch andere Symbole in Bezeichnern erlaubt habe, ist einfach: Je
weniger Symbole in Bezeichnern verwendet werden d"urfen, desto mehr
hat man f"ur andere Zwecke frei.
                            
\section{Weiterf\protect{"uhrende} Projekte}
\label{langext}

Keine Programmiersprache ist perfekt, und schon gar nicht SEL,
f"ur dessen jetzige Form ich bewu"st die Versionsnummer 0.3
gew"ahlt habe. Leider war es mir aus Zeit- und Platzgr"unden
nicht m"oglich, auch nur einen Teil der von mir geplanten
Erweiterungen in diese Arbeit einzubringen, soda"s ich sie in
diesem Kapitel kurz zusammenfassen m"ochte, einerseits, um zu
zeigen, was SEL noch alles fehlt und andererseits, um den Leser
anzuregen, sich mit der Materie eingehender auseinanderzusetzen.

\begin{description}
\item[Zeiger auf Funktionen] stellen ein sehr praktisches
Sprachkonzept dar. Sie erm"oglichen es z.~B., in einer
graphischen Benutzeroberfl"ache, wie etwa Windows NT oder meinen
Schani-Tools, jedem Bildschirmfenster eine Funktion zuzuweisen,
die die Verwaltung desselben "ubernimmt. Sprachlich gesehen liegt
die einzige Schwierigkeit bei der Typangabe eines solchen
Funktionszeigers, denn der Aufruf einer Funktion durch einen
Zeiger kann genauso wie bei einer "`richtigen"' Funktion durch
den Namen der Zeigervariablen, gefolgt von einer in Klammern
gesetzten Parameterliste erfolgen, und die Adresse einer Funktion
kann mittels deren Namen, ohne nachfolgende Parameterliste in
Klammern, ermittelt werden. Doch auch der Typ f"ur einen
Funktionszeiger kann relativ leicht ohne Zweideutigkeiten in die
Grammatik eingebaut werden. So k"onnte etwa \fu{char(int,float)}
f"ur einen Zeiger auf eine Funktion, die einen \tint\ und einen
\tfloat\ als Parameter erwartet und einen \tchar\ zur"uckliefert,
stehen.\footnote{Tats"achlich f"uhrt eine Erweiterung wie oben
beschrieben zu einem kleinen Konflikt in der Grammatik, der
jedoch leicht beseitigt werden kann.}

\item[Variante Strukturen] sind Strukturen, die zwar mehrere
Elemente haben, diese jedoch denselben Speicherbereich
einnehmen. So etwas ist immer dann von Vorteil, wenn eine
Struktur ben"otigt wird, die Objekte beschreibt, die zwar in
gewisser Hinsicht "ahnlich sind, jedoch andere Parameter haben.
Ein Beispiel hierf"ur ist die in Kapitel \ref{intexprs}
besprochene Datenstruktur f"ur SEL-Ausdr"ucke.

\item[Makros mit Parametern] w"urden, wenn sie so flexibel
aufgebaut sind, da"s neue Befehle durch sie
"`vorget"auscht"' werden k"onnen, perfekt in eine Makrosprache
passen. Ein gutes Beispiel f"ur solche Makros sind die sog. User
Defines Commands (UDC's) von Clipper.

\item[Objektorientierte Programmierung] ist zwar ein edles Ziel,
jedoch nicht leicht zu verwirklichen.
\end{description}

\chapter{Der Interpreter}

Das Ziel dieses zweiten Kapitels ist es, f"ur die im ersten
Kapitel erkl"arte Programmiersprache SEL einen Interpreter zu
entwerfen. Der Titel dieser Arbeit deutet darauf hin, da"s dieser
Entwurf theoretisch ist, also in dieser Arbeit keine
Implementation eines Interpreters zu finden ist. Die Gr"unde
hierf"ur sind verschieden. Schon lange bevor ich "uberhaupt die
Idee zu dieser Fachbereichsarbeit hatte, wollte ich einen
Interpreter schreiben. Als ich mit Herrn Prof. Freiler "uber das
Thema dieser Arbeit diskutierte, legte er mir die Beschr"ankung
auf, da"s, wenn das Thema dieser Arbeit ein Programm w"are,
dieses in \pascal\ geschrieben sein m"u"ste. Da ich aber den
Interpreter in erster Linie f"ur mich und nicht f"ur diese Arbeit
programmieren wollte, hatte ich vor, ihn in C zu schreiben. Aus
diesem Grund einigte ich mich mit Herrn Prof. Freiler darauf,
da"s ich den Interpreter nur theoretisch, also unabh"angig von
irgendeiner Implementation, behandeln w"urde. Im Nachhinein
betrachtet bin ich froh "uber diese Entscheidung, da es kaum
m"oglich gewesen w"are, in einer im Umfang derma"sen beschr"ankten
Arbeit einen ganzen Interpreter und eine angemessene Erkl"arung
dazu unterzubringen.

Doch auch mit der Beschr"ankung auf eine theoretische Abhandlung
ist der Umfang dieser Arbeit zu gering, um ausf"uhrlich auf alle
Details einzugehen. Aus diesem Grund m"ochte ich im Gegensatz 
zum ersten Kapitel, in dem ich alle meine Entscheidungen 
begr"undet und evtl. auch andere L"osungsans"atze beschrieben 
habe, in diesem nur den Entwurf pr"asentieren und lediglich in
einigen seltenen F"allen andere m"ogliche L"osungen diskutieren.

Dieses Kapitel beschreibt alle wichtigen Algorithmen und
Datenstrukturen des Interpreters, geht jedoch nicht n"aher auf
grundlegende Datenstrukturen wie etwa dynamische Listen ein.
Sollte der Leser mit diesen Strukturen nicht vertraut sein,
empfehle ich ihm die Lekt"ure von \cite{sedge}, Kapitel 3--5 und
16.

\section{Der Begriff des Interpreters}
\label{whatisit}

Um ein Programm einer h"oheren Programmiersprache auf einem
Computer auszuf"uhren, kann man sich zweier verschiedener
Werkzeuge bedienen: Interpreter und Compiler. Der Begriff
"`Compiler"' wird vom Duden Informatik wie folgt definiert:

\begin{quotation}
Programm, das Programme aus einer Programmiersprache A in eine
Programmiersprache B "ubersetzt.\footnote{siehe \cite{duden},
Seite 745}
\end{quotation}

Das bedeutet, da"s ein Compiler ein Programm gar nicht direkt
ausf"uhrt, sondern es lediglich in eine andere Programmiersprache
"ubersetzt. In der Regel ist diese Sprache eine Assemler- oder
Maschinensprache, die direkt vom Computer "`verstanden"' wird. Im
Gegensatz dazu ist ein Interpreter

\begin{quotation}
Ein Programm, welches ein Programm einer anderen
Programmiersprache nach den notwendigen syntaktischen
"Uberpr"ufungen sofort ausf"uhrt. Im Gegensatz zu einem
"Ubersetzer\footnote{andere Bezeichnung f"ur Compiler} mu"s das
Quellprogramm nicht erst in eine andere Programmiersprache
"ubersetzt werden, sondern der Interpreter analysiert
nacheinander jede Anweisung und Deklaration des Quellprogramms
und f"uhrt diese unmittelbar aus.\footnote{siehe \cite{duden},
Seite 322}
\end{quotation}

In der Praxis jedoch wird kaum ein Interpreter ein Programm
ausf"uhren, ohne es vorher in eine geeignetere Form umgewandelt
zu haben. Man kann sich leicht vorstellen, da"s ein Interpreter,
der etwa bei einem Funktionsaufruf alle Module nach der Deklaration
dieser Funktion durchsuchen mu"s, zu langsam ist, um sinnvoll
eingesetzt werden zu k"onnen. Auch das wiederholte Lesen und
Erkennen der Schl"usselw"orter, die im Quelltext verwendet
werden, ist zu zeitaufwendig, als da"s es w"ahrend der
Ausf"uhrung des Programms durchgef"uhrt werden k"onnte. In den
folgenden Kapiteln werden wir deshalb besprechen, wie sich ein
Interpreter die Arbeit so einteilen kann, da"s eine m"oglichst
hohe Ausf"uhrungsgeschwindigkeit erreicht wird.

Ich m"ochte anmerken, da"s das in diesem Kapitel besprochene
Design f"ur einen Interpreter zwar theoretisch, jedoch keineswegs
realit"atsfern ist. Ich habe selbst einen SEL-Interpreter auf
diese Weise implementiert.

\section{Einsatzgebiete eines Interpreters}

Wenn man bedenkt, da"s die Ausf"uhrung eines Programms durch
einen Interpreter oft mehrere hundert Male l"anger dauert als die
eines mit einem Compiler erzeugten Maschinenspracheprogramms, fragt
man sich, welchen Zweck ein Interpreter eigentlich erf"ullen
kann. Ein Haupteinsatzgebiet f"ur Interpreter sind Computer, die
f"ur die sinnvolle Verwendung eines Compilers zu langsam sind
und/oder zu wenig Speicher bieten. Beispiele hierf"ur sind alte
Home-Computer wie z.~B. der Commodore~C-64 aber auch
Taschenrechner wie etwa der Texas Instruments~TI-85 oder der
Hewlett Packard~HP-48. Oft werden Interpreter auch zum
Entwickeln von Programmen eingesetzt, da sie den Vorteil haben,
Programme ohne Zeitverz"ogerung ausf"uhren zu k"onnen. Ein
anderes wichtiges Einsatzgebiet f"ur Interpreter ist die
Steuerung von Me"sger"aten und Robotern. Dort stellt die langsamere
Ausf"uhrung von Programmen kaum einen Nachteil dar, andererseits
ist es ein gro"ser Vorteil, da"s Programme ge"andert werden
k"onnen und sofort wieder einsatzbereit sind.

\section{Bestandteile eines Interpreters}

Bevor wir uns daranmachen, uns dar"uber zu unterhalten, wie der
Interpreter aussehen und welche Leistung er bieten soll, m"ussen
wir "uber den groben Aufbau eines Interpreters Bescheid wissen.
Wie bereits oben erw"ahnt wird in der Praxis das Quellprogramm
vor dem Interpretieren in eine f"ur den Interpreter geeignete
Form umgewandelt ({\em Tokenizing}). Dieser Vorgang wird in der
Regel von der lexikalischen Analyse erledigt. Dabei wird der
Quellcode sequentiell abgearbeitet und Schl"usselw"orter und
Bezeichner codiert. Zus"atzlich werden die Positionen von wichtigen
Schl"usselw"ortern wie z.~B. \rw{function}, \rw{global} und
\rw{local} gespeichert. Im Fall von SEL wird gleichzeitig
zur lexikalischen Analyse der Pr"aprozessor ausgef"uhrt. Er hat die
Aufgabe, Include-Dateien einzubinden und Makros in den
Programmtext einzuf"ugen.

Im Anschlu"s an die lexikalische Analyse und das Tokenizing kann
mit der Ausf"uhrung des Programms begonnen werden. Bei einer
Programmiersprache wie SEL m"ussen zuerst nat"urlich spezielle
Initialisierungen vorgenommen werden, n"amlich das Einrichten der
modul- und programmglobalen Variablen und das Analysieren der
Funktionsk"opfe. Zus"atlich speichert der Interpreter die
Stellen, an denen die jeweiligen Funktionen beginnen, soda"s
nicht das gesamte Programm durchsucht werden mu"s. Jetzt wird das
Programm durch einen Aufruf der \fu{main}-Funktion gestartet. Die
einzelnen Anweisungen und Ausdr"ucke werden vom Parser analysiert
und sofort ausgef"uhrt. Dabei mu"s nat"urlich auf spezielle
Routinen, die z.~B. f"ur die Verwaltung der Variablen und
"ahnliches verantwortlich sind, zur"uckgegriffen werden.

Zusammenfassend l"a"st sich also sagen, da"s ein Interpreter aus
drei wichtigen Teilen besteht: Der lexikalischen Analyse, die das
Tokenizing vornimmt, dem Parser, der das Programm syntaktisch
analysiert und den Verwaltungsroutinen, die die Hintergrundarbeit
leisten.

\section{Zielsetzung f\protect{"ur} das Design}
\label{intgoals}

Wie schon im ersten Kapitel erw"ahnt, ist SEL mit dem Ziel, sie
als Makrosprache zu verwenden, entworfen worden. Es w"are also
angebracht, den Interpreter von Beginn an so zu entwerfen, da"s
er nicht ein eigenst"andiges Programm darstellt, sondern eher
eine Art Modul, das man in andere Programme aufnehmen kann. Das
setzt nat"urlich voraus, da"s er nicht eine Unmenge an Speicher
verbraucht und auch, da"s er den Speicher, den er verbraucht,
nach der Ausf"uhrung des SEL-Programms wieder freigibt, soda"s er
wieder dem Hauptprogramm zur Verf"ugung steht.

Der Interpreter sollte nat"urlich, um ohne Bedenken in einem
Programm eingesetzt werden zu k"onnen, Programmierfehler so
schnell wie m"oglich erkennen k"onnen, den Ausf"uhrungsproze"s
angemessen stoppen, allen Speicher, den das SEL-Programm
angefordert hat, und auch den, den der Interpreter im Zuge der
Ausf"uhrung des Programms benutzt hat, freigeben und den
Fehler dem Hauptprogramm melden.

Wie auch die Programmiersprache SEL sollte auch der Interpreter
erweiterungsf"ahig sein, d.~h., er sollte dem Hauptprogramm
erlauben, bestimmte Funktionen auch SEL-Programmen zug"anglich zu
machen, wie beispielsweise f"ur das Zeichnen einer Linie in einem
Zeichenprogramm.

Da ich mich in letzter Zeit mit dem Multitasking-Betriebssystem
Windows NT auseinandergesetzt habe, wollte ich auch meinen
Interpreter multitaskingf"ahig machen, soda"s es in einem
Programm m"oglich ist, mehrere SEL-Programme gleichzeitig
auszuf"uhren. Um mich nicht falsch zu verstehen: Das Multitasking
wird in keiner Weise vom Interpreter in die Hand genommen,
sondern obliegt einzig und allein dem Betriebssystem. Der
Interpreter allerdings mu"s so entworfen werden, da"s er das
Multitasking nicht erschwert oder gar unm"oglich macht, was, wie
sich zeigen wird, nicht sehr schwer ist.

Im Laufe dieses Kapitels werde ich auf einige, aber nicht auf alle
Forderungen dieses Unterkapitels eingehen, weil sie zu trivial
sind. So ist z.~B. die Freigabe des benutzten Speichers nur eine
Formsache und mu"s nicht extra von mir erl"autert werden.

\section{Die lexikalische Analyse}
\label{scanner}

Die lexikalische Analyse hat die Aufgabe, den Zeichenstrom des
Quelltextes in einen Tokenstrom zu verwandeln. Als Token
bezeichnet man die kleinste Einheit, die vom Parser bearbeitet
wird. Token sind z.~B. die Schl"usselw"orter, Bezeichner,
Konstanten u.~s.~w. Der Grund, warum man die Erkennung von
Schl"usselw"ortern etc. nicht auch dem Parser "uberl"a"st, was
durchaus m"oglich w"are, ist der, da"s ein Parser diese Aufgabe
bei weitem nicht so schnell ausf"uhren kann wie ein eigens daf"ur
angefertigter {\em Scanner}.\footnote{Begriff f"ur ein Programm
oder einen Programmteil, der die lexikalische Analyse
durchf"uhrt} Au"serdem wird der Bau eines Parsers wesentlich
vereinfacht, wenn man einen Scanner verwendet.

Um die Arbeitsweise eines Scanners zu beschreiben, m"ochte ich
zun"achst ein Beispiel heranziehen. Nehmen wir an, der Scanner
h"atte den Zeichenstrom \fu{for (x=1 to 10)} zu bearbeiten. Er
untersucht nun das erste Zeichen und merkt, da es sich um einen
Buchstaben handelt, da"s er es entweder mit einem Bezeichner oder
mit einem Schl"usselwort zu tun hat. Die einfachste Methode,
dieses Token zu analysieren ist, so lange Buchstaben, Ziffern und
Unterstriche zu lesen, bis er an einem Zeichen ankommt, das
nicht mehr aus einer der genannten Kategorien stammt. Der
Scanner w"urde also zuerst das \fu{f}, dann das \fu{o} und
anschlie"send das \fu{r} lesen. Da das n"achste Zeichen weder
Buchstabe, Ziffer noch Unterstrich ist, stoppt er hier. Die
einfachste Methode, herauszufinden, ob man es hier mit einem
Schl"usselwort oder einem Bezeichner zu tun hat, ist, eine
Tabelle anzulegen, die alle Schl"usselw"orter der Sprache
umfa"st. Findet der Scanner das Wort in der Tabelle, ist es ein
Schl"usselwort, andernfalls ein Bezeichner. Das n"achste Zeichen
ist ein Leerzeichen. In einer formlosen Programmiersprache wie
SEL werden Leerzeichen und Tabulatoren "ubersprungen, soda"s der
Scanner nichts anderes zu tun hat, als zum n"achsten Zeichen
vorzur"ucken. Zwar sind Zeilenumbr"uche f"ur den Parser nicht
relevant, doch f"ur den Pr"aprozessor sind sie wichtig, soda"s
sie gemeldet werden sollten.

Da es nur ein Symbol in SEL gibt,
das mit einer "offnenden, runden Klammer beginnt, n"amlich dieses
selbst, wei"s der Scanner sofort, mit welchen Token er es zu tun
hat und kann vorr"ucken. Interessant wird es nun bei der Analyse
des Symbols \fu{=}, denn hier kann der Scanner nicht sofort davon
ausgehen, da"s es sich um den Zuweisungsoperator handelt, da es
noch ein anderes Symbol, n"amlich \fu{==} gibt, das mit einem
\fu{=} anf"angt. So mu"s der Scanner auch hier das n"achste
Zeichen betrachten um entscheiden zu k"onnen, welches Symbol er
vor sich hat.

Wie dieses Beispiel zeigt, ist es f"ur einen
Scanner vital zu wissen, welches Zeichen als n"achstes vorliegt,
wobei er sich dieses aber auch f"ur das n"achste Token aufheben
k"onnen mu"s. Die L"osung ist einfach: Man arbeitet mit einem
Zeichen im voraus, d.~h. es wird immer das kommende Zeichen
betrachtet, nicht das aktuelle. Wird es akzeptiert, liest man das
n"achste Zeichen ein, wenn nicht, kehrt man mit dem erkannten
Token zur aufrufenden Funktion (in Falle von SEL der
Pr"aprozessor) zur"uck. Beim n"achsten Aufruf des Scanners wird
mit dem vorher zur"uckgestellten Zeichen weitergearbeitet.

Um den Proze"s der lexikalischen Analyse zus"atzlich zu
beschleunigen, sollte die Eingabe gepuffert werden. Dabei wird
ein Teil der Eingabe in einen Zwischenspeicher (Puffer) gelesen,
um von dort aus vom Scanner abgearbeitet zu werden. Wird das Ende
des Puffers erreicht, mu"s abermals ein Teil der Eingabe gelesen
werden, so lange, bis das Ende der Datei erreicht ist. In diesem
Fall sollte der Scanner der "ubergeordneten Funktion Bescheid
geben.

Bei bestimmten Tokentypen gen"ugt es nicht, zu melden, da"s
dieses Token erkannt wurde, sondern es mu"s auch ein Attribut
zur"uckgeliefert werden, im Falle eines Bezeichners sein Name, im
Fall einer \tint -Konstante ihr Wert u.~s.~w., es sollte also
f"ur ein Token eine geeignete Datenstruktur verwendet werden, in
der alle n"otigen Attribute Platz finden. F"ur unseren
Interpreter k"onnte diese Struktur etwa so aussehen:

\newpage
\begin{verbatim}
struct Token
  iType  : int;
  pcName : *char;
  cChar  : char;
  iInt   : int;
  fFloat : float;
end
\end{verbatim}

Ist das Token ein Bezeichner, enth"alt \fu{pcName}
dessen Namen, im Falle einer \tchar -Konstanten enth"alt
\fu{cChar} ihren Wert u.~s.~w.

\subsection{Der Pr\protect{"aprozessor}}

Der Pr"aprozessor ist in gewisser Weise eine Zwischenstation von
Scanner und Parser. In SEL hat er die Aufgabe, den Inhalt von
Include-Dateien in den Programmtext einzuf"ugen, Kommentare zu
entfernen und Makros zu bearbeiten.

Das Entfernen von Kommentaren ist im Grunde kein Problem. Wird
vom Scanner das Symbol \fu{//} erkannt, dann werden so lange
Token gelesen, bis ein Zeilenumbruch gemeldet wird, wonach der
Kommentar beendet ist. Zus"atzlich zu den Kommentaren sollten
auch alle Zeilenumbr"uche herausgefiltert werden, da sie den
Parser st"oren w"urden.

Etwas komplizierter ist das Einf"ugen von Include-Dateien. Da es
m"oglich sein sollte, Include-Dateien zu verschachteln, gibt es
zwei M"oglichkeiten zur Implementation. Entweder, der
Pr"aprozessor arbeitet rekursiv, soda"s er sich selbst nochmals
aufruft, wenn eine Include-Datei zu bearbeiten ist, oder aber er
benutzt einen Stack zur Zwischenspeicherung der einzelnen
Include-Level, was am Ende eigentlich auf dasselbe hinausl"auft.
Ein kleines Problem ergibt sich mit dem Puffer f"ur den
Scanner: Wenn der Pr"aprozessor nur einen einzigen Puffer f"ur alle
Include-Level zusammen benutzt, mu"s er beim Eintritt in eine
Include-Datei geleert und mit
der Eingabe aus der Include-Datei gef"ullt werden. Ist diese zu Ende,
mu"s der Puffer abermals geleert werden und mit dem
"ubergeordneten Level gef"ullt werden. Das Problem dabei ist,
da"s man in der Datei dabei ein wenig zur"ucksetzen mu"s, weil
ein Teil des Puffers gel"oscht wurde. Ist das aus
irgendwelchen Gr"unden nicht m"oglich oder zu ineffizient,
mu"s man f"ur jeden einzelnen Include-Level einen Puffer
benutzen, was nat"urlich mehr Speicher verbraucht.

Zum Bearbeiten von Makros mu"s eine Datenstruktur
angelegt werden, in der alle Makros aufgezeichnet und auch schnell wieder
gefunden werden k"onnen. Ein bin"arer Baum und eine Hash-Tabelle
bieten sich hierzu an. F"ur jeden einzelnen Eintrag in dieser
Struktur mu"s zus"atzlich noch gespeichert werden, durch welchen
Text er zu ersetzen ist. Da der Pr"aprozessor auf den Scanner
aufbaut, ist es einfacher und effizienter, nicht den Zeichenstrom
zu speichern, sondern den Tokenstrom, der vom Scanner geliefert
wird. Dazu verwendet man am einfachsten eine einfach verkettete
Liste aus Token:

\begin{verbatim}
struct TokenList
  tokToken : Token;
  Next     : *TokenList;
end
\end{verbatim}

Trifft der Pr"aprozessor nun auf einen Bezeichner, mu"s er zun"achst
"uberpr"ufen, ob nicht vielleicht ein Makro unter diesem Namen
exisitiert. Dazu sucht er
nach ihm in der Datenstruktur. Ist er dort vorhanden,
setzt der Pr"aprozessor anstatt seiner die Tokenliste ein,
andernfalls wird er wie ein normaler Bezeichner behandelt.

Nachdem es m"oglich sein soll, innerhalb eines Makros ein anderes
zu benutzen, m"ussen wir noch etwas einbauen. Die einfachste
L"osung ist, da"s wir zum Einf"ugen eines Tokens in unseren
Zwischencode eine eigene Funktion verwenden, die, falls sie auf
einen Bezeichner trifft, der als Makro definiert ist, sich f"ur
jedes der Token aus der Liste selbst wieder aufruft. Dadurch ist
garantiert, da"s, falls in dieser Liste ein Makro auftaucht,
auch dieses entsprechend behandelt wird.
                              
\subsection{Der Zwischencode}
\label{tokencode}

Wie schon oben erw"ahnt, verwenden wir eine Art Zwischencode, um
dem Interpreter eine schnelle Ausf"uhrung zu erm"oglichen. Dabei
werden Schl"usselw"orter und Symbole mit einer eindeutigen Nummer
codiert, soda"s der Tokentyp sofort ersichtlich ist. Bei
bestimmten Tokentypen m"ussen zus"atzliche Informationen
gespeichert werden, wie z.~B. bei Konstanten oder Bezeichnern.

Um die Gr"o"se des Zwischencodes nicht zu sehr aufzubl"ahen,
werden auch die Namen der Bezeichner mit einer Nummer
codiert. Dazu wird eine Tabelle aller im Modul vorkommenden
Bezeichner angelegt und jedem Eintrag eine eindeutige Nummer
zugewiesen. Diese Nummer wird anschlie"send im Zwischencode
verwendet, um Bezeichner zu charakterisieren. Zu beachten ist,
da"s nat"urlich auch der Parser "uber diese Tabelle verf"ugen
mu"s, da er sonst nur die Nummern der Bezeichner wei"s. Wird
allerdings in allen Modulen des Programms dieselbe Tabelle
verwendet, so kann der Interpreter das Programm auch ohne das
Wissen um die Namen der Bezeichner ausf"uhren.

Eine besondere Hilfeleistung erstattet der Zwischencode dem
Interpreter bei Schleifen und Kontrollanweisungen: F"ur jedes \rw{if}
wird zus"atzlich die Position im Zwischencode abgespeichert, an der das
zugeh"orige \rw{end} bzw. \rw{else} steht. Das hat den Vorteil,
da"s der Interpreter, wenn er z~B. auf eine \rw{if}-Anweisung
trifft, deren Bedingung logisch falsch ist, ohne lange zu suchen
zum zugeh"origen \rw{end} bzw. \rw{else} springen kann, was
verst"andlicherweise die Ausf"uhrungsgeschwindigkeit erheblich
steigert, besonders bei Schleifen.

Wie bereits erw"ahnt, werden dem Interpreter zus"atzlich die
Position aller Vorkommnisse der Schl"usselw"orter \rw{global},
\rw{local} und \rw{function} mitgeteilt, soda"s er ohne zu suchen die
Initialisierung der programm- und modulglobalen Variablen und der
Funktionsk"opfe vornehmen kann.

Neben der erheblichen Geschwindigkeitssteigerung, die der
Zwischencode erm"oglicht, hat er noch den Vorteil, da"s er
weniger Speicherplatz als der Quellcode ben"otigt, bei etwas
gr"osseren Programmen etwa die H"alfte, was besonders unter
MS-DOS entgegenkommend ist.

Ich m"ochte hier kurz noch darauf eingehen, wie die oben
angesprochenen Verweise von Anweisungen auf die zugeh"origen
\rw{end}s bzw. \rw{else}s erzeugt werden k"onnen. Diese auf den
ersten Blick umst"andliche Aufgabe l"a"st sich relativ leicht
l"osen, indem man einen Stack zur Speicherung der
Verschachtelungen der Anweisungen verwendet. Trifft der Scanner z.~B.
auf ein \rw{if}, dann pusht er die Position, an der er sich
gerade im Zwischencode befindet, auf den Stack und l"a"st Platz
f"ur den sp"ater folgenden Verweis frei. Als n"achstes trifft
er z.~B. auf ein \rw{for} und tut dasselbe. Irgendwann mu"s
er jedoch zu dem \rw{end} gelangen, da"s die jeweils letzte
Anweisung schlie"st. In diesem Fall speichert er das \rw{end} im
Zwischencode, poppt vom Stack die Position der letzten
Anweisung und f"ugt an dieser Stelle die Position ein, an der
er soeben das \rw{end} gespeichert hat. Diese Technik wird
"ubrigens {\em Backpatching} genannt und oft im Compilerbau
eingesetzt.\footnote{Tats"achlich ist unser Scanner, der
einen SEL-Quelltext in Zwischencode umwandelt, streng gesehen
auch ein Compiler.}

\bigskip

Eine effiziente und einfache Darstellung des Zwischencodes ist
die Verwendung jeweils eines Bytes f"ur ein Token, wobei bei
bestimmten Tokentypen, wie z.~B. dem \rw{if}-Token, ein
zus"atzlicher Verweis auf ein anderes Token eingebaut wird. Wie
viele Bytes f"ur diesen Verweis verwendet werden, h"angt von den
Anforderungen an den Interpreter ab. Ich verwende hierzu in
meinem Interpreter 32 Bits, soda"s der Zwischencode theoretisch
bis zu 4 GByte gro"s sein kann, was sicherlich f"ur alle Anwendungen
mehr als ausreichend ist. F"ur die meisten
Anwendungen w"urden auch 16 Bits reichen, da Module in der Regel
sowieso klein gehalten werden. Bei anderen Token wie etwa einer
String-Konstante mu"s au"ser dem Tokentyp noch ein Wert
gespeichert werden, wie in diesem Fall der String.

\section{Der Parser}
\label{parser}

Der wohl komplizierteste Teil des Interpreters ist der
Parser. Er hat zur Aufgabe, aus dem Tokenstrom des Zwischencodes
die einzelnen Anweisungen zu erkennen. Sehr wichtig f"ur den Bau
eines Parsers ist die Grammatik, denn sie gibt an, was der Parser
eigentlich erkennen soll. Aus verschiedenen Gr"unden werden wir
Parser schreiben, die Teile der Sprache SEL erkennen k"onnen,
n"amlich Typen, Ausdr"ucke und Folgen von Anweisungen. Die Parser
f"ur Funktionsk"opfe und \rw{global}- bzw. \rw{local}-Bl"ocke
sind trivial, soda"s ich hier nicht speziell auf sie eingehen
werde. Ein Grund f"ur diese Trennung ist, da"s der
SEL-Interpreter vor allem f"ur die Verwendung in anderen
Programmen gedacht ist und dort ist es in vielen Situationen von
Vorteil, wenn der Benutzer Ausdr"ucke eingeben kann, z.~B. als
Filterbedingung in einer Datenbank.
                              
\subsection{Der Typen-Parser}

Beim Design eines Parsers ist es zun"achst wichtig, "uber die
Grammatik Bescheid zu wissen, die der Parser erkennen soll. In
diesem Unterkapitel besch"aftigen wir uns mit einem sehr
einfachen Teil der Grammatik von SEL, n"amlich den Typen. Das
gibt uns die M"oglichkeit, uns "uber die verschiedenen Arten von
Parsern und "uber die Zusammenh"ange zwischen Grammatik und
Parser zu informieren. Der "Ubersichtlichkeit halber sei an
dieser Stelle die Grammatik dargestellt:

\begin{verbatim}
type:          `char'
             | `int'
             | `float'
             | `void'
             | `[' INT_CONSTANT `]' type
             | `*' type
             | IDENTIFIER
             ;
\end{verbatim}

Diese Grammatik hat sieben Produktionen, das Startsymbol
ist \fu{type}.

\subsubsection{Top-Down und Bottom-Up Parser}
\label{topdbotu}

Im Interpreter- und Compilerbau werden haupts"achlich zwei
verschiedene Arten von Parsern verwendet: Top-Down und Bottom-Up
Parser. W"ahrend Top-Down Parser beim Startsymbol beginnen und
immer weiter hinabsteigen, bis sie auf die Token treffen, fangen
Bottom-Up Parser bei den Token an und arbeiten sich bis zum
Startsymbol hinauf. Das l"a"st sich am besten
verdeutlichen, indem man versucht, einen Parse-Tree zu
konstruieren, also einen Baum, der f"ur einen g"ultigen
Ausdruck der Sprache, in diesem Fall etwa \fu{*[10]Test}, den
Weg vom Startsymbol bis hinunter zu den einzelnen Token,
verdeutlicht. Der Parse-Tree f"ur den Beispieltyp ist in
Abbildung \ref{parsetree} dargestellt.

\begin{figure}

\centerline{\includegraphics[height=4cm]{partree1.eps}}

\caption{Der Parse-Tree f"ur \fu{*[10]Test}}
\label{parsetree}
\end{figure}

Die Wurzel des Baumes ist immer das Startsymbol, in diesem Fall
also \fu{type}. Dieser \fu{type} kann auf der ersten Ebene
zerlegt werden in das Sternsymbol und einen zweiten \fu{type},
gem"a"s der Produktion \fu{`*' type}. Dieser zweite \fu{type}
kann zelegt werden in eine "offnende und eine schlie"sende
Klammer und eine \fu{INT\_CONSTANT} dazwischen, gefolgt von noch
einem \fu{type}. Dieser wiederum ist ein \fu{IDENTIFIER}, n"amlich
\fu{Test}.

Nun zu den Vorgehensweisen der beiden Parsertypen. Ein Bottom-Up
Parser arbeitet sich wie gesagt von den einzelnen Token zum
Startsymbol hinauf. Dazu benutzt er einen Stapel, auf dem er zwei
verschiedene Aktionen durchf"uhren kann: {\em Shift} und
{\em Reduce}. Aus diesem Grund werden Bottom-Up Parser auch oft
als Shift-Reduce Parser bezeichnet. Die Analyse des
Beispieltyps f"uhrt der Parser so durch: Zun"achst liest er den
Stern, den er shiftet. Auch die beiden Klammern und die
\fu{INT\_CONSTANT} werden von ihm geshiftet. Nun liest er einen
Bezeichner und entscheidet, diesen zu einem Typen zu reducen,
gem"a"s der Produktion \fu{IDENTIFIER}. M"u"ste er einen
Parse-Tree erstellen, h"atte er an dieser Stelle einen Baum
mit der Wurzel 3\footnote{siehe Abbildung \ref{parsetree}}
erzeugt. Auf
dem Stack des Parsers liegen nun die Token \fu{`*'}, \fu{`['},
\fu{INT\_CONSTANT}, \fu{`]'} und das Nonterminal \fu{type}, in der
angegebenen Reihenfolge. An dieser Stelle entscheidet der Parser,
die letzten vier Symbole zu reducen. Um einen Parse-Tree zu
erhalten, m"u"ste er nun Baum 3 verwenden, um den gr"o"seren
Baum 2 zu konstruieren. Zuletzt w"urde der Parser nochmals
reducen und den Baum 1 erstellen. Der Parser hat nun
seine Analyse beendet, da das einzige, was auf seinem Stack
liegt, das Startsymbol ist. Die Hauptaufgabe bei der Erstellung
eines Bottom-Up Parsers ist, zu entscheiden, wann er zu shiften
und wann zu reducen hat. Dies ist eine schwere Aufgabe und wird
in der Praxis von einem Parser-Generator (auch Compiler-Compiler
genannt) wie etwa YACC erledigt. Da wir jedoch unseren Parser
selbst programmieren wollen, beschr"anken wir uns auf die
wesentlich einfacheren Top-Down Parser.

Ein Top-Down Parser geht bei der Analyse von \fu{*[10]Test}
folgenderma"sen vor: Er erkennt das Sternsymbol und wei"s nun,
da"s er nach der Produktion \fu{`*' type} vorzugehen hat. An
dieser Stelle w"urde er den Baum 1 erzeugen. Er liest nun den
Stern ein und mu"s nun einen weiteren Typ analysieren. Hier wird
klar, da"s ein Top-Down Parser in der Regel rekursiv
implementiert wird. Zwar ist es m"oglich, diese Rekursion mit
einem Stack zu beseitigen, doch ein solcher Parser ist schwerer
zu programmieren und zu warten, soda"s ich eine rekursive
Implementation gew"ahlt habe.

Der Parser ruft sich nun selbst auf und erkennt anhand der
"offnenden eckigen Klammer, da"s er die Produktion
\fu{`[' INT\_CONSTANT `]' type} w"ahlen mu"s und erzeugt
gleichzeitig Baum 2. Nacheinander liest er nun die drei folgenden
Token und ruft sich abermals selbst auf. Nun wendet er die
Produktion \fu{IDENTIFIER} an und erstellt Baum 3. Damit ist die
Analyse beendet.

\subsubsection{Look-Aheads}

Wie sich in diesem Beispiel gezeigt hat, ist es f"ur einen Parser wichtig, zu
wissen, welches Token er als n"achstes zu bearbeiten hat, denn
danach hat er zu entscheiden, welche Produktion er ausw"ahlt. Zu
diesem Zweck verwendet man sog. {\em Look-Aheads}. F"ur einen
SEL-Parser und f"ur fast alle Programmiersprachen ist ein Look-Ahead
v"ollig ausreichend, es gen"ugt also, ein Token
vorauszusehen. Die Implementation eines Look-Ahead ist simpel:
Anstatt nur des aktuellen Tokens verwendet der Parser zus"atzlich das
Look-Ahead. Soll nun das aktuelle Token "ubersprungen werden,
weist er einfach den Wert des Look-Aheads dem aktuellen Token
zu und liest einen neuen Look-Ahead aus dem Zwischencode ein.

\subsubsection{Die Konstruktion des Parsers}

Mit diesen Informationen ist es nun ein Leichtes, einen
rekursiven Top-Down Parser f"ur Typen zu bauen. Er besteht aus
einer einzigen Funktion, die durch Betrachten des Look-Ahead
entscheidet, was passieren soll. Ist der Look-Ahead eines der
Token \fu{`char'}, \fu{`int'}, \fu{`float'}, \fu{`void'} oder
\fu{IDENTIFIER}, mu"s der Parser dieses Token nur lesen und die Analyse
ist beendet. Im Fall einer "offnenden eckigen Klammer liest er
drei Token und ruft sich anschlie"send selbst auf.\footnote{Es
ist leicht m"oglich, eine am Ende stehende Rekursion durch einen
Sprung an den Beginn der Funktion, z.~B. durch eine
\rw{repeat}-Schleife, evtl. durch einen Stack unterst"utzt, zu
ersetzen, was im Normalfall die bessere L"osung ist.} Analog dazu
ist die Produktion \fu{`*' type} zu implementieren. Sollte der
Look-Ahead keinem der Token entsprechen, hat es der Parser
offensichtlich mit einem Fehler zu tun, den er entsprechend
behandeln mu"s. Ein Interpreter wird in diesem Fall eine
Fehlermeldung ausgeben und die Ausf"uhrung beenden.

Nat"urlich ist klar, da"s unser Parser nicht nur die Aufgabe hat,
einen Typ zu erkennen, sondern auch, mitzuteilen, um welchen Typ
es sich dabei handelt. Das jedoch f"allt unter die semantische
Analyse und wird in Kapitel \ref{inttypes} behandelt.

\subsection{Der Ausdruck-Parser}
\label{exprparser}

Nachdem wir nun die grundlegenden Funktionen eines Parsers
kennen, wollen wir einen etwas komplexeren Parser konstruieren:
Den Parser f"ur SEL-Ausdr"ucke. Die Grammatik hierzu kann im
Anhang \ref{grammar} oder in \cite{selkr}, Seite 14 nachgelesen
werden. Ausdr"ucke in SEL werden in der Grammatik durch die
beiden Nonterminals \fu{expr} und \fu{primary} repr"asentiert. In
den Produktionen f"ur \fu{primary} taucht noch das dritte
Nonterminal \fu{type} auf, doch dieses stellt hier kein
Problem mehr dar, da der entsprechende Parser bereits im letzten
Unterkapitel skizziert wurde.

Besonders interessant sind die Produktionen f"ur \fu{expr}, denn
sie bestehen gr"o"stenteils aus dem Nonterminal \fu{expr} selbst.
Wie sich leicht ersehen l"a"st, ist diese Grammatik auf keinen
Fall eindeutig, was die beiden Parse-Trees f"ur den Ausdruck
\fu{a=b+1} in Abbildung \ref{twopartrees}
verdeutlichen.\footnote{In einer eindeutigen Grammatik darf es
f"ur einen korrekten Satz nur einen einzigen Parse-Tree geben.}

\begin{figure}

\centerline{\includegraphics[height=3cm]{partree2.eps}}

\caption{Zwei Parse-Trees f"ur \fu{a=b+1}}
\label{twopartrees}
\end{figure}

Alle diese Konflikte jedoch lassen sich durch die in
\cite{selkr},~Seite~21 angegebene Priorit"atstabelle l"osen. Die
Frage, die sich nun stellt, ist, wie man diese Rangfolge in den
Parser einbringen kann. Betrachten wir der Einfachheit halber
nur die folgenden drei Produktionen von \fu{expr}:

\begin{verbatim}
expr:   expr `=' expr
      | expr `+' expr
      | expr `*' expr
      | primary
      ;
\end{verbatim}

Die Priorit"atstabelle sagt, da"s die h"ochste Priorit"at der
Operator \fu{*} hat, danach folgen \fu{+} und~\fu{=}. Der
Konflikt kann nun gel"ost werden, indem man f"ur jede Priorit"atsstufe ein
eigenes Nonterminal einf"uhrt, wobei die Produktionen f"ur die
Operatoren niedriger Priorit"at auf die Nonterminals der
Operatoren h"oherer Priorit"at zur"uckgreifen, wie hier
dargestellt:

\begin{verbatim}
expr3:   expr2 `=' expr2
       | expr2
       ;
expr2:   expr1 `+' expr1
       | expr1
       ;
expr1:   primary `*' primary
       | primary
       ;
\end{verbatim}

Diese Grammatik kann, angewendet auf den Ausdruck \fu{a=b+1},
nur einen Parse-Tree erzeugen, n"amlich den, der der
Klammerung des Ausdrucks \fu{a=(b+1)} entspricht.
Ungl"ucklicherweise geht mit dieser Grammatik die F"ahigkeit
verloren, einen Ausdruck wie \fu{a+b+c} darzustellen. Um auch
diesen Mangel zu beseitigen, ersetzt man in jeder Produktion mit
einem Operator die linke bzw. rechte Seite mit dem Nonterminal
des jeweiligen Levels, je nachdem, ob der Operator links- oder
rechtsassoziativ ist. Daraus resultiert folgende Grammatik:

\begin{verbatim}
expr3:   expr2 `=' expr3
       | expr2
       ;
expr2:   expr2 `+' expr1
       | expr1
       ;
expr1:   expr1 `*' primary
       | primary
       ;
\end{verbatim}

Diese Grammatik ist eindeutig und erzeugt sowohl f"ur den
Ausdruck \fu{a=b+1} als auch f"ur \fu{a+b+c} die richtigen
Parse-Trees, man kann also mit dieser Methode die "`gro"se"'
Grammatik f"ur den Bau des Anweisungs-Parsers umwandeln. Ich m"ochte die
sich daraus ergebende Grammatik in dieser Arbeit nicht
auff"uhren, da sie sehr leicht hergeleitet werden kann und zu
viel Platz beanspruchen w"urde.

Mit den bisher besprochenen Techniken kann man bereits einen
Parser f"ur einen Teil der Grammatik konstruieren, sp"atestens
bei den Linksrekursionen, wie z.~B. in den
Produktionen f"ur \fu{expr2} in der obigen Grammatik, wird man
allerdings Schwierigkeiten bekommen. Eine Rechtsrekursion wie etwa bei
\fu{expr3} ist keine Schwierigkeit: Man kann sie entweder
durch eine echte Rekursion implementieren oder durch eine
Schleife, unterst"utzt durch einen Stack. F"ur die erste Produktion
von \fu{expr2} m"u"ste man jedoch eine Funktion programmieren,
die als erstes sich selbst aufruft, was nat"urlich in einer
endlosen Rekursion m"unden w"urde.

Um dieses Problem zu l"osen, versuchen wir zu ergr"unden, welche
Menge von Ausdr"ucken eigentlich die beiden Produktionen von
\fu{expr2} beschreiben. Der einfachste Ausdruck, den die
Grammatik beschreiben kann, ist offensichtlich \fu{expr1}. Er
l"a"st sich mit Hilfe der beiden Produktionen nicht weiter
expandieren. Benutzt man als Basis nun die erste Produktion,
kann man den ersten \fu{expr2} expandieren, entweder
durch die erste oder die zweite Produktion. Benutzt man die
zweite Produktion, bekommt man als Ergebnis den Ausdruck
\fu{expr1 `+' expr1}, der sich abermals nicht weiter expandieren
l"a"st. F"uhrt man diesen Proze"s fort, wird man bald merken,
da"s diese beiden Produktionen eine Abfolge von beliebig vielen
\fu{expr1}, verbunden jeweils durch ein \fu{`+'}, beschreiben.
Zwar beschreiben die Produktionen f"ur \fu{expr3} im Grunde
dieselbe Menge, doch ist der Operator \fu{=} rechts-assoziativ,
w"ahrend \fu{+} links-assoziativ ist, was sich z.~B. beim
Erzeugen eines Parse-Trees zeigt.

Mit diesem Wissen k"onnen wir nun einen Parser konstruieren, der
einen \fu{expr2} akzeptiert, m"ussen aber wie gesagt die
Links-Assoziativit"at beachten. Um nicht unn"otig viele Worte
machen zu m"ussen, folgt ein Beispiel einer Implementation, die
einen (primitiven) Parse-Tree erzeugt. F"ur den Parse-Tree wird
folgende Struktur verwendet, die keiner Erkl"arung bedarf:

\begin{verbatim}
struct ParseTree
  Operator : int;
  Value    : int;
  Left     : *ParseTree;
  Right    : *ParseTree;
end
\end{verbatim}

Die Funktion \fu{expr2} erzeugt einen Parse-Tree f"ur einen
\fu{expr2} und greift dabei auf die Funktion \fu{expr1} zur"uck,
die dasselbe f"ur einen \fu{expr1} tut. Zus"atzlich werden die
Funktionen \fu{LookAhead} und \fu{Accept} verwendet, die ich
nicht n"aher erl"autern mu"s:

\begin{verbatim}
function expr2 () : *ParseTree

  Tree : *ParseTree;
  Temp : *ParseTree;

begin
  Tree = expr1();
  while (LookAhead() == PLUS)
    Accept(PLUS);
    Temp = new ParseTree;
    Temp->Operator = PLUS;
    Temp->Left = Tree;
    Temp->Right = expr1();
    Tree = Temp;
  end
  return Tree;
end
\end{verbatim}

Auf diese Art und Weise k"onnen Funktionen f"ur jeden der
Priorit"atsgrade erstellt werden. Im wesentlichen mit derselben
Methode werden auch die Links-Rekursionen in den Produktionen von
\fu{primary} beseitigt. Eine Kleinigkeit in der Implementation
von \fu{primary} m"ochte ich noch besprechen: Die Produktion
\fu{IDENTIFIER `(' [expr \{`,' expr\}] `)'} enth"alt ein
optionales Element und darin eingeschlossen ein Element, da"s
sich beliebig oft wiederholen kann. Es ist zwar m"oglich, solche
Konstruktionen durch Umformung der Grammatik mit den oben
besprochenen Methoden zu analysieren, doch die direkte
Implementation ist nicht schwer zu realisieren und wird deshalb
von mir vorgezogen.

Die beiden eckigen Klammern in der Produktion bedeuten, da"s auf
die "offnende runde Klammer, die nach dem \fu{IDENTIFIER} zu
stehen hat, entweder eine schlie"sende runde Klammer oder ein
\fu{expr} folgt. Welche der beiden M"oglichkeiten zutrifft,
kann der Parser dadurch erkennen, da"s er den Look-Ahead
betrachtet.
Da ein \fu{expr} niemals mit einer schlie"senden Klammer beginnen
kann, ist klar, da"s, wenn der Look-Ahead die schlie"sende
Klammer ist, die erste Variante zutrifft, andernfalls die zweite.
In letzterem Fall mu"s die Funktion \fu{expr} aufgerufen werden, um
den ersten \fu{expr} zu erkennen. Nun steht der Parser vor einem
"ahnlichen Problem. Entweder, es folgt ein Beistrich und noch ein
\fu{expr}, dann vielleicht noch ein Beistrich u.~s.~w., oder es
folgt die schlie"sende Klammer. Wie oben kann auch dieses
Problem mit dem Look-Ahead gel"ost werden. Entspricht er einem Beistrich,
mu"s noch ein \fu{expr} analysiert werden, wonach der Parser wieder
auf einen Beistrich testen mu"s. Ist der Look-Ahead die
schlie"sende Klammer, mu"s er diese nur noch lesen und ist
mit der Analyse des \fu{primary} fertig. Wenn jedoch der
Look-Ahead keinem der beiden Symbole entspricht, hat er
offensichtlich einen Fehler gefunden, der dementsprechend zu
behandeln ist.

\subsection{Der Anweisungsparser}

Der Parser f"ur Anweisungen ist relativ einfach zu realisieren,
da er weder Linksrekursionen enth"alt, noch irgendwelche
Priorit"atsregeln zu befolgen sind. Die Anweisungen von SEL sind
so konzipiert, da"s der Parser schon anhand des ersten Tokens
entscheiden kann, um welche Anweisung es sich handelt. Er
unterscheidet sich jedoch insofern von den bis jetzt besprochenen
Parsern, als da"s bei bestimmten Anweisungen nur Teile dieser
Anweisung geparst werden m"ussen, wie z.~B. bei \rw{if}, wo der Wert der
Bedingung bestimmt, ob der \rw{if}- oder der \rw{else}-Zweig,
wenn vorhanden, ausgef"uhrt werden soll.

Betrachten wir als Beispiel die zitierte \rw{if}-Anweisung. Der
Parser erkennt anhand des Look-Ahead, da"s es sich um eine
\rw{if}-Anweisung handelt. Er "uberliest also die beiden Token
\rw{if} und \fu{(} und ruft anschlie"send den im vorigen
Unterkapitel besprochenen Ausdruck-Parser auf, um den \fu{expr}
zu lesen. Wie sich in Kapitel \ref{semantics} zeigen wird,
berechnet dieser Parser den Wert des Ausdrucks, der nun
benutzt werden kann, um zu entscheiden, ob die dem \rw{if} folgenden
Anweisungen ausgef"uhrt werden sollen. Ist der Wert ungleich \fu{0}, ist
die Sache leicht. In diesem Fall mu"s sich der Anweisungs-Parser
nur selbst aufrufen, um die folgenden Anweisungen bis zum
\rw{end} bzw. \rw{else} zu parsen. Andernfalls benutzt er die
zus"atzliche Information, die ihm der in Kapitel \ref{tokencode}
besprochene Zwischencode zum \rw{if}-Token liefert. Sie gibt die
Position an, an der das zugeh"orige \rw{end} bzw. \rw{else}
steht. So m"ussen nicht langwierig alle folgenden Anweisungen bis
zum jeweiligen Token durchsucht werden, sondern es kann direkt an
die entsprechende Position gesprungen werden. Steht an dieser
Stelle ein \rw{end}, ist die Analyse der \rw{if}-Anweisung
beendet, andernfalls mu"s der \rw{else}-Zweig ausgef"uhrt werden.

Noch etwas komplizierter verh"alt es sich mit der
\rw{switch}-Anweisung. Dort mu"s das \rw{switch}-Token die
Position des ersten zugeh"origen \rw{case} enthalten, dieses
wiederum die des n"achsten \rw{case} bzw. \rw{default}
u.~s.~w. Auf diese Weise ist es m"oglich, einen Zweig nach
dem anderen auf "Ubereinstimmung zu untersuchen und, falls eine
"ubereinstimmende Konstante gefunden wurde, diesen Zweig
auszuf"uhren, indem die Funktion sich selbst aufruft. Damit die
\rw{switch}-Anweisung auch wie vorgeschrieben ausgef"uhrt wird,
mu"s der Anweisungs-Parser beim normalen Parsen ein
\rw{case}-Schl"usselwort einfach "uberspringen und bei einem
\rw{break} zur "ubergeordneten Funktion zur"uckkehren.

Ein Problemfall ist die \rw{return}-Anweisung. Wenn der
Parser auf sie trifft, mu"s er auch zur "ubergeordneten Funktion
zur"uckkehren, doch er mu"s das im Gegensatz zur
\rw{break}-Anweisung so lange tun, bis der Anweisungs-Parser
komplett verlassen wurde. Wir m"ussen uns also "uberlegen, wie
wir es anstellen k"onnten, da"s der Parser wei"s, wann die
untergeordnete Instanz ein \rw{return} entdeckt hat. Die
einfachste L"osung ist, das mit dem R"uckgabewert zu tun. Im
n"achsten Kapitel wird sich zeigen, da"s SEL-Werte vom
Interpreter in Strukturen gehalten werden. Am sinnvollsten ist
es, als R"uckgabetyp f"ur den Parser einen Zeiger auf eine solche
Struktur zu benutzen. Trifft er nun auf ein \rw{break}, gibt er
einen Nullzeiger zur"uck und der "ubergeordnete Parser wei"s, das
alles in Ordunung ist. Wurde aber eine \rw{return}-Anweisung
entdeckt, wird ein Zeiger auf den zugeh"origen Wert
zur"uckgegeben. Die "ubergeordnete Funktion wei"s nun, da"s auch
sie beendet werden mu"s und gibt ihrerseits den Wert weiter. Wird
eine Funktion geparst, die als R"uckgabetyp \tvoid\ hat, mu"s
ein Zeiger auf einen Dummy-Wert zur"uckgegeben werden.

Der Anweisungs-Parser mu"s auch mit einem Null-Zeiger verlassen
werden, wenn er auf ein \rw{end} oder \rw{else} trifft, die zu
keiner von ihm behandelten Anweisungen geh"ort, denn dann f"allt
es in den Zust"andigkeitsbereich der "ubergeordneten
Rekursionsebene.

Hier wird auch klar, warum ich auf eine
\rw{goto}-Anweisung\footnote{siehe Kapitel \ref{goto}} verzichtet
habe: In einem rekursiven Parser w"are es f"ur einen Sprung
n"otig, zun"achst zu ermitteln, auf welcher Rekursionsebene das
Label liegt, und ob es vielleicht innerhalb einer anderen Anweisung
liegt, was einen besonders schlimmen Fall darstellt. Betrachten wir
dazu folgendes (theoretische) Programmfragment:

\begin{verbatim}
if (x == 5)
  goto label;
end
if (y == 3)
  do_something();
  if (z == 4)
    label:
    do_something_else();
  end
end
\end{verbatim}

In diesem Fall m"u"ste der Interpreter, wenn er die
\rw{goto}-Anweisung ausf"uhrt, zun"achst eine Rekursionsebene
verlassen und dann in zwei andere einsteigen, was nat"urlich sehr
schwer zu realisieren ist. "Uberdies m"u"ste der Parser sehr
schnell herausfinden k"onnen, wieviel Ebenen zu verlassen und in
welche er einzusteigen hat. Zu diesem Problem ist mir bis jetzt
noch keine zufriedenstellende L"osung eingefallen.

\section{Die Verwaltung}
\label{semantics}

In den letzten Unterkapiteln wurde behandelt, wie der Interpreter
den Quelltext (genauer gesagt, den Zwischencode) lexikalisch
analysiert, in Zwischencode verwandelt und diesen syntaktisch
analysiert, jedoch wurde noch nicht gesagt, wie der Interpreter
die durch diese Analyse gewonnenen Informationen dazu verwenden
kann, das Programm auszuf"uhren. Das ist die Aufgabe dieses
Unterkapitels. Hier wird besprochen, was der Interpreter bei der
Ausf"uhrung des Programms tut, und wie dies implementiert werden
kann.

Die meisten der in diesem Unterkapitel besprochenen Algorithmen
werden w"ahrend der Syntaxanalyse, also des im letzten
Unterkapitel besprochenen Parsing, ausgef"uhrt, was auch der
Definition des Interpreters\footnote{siehe Kapitel
\ref{whatisit}} entspricht.
        
\subsection{Module}
\label{modules}

Da SEL es erlaubt, modulglobale Variablen und Funktionen zu
benutzen, ist es vonn"oten, ein SEL-Programm auch im Speicher in
Module aufzuteilen. Dazu m"ussen f"ur jedes Modul mindestens der
Zwischencode und die modulglobalen Variablen und Funktionen
gespeichert werden. Im n"achsten Unterkapitel werden sich dazu
noch einige zus"atzliche Informationen gesellen, die bei der
Fehlerbehandlung helfen.

Die programmglobalen Funktionen und Variablen werden au"serhalb
der Module f"ur das gesamte Programm gemeinsam verwaltet, was
wesentlich effizienter und auch einsichtiger ist (und es nebenbei
einfacher und schneller macht, Namenskonflike zu bemerken).
Weiters gibt es zwei Variablen, die f"ur das gesamte Programm
gelten. Sie geben die Position an, an der sich der Parser
gerade mit der Ausf"uhrung aufh"alt. Die erste dieser Variablen
nenne ich \fu{MP} (Module Pointer). Sie enth"alt einen Wert, der
das Modul, in dem sich die gerade abgearbeiteten Anweisungen
befinden, eindeutig identifiziert. Am effizientesten ist es, ein
Array aus Modulbeschreibungen anzulegen, das durch diese
Variable indiziert wird. Die zweite Variable, \fu{IP}
(Instruction Pointer), gibt die Position innerhalb des
Zwischencodes des aktuellen Moduls an.\footnote{Genauerweise
m"u"ste diese Variable den Namen \fu{TP} f"ur Token Pointer
tragen, da sie sich nicht auf die aktuelle Anweisung, sondern auf
das aktuelle Token bezieht.}

\bigskip

In Kapitel \ref{intgoals} wurde als Ziel f"ur den Interpreter
festgelegt, da"s er Multitasking unterst"utzt. Um die dazu
notwendigen Schritte anzugeben, mu"s erst gekl"art werden, wie
Multitasking funktioniert. Zun"achst m"ochte ich feststellen, da"s
ich mit Multitasking nicht die F"ahigkeit eines Betriebssystems
meine, mehrere Programme gleichzeitig zu bearbeiten, sondern die
F"ahigkeit, innerhalb eines Programms mehrere Funktionen zur
gleichen Zeit auszuf"uhren. Unter den Betriebssystemen OS/2 und
Windows NT wird daf"ur der Begriff "`Multithreading"' gepr"agt.
Im Fall eines SEL-Interpreters l"auft dieses Multitasking darauf
hinaus, da"s durch diesen Interpreter (der ja ein Programm
darstellt) mehrere SEL-Programme gleichzeitig ausgef"uhrt werden
k"onnen.

Zwei Funktionen, die gleichzeitig ausgef"uhrt werden, k"onnen
zwar auf dieselben globalen Variablen zugreifen, soda"s eine
gemeinsame Benutzung eines Datenpools, etwa einer
Funktionssammlung f"ur ein Mathematikprogramm, gew"ahrleistet
ist, jedoch besitzt jede der beiden Funktionen einen eigenen
Stack, was garantiert, da"s sie ungest"ort von den anderen
Funktionen ihre eigenen lokalen Variablen verwenden kann.

Wenn nun der Interpreter zwei verschiedene SEL-Programme
gleichzeitig ausf"uhren soll, tut er das, indem er zwei Tasks
startet, die jeweils ein Programm bearbeiten. Um zu
gew"ahrleisten, da"s sich diese beiden Funktionen dabei nicht in
die Quere kommen, d"urfen die Daten, die das aktuelle Programm
betreffen, nat"urlich nicht in globalen Variablen gespeichert
werden, sondern m"ussen der Funktion, die f"ur die Ausf"uhrung
zust"andig ist, als Parameter mitgegeben werden. F"ur unseren
rekursiven Parser bedeutet das, da"s diese Daten sehr oft als
Parameter "ubergeben werden, was nicht nur den Stack belastet,
sondern auch die Ausf"uhrungsgeschwindigkeit. Um diese Verluste
gering zu halten, empfielt es sich, nicht direkt die Daten,
sondern Zeiger auf sie zu verwenden, was aber in den meisten
Programmiersprachen ohnedies die einzige M"oglichkeit ist.

\subsection{Typen}
\label{inttypes}

Eine f"ur die Ausf"uhrungsgeschwindigkeit sehr wichtige
Entscheidung bei einem Interpreter f"ur eine Programmiersprache,
die ein im Vergleich zu anderen Interpretersprachen so komplexes
Typkonzept verwendet, wie SEL, ist, wie diese Typen
repr"asentiert werden sollen. In Compilern werden hierf"ur
h"aufig verkettete Listen verwendet, die jedoch f"ur einen
Interpreter nicht sehr brauchbar sind, da sie wegen ihres
Aufbaus, der auf dynamischer Speicherverwaltung basiert, zu
langsam sind und au"serdem zu viel Speicher verbrauchen.

Das von mir verwendete Konzept basiert darauf, da"s Typen durch
Byte-Arrays mit begrenzter L"ange repr"asentiert werden. Da es
in SEL nur eine geringe Zahl an Grundtypen und auch nicht viele
M"oglichkeiten, diese zu verkn"upfen, gibt, kann ein Byte
jeweils eine Zusammensetzungsstufe oder einen Grundtypen
enthalten. Ein Beispiel: Der Typ \type{\poi \arr{10}int} wird
durch ein Array mit der Gr"o"se von 5 Bytes repr"asentiert,
dessen erstes Byte angibt, da"s es sich um einen Zeiger handelt.
Das zweite Element gibt an, da"s ein Array im Spiel ist. F"ur
dieses Array werden auch die n"achsten zwei Bytes ben"otigt, die
die Gr"o"se angeben. Das f"unfte Element schlie"slich steht f"ur
den Grundtypen \tint\ und ist gleichzeitig die Ende-Markierung,
da es sich dabei um einen Grundtypen handelt. In gewisser Weise
"ahnlich wie \tint\ werden Strukturen behandelt, die ja auch
immer am Ende des Typs stehen. Bei ihnen mu"s, "ahnlich dem
Array, ein Wert angegeben werden, der den Typ der Struktur
bestimmt. Zu diesem Zweck vergibt der Interpreter f"ur die Strukturen
Nummern, wobei die Werte von 0 bis 9999 f"ur programmglobale
Strukturen reserviert sind. Die weiteren Zuordungen richten sich
danach, wie viel Strukturen in den einzelnen Modulen definiert
sind.

Wie schon oben erw"ahnt, besitzen die Typenarrays eine begrenzte
L"ange, was es erm"oglicht, sie schon vor ihrer Verwendung
bereitzustellen, soda"s sie nicht w"ahrend der Ausf"uhrung des
SEL-Programms dynamisch allokiert werden m"ussen. In meiner
Implementation habe ich mich f"ur eine L"ange von 20 Bytes
entschieden. Das mag nicht viel erscheinen, vor allem in
Anbetracht der Tatsache, da"s es in SEL m"oglich ist, Typen
beliebig zu verschachteln, doch ist es m"oglich, mit diesen 20
Bytes einen Zeiger auf einen Zeiger auf einen Zeiger auf ein Array
aus Arrays aus Strukturen zu repr"asentieren, wobei noch 8 Byte
frei bleiben. Dieses Beispiel soll zeigen, da"s 20 Bytes f"ur die
Praxis eine absolut ausreichende Gr"o"se ist.

Beim Analysieren von Ausdr"ucken ist es in einem SEL-Interpreter
sehr oft n"otig, Typen bereitzustellen, von denen aber viele
schon im n"achsten Moment wieder freigegeben werden. Ist es nicht
m"oglich, diese Operationen effizient auszuf"uhren, leidet
darunter die Geschwindigkeit des Interpreters sehr. Ich habe zu
diesem Zweck einen Algorithmus entwickelt, der es erm"oglicht, in
konstanter Zeit, also unabh"angig davon, wieviele von den
bereitgestellten Typenarrays bereits vergeben sind, Typenarrays
zu allokieren und freizugeben:

\bigskip

Sei $n$ die Anzahl der maximal bereitzustellenden Typen und $s$
die Gr"o"se eines Typenarrays, dann wird vor der Ausf"uhrung des
SEL-Programms ein Byte-Array mit der Gr"o"se $n \cdot s$ erzeugt,
wobei $t$ die Speicheradresse sei, an der das Array beginnt.
Weiters sei $r$ die Speicheradresse, an der ein Array der Gr"o"se
$n$ beginnt, dessen Elemente vern"unftigerweise Maschinenworte
sind. Dieses Array wird initialisiert, soda"s $r[x] = x$.
Zus"atzlich sei $i$ eine numerische Variable, die mit
dem Wert $0$ initialisiert werden.

Soll ein Typenarray allokiert werden, wird zun"achst "uberpr"uft,
ob $i = n$. In diesem Fall sind bereits alle Arrays vergeben und
eine entsprechende Fehlermeldung wird ausgegeben. Andernfalls ist
$t + r[i] \cdot s$ die Adresse, bei der das angeforderte Array
beginnt. Zus"atzlich ist $i = i + 1$ zu setzen.

Soll das Typenarray mit der Speicheradresse $f$ freigegeben
werden, so ist zuerst $i = i - 1$ und dann $r[i] = (f - t) / s$
zu setzen.

Der angegebene Algorithmus verl"a"st sich darauf, da"s nicht
Arrays freigegeben werden sollen, die nicht allokiert sind, und
da"s jeweils maximal $s$ Bytes eines Arrays genutzt werden.

\bigskip

Dieser Algorithmus verwendet zwei Datenstrukturen. Das Array $t$
ist sozusagen ein Zusammenschlu"s aller Arrays, die zur
Verf"ugung gestellt werden k"onnen, was sich nat"urlich auch
mittels eines zweidimensionalen Arrays realisieren lie"se. Das
Array $r$ ist ein Puffer, der die Nummern (0 bis $n - 1$)
aller noch freien Typenarrays enth"alt. Die Variable $i$ gibt den
Index des n"achsten zu vergebenden Typenarrays an. Wird dieses
angefordert, wird $i$ um 1 erh"oht. Umgekehrt wird $i$ um 1
vermindert, wenn ein Array freizugeben ist. Dessen Nummer wird
anschlie"send in $r[i]$ gespeichert, soda"s es bei der n"achsten
Anforderung sofort wieder vergeben wird. Wie schon gesagt, nimmt
der Algorithmus an, da"s keine Fehler vom Programmierer begangen
werden, wie z~B. die Freigabe eines bereits freigegebenen Arrays.

In meiner Implementation von SEL habe ich f"ur $n$ den Wert 400
benutzt, f"ur $s$ wie gesagt 20. Diese Werte sind f"ur die
meisten Anwendungen ausreichend, und wenn n"otig, k"onnten sie
vom Benutzer entsprechend ge"andert werden. Die
Speicheranforderungen dieses Algorithmus sind durchaus
vertretbar. Bei den angegebenen Werten erw"achst f"ur die beiden
Arrays unter Windows NT ein Speicherbedarf von genau 9600 Byte,
was bei den heute "ublichen Speicherkapazit"aten kaum ein Problem
darstellt.

\subsection{Ausdr"ucke}
\label{intexprs}

"Ahnlich wichtig wie die Typen sind in SEL die Ausdr"ucke, die
nat"urlich auch verwaltet werden m"ussen. Dabei werden wir jedoch
nicht ganze Ausdr"ucke verwenden, sondern uns lediglich auf die
Werte der Ausdr"ucke beschr"anken, die im Laufe des
Parsevorganges errechnet werden.

Um zu zeigen, wie wichtig die effiziente Verwaltung von
Ausdr"ucken ist, hier ein konkretes Beispiel: Ein
QuickSort-Programm f"ur 500 zuf"allig angeordnete Zahlen hat auf
meinem Computer eine Laufzeit von etwa 25 Sekunden, wobei w"ahrend
dieser Zeit ungef"ahr 60000 Ausdr"ucke allokiert und wieder
freigegeben werden m"ussen.

Zum Allokieren und Freigeben von Ausdr"ucken verwende ich aus
diesem Grund denselben Algorithmus wie f"ur Typen, wobei ich von
Haus aus maximal 200 Ausdr"ucke zur Verf"ugung stelle, was in den
meisten F"allen ausreichend ist.

Sehr wichtig f"ur eine effiziente Verarbeitung ist nat"urlich
auch der Aufbau eines Ausdrucks. Ich verwende hierzu eine
Struktur. Ein obligater Bestandteil eines Ausdrucks ist dessen
Typ, soda"s ein Typenarray Teil dieser Struktur ist. Weiters gibt
es je eine Variable f"ur die Basistypen \tchar , \tint\ und
\tfloat . Hierzu wird es aus Effizienzgr"unden am besten sein,
passende Typen der Implementationssprache zu verwenden. F"ur den
Fall, da"s der Ausdruck ein Array, ein Zeiger oder eine Struktur
ist, habe ich zus"atzlich einen typenlosen Zeiger aufgenommen,
der, falls die Implementationssprache keine typenlosen Zeiger
zul"a"st, durch einen Zeiger auf ein Byte ersetzt werden kann.
Im Falle eines Funktionsaufrufs ist es n"otig, mehrere Ausdr"ucke
zu verketten, um die Parameter der Funktion anzugeben, soda"s
zus"atzlich ein Zeiger auf einen Ausdruck in die Struktur
eingebaut werden sollte, soda"s sie sich zu einer verketteten
Liste erweitern l"a"st.

Soviel zur Verwaltung der Ausdr"ucke im Speicher. Eine weit
wichtigere Frage jedoch ist, was der Interpreter mit diesen
Ausdr"ucken tut. Nun, zun"achst einmal mu"s jede der in
Kapitel \ref{exprparser} besprochenen Produktionsfunktionen so
konzipiert werden, da"s sie der "ubergeordneten Funktion einen Ausdruck
mit dem Wert des von ihr geparsten Ausdrucks zur"uckgibt. Diese
verwendet ihn ihrerseits, um den Wert ihres Ausdrucks zu
berechnen u.~s.~w. In der Kette der Funktionsaufrufe gibt es dann
irgendwann eine Funktion, die den Wert des Ausdrucks benutzt, um
damit etwas zu entscheiden, wie z.~B., ob der Zweig einer
\rw{if}-Anweisung ausgef"uhrt werden soll.

Solange man lediglich Berechnungen mit den Ausdr"ucken
anstellt, ergeben sich mit der oben skizzierten Methode
keinerlei Probleme, wie ist es jedoch m"oglich, den
Zuweisungsoperator zu implementieren? Bauen wir die Funktion zum
Parsen eines \fu{primary} so auf, da"s sie f"ur eine Variable
lediglich deren Wert zur"uckgibt, werden wir Probleme haben, denn
dann wei"s der Interpreter nicht, wo diese Variable im Speicher liegt und
kann ihr so auch keinen Wert zuweisen. Um dieses Problem zu
"uberwinden, wird intern ein zus"atzlicher Typ,
n"amlich der Referenztyp, der einem Zeiger sehr "ahnlich ist,
verwendet.
Trifft die Funktion zum Parsen eines \fu{primary} nun im Zwischencode auf eine Variable des Typs
\tint , dann gibt sie nicht deren Wert zur"uck, sondern ihre
Adresse und vermerkt, da"s es sich dabei um eine Referenz auf
einen \tint\ handelt, der Einfachheit halber von mir mit
\type{\reference int} bezeichnet.
Mu"s nun mit diesem Ausdruck gerechnet werden, gen"ugt es, ihn zu
dereferenzieren, soda"s ein Ausdruck mit dem Typ \tint\ entsteht.
Wird ihm ein Wert zugewiesen, mu"s dieser lediglich an die
Speicheradresse geschrieben werden, die er angibt. Soll der
un"are Operator \fu{\&} auf ihn angewendet werden, dann ist es
lediglich n"otig, das erste Byte im Typenarray so zu "andern,
da"s daraus ein Zeiger wird.
                
\subsection{Variablen}
\label{maintvars}

Variablen sind von den oben besprochenen Ausdr"ucken und Typen
insofern verschieden, als da"s sie nicht intern vom Interpreter
verwendet werden, sondern direkt mit dem Quellprogramm in
Verbindung stehen. Aus dem Modulkonzept und der Tatsache, da"s
Funktionsvariablen nur in der jeweiligen Funktion verf"ugbar
sind, folgt, da"s die Verwaltung von Variablen auf mehrern Stufen
erfolgen mu"s. Die umfassendste Stufe, allerdings auch diejenige
mit der niedrigsten Priorit"at, ist die der programmglobalen
Variablen. Auf sie folgt die Stufe der modulglobalen Variablen.
Die h"ochste Priorit"at besitzt die Stufe der Funktionsvariablen,
die allerdings in der Regel auch die wenigsten Variablen
aufnimmt.

Wird nun im Programm eine Variable verwendet, mu"s zun"achst
gepr"uft werden, ob es eine Funktionsvariable ihres Namens gibt.
Ist das nicht der Fall, kann es sich nur mehr um eine globale
Variable handeln. Gibt es eine modulglobale Variable, die den
gesuchten Namen hat, wird sie verwendet, andernfalls mu"s es eine
programmglobale Variable geben. Ist auch das nicht zutreffend,
hat es der Interpreter mit einem Programmfehler zu tun (siehe Kapitel
\ref{errortrapping}).

Die komplizierteste Verwaltungsstufe ist die der
Funktionsvariablen (es ist am einfachsten, bei der Verwaltung der
Variablen nicht zwischen Parametern und direkten Funktionsvariablen
zu unterscheiden), was daran liegt, da"s nicht alle
Funktionsvariablen in einen "`Topf"' geworfen werden k"onnen,
sondern da"s f"ur jede Funktion eine eigene Variablenliste
angelegt werden mu"s. Die einfachste L"osung f"ur dieses Problem
ist, einen Stack zu verwenden, auf dem f"ur jeden
Funktionsaufruf eine Liste aller Funktionsvariablen angelegt
wird, "ahnlich dem Code, den ein Compiler erzeugen w"urde. Analog
dazu werden bei der R"uckkehr aus einer Funktion die vorher
erzeugten Variablen wieder entfernt. Um nun
nach einer Funktionsvariablen zu suchen, mu"s nur der oberste
Stacklevel in Betracht gezogen werden, da er alle Variablen
beinhaltet, die zu der gerade auszuf"uhrenden Funktion geh"oren.

Wie schon in Kapitel \ref{modules} angesprochen, stellt der
Interpreter f"ur jedes Modul eine eigene Datenstruktur zur
Verf"ugung, die unter anderem alle modulglobalen Variablen dieses
Moduls verwaltet. Ist also nach einer modulglobalen Variablen zu
suchen, kann man die Suche auf das durch \fu{MP} angegebene Modul
beschr"anken.

Der einfachste Fall liegt bei programmglobalen Variablen vor.
F"ur ihre Verwaltung gen"ugt es, eine einzige Liste
bereitzustellen, die die programmglobalen Variablen aller Module
bereitstellt, ohne Unterscheidung, in welchem Modul diese
definiert sind.

\bigskip

Bis jetzt habe ich, wenn es um die Details der
Variablenverwaltung ging, lediglich von "`Listen"' gesprochen,
wobei ich damit weniger dynamische Listen, als vielmehr
geeignete Datenstrukturen f"ur das schnelle Auffinden von
Variablen gemeint habe. Diese Datenstrukturen m"ochte ich jetzt
genauer besprechen.

Das Problem, auf das wir hier treffen ist, eine
Datenstruktur (und nat"urlich auch die dazugeh"origen
Algorithmen) aussuchen zu m"ussen, die es erm"oglicht, einen
Datensatz in einer Menge von Datens"atzen m"oglichst
schnell zu finden. Dieses Problem ist ein gutes Beispiel f"ur das
Grundproblem eines Programmdesigners: F"ur eine bestimmte Aufgabe
den besten Algorithmus bzw. die beste Datenstruktur (was in den
meisten F"allen Hand in Hand geht) auszuw"ahlen. G"abe es eine
absolut beste L"osung, dann h"atte der Designer seine Aufgabe
schnell erledigt, doch leider gibt es eine solche nur in den
seltensten F"allen.

Das wichtigste Entscheidungskriterium f"ur unsere Datenstruktur
ist, wieviel Datens"atze zu verwalten sind. F"ur die
Variablenverwaltung eines Interpreters bieten sich grunds"atzlich
zwei Strukturen an: Bin"are Suchb"aume und Hash-Tabellen. Um
entscheiden zu k"onnen, welche der beiden Strukturen f"ur unsere
Anwendungen die bessere ist, ist es n"otig, ein wenig zu rechnen.
Der Einfachheit halber gehen wir vom g"unstigsten Fall aus (die
Ergebnisse sind von denen des durchschnittlichen Falls nicht sehr
verschieden). Dieser ist, da"s der Benutzer die Variablennamen so
gew"ahlt hat, da"s sich eine gleichm"a"sige Anordung in der
Hash-Tabelle ergibt bzw. der Suchbaum ausgeglichen ist.

In einer Hash-Tabelle, bei der die Eintr"age f"ur gleiche
Hash-Werte linear verkettet sind, werden f"ur eine Suche nach
einem Datensatz im Mittel $\max (N / 2M + 1/2, 1)$ Vergleiche
ben"otigt, wohingegen f"ur die Suche in einem Baum etwa
$\log_2N - 1$ Vergleiche notwendig sind, wobei $N$ die Anzahl der
Datens"atze und $M$ die Gr"o"se der Hash-Tabelle angibt. Auf den
ersten Blick scheint der Baum die bessere L"osung zu sein, da die
Anzahl der ben"otigten Vergleiche f"ur gro"se $N$ in etwa
proportional zum bin"aren Logarithmus der Anzahl der Datens"atze
ist, bei einer Hash-Tabelle allerdings die Anzahl der Vergleiche
direkt proportional zur Anzahl der Datens"atze ist, was f"ur
gro"se $N$ nat"urlich bedeutend schlechter ist, als die
Leistungsdaten des Baums. Nehmen wir allerdings an, wir benutzen
eine Hash-Tabelle mit 37 Eintr"agen, was in Bezug auf den
Speicherplatz "uberhaupt kein Problem darstellt. In diesem Fall
sind f"ur eine Suche in einer Hash-Tabelle mit 100 Datens"atzen
im Durchschnitt weniger als zwei Vergleiche n"otig, wohingegen in
einem Baum mehr als sechs Vergleiche ben"otigt werden. Erst wenn
mehr als 550 Datens"atze vorhanden sind, ist die Suche in einem
Baum g"unstiger als in einer Hash-Tabelle.

Es ist kaum anzunehmen, da"s jemand f"ur einen Interpreter ein
Programm schreibt, welches mehr als 300 globale Variablen verwendet,
soda"s es f"ur unsere Zwecke die bei weitem bessere L"osung ist,
f"ur die Verwaltung programm- und modulglobaler Variablen
Hash-Tabellen zu verwenden, wobei eine Gr"o"se von 37 durchaus
praktikabel ist, sowohl in Bezug auf Schnelligkeit als auch auf
die vern"unftige Ausnutzung des Speicherplatzes. F"ur
Funktionsvariablen sind 37 Eintr"age wohl "uberdimensioniert.
Hierf"ur bieten sich die Primzahlen 7 und 11 an.

Zuletzt stellt sich noch die Frage nach der Hash-Funktion. Ich
verwende sowohl der Einfachheit als auch der Effizienz halber die
Modulo-Funktion, d. h., ich betrachte eine Zeichenkette als lange
Ganzzahl, deren Rest bez"uglich des Moduls 37 den Index in der
Hash-Tabelle angibt. Die folgende Implementation geht davon aus, 
da"s ein \tchar\ 256 verschiedene Werte annehmen kann.

\begin{verbatim}
function hash (s : *char, m : int) : int

  h : int = 0;

begin
  while (*s)
    h = (h * 256 + *s) % m;
    s = s + 1;
  end
  return h;
end
\end{verbatim}

\subsection{Funktionen}

Die Funktionsverwaltung hat in erster Linie die Aufgabe, eine
Funktion eines bestimmten Namens zu finden und deren Adresse im
Zwischencode zu ermitteln. Weiters ist es wichtig, Namen und
Typen der Funktionsparameter zu erfahren, um erstens "uberpr"ufen
zu k"onnen, ob die Funktion mit der richtigen Anzahl und den
richtigen Typen der Parameter aufgerufen worden ist, und zweitens,
um die Parameter auf dem Funktionsstack anlegen zu
k"onnen.\footnote{siehe Kapitel \ref{maintvars}}

Zus"atzlich sollte es m"oglich sein, nicht nur Funktionen, die in SEL
implementiert sind, aufzurufen, sondern auch solche, die in der
Implementationssprache des Interpreters geschrieben sind,
da eine Implementation der Grundfunktionen von SEL in
SEL selbst f"ur einen Interpreter viel zu ineffizient w"are.
Deshalb schreibt man Funktionen wie etwa zum Kopieren von
Strings in der Implementationssprache des Interpreters (in
meinem Falle C) und ruft diese auf. Um das nahtlos in die
Verwaltung einbinden zu k"onnen, ist es notwendig, da"s die
Implementationssprache Zeiger auf Funktionen unterst"utzt. Ist
diese M"oglichkeit nicht vorhanden (wie etwa im Falle von
\pascal ), dann mu"s jede einzelne Funktion im Programmcode extra
behandelt werden, was sowohl umst"andlicher als auch langsamer
ist.\footnote{In Turbo Pascal ab der Version 5.5 ergibt sich die
M"oglichkeit, statt eines Zeigers auf eine Funktion einen Zeiger
auf ein Objekt zu verwenden, welches eine Methode enth"alt, die
die gew"unschte Funktion ersetzt.}

Um eine Funktion eines bestimmten Namens zu finden, werden
Hash-Tabellen, wie im letzten Unterkapitel beschrieben,
verwendet: Eine f"ur programmglobale Funktionen und eine
pro Modul f"ur modulglobale Funktionen, der Einfachheit halber
jeweils mit 37 Eintr"agen, was bei einem Programm mit 10 Modulen
unter Windows NT einen Speicherplatz von 1628 Bytes beansprucht
(was ohne schlechtes Gewissen als "`unbedenklich"' eingestuft
werden kann). Eintr"age f"ur denselben Hash-Wert werden
unsortiert linear verkettet.\footnote{Eine Sortierung h"atte bei
einer linearen Verkettung auch nicht den geringsten Sinn, da
ohnedies alle Eintr"age bis zum gew"unschten durchlaufen werden
m"ussen, unabh"angig davon, in welcher Reihenfolge die Eintr"age
gespeichert sind.}

Die Datenstruktur, die eine einzelne Funktion beschreibt, mu"s
neben deren Namen und R"uckgabetyp auch eine Angabe dar"uber
enthalten, wo der Code der Funktion zu finden ist. Hierf"ur wird
zun"achst eine Variable ben"otigt, die Auskunft dar"uber gibt, ob
die Funktion in SEL geschrieben ist oder in der
Implementationssprache des Interpreters. Im ersten Fall mu"s die
Datenstruktur Auskunft dar"uber geben k"onnen, in welchem Modul
und an welcher Stelle der Code beginnt, andernfalls wird ein
Zeiger auf die jeweilige Funktion ben"otigt. Zus"atzlich m"ussen
Informationen "uber Namen und Typen der Parameter gespeichert
sein. Ich verwende hierzu eine einfach verkettete Liste.

\bigskip

Interessant ist die Realisierung von Funktionsaufrufen. Dazu ist
zun"achst zu bedenken, da"s der Parser rekursiv implementiert ist
und der Funktionsaufruf ein \fu{primary} ist, sich also auf der
untersten Grammatikebene der Ausdr"ucke befindet. Aus diesem
Grund wird man schwer umhinkommen, auch Funktionsaufrufe rekursiv
zu behandeln, was nat"urlich bei einem rekursiven SEL-Programm zu
einer starken Belastung des Prozessorstacks f"uhren kann, weshalb
dieser m"oglichst gro"s dimensioniert werden sollte. Mu"s nun
innerhalb eines \fu{primary} eine Funktion aufgerufen werden,
wird in den oben besprochenen Hash-Tabellen die Funktion gesucht
und, wenn vorhanden, wird an die angegebene Stelle im
Zwischencode gesprungen. Dort werden zun"achst der Funktionskopf
und die Deklaration der direkten Funktionsvariablen interpretiert
und anschlie"send der Code der Funktion mittels des
Anweisungs-Parsers ausgef"uhrt. Ist die Ausf"uhrung der Funktion
beendet, mu"s zur aufrufenden Funktion zur"uckgekehrt, also die
Werte der Variablen \fu{MP} und \fu{IP} restauriert werden.
Hierzu m"ussen klarerweise die alten Werte irgendwo
zwischengespeichert werden. Weiters m"ussen die lokalen Variablen
der Funktion abgelegt werden, und zwar so, da"s auch mehrere
Instanzen dieser Funktion existieren k"onnen und jede ihre
eigenen Variablen zur Verf"ugung hat. Wie schon im letzten
Kapitel angedeutet, wird hierzu ein Stack verwendet, in Anlehnung
an den Prozessorstack, der bei compilierten Programmen dieselben
Aufgaben wahrnimmt. Relevant ist jeweils nur der oberste Level
des Stacks, der die R"ucksprungadresse und die lokalen Variablen
der jeweils aktiven Funktion enth"alt. Wird aus dieser eine
andere Funktion aufgerufen, kommt ein Level zum Stack hinzu,
andernfalls wird ein Level entfernt. Um sich nicht auf die
relativ langsame dynamische Allokation verlassen zu m"ussen, ist
es am besten, diesen Stack als Array zu
konzipieren.\footnote{Eine Konzeption als dynamische Liste hat
auf den ersten Blick den Vorteil, da"s die Rekursionstiefe
lediglich vom vorhandenen Speicherplatz abh"angig ist, was jedoch
bei einem rekursiven Interpreter nicht zutrifft: dort ist der
Prozessorstack das Ma"s aller Dinge.}

\section{Fehlerbehandlung}
\label{errortrapping}

Da nicht davon auszugehen ist, da"s das Programm, das
interpretiert werden soll, fehlerfrei ist, mu"s der Interpreter
auch die M"oglichkeit haben, Fehler zu entdecken und sie
entsprechend zu behandeln. Im Gegensatz zu Compilern, die wenn
m"oglich nach dem Erkennen eines Fehlers mit dem "Ubersetzen
fortfahren sollten, um evtl. weitere Fehler zu entdecken, ist das
bei einem Interpreter nicht n"otig und auch wenig sinnvoll, da
das auszuf"uhrende Programm in seiner Stabilit"at gest"ort und
vermutlich abst"urzen w"urde. Eine sogenannte "`{\em Error
Recovery}"' bleibt uns also erspart. Die Ziele f"ur die
Fehlerbehandlung sind also das Erkennen von Fehlern und die
Ausgabe einer f"ur den Benutzer m"oglichst hilfreichen
Fehlermeldung.

Grunds"atzlich mu"s ein Interpreter zwei Arten von Fehlern
behandeln: syntaktische und semantische Fehler. Syntaktische
Fehler treten dann auf, wenn die im Anhang \ref{grammar}
angegebene Grammatik von SEL durch das Programm verletzt wird.
Diese Art von Fehler ist in der Regel leicht zu erkennen und kann
dem Benutzer relativ verst"andlich erkl"art werden. Ein Beispiel:
Der Interpreter ist gerade besch"aftigt, die Anweisung
\fu{while [x != 10)} auszuf"uhren. Der Parser erkennt das
Schl"usselwort \rw{while}, dem laut Grammatik immer
eine "offnende runde Klammer folgen mu"s. Offenbar hat sich der Benutzer
jedoch vertippt und statt der runden eine "offnende eckige
Klammer eingegeben. Es ist leicht, den Parser so zu erweitern,
da"s er an dieser Stelle den Fehler erkennt. Die entsprechende
Fehlermeldung sollte lauten: \fu{'(' expected}.

Semantische Fehler treten auf, wenn das Programm versucht,
unerlaubte Operationen auszuf"uhren, wie etwa das Casten eines
\tfloat\ in einen Zeiger, die Verwendung einer Variablen,
die gar nicht existiert, oder die Anwendung eines Operators auf
einen nicht erlaubten Typen. Tritt ein solcher Fall auf,
l"a"st sich das in der Regel ohnedies leicht feststellen und auch
eine passende Fehlermeldung ist nicht schwer zu schreiben.

F"ur den Benutzer des SEL-Interpreters ist es klarerweise nicht
nur wichtig, da"s ein Fehler aufgetreten und welcher Art
dieser ist, sondern auch, in welchem Modul in welcher Zeile der
Fehler steckt. Da der Interpreter weder syntaktische noch
semantische Fehler erkennen kann, w"ahrend er den Quellcode
untersucht, und der Zwischencode aus Effizienzgr"unden keine
Zeilenumbr"uche enth"alt, m"ussen die Zeilennummern gesondert
gespeichert werden. Ich verwende hierzu ein durch die
Zeilennummer indiziertes Array, das f"ur jede
Zeile eines Moduls die Anzahl der Bytes angibt, die diese Zeile
im Zwischencode beansprucht. Da der Zwischencode in der Regel
k"urzer als der Quellcode und kaum eine Programmzeile l"anger
als 100 Zeichen ist, gen"ugt es, pro Arrayelement ein Byte zur
Verf"ugung zu stellen. Tritt nun an der Stelle \fu{MP==3} und
\fu{IP==1364} ein Fehler auf, dann mu"s im Modul Nr. 3 das
Zeilennummernarray so lange durchlaufen werden, bis die Summe der
Elemente gr"o"ser als 1364 ist. Der folgenden Funktion, die die
gesuchte Zeilennummer zur"uckgibt, ist
das Zeilennummernarray als Zeiger und der Wert von \fu{IP} zu
"ubergeben:

\begin{verbatim}
function line (a : *char, ip : int) : int

  l : int = 0;
  c : int = 0;
  
begin
  repeat
    c = c + a[l];
    l = l + 1;
  until (c > ip)
  return l;
end
\end{verbatim}

In einem rekursiven Interpreter ist neben dem Erkennen und
Anzeigen eines Fehlers noch ein weiteres Problem zu "uberwinden:
das au"serordentliche Beenden des Interpreters. Es ist davon
auszugehen, da"s der Interpreter schon eine gewisse
Rekursionsebene erreicht hat, wenn der Fehler auftritt, soda"s es
nicht so einfach ist, ihn zu beenden und evtl. zu einer (den
Interpreter) aufrufenden Funktion zur"uckzukehren. Die
ANSI-Bibliothek von C bietet hierzu das praktische Konzept der
"`{\em Long-Jumps}"' an, das ungef"ahr einem \rw{goto}
gleichkommt, das nicht nur Spr"unge innerhalb einer Funktion
erlaubt, sondern auch aus einer Funktion in eine andere,
"ubergeordnete Funktion hinein. In C++ z.~B. ist es nicht einmal
mehr n"otig, sich auf Bibliotheksfunktionen zu st"utzen, da ein
"ahnliches Konzept (das sog. "`{\em Exception-Handling}"') in der
Programmiersprache selbst eingebaut ist. Leider besitzen nicht
alle Sprachen einen derartigen Luxus. In \pascal\ etwa mu"s man
sich einer anderen Methode bedienen. Eine einfache und
naheliegende L"osung ist, eine Variable zu definieren, die
angibt, ob bei der Ausf"uhrung des Programms bisher ein Fehler
aufgetreten ist. Eine Funktion, die einen Fehler entdeckt, hat
diesen zu melden und die Variable auf einen Wert zu setzen, der
anzeigt, da"s ein Fehler gefunden wurde. Kehrt diese Funktion zu
der aufrufenden Funktion zur"uck, mu"s diese "uberpr"ufen, ob die
Variable gesetzt ist und in diesem Fall zu ihrer aufrufenden
Funktion zur"uckkehren, so lange, bis man an der Funktion
angelangt ist, die die Ausf"uhrung des Programms initiiert hat.
Zwar ist diese Methode simpel, doch hat sie einen gewissen
Overhead zur Folge, der nicht nur den Code gr"o"ser, sondern auch
un"ubersichtlicher macht.
                                   
\newpage
\section{Weiterf\protect{"uhrende} Projekte}

Ebenso wie im ersten Kapitel m"ochte ich auch hier einige meiner
Ideen darlegen, die leider zu umfangreich sind, als da"s ich sie
ausf"uhrlich in dieser Arbeit h"atte behandeln k"onnen.

\begin{description}
\item[Verwendung einer Baumstruktur] anstatt des Zwischencodes
k"onnte die Geschwindigkeit des Interpreters um einiges erh"ohen.
Diese Baumstruktur w"urde in gewisser Weise einem
Parse-Tree\footnote{siehe Kapitel \ref{topdbotu}} "ahneln, bei
dem jedoch die unn"otigen Bl"atter entfernt sind. Zus"atzlich
k"onnte vor der Ausf"uhrung des Programms der Baum optimiert
werden, z.~B. durch Vereinfachung von konstanten Ausdr"ucken
(\fu{2*3+4} kann ohne Ver"anderung der Laufzeiteigenschaften zu
\fu{10} umgeformt werden). Eine Darstellung als Baum hat f"ur den
Interpreter den gro"sen Vorteil, da"s die lexikalische Analyse
und der Parsevorgang nur einmal ausgef"uhrt werden m"ussen. Ist die
Baumstruktur gut konzipiert, k"onnte der Parser sogar als
Front-End f"ur einen Compiler dienen.

\item[Ein Compiler f"ur SEL] ist mein n"achstes Ziel. Leider sind
jedoch die meisten Maschinensprachen entweder zu komplex (wie
etwa die des Intel 80386) oder zu umst"andlich (wie die des Intel
8086), als da"s es Spa"s machen k"onnte, f"ur diese Prozessoren
einen Compiler zu schreiben. Aus diesem Grund erw"age ich zwei
M"oglichkeiten, um diese Schwierigkeiten zu umgehen.
Die einfachere der beiden ist, ein SEL-Programm anstatt in
Maschinencode in eine andere h"ohere Programmiersprache,
vorzugsweise C, zu "ubersetzen. Besonders interessant w"are es,
diesen Compiler selbst in SEL zu schreiben und ihn dann durch sich
selbst "ubersetzen zu lassen (mit Hilfe des Interpreters).

Die zweite M"oglichkeit ist,
anstatt f"ur einen real existierenden Prozessor Maschinencode
f"ur einen einfachen, auf die Programmiersprache zugeschnittenen,
hypothetischen Prozessor zu erzeugen, wobei dieser Prozessor
nat"urlich zun"achst entworfen und auf einem existierenden
Prozessor emuliert werden m"u"ste. Obwohl auch hier ein Interpreter
im Spiel ist, ist die Ausf"uhrungsgeschwindigkeit eines
Pseudo-Code weit h"oher als die eines Interpreters, der den
Quellcode (oder auch einen Zwischencode) der Programmiersprache
interpretiert. Der Vorteil dieser beiden M"oglichkeiten ist, da"s
sie gr"o"stenteils prozessorunabh"angig sind. So existiert heute
beinahe f"ur jeden Prozessor eine ausreichende Menge an
C-Compilern und f"ur den P-Code m"u"ste lediglich der Interpreter
neu geschrieben, oder, falls auch er in einer h"oheren
Programmiersprache geschrieben ist, lediglich neu compiliert
werden. Aus Effizienzgr"unden ist es jedoch ratsam, einen
Interpreter f"ur P-Code in Assembler zu implementieren.

\item[Ein Pretty-Printer] ist ein Programm, das einen Quelltext
einer bestimmten Programmiersprache in eine ansprechende Form
bringt, also Einr"uckungen vornimmt, wenn n"otig, und Kommentare
einheitlich gestaltet. F"ur C existiert eine Reihe solcher
Programme, das bekannteste ist wohl C-Indent. Es ist eine
interessante, wenn auch nicht "uberm"a"sig schwere Aufgabe, einen
solchen Pretty-Printer f"ur SEL zu schreiben.
\end{description}

\begin{appendix}

\chapter{Die Grammatik von SEL}
\label{grammar}

\begin{verbatim}
module:        {block}
             ;
block:         function
             | local
             | global
             ;
function:      `function' IDENTIFIER `(' [formal_params] `)' `:' type
                 {func_var}
               `begin'
                 {statment}
               `end'
             ;
formal_params: IDENTIFIER `:' [`&']  type [`,' formal_params]
             ;
func_var:      IDENTIFIER `:' type [`=' expr] `;'
             ;
func_proto:    `function' IDENTIFIER
               `(' [[`&' type {`,' [`&'] type}] `)' `:' type
             ;
local:         `local' {local_decl} `end'
             ;
local_decl:    struct
             | func_proto
             | var
             ;
global:        `global' {global_decl} `end'
             ;
global_decl:   struct
             | func_proto
             | var
             | extern
             ;
\end{verbatim}
\newpage
\begin{verbatim}
statement:     [expr] `;'
             | if
             | while
             | for
             | repeat
             | switch
             | prefix
             | `break' `;'
             | `return' [expr] `;'
             ;
if:            `if' `(' expr `)' {statement} [`else' {statement}] `end'
             ;
switch:        `switch' `(' expr `)' {statement} `end'
             ;
prefix:        `case' constant `:'
             | `default' `:'
             ;
while:         `while' `(' expr `)' {statement} `end'
             ;
repeat:        `repeat' {statement} `until' `(' expr `)'
             ;
for:           `for' `(' primary `=' expr (`to' | `downto') expr
                         [`step' expr] `)'
                 {statement}
               `end'
             ;
expr:          expr (`=' | `&&' | `||' | `^^' | `==' | `<' | `>'
                   | `<=' | `>=' | `!=' | `&' | `|' | `^' | `<<'
                   | `>>' | `+' | `-' | `*' | `/' | `%') expr
             | (`!' | `~' | `-' | `*' | `&') expr
             | `delete' expr
             | primary
             ;
primary:       IDENTIFIER
             | `(' expr `)'
             | CHAR_CONSTANT
             | INT_CONSTANT
             | FLOAT_CONSTANT
             | STRING_CONSTANT
             | primary `->' IDENTIFIER
             | primary `:' type
             | primary `[' expr `]'
             | IDENTIFIER `(' [expr {`,' expr}] `)'
             | `sizeof' type
             | `new' type
             ;
var:           IDENTIFIER `:' type [`=' constant] `;'
             ;
extern:        `extern' IDENTIFIER `:' type `;'
             ;
\end{verbatim}
\newpage
\begin{verbatim}
struct:        `struct' IDENTIFIER {IDENTIFIER `:' type `;'} `end'
             ;
type:          `char'
             | `int'
             | `float'
             | `void'
             | `[' INT_CONSTANT `]' type
             | `*' type
             | IDENTIFIER
             ;
constant:      CHAR_CONSTANT
             | [`-'] INT_CONSTANT
             | [`-'] FLOAT_CONSTANT
             | STRING_CONSTANT
             ;
INT_CONSTANT:    [0-9]+
CHAR_CONTANT:    '(.|\.)'
FLOAT_CONSTANT:  [0-9]*(`.'[0-9]*)?((e|E)(+|-)?[0-9]+)?
STRING_CONSTANT: "([^"\n]|\.)*"
IDENTIFIER:      [A-Za-z_][A-Za-z0-9_]*
\end{verbatim}

\end{appendix}

\begin{thebibliography}{9}
\addcontentsline{toc}{chapter}{Literaturverzeichnis}
\bibitem{selkr} Mark Probst. {\em SEL-Kurzreferenz}, Version~0.3.
Beilage.
\bibitem{kandr} Brian W.~Kernighan, Dennis M.~Ritchie. {\em
Programmieren in C}, deutsche Ausgabe bearbeitet
von A.~T.~Schreiner, Ernst Janich. Carl Hanser Verlag, Wien,~1983.
ISBN~3-446-13878-1.
\bibitem{sedge} Robert Sedgewick. {\em Algorithmen}, deutsche
Ausgabe. Addison-Wesley~GmbH, Deutschland, 1.~Auflage~1991.
ISBN~3-89319-301-4.
\bibitem{duden} {\em Duden Informatik}, 2.~Auflage. Dudenverlag,
Mannheim,~1993. ISBN~3-411-05232-5.
\end{thebibliography}

% \newpage
% \addcontentsline{toc}{chapter}{Index}
% \printindex

\end{document}

\begin{document}

\title{Design einer Programmiersprache\\
       und theoretische Besprechung\\
       eines Interpreters anhand der\\
       Programmiersprache SEL}
\author{Mark Probst}

\maketitle

\tableofcontents

\chapter{Die Programmiersprache}

Dieses Kapitel hat das Ziel, das Design der Programmiersprache
SEL, wie in \cite{selkr} beschrieben, zu erkl"aren bzw. zu
begr�nden, und zwar in Hinblick auf den Bau eines Interpreters
und eines Compilers. Aus diesem Grund setze ich eine gewisse
Kenntnis sowohl einer strukturierten Programmiersprache,
insbesondere nat"urlich der Sprache SEL, als auch der
grundlegenden Funktionsweise eines Interpreters bzw. eines
Compilers beim Leser voraus. Der Leser sollte also zumindest
\cite{selkr} gelesen haben, da"s, obwohl es als Kurzreferenz
konzipiert ist, durchaus auch geeignet ist, einen "Uberblick
"uber die Programmiersprache SEL zu vermitteln.

\section{Der Begriff der Programmiersprache}

Grunds"atzlich kann eine Programmiersprache definiert werden als

\begin{quotation}
Sprache zur Formulierung von Rechenvorschriften, d. h. von
Datenstrukturen und Algorithmen, die von einem Computer
ausgef"uhrt werden k"onnen.\footnote{siehe \cite{duden},
Seite~544}
\end{quotation}

Die in dieser Arbeit behandelte Programmiersprache SEL geh"ort
zur Gruppe der imperativen Programmiersprachen, deren wichtigste
Untergruppe wohl die der strukturierten Programmiersprachen ist.
Imperative Programmiersprachen werden im Duden Informatik
wie folgt definiert:

\begin{quotation}
Bei diesen Sprachen besteht ein Programm aus einer Folge von
Befehlen an den Computer, wie z.~B. "`Schreibe in die Variable a
den Wert 3"', "`Springe an die Stelle x im
Programm"'.\footnote{siehe \cite{duden}, Seite~545}
\end{quotation}

Strukturierte Programmiersprachen stellen f"ur den Programmierer
zus"atzlich Hilfen in Form von Kontrollanweisungen, Schleifen und
benutzerdefinierten Funktionen zur Verf"ugung. Diese machen
sowohl das Schreiben der Programme einfacher und schneller wie
auch die Programmtexte leichter lesbar. Zur Zeit sind sie die am
meisten verwendete Art der Programmiersprachen.

\section{Ziele f\protect{"ur} das Design von SEL}

Bevor man an die Arbeit an einem so komplexen Projekt wie das Design
einer Programmiersprache herangeht, sollte man sich im Klaren
dar"uber sein, welche Ziele man eigentlich verfolgt. So habe ich
auch an meine zuk"unftige Programmiersprache einige Anforderungen
gestellt, von denen ich die meisten in der gegenw"artigen Version
der Sprache verwirklicht sehe.

Die wichtigste Entscheidung bei einem Projekt ist die nach dem
Zweck desselben. Im Fall von SEL wollte ich weniger eine
Programmiersprache f"ur das Schreiben von kommerzieller Software
schaffen, als vielmehr eine kleine, feine Sprache, die ich in
vielen Situationen benutzen k"onnte. Haupts"achlich schwebte mir
dabei die Verwendung als Makrosprache in diversen Programmen vor.
Eine solche Makrosprache sollte vor allem leicht erlernbar sein
und keine gro"se Umstellung von Seiten des Benutzers verlangen,
sollte aber andererseits nicht so primitiv aufgebaut sein, da"s
sie eine vern"unftige Verwendung nicht mehr m"oglich macht, was
die meisten Makrosprachen leider sind.

Um eine geringe Umstellungszeit m"oglich zu machen, mu"s die
Programmiersprache an eine bereits bestehende, weit verbreitete
Sprache angelehnt sein. Im Fall von SEL habe ich mich
haupts"achlich an C und \pascal\ gehalten, weil mir diese am
gel"aufigsten und vermutlich die am weitest verbreiteten
Sprachen sind.

Da ich ein gro"ser Freund von C und ihrer pr"agnanten und kurzen
Ausdrucksweise bin, habe ich mir zum Ziel gesetzt, auch SEL zu
einer effizienten, kurzen Ausdrucksweise zu verhelfen. Eine
Konsequenz aus dieser Entscheidung ist es, anstatt von
Schl"usselw"ortern, so wie in \pascal , Symbole wie in C zu
verwenden, soferne dies nicht die "Ubersichtlichkeit und
Konsistenz innerhalb der Sprache zu sehr gef"ahrdet. Die zweite
Konsequenz ist es, redundante Sprachteile zu entfernen, solange
auch das nicht die Lesbarkeit der Programmtexte verletzt.

Ein weiteres wichtiges Kriterium bei einer Makrosprache ist die
Erweiterungsf"ahigkeit, da Makrosprachen in so verschiedenen
Anwendungen wie z.~B. Textverarbeitungen und Zeichenprogrammen
eingesetzt werden, da"s es schlichtweg unm"oglich ist, eine
Sprache zu konstruieren, die f"ur alle diese Anwendungsgebiete
passend ist. Diese Erweiterungsf"ahigkeit sollte einerseits durch
die Erweiterung mit Hilfe von Funktionen wie z.~B.
f"ur das Zeichnen von Linien oder Kreisen und andererseits durch
die M"oglichkeit der "Anderung in der Grammatik, wie etwa
f"ur das flexiblere Handling von Vektoren und Matrizen, gegeben
sein. Einige Vorschl"age f"ur "Anderungen der letzteren Art sind
in Kapitel \ref{langext} gegeben.

Zuletzt soll die Programmiersprache SEL so aufgebaut sein, da"s
es nicht nur m"oglich ist, einen Interpreter, sondern auch mit
vertretbarem Aufwand einen Compiler zu implementieren. Nicht nur,
weil das bei einer Programmiersprache immer eine gute Idee ist,
sondern auch in Hinblick darauf, da"s ich selbst einmal einen
Compiler schreiben m"ochte.

Als Abschlu"s dieses Unterkapitels seien nochmals alle Forderungen
kurz zusammengefa"st, um in sp"ateren Abschnitten darauf
zur"uckgreifen zu k"onnen.

\begin{itemize}
\item Geringer Sprachumfang, jedoch auch f"ur mittelgro"se
Programme geeignet
\item Anlehnung an C und \pascal\
\item Verwendung von Symbolen anstatt von Schl"usselw"ortern
(wenn m"oglich)
\item Vermeidung von Redundanzen
\item Konsistenz innerhalb der Sprache
\item Erweiterungsf"ahigkeit
\item Sowohl f"ur Interpreter als auch f"ur Compiler geeignet
\end{itemize}

\section{Das Programm}

Das gr"o"ste in sich abgeschlossene Objekt, da"s ein Compiler bzw.
Interpreter zu behandeln hat, ist das Programm. Uns stellt sich nun die
Frage, wie dieses Programm aufgebaut sein soll.

Zun"achst ist zu
entscheiden, ob das Programm in Module eingeteilt werden soll oder
nicht. Eine solche Einteilung hat den Vorteil, da"s der Programmierer
Implementationsdetails eines Programmteils oder einer Funktionssammlung
vom Rest des Programms abschirmen kann und den Anwender dieser
Funktionssammlung gar nicht erst in Versuchung bringt, in die internen
Abl"aufe einzugreifen. Nat"urlich ist das nicht der einzige Vorteil
dieser Technik. So kann es z.~B. bei einer Modularisierung des Programms
auch kaum zu Namens"uberschneidungen bei Funktionen und Variablen
kommen. Last but not least wird auch das Programm leichter
"uberschaubar, weil man sich jeweils nur auf ein Modul konzentrieren
mu"s. Ich habe mich f"ur eine Modularisierung entschieden. Erstens, weil
es die M"oglichkeiten des Programmierers um einen nicht zu
untersch"atzenden Faktor erweitert und zweitens, weil es nicht schwer
ist, ein solches Modulkonzept in einen Interpreter oder einen Compiler
einzubauen---jedenfalls nicht das hier verwendete.

In SEL wird ein Modul in mehrere Bl"ocke eingeteilt, die in beliebiger
Reihenfolge auftreten d"urfen, wobei es drei Arten
von Bl"ocken zu unterscheiden gibt:

\begin{itemize}
\item Die \rw{global}-Bl"ocke werden verwendet, um programmglobale Variablen,
Funktionen oder Datentypen
zu deklarieren, also solche, die im ganzen Programm bekannt
sind.
\item Die \rw{local}-Bl"ocke werden analog zu \rw{global}-Bl"ocken
verwendet, um
modulglobale Variablen, Funktionen oder Datentypen zu deklarieren.
\item Funktionen werden "ahnlich wie in C und \pascal\ definiert.
\end{itemize}

Der Hauptgrund, warum ich mich f"ur diese Art der Modulstruktur
entschieden habe ist der, da"s es f"ur einen Interpreter weit
praktikabler ist, zur Initialisierung der Variablen vor dem Start eines
Programms einfach nach den Schl"usselw"ortern \rw{global} und
\rw{local} zu suchen, als zwischen den Funktionsdefinitionen. Dieses
Problem h"atte nat"urlich auch eine Programmstruktur wie die in \pascal ,
wo Variablen- und Typendeklarationen ihren festen Platz haben, gel"ost,
doch erscheint mir dieser Ansatz als zu starr, vermutlich, weil ich die
Freiheiten von C gewohnt bin.

Noch ein kleines Wort zur allgemeinen Struktur eines Moduls: Wie
bereits oben erw"ahnt, besteht ein Modul aus einer beliebigen
Abfolge von Bl"ocken, wobei es durchaus m"oglich ist, da"s
mehrere Bl"ocke ein und desselben Typs vorkommen. F"ur das
gesamte Modul gilt jedoch eine Regel: Ein Objekt ist erst ab der
Zeile bekannt, in der es deklariert wurde, das hei"st, da"s es
z.~B. nicht m"oglich ist, eine Variable in Zeile 154 zu
verwenden, wenn sie erst in Zeile 186 deklariert wird.
Diese Regelung macht es einem Compiler m"oglich, ein Programm in
einem Durchgang zu "ubersetzen ({\em One-Pass Compiler}). Eine
Ausnahme bilden die Strukturen (siehe Kapitel \ref{structs}).

\subsection{global-Bl\protect{"ocke}}

Ein \rw{global}-Block\footnote{siehe Anhang \ref{grammar} oder
\cite{selkr}, Seite 9} wird durch das Schl"usselwort \rw{global}
eingeleitet und durch \rw{end} beendet. Dazwischen befindet sich
eine beliebige Abfolge von Variablen-, Funktions- und
Strukturdeklarationen. Alle in diesem Abschnitt deklarierten
Objekte sind dazu bestimmt, auch in anderen Modulen des Programms
verwendet zu werden, was jedoch nicht unbedingt bedeutet, da"s
sie in diesen auch ohne weiteres verf"ugbar sind. Soll ein
programmglobales Objekt auch in einem anderen Modul verwendet
werden, so ist es n"otig, es auch dort zu deklarieren und---im
Falle von Variablen---als \rw{extern} zu bezeichnen. Der Grund
f"ur diese Vorgangsweise ist, da"s sich ein Compiler, der die
Deklarationen aller anderen Module verwenden kann, nicht
konstruieren l"a"st, ohne nicht ein eigenes Objektformat zu
benutzen, wie das z.~B. bei Turbo \pascal\ der Fall ist. Das
Standard-Objektformat von MS-DOS und Windows NT ist dazu nicht in
der Lage. Au"serdem w"are es erheblich schwerer, Module
zu schreiben, die gegenseitig voneinander abh"angig sind, soda"s
ich von dieser L"osung Abstand genommen habe.

Auffallend an den drei Arten der Deklarationen, die innerhalb
eines \rw{global}-Blocks vorkommen d"urfen, ist, da"s die
Variablendeklaration als einzige mit einem Semikolon
abgeschlossen wird, und das, obwohl es ohne Gef"ahrdung der
Eindeutigkeit m"oglich w"are, dieses wegzulassen. Der Grund
hierf"ur ist einerseits die Einhaltung der Konsistenz der
Sprache\footnote{siehe auch Kapitel \ref{funcs}} und andererseits
die Offenhaltung f"ur Erweiterungen der Sprache. So ist es im
Moment zwar lediglich m"oglich, eine Variable mit einer \tchar -,
\tint -, \tfloat - oder Stringkonstante zu initialisieren, jedoch
ist es durchaus denkbar, dies in zuk"unftigen Versionen auf
Arrays, Strukturen und Zeiger auszuweiten, wo die Eindeutigkeit
nicht mehr unbedingt gegeben sein mu"s.

Der Grund, warum Variablen, die in einem anderen Modul definiert
worden sind und in einem bestimmten Modul verwendet werden, als
\rw{extern} zu deklarieren sind und dies bei Strukturen und
Funktionen nicht n"otig ist, ist einfach: Bei einer
Strukturdeklaration wird dem Compiler bzw. Interpreter lediglich
"`mitgeteilt"', wie diese aufgebaut ist. Es wird in diesem Fall
kein vom Programm verwendetes Objekt im Speicher erzeugt. Bei
einer Variablendefinition hingegen wird dem Compiler sowohl der
Typ der Variable mitgeteilt als auch der Befehl, sie im Speicher
zu erzeugen. Ohne Unterscheidung zwischen Variablendefinition und
\rw{extern}-Deklaration w"u"ste der Compiler nicht, ob er nun
eine Variable im Speicher zu erzeugen h"atte oder sich lediglich
vormerken sollte, da"s diese Variable in einem anderen Modul
existiert und in diesem verwendet wird. Bei Funktionen ist es
kaum anders: Eine Deklaration innerhalb eines \rw{local}- oder
\rw{global}-Blocks sagt dem Compiler lediglich, da"s es
"`irgendwo anders"' eine solche Funktion gibt und da"s sie hier
verwendet wird. Eine Funktionsdeklaration innerhalb eines
\rw{global}-Blocks dient weiterhin dazu, zu kennzeichnen, da"s
diese Funktion, falls sie in diesem Modul implementiert ist, f"ur
die Verwendung in anderen Modulen bestimmt ist. Eine zweite
M"oglichkeit, dies zu erzielen, w"are, dem Funktionskopf bei der
Implementation noch ein spezielles Schl"usselwort hinzuzuf"ugen,
doch erscheint mir die in SEL verwendete L"osung als die bessere
und auf jeden Fall als die "ubersichtlichere.

Variablen d"urfen in \rw{local} bzw. \rw{global}-Bl"ocken
lediglich mit Konstanten, in
Funktionen hingegen auch mit Ausdr"ucken initialisiert werden.
Dies ist einerseits der Einfachheit halber so, andererseits, weil
es immer eine riskante Sache ist, Werte zu berechnen, noch bevor
die einzelnen Teile des Programms initialisiert sind. Dar"uber hinaus w"are
es notwendig, eine Regelung zu erstellen, die angibt, welche
Variablen als erste initialisiert werden und welches Modul einem
anderen vorangeht. Dieser Aufwand erscheint mir allerdings etwas
"ubertrieben, nicht zuletzt deshalb, weil eine solche
Initialisierung leicht in der \fu{main}-Funktion\footnote{siehe
Kapitel \ref{main}} nachgeholt werden kann.

\subsection{local-Bl\protect{"ocke}}

Ein \rw{local}-Block\footnote{siehe Anhang \ref{grammar} oder
\cite{selkr}, Seite 9} unterscheidet sich von einem
\rw{global}-Block, abgesehen davon, da"s er durch ein anderes
Schl"usselwort eingeleitet wird, dadurch, da"s eine
\rw{extern}-Deklaration nicht erlaubt ist. Dies ist auch ganz
verst"andlich, denn es ist ein Widerspruch in sich, eine
modulglobale Variable, die in einem anderen Modul definiert
ist, zu verwenden.

Die Tatsache, da"s Funktionen, die nicht explizit in einem
\rw{global}-Block aufscheinen, automatisch modulglobal sind,
verleitet zum Schlu"s, da"s Funktionsdeklarationen in
\rw{local}-Bl"ocken redundant sind. Es k"onnte jedoch ohne
weiteres die Notwendigkeit entstehen, da"s zwei Funktionen sich
gegenseitig aufrufen. In diesem Fall w"are jedoch ein
Aufruf illegal, da Funktionen erst ab dem Punkt ihrer
Deklaration im Modul bekannt sind. Dieser Konflikt ist sehr
einfach zu l"osen, indem man eine der Funktionen (oder der
"Ubersichtlichkeit halber beide) in einem \rw{local}-Block
deklariert und diese erst nach der anderen implementiert. Ein
Beispiel: Die Funktion \fu{A} und \fu{B} rufen sich gegenseitig
auf. Nun wird vor der Implementation die Funktion \fu{A} als
\rw{local} deklariert, soda"s sie im Rest des Moduls bekannt ist.
Nun wird die Funktion \fu{B} implementiert. Der Aufruf von \fu{A}
stellt kein Problem dar, der oben geschilderten Deklaration
wegen. Die Implementation von \fu{A} ist nun auch kein
Problem mehr, da an diesem Punkt schon beide Funktionen bekannt
sind.

\subsection{Funktionen}
\label{funcs}

Im Gegensatz zu \pascal\ wird in SEL, so wie in C, nicht zwischen
Funktionen und Prozeduren unterschieden. Dies hat---anders als es
zu erwarten w"are---fast nur Vorteile. Auf der einen Seite ist es
m"oglich, Prozeduren durch Funktionen mit dem R"uckgabetyp
\tvoid \footnote{siehe Kapitel \ref{elemtypes}} nachzubilden,
andererseits vereinfacht es den Bau eines
Interpreters bzw. eines Compilers. Der einzige Nachteil ist ein
geringer Verlust an Lesbarkeit.

Eine Funktion\footnote{siehe Kapitel \ref{exprfuncs} und
Anhang \ref{grammar} oder
\cite{selkr}, Seiten 7--9 und 17} gliedert sich augenscheinlich in
drei gro"se Teile:
Den Funktionskopf, die Definitionen der direkten
Funktionsvariablen\footnote{siehe Kapitel \ref{variables}}
und den Funktionsk"orper.

Der Funktionskopf besteht aus dem Schl"usselwort \rw{function},
dem Namen der Funktion, einer in Klammern eingeschlossenen Liste
der Parameter inklusive deren Typen und dem R"uckgabetyp. Der
Doppelpunkt vor dem Typ soll andeuten, da"s eine Funktion in
gewisser Weise einer Variablen eines bestimmten Typs "ahnelt.
Es sei angemerkt, da"s dieser Doppelpunkt ohne weiteres
ausgelassen werden k"onnte, ohne die Eindeutigkeit der Sprache zu
gef"ahrden.\footnote{Tats"achlich k"onnte sogar die schlie"sende
Klammer ausgelassen werden, doch diente das sicherlich nicht der
Lesbarkeit.}

Bei genauerer Betrachtung des Funktionskopfes und der
Funktionsdeklaration in einem \rw{global} bzw. \rw{local}-Block
f"allt auf, da"s in einem Funktionskopf die Namen der
Parameter aufscheinen, in den beiden letztgenannten Bl"ocken dies
jedoch nicht der Fall ist. Betrachten wir zur Begr"undung dieser
Entscheidung einen Interpreter und einen Compiler. Der
Interpreter kann alle Module des Programms auf einmal betrachten,
er kennt also die Namen der Parameter auch dann, wenn sie in
einem anderen Modul erw"ahnt sind als in dem, von dem aus die
Funktion aufgerufen wird. Ein Compiler jedoch besitzt diesen
Vorteil nicht: Er beschr"ankt sich jeweils auf ein einzelnes
Modul. Bei einem Funktionsaufruf ist aber das Einzige, was der
Compiler zu tun hat, zu "uberpr"ufen, ob die Parameter die
richtigen Typen haben, um sie dann in der richtigen Reihenfolge
auf dem Stack abzulegen. Wirkliche Bedeutung haben die Namen der
Parameter nur f"ur die eigentliche Funktion und insofern stellt
die Beschr"ankung auf die Typen keine Einschr"ankung dar, es ist
jedoch durchaus denkbar, da"s es der "Ubersichtlichkeit wegen in
einer zuk"unftigen Version von SEL m"oglich sein wird, auch die
Namen anzugeben, wobei diese Informationen aber vom Interpreter
bzw. Compiler ignoriert werden.

Dem Funktionskopf folgen die Definitionen der direkten
Funktionsvariablen.
Diese werden bei jedem Funktionsaufruf neu auf dem Stack angelegt
und belegen daher bei einer rekursiven Funktion in den
verschiedenen Instanzen nicht denselben Speicherbereich---erst
diese Tatsache macht Rekursionen so effektiv. Es ist m"oglich,
die Variablen mit Werten zu initialisieren, die sich aus Ausdr"ucken ergeben.
Dies stellt weder einen Compiler noch einen Interpreter vor
Probleme, soda"s ich hier nicht n"aher darauf eingehen m"ochte.
Das Semikolon am Ende jeder Variablendefinition w"are nicht
unbedingt notwendig, ergibt jedoch sowohl in Hinsicht auf die
"Ubersichtlichkeit als auch auf die Erweiterungsf"ahigkeit der
Sprache durchaus Sinn.

Der Funktionsk"orper wird durch das Schl"usselwort \rw{begin}
eingeleitet und durch \rw{end} beendet. Dazwischen befindet sich
eine beliebige Anzahl von Anweisungen, die bei Aufruf der
Funktion ausgef"uhrt werden.

\subsubsection{Die main-Funktion}
\label{main}

Jedes Programm mu"s einen Punkt haben, an dem es beginnt, soda"s
nicht unkontrolliert irgendwelche Funktionen ausgef"uhrt werden.
Hierf"ur gibt es grunds"atzlich zwei M"oglichkeiten. Entweder,
man f"uhrt bestimmte Konstrukte ein, um diesen Hauptteil des
Programms zu kennzeichnen, so wie es z.~B. in \pascal\
praktiziert wird, oder man verwendet eine bestimmte Funktion
dazu, wie in C. Ich habe mich f"ur letzteren
L"osungsansatz entschieden, weil er es erlaubt, das gesamte
Programm als eine Art Funktion zu sehen, die vom Anwender
aufgerufen werden kann. Es wird auch m"oglich, die
Kommandozeilenparameter als einfache Funktionsparameter zu
"ubergeben---wenn dies auch in der aktuellen Version von SEL
noch nicht eingef"uhrt ist. Weiters kann durch dieses Prinzip dem
Betriebssystem ein Wert zur"uckgegeben werden---der sogenannte
{\em Exit-Code}. In \pascal\ sind hierf"ur bestimmte Funktionen
bzw. Prozeduren notwendig.

Der Einfachheit halber und wegen der Verwandschaft mit C wird die
Hauptfunktion des Programms \fu{main} genannt. Sie nimmt keine
Parameter an und hat entweder den R"uckgabetyp \tvoid\ oder
\tint . Ist ersteres der Fall, wird immer der Exit-Code
\fu{0} zur"uckgegeben, andernfalls wird dieser durch die
\rw{return}-Anweisung bestimmt.

\section{Anweisungen}
\label{stmts}

Anweisungen kontrollieren das Geschehen innerhalb eines
Programms. Sie befehlen\footnote{daher die Bezeichnung
"`imperative Programmiersprache"'} dem Computer, was er tun soll.
In SEL wird grunds"atzlich zwischen drei Arten von Anweisungen
unterschieden: Kontrollanweisungen, Schleifen und Ausdr"ucken.
W"ahrend die ersten beiden in diesem Unterkapitel besprochen
werden, widmet sich das gesamte n"achste Unterkapitel den
Ausdr"ucken. Der Korrektheit wegen sei hier gesagt, da"s es
eigentlich noch eine vierte Art von Anweisungen gibt: Die {\em
Nullanweisung}, die jedoch, da sie keinen bestimmten Zweck
erf"ullt, au"ser vielleicht die Lesbarkeit des Programms zu
erh"ohen, f"ur unsere Betrachtungen nicht weiter wichtig ist.

Es ist auffallend, da"s SEL keine Anweisungen zur Verf"ugung
stellt, die in irgendeiner Weise mit der Hardware des Computers
kommunizieren k"onnen, wie z.~B. Anweisungen zum Einlesen von der
Tastatur, oder zur Ausgabe von Zeichen auf den Bildschirm oder in
eine Datei. Die Gr"unde hierf"ur sind vielf"altig. Erstens haben
wir es uns zum Ziel gesetzt, den Sprachumfang von SEL m"oglichst
gering zu halten und Ein/Ausgabeanweisungen alleine
w"urden---sollten sie einigerma"sen vern"unftig sein---die
Grammatik von SEL betr"achtlich aufbl"ahen. Weiters habe ich SEL
mit dem Ziel vor Augen entworfen, sie einmal als Makrosprache
einzusetzen. Makrosprachen allerdings m"ussen in vielen
verschiedenen Umgebungen sinnvoll einzusetzen sein, wie z.~B. in
einer Textverarbeitung oder in einem Zeichenprogramm, das unter
einer graphischen Benutzeroberfl"ache l"auft. Nun sind aber die
M"oglichkeiten zur Textausgabe in einer graphischen
Benutzeroberfl"ache grundverschieden zu denen von z.~B. DOS.
Ausgabeanweisungen zu entwerfen, die f"ur alle diese
M"oglichkeiten halbwegs geeignet sind, erscheint mir fast
unm"oglich. Deshalb ist es unumg"anglich, jede Implementation von
SEL mit verschiedenen Grundfunktionen auszustatten, die auf die
jeweilige Umgebung zugeschnitten sind. Eine Standardisierung
dieser Grundfunktionen ist immer eine gute Idee, doch sicher sind
diese Standardfunktionen nicht "uberall anwendbar.\footnote{Ein
Beispiel hierf"ur w"are der ANSI-Standard von C, dessen
Ein/Ausgabefunktionen auf eine textorientierte Oberfl"ache
zugeschnitten sind. Die \fu{printf}-Funktion z.~B. ist unter der
graphischen Benutzeroberfl"ache Windows nutzlos.}

\bigskip

An dieser Stelle scheint es mir wichtig, ein Konzept aus anderen
strukturierten Programmiersprachen zu erw"ahnen. Das
Blockkonzept\footnote{nicht zu verwechseln mit den Bl"ocken in
SEL} erm"oglicht es, eine Anweisung entweder als eben nur eine
Anweisung oder als eine von zwei Schl"usselw"ortern oder Symbolen
begrenzte, beliebige Abfolge von Anweisungen zu sehen. So verlangt
eine \rw{if}-Anweisung in C oder \pascal\ nur eine Anweisung als
Zweig, es ist jedoch kein Problem, mit Hilfe der Blockstruktur
auch mehrere Anweisungen ausf"uhren zu lassen. Dieses
Blockkonzept ist in manchen Situationen sehr praktisch, aber es
bl"aht die Grammatik einer Programmiersprache auch etwas auf und
macht es so nat"urlich schwerer, einen Interpreter oder Compiler
zu schreiben. Aus diesem Grund habe ich mich entschieden, es in
SEL nicht einzubauen. Die Quelltexte sind so zwar etwas l"anger,
doch es entsteht keine wirkliche Beeintr"achtigung.

\subsection{Kontrollanweisungen}

Kontrollanweisungen haben die Aufgabe, zu kontrollieren, ob und
was ausgef"uhrt wird. In SEL gibt es drei Anweisungen dieser Art.
Die \rw{if}-Anweisung "uberpr"uft eine Aussage und f"uhrt, wenn
sie einen positiven Wahrheitswert ergibt, bestimmte Anweisungen
aus. Auf Wunsch werden auch Anweisungen ausgef"uhrt, wenn das
Ergebnis der "Uberpr"ufung negativ ist. Die \rw{switch}-Anweisung
f"uhrt f"ur bestimmte Werte eines Ausdrucks verschiedene
Anweisungen aus. Die \rw{return}-Anweisung springt zur
aufrufenden Funktion zur"uck.

\subsubsection{Die if-Anweisung}

Die \rw{if}-Anweisung\footnote{siehe Anhang \ref{grammar} oder
\cite{selkr}, Seite 10} wird durch das Schl"usselwort \rw{if}
eingeleitet, auf welches ein in Klammern eingeschlossener Ausdruck
folgt. Die Alternative zu dieser L"osung w"are gewesen, den
Ausdruck sofort auf das \rw{if} folgen zu lassen und durch ein
bestimmtes Schl"usselwort (z.~B. \rw{then}) oder ein Symbol zu
beenden. Der Doppelpunkt jedoch w"urde Mehrdeutigkeiten
hervorrufen und ein Semikolon erscheint mir als unpassend, und da
wir uns entschienden haben, Symbole zu verwenden, wo nicht
unbedingt Schl"usselw"orter erforderlich sind, f"allt auch das
\rw{then} weg. Die einfachste und unauff"alligste L"osung in
diesem Fall ist die Klammerung.

Auf die schlie"sende Klammer folgt eine beliebige Anzahl von
Anweisungen, beendet entweder durch \rw{end} oder \rw{else}. Da
auch \rw{if} eine Anweisung ist, ist es m"oglich, \rw{if}s zu
verschachteln. Im zweiten Kapitel wird sich zeigen, da"s es f"ur
einen Interpreter bzw. Compiler kein Problem ist, mit einer
solchen Schachtelung fertigzuwerden. Folgt auf die Anweisungen ein
\rw{end}, so ist die \rw{if}-Anweisung beendet. Es w"are
m"oglich, auch das \rw{end} durch ein Symbol zu ersetzen,
doch alle sinnvollen Symbole, bis auf die Klammern, sind bereits
f"ur andere Zwecke vergeben. Nat"urlich k"onnte man auch eine
schlie"sende Klammer benutzen, doch um ein solches Konstrukt
sinnvoll zu machen, ben"otigt man auch eine "offnende Klammer.
Eine runde Klammer w"urde den falschen Eindruck erwecken, eckige
Klammern sind bereits vergeben, es bleiben die geschwungenen
Klammern, die auch in C zu diesem Zweck verwendet werden. Da ich
aber auch in allen anderen F"allen einem \rw{end} bzw. einem
\rw{begin}-\rw{end} Paar den Vorrang gegeben habe, f"uhrte eine
Verwendung hier h"ochstens zu einer Verwirrung des Benutzers.

Folgt auf die Anweisungen ein \rw{else}, dann folgen auf
dieses weitere Anweisungen, durch ein \rw{end} begrenzt.
Nat"urlich k"onnte man auch anstelle des \rw{else} ein Symbol
verwenden, doch das w"urde den Zweck der Konstruktion zu sehr
verschleiern:

\begin{verbatim}
if x == 1
{
  do_what_you_want_to();
::
  do_something_else();
}
\end{verbatim}

Ich glaube, bei Betrachtung dieses Beispiels zeigt sich, warum
ich die Schl"usselw"orter den Symbolen vorgezogen habe.

Manche Programmiersprachen, wie z.~B. \modula , stellen
zus"atzlich zum normalen \rw{else} ein sogenannte \rw{elseif}
oder auch \rw{elsif} bereit. Es erlaubt, sollte die
\rw{if}-Bedingung negativ sein, noch eine zweite Bedingung zu
"uberpr"ufen und, sollte auch diese logisch falsch sein, eine
dritte u. s. w. Ich habe mich entschieden, ein solches
\rw{elseif} nicht in SEL einzubauen, aus zwei Gr"unden.
Erstens l"a"st sich diese \rw{if}-\rw{elseif} Kombination in
vielen F"allen durch ein besser passendes
\rw{switch}\footnote{siehe Kapitel \ref{switch}} ersetzen und
zweitens ist es m"oglich, den \rw{elseif}-Effekt durch ein
verschachteltes \rw{if} zu erzeugen---\rw{elseif} selbst ist
nichts anderes als eine Vereinfachung dieser Konstruktion.

\subsubsection{Die switch-Anweisung}
\label{switch}

Die \rw{switch}-Anweisung\footnote{siehe Anhang \ref{grammar} oder
\cite{selkr}, Seite 11} erlaubt es, abh"angig vom Wert eines
Ausdrucks verschiedene Anweisungen auszuf"uhren. Dazu wird
zun"achst der Ausdruck evaluiert und ein \rw{case} gesucht,
dessen Wert mit dem des Ausdrucks "aquivalent ist. Ist kein
passendes \rw{case} vorhanden, wird zu einem evtl. vorhandenen
\rw{default} gesprungen. Dort wird mit der Ausf"uhrung
fortgefahren und erst bei einem \rw{break} oder dem zum
\rw{switch} geh"orenden \rw{end} beendet, also nach dem \rw{end}
fortgesetzt. Ist erst einmal ein bestimmter \rw{case}- oder
\rw{default}-Zweig eingeschlagen, werden alle weiteren
\rw{case}'s und \rw{default}s ignoriert. Das erm"oglicht es, f"ur
mehrere Werte ein und dieselben Anweisungen auff"uhren zu lassen,
indem man die einzelnen \rw{case}s hintereinander anreiht.

Die \rw{case}- und \rw{default}-Schl"usselw"orter werden in SEL,
ebenso wie in C, als Pr"afixe f"ur Anweisungen behandelt: bei
normaler Ausf"uhrung werden sie einfach ignoriert. Sie sind also
im Grunde nichts anderes als eine Art Sprunglabel, so wie sie in
vielen Programmiersprachen f"ur die
\rw{goto}-Anweisung\footnote{siehe Kapitel \ref{goto}}
verwendet werden. Eine andere---auf den ersten Blick
nat"urlichere---L"osung w"are, auf die \rw{switch}-Anweisung eine
Liste von \rw{case}'s oder \rw{default}s folgen zu lassen, die
ihrerseits jeweils von Anweisungen gefolgt sind. Um es m"oglich
zu machen, da"s eine Anweisungsfolge f"ur mehrere verschiedene
Werte ausgef"uhrt wird, m"u"ste man statt einem \rw{case} jeweils
mehrere \rw{case}'s erwarten. Welche der beiden L"osungen man
bevorzugt, ist Geschmackssache. Ich habe mich f"ur die
Pr"afix-L"osung entschieden, weil ich sie erstens von C gewohnt
bin und sie zweitens ein wenig leistungsf"ahiger ist.

Das Schl"usselwort \rw{case} vor dem jeweiligen Wert ist
notwendig, um Zweideutigkeiten in der Sprache zu vermeiden. Um
dieses Schl"usselwort auslassen zu k"onnen, m"u"sten wir die oben
besprochene zweite L"osung w"ahlen, jedoch m"u"sten die einzelnen
Anweisungsfolgen wiederum durch ein anderes Schl"usselwort, z.~B.
\rw{end} beendet werden, soda"s sich keine wirkliche Ersparnis
erg"abe.

Auch daf"ur, da"s in \rw{case}-Zweigen nur Konstanten erlaubt
sind, gibt es Gr"unde. Erstens h"atten wir durch den Doppelpunkt
Zweideutigkeiten---er w"urde als Typecast\footnote{siehe Kapitel
\ref{casts}} interpretiert werden---doch dieses Problem ist durch
Wahl eines anderen Symbols ohne weiteres zu "ubergehen. Der weit
gewichtigere Grund ist, da"s die f"ur die \rw{switch}-Anweisung
typische Effizienz verloren ginge: Bis der Interpreter zum
richtigen Zweig gelangt, m"u"ste er eine Vielzahl von Ausdr"ucken
auswerten, was sehr viel mehr Zeit in Anspruch nimmt, als der
simple Vergleich mit einer Konstante. Au"serdem w"urde diese
Entscheidung das Schreiben eines Compilers etwas verkomplizieren.
In F"allen, in denen eine solche \rw{switch}-Anweisung n"otig
w"are, sollte man auf \rw{if}-\rw{else}-Verschachtelungen umsteigen.

\subsubsection{Die return-Anweisung}

Die \rw{return}-Anweisung hat die Aufgabe, zur aufrufenden
Funktion zur"uckzukehren und evtl. einen Wert zu "ubermitteln,
den sog. {\em Return-Code}. Eine andere M"oglichkeit w"are, den
return-code durch eine Zuweisung an eine Variable, die den Namen
der Funktion hat, zu bestimmen und mittels einer anderen
Anweisung bei Bedarf zur aufrufenden Funktion zur"uckzuspringen.
Da man im Normalfall beides an der gleichen Stelle tut, habe ich
mich f"ur die \rw{return}-L"osung entschieden.

Die einfachste Methode w"are jedoch, "uberhaupt keine Anweisung
bereitzustellen, mit der in die aufrufende Funktion
zur"uckgekehrt werden kann und den R"uckgabewert am Ende der
Funktion festzulegen. F"ur einen rekursiven Interpreter ist dies
die beste L"osung, doch wie wir im n"achsten Kapitel sehen
werden, stellt uns auch der \rw{return}-Ansatz vor keine gro"sen
Probleme.

\subsection{Schleifen}

Schleifen kontrollieren nicht, was, sondern wie oft etwas
ausgef"uhrt wird. Die drei Arten der Schleifen, die in SEL
erlaubt sind, entsprechen den in den meisten strukturierten
Programmiersprachen vorhandenen Schleifen.

\subsubsection{Die repeat-Schleife}

Die \rw{repeat}-Schleife in SEL entspricht der
\rw{repeat}-Schleife in \pascal , bis auf das Detail, da"s die
\rw{until}-Bedingung in Klammern zu setzen ist und kein Semikolon
danach erforderlich ist. Diese Konzeption deckt sich mit den
restlichen Anweisungen von SEL und wurde daher von mir bevorzugt.

In manchen Programmiersprachen, wie z.~B. C, existiert die
\rw{repeat}-Schleife in einer abgewandelten Form, n"amlich der
der \rw{do-while}-Schleife. Der einzige Unterschied
besteht darin, da"s die \rw{repeat}-Schleife eine
Abbruchbedingung erwartet, die \rw{do-while}-Schleife jedoch eine
Fortf"uhrungsbedingung, das hei"st, eine Schleife kann durch die
andere ersetzt werden, wenn die Bedingung negiert wird. Welche
Schleife die bessere ist, kommt auf die Situation an, wirkliche
Vor- oder Nachteile gibt es nicht, die Auswahl ist also reine
Geschmackssache.

\subsubsection{Die while-Schleife}

Die \rw{while}-Schleife ist eng verwandt mit der
\rw{repeat}-Schleife. Der einzige Unterschied besteht darin, da"s
die \rw{repeat}-Schleife mindestens einmal durchlaufen wird, die
\rw{while}-Schleife jedoch nicht unbedingt. Syntaktisch gesehen
ist die \rw{while}-Schleife sehr "ahnlich der von C, die
Unterschiede sind in der in Kapitel \ref{stmts} angesprochenen
Blockstruktur von C begr"undet. So folgt auf die schlie"sende
Klammer der Schleifenbedingung in C in der Regel eine "offnende
geschwungene Klammer, in SEL jedoch ist diese ausgespart, da
nicht notwendig. Die schlie"sende geschwungene Klammer, die in C
die Schleife beendet, ist mit dem Schl"usselwort \rw{end} in SEL
zu vergleichen. Die Gr"unde, warum ich die Konstruktion mittels
Schl"usselw"ortern vorgezogen habe, ist bereits in vorhergehenden
Kapiteln er"ortert, ich gehe deshalb hier nicht n"aher darauf
ein.

\subsubsection{Die for-Schleife}

Die \rw{for}-Schleife in SEL erlaubt es, Anweisungen in einer
Anzahl, die schon im voraus bekannt ist, zu wiederholen.
Damit unterscheidet sie sich signifikant von den \rw{while} und
der \rw{repeat}-Schleifen, bei denen die Anzahl der Wiederholungen
nicht unbedingt bekannt sind. Das verhilft der
\rw{for}-Schleife---zumindest auf einem Interpreter---zu einer
sehr schnellen Ausf"uhrung, da die Bedingungen f"ur die Schleife
nur einmal berechnet werden m"ussen, w"ahrenddessen dies bei den
beiden anderen Schleifenarten bei jedem Durchlauf geschehen mu"s.

Dieses eine Merkmal der \rw{for}-Schleife ist vielen
Programmiersprachen gemeinsam, doch eine bekannte Sprache macht
eine Ausnahme: C. In C werden nicht eine Z"ahlervariable und
deren untere und obere Grenzen und evtl. noch die Schrittweite
angegeben, sondern es werden drei Ausdr"ucke verwendet: Der erste
wird lediglich beim Start der Schleife evaluiert, der dritte nach
jedem Durchlauf. Der zweite Ausdruck schlie"slich gibt die
Bedingung an, bei der die Schleife nochmals ausgef"uhrt werden
soll---ist sie logisch falsch, wird die Schleife abgebrochen. In
C ist also eine \rw{for}-Schleife nichts anderes als eine
erweiterte \rw{while}-Schleife und hat den Vorteil, da"s mit
ihrer Hilfe auch verkettete Listen und "ahnliches leicht
durchlaufen werden k"onnen. Auch ein "`klassische"'
\rw{for}-Schleife l"a"st sich so leicht nachbilden, doch diese
Art der Schleifenkonstruktion hat f"ur einen Interpreter den
gro"sen Nachteil, da"s bei jedem Schleifendurchlauf gleich zwei
Ausdr"ucke evaluiert werden m"ussen.

Ungl"ucklicherweise habe ich diesen Nachteil nicht sofort bemerkt und
f"uhrte auch in SEL die von C gewohnte \rw{for}-Konstruktion ein. Erst
sp"ater merkte ich, um wie viel effizienter diese Schleife f"ur einen
Interpreter w"are, w"urde ich sie "ahnlich wie in \pascal\ aufbauen. Die
"Anderungen am Interpreter und an der Dokumentation lie"sen sich zum
Gl"uck in weniger als einer Stunde durchf"uhren, soda"s Verz"ogerungen
f"ur den Zeitplan dieser Arbeit ausblieben und ich zur Erkenntnis
gelangte, da"s man aus Fehlern doch am besten lernt.

Noch ein kleines Wort zur Syntax: Die Bedingungen der Schleife sind in
Klammern gesetzt, weil das in SEL auch bei anderen Anweisungen der Fall
ist und dadurch die Konsistenz gewahrt wird, und die Grammatik ohne
Begrenzung nicht eindeutig w"are. Die Z"ahlervariable wird als
\fu{primary} angegeben, weil es auch m"oglich sein sollte, Array- und
Strukturelemente als Z"ahler zu benutzen, und nicht als \fu{expr}, weil
dann der folgende Zuweisungsoperator als ein Teil dieses Ausdrucks
interpretiert w"urde. Sollte es wirklich notwendig sein, einen \fu{expr}
zu benutzen, so kann dieser in Klammer gesetzt werden, und entspricht
damit einem \fu{primary}.\footnote{siehe Anhang \ref{grammar}} Anstatt
der Schl"usselw"orter \rw{to} bzw. \rw{downto} Symbole zu verwenden,
w"are erstens un"ubersichtlich und w"urde zweitens---au"ser bei
Verwendung neuer Symbole---die Eindeutigkeit zunichte machen, weshalb
ich die oben angesprochenen Schl"usselw"orter verwendete. Gleiches gilt
f"ur \rw{step}.

\subsection{Die goto-Anweisung}
\label{goto}

Die "Uberschrift dieses Absatzes mag darauf deuten, da"s es in SEL eine
\rw{goto}-Anweisung gibt, doch der Schein tr"ugt: Aus Gr"unden, die in
Kapitel \ref{parser} aufgef"uhrt sind und aus dem Grund, da"s diese
Anweisung in dem meisten F"allen ohnedies nicht notwendig ist und nur zu
einem schlechten Programmierstil verf"uhrt, habe ich mich gegen eine
Einbindung in SEL entschieden.    

\section{Ausdr\protect{"ucke}}

Ausdr"ucke sind ein Konzept, das in fast allen strukturierten
Programmiersprachen verwendet wird und das nicht ohne Grund: Es
vereinfacht die Grammatik der Sprache und macht damit das Schreiben
eines Interpreters bzw. eines Compilers einfacher. Au"serdem erleichtert
es dem Anwender der Sprache das Schreiben von Programmen sehr und
unterst"utzt die Lesbarkeit. So k"onnen Ausdr"ucke in vielen
verschiedenen Situationen verwendet werden: Als Bedingungen in
Schleifen, als Funktionsaufrufe oder als die rechten Seiten von
Zuweisungen.\footnote{Tats"achlich ist in SEL sogar die Zuweisung ein
Ausdruck, im Gegensatz zu \pascal\ oder \modula , doch dazu sp"ater
mehr.}

Ein Ausdruck in SEL kann "ahnlich einem Ausdruck im
mathematischen Sinn aufgefa"st werden, er ist eine Verkn"upfung
von Operanden durch Operatoren. Jeder Operand f"ur sich ist
selbst ein Ausdruck\footnote{bis auf die Operanden der Operatoren
\fu{new} und \fu{sizeof}, die Typen sind.} und das Ergebnis der
Verkn"upfung mit einem Operator ist wieder ein Ausdruck. So ist
es in SEL m"oglich, Ausdr"ucke beliebig tief zu schachteln. Diese
Eigenschaften der Ausdr"ucke legen es nahe, einen Ausdruck als
Baum\footnote{siehe \cite{sedge}, Seiten 57--72} darzustellen.
Dabei stellt jeder innere Knoten sowie die Wurzel einen Operator
dar und jedes Blatt einen Wert. Betrachtet man nun jeden
einzelnen Knoten (und dessen Nachfolger), so gelangt man zum
Ergebnis, da"s dieser auch wieder ein Ausdruck f"ur sich ist. Ein
Beispiel f"ur diese Darstellung gibt Sedgewick in \cite{sedge},
Seite 63, Abbildung 4.4.

Jeder Ausdruck in SEL hat nicht nur einen Wert, sondern auch
einen Typ\footnote{siehe auch Kapitel \ref{types}}, mit einer Ausnahme:
Ein Ausdruck des Typs \tvoid\ hat keinen Wert, was auch gar keinen Sinn
h"atte.

\subsection{Operatoren}

Grunds"atzlich gibt es vier Arten, in SEL Operatoren zu
klassifizieren\footnote{genauer gesagt vier "`vern"unftige"' Arten}:
nach der Anzahl ihrer Argumente, nach ihrer funktionellen
Zugeh"origkeit, nach ihrer Priorit"at und nach ihrem Ergebnistyp.

\paragraph{Klassifikation nach der Anzahl der Argumente}

Auf den ersten Blick gesehen gibt es in SEL nach dem Gesichtspunkt der
Zahl der Argumente zwei Arten von Operatoren: un"are und bin"are. Zu den
un"aren z"ahlen z.~B. die Operatoren \fu{+} und \fu{<<}, zu den bin"aren
z"ahlen \fu{!} und \fu{-}. Der letztgenannte Operator existiert auch in
einer bin"aren Version, man k"onnte also sagen, da"s es Operatoren gibt,
die sowohl un"ar als auch bin"ar sind. Ich pers"onlich betrachte diese
beiden "`Versionen"' als zwei verschiedene Operatoren, da auch ihre
Funktionen verschieden sind. Betrachtet man Ausdr"ucke in SEL genauer,
wird man bemerken, da"s ein
Funktionsaufruf durchaus mehrere Argumente erwarten kann, jedoch stellt
sich hier die Frage, ob man einen Funktionsaufruf (besser: eine Funktion) als
Operator auffassen sollte. Ich bin der Meinung, eine Funktion ist
weniger ein Operator als ein Wert, und habe ihr deshalb ein eigenes
Unterkapitel gewidmet.\footnote{siehe Kapitel \ref{exprfuncs}}

\paragraph{Klassifikation nach der funktionellen Zugeh"origkeit}

Auf diese Weise kann man die Operatoren von SEL in mehrere
Gruppen einteilen, wobei die Grenzen eher flie"send sind. Mir
erscheint folgende Einteilung am vern"unftigsten: arithmetische
Operatoren wie \fu{+} oder \fu{<<}, logische Operatoren wie
\fu{\&\&} oder \fu{!}, Vergleichsoperatoren wie \fu{>=} oder
\fu{==}, Operatoren zur Speicherverwaltung wie \fu{new} oder
\fu{[]} und als zwei eigene Gruppen der Zuweisungsoperator \fu{=}
und der Typecasting-Operator \fu{:}. Wie gesagt sind die
"Uberg"ange zwischen diesen Gruppen sehr flie"send. So k"onnte
man beispielsweise meinen, der Operator \fu{[]} geh"ore zur
Gruppe der arithmetischen Operatoren, da er ja in gewisser Weise
Vektorarithmetik m"oglich macht. Auch w"are es m"oglich, den
Typecasting-Operator zur Speicherverwaltung z"ahlen, da er auch
den Speicher modifiziert.

\paragraph{Klassifikation nach der Priorit"at}

Operatoren in SEL besitzen verschiedene
Priorit"atsstufen\footnote{siehe \cite{selkr}, Seite 21}, um eine
"uberm"a"sige Klammerung zu vermeiden. Die Priorit"aten in SEL sind
stark an die von C angelehnt\footnote{siehe \cite{kandr}, Seite 54},
weichen jedoch teilweise von diesen ab. Meistens im selben Atemzug mit
der Priorit"at der Operatoren wird die Richtung der Auswertung genannt.
Diese gibt an, ob der Operator links- oder rechts-assoziativ ist. Ein
Beispiel f"ur einen rechts-assoziativen Operator ist \fu{=}: Der
Ausdruck \fu{a=b=0} ist "aquivalent zu \fu{a=(b=0)}. So ist es ohne
Klammerung m"oglich, mehreren Variablen in einer Anweisung denselben
Wert zuzuweisen. Ein
links-assoziativer Operator ist das bin"are \fu{-}. W"are er das nicht,
w"urde der Ausdruck \fu{5-3-1} nicht zu \fu{(5-3)-1}, sondern zu
\fu{5-(3-1)} "aquivalent sein und so das falsche Ergebnis, n"amlich
\fu{3} anstatt \fu{1}, liefern.

\paragraph{Klassifikation nach dem Ergebnistyp}

Diese Einteilung ist wohl die komplizierteste und es w"are wohl besser,
man w"urde als Kriterium nicht den Ergebnistyp selbst, sondern das, 
was den Ergebnistyp bestimmt, heranziehen. 
So sind dies etwa beim Operator \fu{+} die
Typen seiner Operanden, beim Operator \rw{new} hingegen ist es der Typ,
mit dem er verbunden wird. Weiters gibt es Operatoren, die immer denselben
Ergebnistyp haben, wie z.~B. die logischen Operatoren oder der Operator
\rw{sizeof}. Es gibt sogar einen Operator, dessen Ergebnistyp immer
\tvoid\ ist, n"amlich \rw{delete}.

\bigskip

In der folgenden Aufstellung der Operatoren verwende ich die
Klassifikation nach der funktionellen Zugeh"origkeit, weil sie
f"ur diesen Zweck am geeignetsten ist.

\subsubsection{Arithmetische Operatoren}

SEL bietet neben den numerischen Operatoren, zu denen die vier
Grundrechnungsoperatoren f"ur Addition, Subtraktion, Multiplikation und
Division z"ahlen, zus"atzlich der Modulo-Operator, der den Rest bei einer
ganzzahligen Division ergibt, und der Negationsoperator, der
einer Umkehr des Vorzeichens gleichkommt, auch
Operatoren zur Bitmanipulation, n"amlich bitweise Links- und
Rechtsverschiebung und bitweise Und-, Oder- und
XOR-Verkn"upfungen. Die bin"aren numerischen Operatoren, also
\fu{+}, \fu{-}, \fu{*}, \fu{/} und \fu{\%}, besitzen eine
niedrigere Priorit"at als alle un"aren Operatoren, jedoch eine
h"ohere Priorit"at als alle anderen bin"aren Operatoren.
Untereinander sind die Additions- und Subtraktionsoperatoren auf
derselben Priorit"atsstufe, die um einen Grad niedriger steht,
als die der anderen numerischen bin"aren Operatoren, in Anlehnung
an die in der Mathematik "ublichen Notation. Die un"aren
Operatoren haben deshalb eine h"ohere Priorit"at, weil das
erstens unseren Erwartungen entspricht---man betrachtet intuitiv einen
un"aren Operator einem Wert als "`n"aherstehender"', als einen
bin"aren---und es zweitens in fast allen F"allen praktischer ist.
So ist z.~B. das Ergebnis des Ausdrucks \fu{-x*y} dem von
\fu{-(x*y)} "aquivalent. Bes"a"sen aber die bin"aren Operatoren
eine h"ohere Priorit"at, m"u"ste man anstatt \fu{-x-y} den
Ausdruck \fu{(-x)-y} schreiben, was etwas unnat"urlich aussieht.
Hingegen k"onnte man den Ausdruck \fu{-(x+y)} zwar bei anderer
Operatorenrangfolge ohne Klammern schreiben, doch selbst ohne die
Rangfolge zu ver"andern, k"onnte man durch "Anderung des \fu{+} in
ein \fu{-} die Klammerung vermeiden: \fu{-x-y}. Die Rangfolge
erweist sich also als nat"urlich und praktisch. Die
Vergleichsoperatoren haben deshalb niedrigere Priorit"at als die
numerischen Operatoren, weil es kaum vorkommt, da"s man mit
Vergleichsergebnissen rechnet. F"ur die logischen Operatoren gilt
dieselbe Begr"undung. Die Festsetzung der Priorit"at der
Bitmanipulations-Operatoren relativ zu den numerischen Operatoren
ist eher willk"urlich. Es w"urde kaum ein Vor- oder Nachteil
entstehen, h"atten die Operatoren zur Bitmanipulation h"ohere
Priorit"at. Der Grund, warum ich diese Operatoren nicht auf
dieselbe Priorit"atsstufe gesetzt habe ist der, da"s sie zu
verschieden sind, als da"s dies dienlich w"are.

Die numerischen Operatoren haben eine besondere Eigenschaft: Sie
stellen sich auf die Typen ihrer Operanden ein und ver"andern
diese sogar, wenn n"otig. So stellt sich zum Beispiel die Frage,
welchen Typ eine Addition eines \tint\ und eines \tfloat\ haben
sollte. Vern"unftigerweise wandelt der Additionsoperator in
diesem Fall den \tint\ in einen \tfloat\ um, und f"uhrt erst
anschlie"send die Addition aus. Auf diese Weise gehen keine
Nachkommastellen des \tfloat\ verloren. Gleiches passiert mit
einem \tchar\ und einem \tint : Der \tchar\ wird in den Typ
\tint\ umgewandelt, da nur so gew"ahrleistet ist, da"s mit dem
richtigen Wert des \tint\ gearbeitet wird. F"ur die Typen bei
einer numerischen Operation gilt in SEL also die Regel \hbox{$
\tchar < \tint < \tfloat $}.\footnote{siehe \cite{selkr}, Seite
18} Zus"atzlich zu diesen Typen k"onnen die Operatoren \fu{+} und
\fu{-} auch mit Zeigern operieren, doch dazu mehr in Kapitel
\ref{memmgrops}.

Die Operatoren zur Bitmanipulation arbeiten in SEL lediglich mit
dem Typ \tint\, und das aus guten Gr"unden: Erstens ist es f"ur
einen Computer zwar leicht m"oglich, mit den Bits eines \tfloat\
zu operieren, doch es ist f"ur einen Menschen ungleich
schwieriger, diese Bits auch richtig zu interpretieren---auf den
Wert der Zahl bezogen. Zweitens sind die "ublichen 8 Bit eines
\tchar\ im Regelfall zu wenig f"ur die Anwendungen der
Bitmanipulation und drittens vereinfacht es nicht nur den Bau
eines Interpreters bzw. Compilers, es macht diese auch schneller.

\subsubsection{Logische Operatoren}

Neben den Operatoren zur logischen Und- und Oder-Verkn"upfung
bietet SEL auch einen Operator zur logischen XOR-Verkn"upfung an.
Ich habe mich dazu entschieden, weil es leicht in einen
Interpreter bzw. Compiler einzubauen ist und andererseits der
Operator in manchen Situationen sehr praktisch ist und logische
Ausdr"ucke verk"urzen kann. Die logischen Operatoren sind auf der
zweitniedrigsten Priorit"atsstufe angesiedelt---unter ihnen
befinden sich lediglich die Operatoren \fu{=} und \fu{delete}.
Diese Stufe habe ich gew"ahlt, weil logische Operatoren
haupts"achlich in Schleifen- und \rw{if}-Bedingungen eingesetzt
werden und es dort am praktischsten ist, wenn sie die niedrigste
Priorit"at haben. Nur selten ist es notwendig, da"s einer
Variablen ein logischer Wert zugewiesen wird, doch auch dieser
Fall ist ohne Klammerung zu bew"altigen, da ja der
Zuweisungsoperator eine noch niedrigere Priorit"at hat. Der
einzige Fall, in dem sich diese Rangfolge als unpraktisch
herausstellen w"urde ist, da"s in einer Schleifen- oder
\rw{if}-Bedingung ein logischer Wert einer Variablen zugewiesen
wird und dieser Wahrheitswert mit einem anderen mittels eines
logischen Operators verkn"upft werden mu"s. Ein Beispiel hierf"ur
w"are \fu{if (x \&\& (y = z))}. Tats"achlich tritt dieser Fall so
selten auf---in meiner bislang 4-j"ahrigen C-Laufbahn bin ich
noch nie in eine solche Situation gelangt---da"s man ihn getrost
au"ser Acht lassen kann.

Der logische Negationsoperator \fu{!} liegt auf derselben
Priorit"atsstufe wie die Operatoren zur bitweisen und
arithmetischen Negation und die beiden Operatoren \fu{*} und
\fu{\&}, die im Kapitel \ref{memmgrops} besprochen werden, welche
um eine Stufe tiefer liegt als die der Operatoren \fu{->}, \fu{:}
und \fu{[]}. Um das zu begr"unden, betrachten wir den Ausdruck
\fu{!x->y}. Offenbar wird hier versucht, aus der Struktur \fu{x}
auf das Element \fu{y} zuzugreifen und dieses zu negieren. H"atte
der Operator \fu{!} eine h"ohere Priorit"at als \fu{->}, so
w"urde sich die Negation auf die Struktur beziehen, und es ist
nicht m"oglich, eine Struktur zu negieren.

Logische Operatoren arbeiten in SEL nur mit dem Typ \tint .
Zwar w"urde es auch gen"ugen, den Typ \tchar\ zu diesem Zweck
heranzuziehen, doch da in kaum einer Anwendung Wahrheitswerte
gespeichert werden und \tint\ normalerweise der Gr"o"se des
Maschinenwortes des jeweiligen Prozessors entspricht, und daher
schneller zu bearbeiten ist, als ein \tchar , ist er die bessere
Wahl. Sollte es jedoch n"otig sein, Wahrheitswerte abzuspeichern,
kann das ohne weiteres mit Hilfe von \tchar s geschehen, die durch
ein Typecasting entstanden sind, oder aber durch Bitfelder, was
noch speicherplatzsparender ist.

\subsubsection{Vergleichsoperatoren}

SEL bietet insgesamt sechs Vergleichsoperatoren an. Diese sind um
eine Priorit"atsstufe h"oher als die logischen Operatoren
angesiedelt, da sie meistens in Schleifen- und
\rw{if}-Bedingungen eingesetzt werden und dort diese Rangfolge
von Vorteil ist. So wird es wahrscheinlich oft vorkommen, da"s
"uberpr"uft werden soll, ob zwei Werte zwei anderen Werten
entsprechen, aber eher selten wird man "uberpr"ufen wollen, ob
die Oder-Verkn"upfungen jeweils zweier Werte dieselben Ergebnisse
haben.

Vergleichsoperatoren k"onnen in SEL lediglich Werte der Typen
\tchar , \tint\ und \tfloat\ vergleichen, wobei das im vorletzten
Unterkapitel vorgestellte "`Gr"o"senverh"altnis"' gilt. Es ist
hingegen nicht m"oglich, z.~B. Strings miteinander zu
vergleichen, weil das, um es vern"unftig zu machen, eine
Ver"anderung im Typkonzept von SEL mit sich z"oge und au"serdem
sehr gut von einer Funktion erledigt werden kann---SEL stellt
alle n"otigen Konstrukte hierf"ur bereit.

\subsubsection{Speicherverwaltungsoperatoren}
\label{memmgrops}

Unter den Begriff "`Speicherverwaltungsoperatoren"' fallen all
jene Operatoren, die mit Zeigern, Arrays und Strukturen arbeiten.
Dazu z"ahlen nicht nur die Operatoren zur Ermittlung der
Speicheradresse einer Variablen und zur Dereferenzierung eines
Operators, \fu{\&} und \fu{*}, zum Zugriff auf ein Array- bzw.
Strukturelement, \fu{[]} und \fu{->}, sondern auch die Operatoren
zum Allokieren und Freigeben von Speicher, \fu{new} und
\fu{delete}, und der Operator zum Ermitteln der Gr"o"se eines
Typen, \fu{sizeof}. Au"serdem fallen unter diesen Begriff die
beiden bin"aren Operatoren \fu{+} und \fu{-}, mit denen
zus"atzlich zur numerischen Arithmetik auch Zeigerarithmetik
betrieben werden kann.

Die h"ochste Priorit"atsstufe aller Operatoren haben \fu{new} und
\fu{sizeof}, auf der n"achsten Stufe liegen \fu{->} und \fu{[]}.
Es h"atte keinen Sinn, diese Rangfolge zu "andern, denn die
Operanden von \fu{new} und \fu{sizeof} sind keine Ausdr"ucke,
sondern Typen, und es ist nicht m"oglich, den Operator \fu{->}
oder \fu{[]} auf einen Typ anzuwenden, soda"s dies schon von Haus
aus unterbunden werden sollte. Die beiden un"aren Operatoren
\fu{*} und \fu{\&} liegen um eine Priorit"atsstufe tiefer,
gemeinsam mit allen anderen un"aren Operatoren. Eine andere Stufe
f"ur diese Operatoren w"are nicht angebracht, aus Gr"unden, die
ich schon genannt habe.

Ein in gewisser Hinsicht au"sergew"ohnlicher Operator---gemessen
an seinen Gegenst"ucken aus C und \pascal ---ist der Operator
\fu{->}. Er wird benutzt, um auf Elemente einer Struktur
zuzugreifen. Derselbe Operator wird in C und \pascal\ durch den
Punkt (\fu{.}) symbolisiert. In C jedoch gibt es zus"atzlich den
Operator~\fu{->}. Er wird benutzt, um auf Elemente aus einem
Zeiger auf eine Struktur zuzugreifen---in \pascal\ mu"s hierf"ur
zuerst der Zeiger dereferenziert werden. Der SEL-Operator ist nun
insofern au"senstehend, als er die Funktionalit"at des
Punktoperators und des Pfeiloperators aus C vereinigt. So
erm"oglicht er es, angewendet auf eine Struktur, auf ein Element
daraus zuzugreifen. Angewendet auf einen Zeiger auf eine Struktur
dereferenziert er ersteren und greift dann erst auf das jeweilige
Element zu. Diese Vereinigung macht insofern Sinn, als da"s
keine Funktionalit"at verlorengeht, denn der Punktoperator
hat, angewendet auf einen Zeiger, keine Funktion. Auch das Design
eines Interpreters bzw. eines Compilers wird dadurch nicht
erschwert, denn diese wissen, wenn sie auf den Pfeiloperator
treffen, mit welchem Typ gearbeitet wird. Ein Compiler m"u"ste
also nur, tr"afe er auf einen Pfeiloperator, dessen linker Operand
ein Zeiger ist, den Code zur Dereferenzierung einbauen.

\subsubsection{Der Zuweisungsoperator}

Der Zuweisungsoperator von SEL ist ein Sonderfall unter den
Operatoren, aus mehreren Gr"unden. Der erste Grund ist, da"s es
in manchen Programmiersprachen, darunter auch \pascal , gar
keinen Zuweisungsoperator in diesem Sinne gibt. In \pascal\ ist
eine Zuweisung eine Anweisung und kein Ausdruck. Das macht es in
\pascal\ auf der einen Seite schwerer, den Zuweisungsoperator
vom Compiler bzw. Interpreter unentdeckt mit dem
Vergleichsoperator zu verwechseln, denn der Zuweisungsoperator
ist in Ausdr"ucken nicht erlaubt und Ausdr"ucke sind in \pascal\
keine g"ultigen Anweisungen. Andererseits jedoch ist es nicht
m"oglich, mehreren Variablen in einer Anweisung denselben Wert
zuzuweisen, wie man es von C oder auch aus der Mathematik gewohnt
ist (\fu{x=y=0}). Au"serdem ist es nicht m"oglich, z.~B. in einer
\rw{if}-Anweisung gleichzeitig eine Zuweisung vorzunehmen und den
Wert der zuzuweisenden Variablen zu "uberpr"ufen (z.~B.
\hbox{\fu{if ((x=y) == 3)}}). Manche legen das sogar als Vorteil
aus, weil es eine unsch"one Programmierung erschwert, und
insofern gebe ich ihnen Recht, doch ich pers"onlich finde solche
Konstrukte in manchen Situationen doch sehr praktisch, und da der
Aufwand f"ur einen Zuweisungsoperator etwa dem einer
Zuweisungsanweisung gleichkommt, oder sogar noch geringer ist,
habe ich ihn in SEL eingebaut.

Der zweite Grund f"ur die Besonderheit des \fu{=}-Operators ist,
da"s er mit allen Typen zusammenarbeitet. Anders als der
Vergleichsoperator, der auf die Grundtypen beschr"ankt ist, kann
der Zuweisunsoperator eine Variable jeden beliebigen Typs einer
anderen zuweisen---Typgleichheit vorausgesetzt. Ich habe diese
"`Allm"achtigkeit"' des Operators deswegen in SEL erm"oglicht, weil
es leicht ist, das in einen Interpreter einzubauen---im Fall der
Konzeption, die wir im zweiten Kapitel besprechen werden, sogar
noch leichter, als sich auf die Grundtypen zu beschr"anken. Auch
f"ur einen Compiler ist das nicht weiter schwer: Den Code f"ur die
Grundtypen kann der Compiler direkt einf"ugen, f"ur Strukturen
oder Arrays k"onnte eine Funktion aus der Laufzeitumgebung
aufgerufen werden. Hier dr"angt sich nat"urlich die Frage auf,
warum ich dieses Konzept nicht auf den Vergleichsoperator
\fu{==} ausgeweitet habe. Die Antwort ist einfach: H"atte ich das
getan, w"are es unausweichlich gewesen, das auch f"ur die
Operatoren \fu{<}, \fu{>}, \fu{<=}, \fu{>=} und \fu{!=} zu tun.
Mit letzterem h"atte ich keine Probleme gehabt, doch wann ist
eine Struktur, die die Adresse einer Person enth"alt, gr"o"ser
als eine andere desselben Typs? Byte-f"ur-Byte-Vergleich w"are
hier sicher nicht angebracht.\footnote{In C++ kann dieses Problem
gel"ost werden, weil es dort m"oglich ist, Operatoren zu
"uberladen, d. h., da"s der Benutzer imstande ist, eine Funktion
zu schreiben, die zwei Strukturen (beispielsweise des Typs
\fu{Adresse}) vergleicht, die dann anstatt des
Vergleichsoperators aufgerufen wird.}

\subsubsection{Der Typecasting-Operator}

Der Typecasting-Operator \fu{:} wird benutzt, um einen Wert
explizit in einen anderen Typ umzuwandeln. Er wird genauer in
Kapitel \ref{casts} beschrieben. Seine Priorit"at ist mit der der
beiden Operatoren \fu{->} und \fu{[]} "aquivalent, erstens, weil
er ein un"arer Operator ist und zweitens, weil es eher selten
vorkommt, da"s gr"o"sere Ausdr"ucke gecastet werden m"ussen. Wenn
n"otig, wird dies sowieso von SEL "ubernommen.

\subsection{Konstanten}
\label{const}

Grunds"atzlich gibt es in SEL vier verschiedene Arten von
Konstanten: \tchar -Konstanten, \tint -Konstanten,
\tfloat -Konstanten und String-Konstanten. Eine \tchar -Konstante
ist ein unter einfache Anf"uhrungsstriche gesetztes Zeichen bzw.
Steuerzeichen und hat---es ist nicht anders zu erwarten---den Typ
\tchar . Eine \tint -Konstante ist eine simple Abfolge von
Ziffern. Zahlenkonstanten im lexikalischen Sinn sind in SEL immer
positiv---negative Zahlen macht erst der un"are Operator \fu{-}
m"oglich, und insofern ergibt sich daraus kein Nachteil. Der
Grund, warum ich Konstanten nicht auch auf negative Zahlen
ausgeweitet habe ist der, da"s sich die lexikalische
Analyse\footnote{siehe Kapitel \ref{scanner}}
dadurch sehr erschweren w"urde. Nehmen wir als Beispiel den
Ausdruck \fu{7-4}. Der Scanner erkennt als erstes die
\tint -Konstante \fu{7}. Danach m"u"ste ein Operator folgen oder
der Ausdruck zu Ende sein, was aber der Scanner nicht wissen
kann, da er vom Parser in dieser Hinsicht abgekapselt ist. Aus
diesem Grund w"urde er anstatt den Operator \fu{-} zu erkennen
eine weitere \tint -Konstante, n"amlich \fu{-4} lesen, was
nat"urlich falsch ist. Das Gesagte gilt auch f"ur
\tfloat -Konstanten.

String-Konstanten sind in gewisser Hinsicht etwas Besonderes,
n"amlich in Anbetracht der Tatsache, da"s es einen Stringtyp in
SEL gar nicht gibt. Normalerweise werden Strings als Arrays von
\tchar s behandelt, die durch ein Byte mit dem Wert \fu{0}
abgeschlossen sind, es w"are allerdings nicht schwer, Funktionen
zu programmieren, die \pascal -Strings verarbeiten. Da jedoch die
erstere Stringart die verbreitetere ist, wollen wir bei dieser
bleiben. Wie wir wissen, kann ein Array in SEL in vielen
Situationen wie ein Zeiger behandelt werden. Was also w"are
falsch daran, einen String als Pointer-Konstante zu sehen?
Tats"achlich wird das nicht nur in SEL, sondern auch in C so
praktiziert, von wo ich dieses Konzept "`importiert"' habe. Dieser
Pointer zeigt auf einen Speicherbereich, in dem der eigentliche
String gespeichert ist. Nat"urlich mu"s dieser Speicherbereich
w"ahrend der gesamten Programmausf"uhrung an derselben
Speicheradresse bleiben, und nicht irgendeiner L"oschaktion zum
Opfer fallen, soda"s die Pointerkonstante im ganzen
Programm zu jeder Zeit "`aktuell"' ist.\footnote{In Wirklichkeit
mu"s der String nicht w"ahrend der gesamten Programmausf"uhrung
im Speicher sein, sondern erst nach seinem ersten Auftreten im
Code, doch der Einfachheit halber wird er die ganze Zeit im
Speicher gehalten.}

\subsection{Variablen}
\label{variables}

In SEL k"onnen zwei Arten von Variablen unterschieden werden, die
ihrerseits in jeweils zwei Untertypen unterteilt werden k"onnen.
Lokale Variablen sind nur innerhalb einer bestimmten Funktion
bekannt und existieren f"ur jede Instanz dieser Funktion
getrennt.\footnote{Mehrere Instanzen einer Funktion sind dann
m"oglich, wenn sich diese direkt oder indirekt selbst aufruft.}
Die Gruppe der lokalen Variablen werden weiter unterteilt in
Parameter und direkte Funktionsvariablen.\footnote{Eine bessere
Bezeichnung als "`direkte Funktionsvariablen"' habe ich nicht
gefunden.} Die Werte der Parameter werden von der aufrufenden
Funktion bestimmt, wohingegen die direkten Funktionsvariablen der
Funktion alleine unterstehen. Lokale Variablen werden auf dem
Stack gespeichert. Wird eine Funktion aufgerufen, werden die
Parameter von der aufrufenden Funktion auf den Stack gelegt und
anschlie"send zum Code der Funktion
gesprungen.\footnote{Nat"urlich wird auch die R"ucksprungadresse
auf dem Stack gespeichert} Die Funktion ihrerseits legt die
direkten Funktionsvariablen auf den Stack und f"uhrt dann ihre
Anweisungen aus. Wird sie beendet, entfernt sie sowohl
direkte Funktionsvariablen als auch Parameter vom Stack
und springt dann zur aufrufenden Funktion zur"uck. Das Konzept,
den Stack als Speicher f"ur lokale Variablen zu benutzen macht
es m"oglich, da"s bei einer rekursiven Funktion jede Instanz ihre
eigenen Variablen zur Verf"ugung hat.

Globale Variablen werden in SEL unterteilt in modulglobale und
programmglobale Variablen. Erstere sind nur in der Quelldatei
(Modul) bekannt, in der sie definiert sind, wohingegen letztere
auch in anderen Modulen benutzt werden k"onnen. Sowohl modul- als
auch programmglobale Variablen werden im statischen Speicher
aufbewahrt. Daf"ur, da"s modulglobale Variablen nur in ihrem
jeweiligen Modul referenziert werden, haben der Interpreter bzw.
der Compiler oder der Linker zu sorgen.

Ich glaube, "uber die Verwendung von Variablen in Ausdr"ucken
mu"s ich hier nicht viele Worte verlieren: das geschieht ganz
einfach---wie in fast allen Programmiersprachen "ublich---durch
die Benutzung des Variablennamens.

\subsection{Funktionen}
\label{exprfuncs}

Die Aufrufkonvention f"ur Funktionen ist einfach: Der
Funktionsname gefolgt von in Klammern gesetzten Parametern,
jeweils durch einen Beistrich getrennt, bilden einen Funktionsaufruf.
Hat die Funktion keine Parameter, sind trotzdem "offnende und
schlie"sende Klammern anzugeben, aus mehreren Gr"unden. Erstens
kann der Interpreter bzw. Compiler daran ohne nachzusehen
erkennen, da"s es sich um einen Funktionsaufruf handelt und nicht
um die Verwendung einer Variablen. Zweitens ist dies in Hinblick
auf die evtl. Erweiterung der Sprache durch
Funktionszeiger\footnote{siehe Kapitel \ref{langext}}, also
Zeigervariablen, die die Adressen von Funktionen aufnehmen,
n"otig. In diesem Fall signalisiert das Klammerpaar, da"s wir
vorhaben, diese Funktion aufzurufen und nicht, sie als Zeiger zu
verwenden.

Die Anzahl und die Typen der Funktionsparameter sind nat"urlich
einzuhalten. Die Regelung bei letzteren ist nicht so streng:
Stimmen die Typen nicht "uberein, versucht SEL mittels impliziten
Typecastings\footnote{siehe Kapitel \ref{casts}}, diese in die
richtigen umzuwandeln, was jedoch nicht unbedingt gelingen mu"s,
da sich der Interpreter bzw. der Compiler an bestimmte Regeln zur
Typumwandlung halten mu"s.

\section{Typen}
\label{types}

In diesem Abschnitt m"ochte ich das Typkonzept von SEL erkl"aren.
Zun"achst stellt sich beim Design einer Programmiersprache die
Frage, welches Typkonzept "uberhaupt zur Anwendung kommen soll.
So gibt es z.~B. in \basic\ einige wenige Typen, aus denen der
Benutzer ausw"ahlen kann, es ist dort jedoch nicht m"oglich,
selber Typen zu definieren. Zwar k"onnen Arrays definiert werden,
doch ist ein Array in \basic\ nichts anderes als eine Menge von
Variablen, deren Namen sich jeweils durch eine Nummer
unterscheiden---ein Array als Ganzes zu behandeln, ist in \basic\
nicht m"oglich.

Dann gibt es Programmiersprachen---\dbase\ ist
ein Beispiel hierf"ur---in denen Variablen zwar einen bestimmten
Typ haben, dieser sich aber jederzeit "andern kann. Dort ist es
z.~B. m"oglich, einer Variablen einen numerischen Wert zuzuweisen
und mit ihr Berechnungen anzustellen. Sp"ater kann man ihr bei
Bedarf einen String zuweisen und diesen am Bildschirm ausgeben.
Zwar bietet dieses Konzept auf den ersten Blick f"ur den Benutzer
Vorteile\footnote{Auch diese Vorteile sind nur scheinbar, denn
sie verf"uhren zu einem sehr schmutzigen Programmierstil} und es
ist f"ur eine reine Interpretersprache eine gute Wahl, doch einem
Compiler macht ein solches Konzept gro"se Probleme: In
Programmiersprachen wie C oder \pascal\ haben Variablen einen
vorgegebenen, konstanten Typ. Vom Compiler werden diese Variablen
an eine bestimmte Speicheradresse gelegt, die sich im Laufe des
Programms nicht ver"andert, was ohne weiteres m"oglich ist, da
sich ja die Speicheranforderungen der Variablen auch nicht
ver"andern. K"onnen sich aber die Typen der Variablen "andern,
"andern sich auch die Gr"o"sen der Variablen, d. h., da"s es
nicht m"oglich ist, sie an eine konstante Stelle im Speicher zu
setzen. W"ahrend ein C-Compiler, wenn er die Variable zu
referenzieren hat, ganz einfach ihre Speicheradresse einf"ugen
kann, mu"s ein compiliertes \dbase -Programm eine Tabelle
verwalten, in der alle Variablen und deren momentane Adressen und
Typen verzeichnet sind. Wird nun der Ausdruck \fu{x+1}
abgearbeitet, mu"s zun"achst in der Tabelle nachgesehen werden,
an welcher Adresse sich die Variable \fu{x} befindet. Zus"atzlich
mu"s auch gepr"uft werden, welchen Typ diese hat, da z.~B. eine
ganzzahlige Addition einen anderen Code verlangt, als eine
Addition mit Gleitkommazahlen. Erst jetzt kann die Addition
ausgef"uhrt werden. Es ist also nur verst"andlich, da"s Compiler
f"ur eine solche Sprache einen Code erzeugen, der zehnmal
langsamer\footnote{Die Zahl 10 mag wie eine "Ubertreibung
klingen, tats"achlich ist es nicht selten, da"s solche Codes noch
um einiges langsamer sind. So ben"otigt z.~B. ein Programm, da"s
einen QuickSort-Algorithmus ausf"uhrt in Clipper (einem
\dbase -Compiler) sage und schreibe 120 mal l"anger als ein
vergleichbares Programm in C} ist als der eines C- oder
\pascal -Compilers.

Eine weitere wichtige Kategorie ist die der typenlosen
Programmiersprachen---AWK f"allt teilweise darunter. Dort ist es
m"oglich, einer Variablen einen String zuzuweisen, der eine Zahl
enth"alt und in der n"achsten Zeile mit dieser Variable zu
rechnen als w"are sie numerischen Typs. Dieses Konzept ist
insofern unpraktisch, als man f"ur verschiedene Typen (obwohl
es diese eigentlich nicht gibt) auch verschiedene Operanden
braucht. Bei einem Ausdruck \fu{x+y} w"are also nicht klar, ob
die beiden Strings zu konkatenieren sind oder ob deren numerische
Werte addiert werden sollen. Auch in Bezug auf den Bau eines
Compilers ist dieses Konzept von Nachteil, aus "ahnlichen
wie im vorigen Absatz dargelegten Gr"unden.

Manche alte Programmiersprachen stellen als einzigen Typ das
Maschinenwort zur Verf"ugung, was einerseits die Gefahr birgt,
da"s in einer solchen Sprache geschriebene Programme sehr vom
Prozessor abh"angig sind und andererseits den Nachteil hat, da"s
viele Aufgaben, die der Interpreter bzw. Compiler ohne 
Effizienzverluste "ubernehmen k"onnte, vom Programmierer 
bestritten werden m"ussen.

Wie sich aus den vorigen Abs"atzen leicht herauslesen l"a"st,
habe ich mich f"ur das Typkonzept entschieden, das auch in C
oder \pascal\ verwendet wird. Dort mu"s jede Variable vor ihrer
Verwendung definiert werden, wobei gleichzeitig ihr Typ
festgelegt wird, der sich im Laufe des Programms nicht "andert.
Au"serdem ist es dem Benutzer m"oglich, aus den Basistypen der
Programmiersprache eigene Typen herzuleiten, wenn auch nicht
unbeschr"ankt.

\subsection{Elementare Typen}
\label{elemtypes}

Wie bereits gesagt, ben"otigt man zur Konstruktion von
komplizierteren Typen eine Menge an Basistypen, mit denen man
mittels der Konstrukte, die SEL bereitstellt, arbeiten kann.
Programmiersprachen wie C oder \pascal\ bieten eine relativ
gro"se Zahl solcher Basistypen an. So gibt es in C zwei
verschiedene Typen, die ein einzelnes Zeichen repr"asentieren,
drei Grundtypen f"ur ganze Zahlen, wobei von jedem dieser Typen
ein vorzeichenloser sowie ein vorzeichenbehafteter Typ
exisitiert, und mehrere Gleitkomma-Typen. Die Vielfalt in
\pascal\ ist "ahnlich gro"s, wobei zus"atzlich noch ein Typ f"ur
logische Werte exisitiert (\type{boolean}). Ich habe mich
entschieden, in SEL nur vier Elementartypen aufzunehmen:
\tchar , \tint , \tfloat\ und \tvoid .

Der Typ \tchar\ repr"asentiert ein einzelnes Zeichen und nimmt
dementsprechend nur ein Byte ein. Arrays aus \tchar s k"onnen
z.~B. als Strings interpretiert werden.

Ein \tint\ repr"asentiert eine ganze Zahl mit Vorzeichen und hat
in der Regel die Gr"o"se eines Maschinenwortes.\footnote{Eine
Ausnahme bildet MS-DOS, wo ein Maschinenwort nur 16 Bits lang
ist, d.~h. der Wertbereich w"urde von etwa -32000 bis 32000
reichen, was zwar f"ur viele Anwendungen reicht, aber ungen"ugend
ist, da es ja nur einen Typ f"ur ganze Zahlen gibt. Auf MS-DOS
Maschinen ist deshalb die Gr"o"se eines \tint\ zwei
Maschinenworte, also 32 Bits.} Zus"atzlich hat \tint\ die
Aufgabe, logische Werte zu repr"asentieren, wobei der Wert \fu{0}
einem "`logisch falsch"' und ein Wert ungleich \fu{0} "`logisch
wahr"' entspricht. Unter anderem hat diese Konzeption den
Vorteil, da"s ein Ausdruck des Typs \tint\ nicht nur den Wert des
\tint\ wiederspiegelt, sondern auch die Aussage, ob dieser
ungleich \fu{0} ist oder nicht. So ist es z.~B. m"oglich, mit
einer sehr kurzen \rw{if}-Anweisung zu pr"ufen, ob ein Zeiger ein
Nullzeiger ist. W"ahrend hierf"ur in \pascal\ die Anweisung
\fu{if p <> nil then} benutzt wird, kann in SEL einfach
\fu{if (p)} geschrieben werden. Zieht man es aus Gr"unden der
"Ubersichtlichkeit vor, trotzdem einen boolee'schen Typ zu haben,
kann man diesen in SEL problemlos mit den Direktiven

\begin{verbatim}
#define boolean        int
#define false          0
#define true           1
\end{verbatim}

\noindent simulieren.\footnote{Die Betonung liegt auf "`simulieren"',
da wir uns mit diesen Direktiven keinen eigenen Typ schaffen sondern
lediglich \fu{boolean} zu einem Pseudonym f"ur \rw{int} machen.}

Der Typ \tfloat\ repr"asentiert eine vorzeichenbehaftete
Gleitkommazahl. Ihre Gr"o"se im Speicher h"angt vom jeweiligen
Prozessor ab.

Einen Spezialfall unter den Typen stellt \tvoid\ dar: Er
repr"asentiert gar nichts! F"ur sein Dasein sprechen zwei
Gr"unde. Erstens macht er es m"oglich, Funktionen ohne
R"uckgabewert zu konstruieren, also Konstrukte, die mit den
Prozeduren in \pascal\ oder \modula\ vergleichbar sind. Au"serdem
ist er der Grundtyp f"ur typenlose Zeiger, also Zeiger, die auf
etwas nicht Bestimmtes zeigen. Verwendung finden solche Zeiger
z.~B. in Funktionen zur Speicherallokation. Diese geben einen
Zeiger auf einen Speicherbereich bestimmter Gr"o"se zur"uck,
wobei aber nicht bekannt ist, f"ur welchen Typ dieser Bereich
reserviert wird.

\subsection{Arrays}

Ein Array ist ein Zusammenschlu"s mehrerer Variablen desselben
Typs, die im Speicher sinnvollerweise aufeinander folgen. Ich
habe diesen Typ in SEL aufgenommen, weil er erstens nicht schwer
zu implementieren und zweitens in sehr vielen Situationen
unverzichtbar ist. Auffallend an der Schreibweise des Arraytyps
ist, da"s die Dimension des Arrays dem Grundtyp nicht nach- sondern
vorgestellt sind. Gem"a"s unserem Grundsatz, Symbole anstatt
Schl"usselw"ortern zu verwenden, wo es nicht st"ort, habe ich hier
auf Ausschm"uckungen wie das in \pascal\ verwendete \fu{array
... of ...} verzichtet und statt dessen eine in eckige Klammern
geschlossene Array-Dimension benutzt. Auf ein Angeben der
Grenzen, wie in \pascal\ "ublich, habe ich verzichtet, das erste
Element besitzt den Index \fu{0}. Diese Vereinfachung hat zwei
Gr"unde: Einerseits ist es nicht schwer, auf speziell indizierte
Arrays zu verzichten, andererseits w"urde sich das auf die
Konsistenz in Bezug auf die Verwendung von Zeigern anstatt Arrays
auswirken: Angenommen, wir h"atten ein Array \fu{a : [5..10]int}.
In SEL ist es nun m"oglich, dieses Array als Zeigervariable zu
betrachten und so zu behandeln, d. h., wir k"onnten mit ihm
Zeigerarithmetik betreiben. Per definitionem ist ein Zeiger
\fu{p} gleich \fu{p + 0}, was bedeutet, da"s wir mit \fu{*(a + 0)} auf das
erste Element des Arrays zugreifen k"onnen. Wir k"onnen auf
dieses Element auch direkt "uber das Array zugreifen, also mit
\fu{a[5]}. Ich denke, hier zeigt sich schon, worauf ich
hinauswill: Die Zahlen, die zur Indizierung verwendet werden,
stimmen nicht "uberein. Um das zu umgehen, m"u"ste man entweder
auf die Gleichsetzung von Array und Zeiger verzichten, was ich
nicht wollte, oder auch in die Zeiger eine komplizierte
Nummerierung einbauen, was mir und sicher auch dem Benutzer der
Programmiersprache zu umst"andlich ist.

Zur"uck zu der Dimension des Arrays, die, wie oben erw"ahnt, dem
Grundtyp vorangestellt ist. Der Grund hierf"ur ist in erster
Linie beim Typecasting\footnote{siehe Kapitel \ref{casts}}
von SEL zu suchen. Nehmen wir als Beispiel den Ausdruck
\fu{a : *int[10]}. Hier w"are es unklar, ob die eckige Klammer noch
zum Typ geh"ort, oder ob versucht wird, auf das zehnte Element
des Zeigers (der als Array behandelt wird) zuzugreifen.
Nat"urlich kann auch dieses Problem durch Klammerung umgangen
werden, doch warum etwas komplizierter machen, als es ist? Die
vorangestellte Arraydimension in SEL hat n"amlich noch einen
Vorteil: Es ist ohne eine Rangordnung oder Klammerung m"oglich,
einen Typ eindeutig zu bestimmen. Als Beispiel soll der C-Typ
\fu{*int[10]} dienen. Ohne den beiden Operatoren eine
unterschiedliche Priorit"at zuzuweisen, ist es unm"oglich zu
sagen, ob der Typ nun \fu{*(int[10])} oder \fu{(*int)[10]} ist.
In SEL ist die Sache eindeutig: \fu{[10]*int} kann nur als Array
aus zehn Zeigern auf \tint\ gedeutet werden.

\subsection{Zeiger}

Zeiger stellen ein wichtiges Konstrukt dar, mit dessen Hilfe es
m"oglich ist, sehr komplexe, verkettete und verzweigte
Datenstrukturen zu bilden, die ohne sie nur sehr schwer zu
verwirklichen w"aren. Der Bau eines Interpreters oder gar eines
Compilers z.~B. w"are ohne die Verwendung von Zeigern kaum zu
bewerkstelligen. Aus diesem Grund habe ich Zeiger in SEL
eingebaut, was f"ur eine Interpretersprache eigentlich sehr
ungew"ohnlich ist.

Sehr eng in Verbindung mit Zeigern steht in SEL der Operator
\fu{*}, der einerseits in der Typdefinition eines Zeigers
verwendet wird und andererseits benutzt wird, einen Zeiger zu
dereferenzieren. Letzteres geschieht, indem dem Zeiger der
Sternoperator vorangestellt wird, anders als in \pascal , wo der
Dereferenzierungsoperator als Suffix verwendet wird. Der Grund,
warum in SEL der Operator vorangestellt wird, ist die Tatsache,
da"s er dasselbe Symbol hat wie der Multiplikationsoperator, und
der Parser\footnote{siehe Kapitel \ref{parser}} nicht w"u"ste,
wie etwa der Ausdruck \fu{p * - x} gemeint ist. Einerseits k"onnte es
die Differenz zwischen der Variablen, auf die \fu{p} zeigt und
\fu{x} sein, andererseits aber auch \fu{p} multipliziert mit
\fu{-x}. Nat"urlich h"atte ich auch ein anderes Symbol f"ur die
Dereferenzierung verwenden k"onnen, doch das wollte ich nicht,
weil ich den Operator \fu{*} von C gewohnt bin, au"serdem fand
ich es schon immer ein wenig verwirrend, da"s bei der
Typdeklaration der Operator vorangestellt, bei der
Dereferenzierung aber nachgestellt sein mu"s.

Der zweite wichtige Operator in Verbindung mit Zeigern ist
\fu{\&}. Wird er einer Variablen vorangestellt, dann ergibt er
einen Zeiger auf diese Variable. Mit seiner Hilfe ist es also
m"oglich, einen Zeiger auf eine bestimmte Variable zeigen zu
lassen. Verst"andlicherweise ist es nicht m"oglich, \fu{\&} in
Verbindung mit einer Konstante zu benutzen, denn eine Konstante
hat keinen bestimmten Platz im Speicher, sie wird vom Interpreter
direkt in einen Wert umgewandelt und von einem Compiler direkt in
den Code eingef"ugt. Zu verwechseln ist dies jedoch nicht mit den
String-Konstanten\footnote{siehe Kapitel \ref{const}}, die
eigentlich nicht konstante Strings sind, sondern konstante Zeiger
auf einen String, der an einer bestimmten Stelle im Speicher
liegt. Doch zur"uck zum Operator \fu{\&}: Sein zweiter
Verwendungszweck ist die Kennzeichnung eines Parameters einer
Funktion, der {\em Call-By-Reference} "ubergeben werden soll. In
\pascal\ z.~B. "ubernimmt diese Aufgabe das Schl"usselwort
\rw{var}, in C gibt es etwas Derartiges nicht, erst in C++ wurden
Referenzparameter eingef"uhrt. Auch dort werden sie mit \fu{\&}
gekennzeichnet. In einer Programmiersprache, in der Zeiger zur
Verf"ugung stehen, sind Referenzparameter nicht zwingend
notwendig: Es k"onnten Zeiger auf die betreffenden Variablen
"ubergeben werden, wie es in C auch praktiziert wird. Ich habe
mich aus zwei Gr"unden entschlossen, Referenzparameter doch in
SEL aufzunehmen. Erstens k"onnen sie in manchen Situationen
Programme um sehr viel "ubersichtlicher und leichter wartbar
machen. Zweitens, und das war f"ur mich wohl der gewichtigere
Grund, arbeitet der im n"achsten Kapitel theoretisch beschriebene
Interpreter intern ohnedies mit Referenzen, soda"s es keine
Schwierigkeit war, Referenzparameter einzubauen.

\subsubsection{Dynamische Variablen}

Eng verbunden mit den Zeigern ist das Konzept der dynamischen
Variablen. Dynamisch deshalb, weil sie erzeugt werden k"onnen,
wenn man sie braucht und aus dem Speicher entfernt, wenn sie
"uberfl"ussig werden. Statische Variablen belegen
Speicherbereiche, die vom Compiler bzw. Interpreter im voraus
bestimmt werden, wohingegen der Speicherbereich f"ur eine
dynamische Variable erst w"ahrend der Programmausf"uhrung
bestimmt werden kann. Das impliziert nat"urlich, da"s dem
Benutzer geeignete Mittel in die Hand gegeben werden m"ussen, um
dynamische Variablen zu erzeugen. Grunds"atzlich gibt es hierf"ur
zwei M"oglichkeiten. Entweder, man baut das Konzept der
dynamischen Speicherverwaltung nicht fest in die
Programmiersprache ein, sondern "uberl"a"st es dem Benutzer bzw.
dem Designer der Laufzeitbibliothek, geeignete Funktionen zu
implementieren, wie es in C der Fall ist, oder man
ber"ucksichtigt dynamische Speicherverwaltung schon beim Design
der Programmiersprache und stellt Konstrukte zur Verf"ugung, um
dynamische Variablen zu erzeugen, wie z.~B. in C++.

Ich habe mich f"ur den zweiten L"osungsansatz entschieden, weil
ich finde, da"s das Konzept der dynamischen Speicherverwaltung zu
wichtig ist, um es beim Design einer Programmiersprache nicht
miteinzubeziehen. Ein daraus resultierender Vorteil ist, da"s
Programme leichter lesbar werden und es leicher ist, sauberer zu
programmieren, als z.~B. in C.\footnote{Was nicht hei"sen will,
da"s es in C nicht m"oglich ist, sauber zu programmieren, es ist
nur nicht so leicht wie etwa in C++} Angelehnt an das Vorbild C++
werden in SEL zur dynamischen Speicherverwaltung die beiden
un"aren Operatoren \rw{new} und \rw{delete}\footnote{siehe
Kapitel \ref{memmgrops}} verwendet. Der Operand von \rw{new} ist ein
Typ, also im Gegensatz zu denen der meisten anderen Operatoren
kein Ausdruck. Der Ergebnistyp ist ein Zeiger auf den
"ubergebenen Typ, der auf den neu allokierten Speicherbereich
zeigt. Dieser Bereich kann nun so lange wie n"otig benutzt
werden. Wird er nicht mehr verwendet, sollte er mit dem Operator
\rw{delete} freigegeben werden, um sp"ater anderen dynamischen
Variablen zur Verf"ugung zu stehen. Der Grund, warum ich aus
\rw{new} und \rw{delete} keine Anweisungen gemacht
habe---besonders bei \rw{delete} w"urde sich das anbieten---ist,
da"s alle anderen Anweisungen in SEL zur Flu"skontrolle im
Programm verwendet werden, soda"s dies gewisserma"sen ein Bruch
in der Konsistenz der Sprache w"are.

\subsection{Strukturen}
\label{structs}

Eine Struktur\footnote{Der Name "`Struktur"' stammt von C, in
\pascal\ wird f"ur dasselbe der Name "`Record"' verwendet} ist
der Zusammenschlu"s mehrerer Variablen, auch verschiedenen Typs,
unter einem Namen. Sie bilden zusammen mit Zeigern und
dynamischen Variablen die Basis f"ur weit komplexere, verzweigte
und verkettete Datenstrukturen. "Uber ihre Definition m"ochte ich
nicht viele Worte verlieren, da sie einfach und zweckm"a"sig ist
und es au"serdem nicht viele andere M"oglichkeiten gibt,
Strukturen zu definieren.  Was jedoch viel interessanter ist, ist
die Benutzung von Zeigern auf eine Struktur in der Struktur
selbst. Ein Beispiel:

\begin{verbatim}
struct List
  i    : int;
  Next : *List;
end
\end{verbatim}

Diese Struktur ist eine m"ogliche Definition einer einfach
verketteten Liste. Neben einem oder mehreren Werten (in diesem
Fall ein \tint ) enth"alt sie einen Zeiger auf eine Struktur
ihres eigenen Typs. Einen Interpreter oder Compiler, der mit nur
einem Pass arbeitet und diese Struktur akzeptiert, zu schreiben,
ist noch leicht: Man k"onnte sich schon beim Lesen des auf das
Schl"usselwort \rw{struct} folgenden Bezeichners \fu{List}
notieren, da"s es diesen Typ von Struktur gibt. Beim Zeiger
\fu{Next} kann diese Notiz benutzt werden, um seinen Typ
festzulegen. Beim Lesen des Schl"usselwortes \rw{end} wird die
Notitz anschlie"send zu einer kompletten Definition erg"anzt. Was
tun wir jedoch, wenn es zwei Strukturen gibt, die jeweils einen
Zeiger auf die andere Struktur enthalten, wie es nicht selten
vorkommt? Eine der beiden Strukturen mu"s zuerst gelesen werden,
und dort treffen wir auf einen unbekannten Typen. In C wird dieses
Problem dadurch gel"ost, da"s man dem Compiler vor der Definition
der beiden Strukturen einen Hinweis darauf gibt, da"s diese
Strukturen irgendwo einmal definiert werden, soda"s er schon vor
der Definition seine Notiz machen kann. In SEL jedoch k"onnen wir
dieses Problem weit einfacher l"osen: Da es nicht m"oglich ist,
einem Typ in SEL einen anderen Namen zu geben, steht, wenn ein
unbekannter Typ gelesen wird, zweifelsfrei fest, da"s es sich
dabei um eine Struktur handeln mu"s, soda"s unser Compiler bzw.
Interpreter bereits hier die Notiz vornehmen kann.

\subsection{Typecasting}
\label{casts}

Bei einer Gegen"uberstellung der Typkonzepte von C und \pascal\
f"allt auf, da"s \pascal\ in der Typ"uberpr"ufung weit strenger
ist als C. So ist es in C ohne weiteres m"oglich, einem \tchar\
einen numerischen Wert zuzuweisen, wohingegen in \pascal\ zu
diesem Zweck ein eigener Typ existiert. In der Sprache \pascal\
ist eine Umwandlung eines Typs in einen anderen gar nicht
vorgesehen,\footnote{au"ser in F"allen, wo die beiden Typen zur
gleichen Gruppe geh"oren, wie z.~B. \tint\ und \type{longint}} in
C wird sie in vielen F"allen sogar vom Compiler "ubernommen, an
anderen Stellen kann sie der Benutzer mittels Typecasting
durchf"uhren. Der Grund, warum ich in SEL ein flexibles
Typkonzept eingebaut und nicht die strenge "Uberpr"ufung von
\pascal\ "ubernommen habe, ist, da"s es bei vier Basistypen sehr
schwer ist, ohne Typumwandlung auszukommen. Au"serdem k"onnen
Programme durch implizites Typecasting um einiges k"urzer und
"ubersichtlicher werden, man denke nur an die Umwandlung zwischen
\tchar\ und \tint , wof"ur in \pascal\ eigene Funktionen
notwendig sind.

Um aber nicht in einem vollst"andigen Chaos zu enden, darf nicht
wahllos jeder Typ implizit in jeden anderen Typen umgewandelt
werden, nur damit im Programm ohne Fehlermeldung fortgefahren
werden kann. Die Regeln, die in SEL f"ur implizites sowie
explizites Typecasting gelten, sind in \cite{selkr}, Seite 24
zusammengefa"st.

Der Grund, warum ich nicht das in C "ubliche Typecasting, einem
Ausdruck den neuen Typ in Klammern voranzustellen, "ubernommen
habe, ist der, da"s ich einfach etwas Neues ausprobieren wollte,
und das Typecasting mittels des \fu{:}-Operators exzellent zu den
"ubrigen Konstrukten in SEL (Variablen- und Strukturdefinition,
Parameterdefinition in Funktionen u.~s.~w.) pa"st. Es zeigt sich,
da"s Programme bei richtiger Verwendung des Typecasting-Operators
weit "ubersichtlicher sind als in C, meine Wahl hat sich also in
keiner Weise als Fehlgriff erwiesen.

\subsection{Mengen, Listen, B\protect{"aume} etc.}

Komplexere Typen, wie z.~B. Mengen, Listen oder B"aume sind in
SEL nicht vorhanden, aus guten Gr"unden. Mengen einerseits
werden relativ selten verwendet und lassen sich leicht
vermeiden. Au"serdem sind sie sehr leicht mit den Mitteln von
SEL zu implementieren, soda"s das Einbeziehen in SEL nur
"au"serst wenig Sinn h"atte. Listen und B"aume hingegen werden
zwar oft verwendet, sind allerdings so komplex und existieren in
einer derartig gro"sen Anzahl von Variationen, da"s es mit
einfachen Mitteln kaum m"oglich ist, sie vern"unftig zu
unterst"utzen.

\section{Bezeichner und Schl\protect{"usselw\protect{"orter}}}

Wie in vielen Programmiersprachen "ublich, sind auch in SEL
Schl"usselw"orter reserviert, d.~h., sie k"onnen nicht als
Bezeichner verwendet werden, aus zwei Gr"unden. Erstens wird
dadurch die Programmierung eines Interpreters bzw. Compilers
erheblich erschwert, und zweitens w"aren dadurch sehr
verwirrende Konstruktionen wie z.~B.

\begin{verbatim}
for (for = to to step step end)
  ...
end
\end{verbatim}

\noindent m"oglich, was nicht im Sinn des Designers sein kann.

Schl"usselw"orter und Bezeichner sind in SEL case-sensitiv,
d.~h., da"s zwei W"orter, die sich nur dadurch unterscheiden,
da"s ein oder mehrere Buchstaben statt gro"s klein geschrieben
sind oder umgekehrt, als unterschiedlich betrachtet werden. Das
hat drei Vorteile: Erstens wird dadurch die Entwicklung der
lexikalischen Analyse\footnote{siehe Kapitel \ref{scanner}}
ein wenig einfacher, zweitens k"onnen dadurch verschiedene Arten
von Bezeichnern (Bezeichner f"ur Funktionen, Variablen, Makros
u.~s.~w.) durch verschiedene Schreibweise gekennzeichnet werden.
Drittens wird durch eine vorgegebene Schreibweise bei
unerfahrenen oder schlampigen Programmierern die
"Ubersichtlichkeit der Quelltexte gesteigert.

Die Syntax der Bezeichner ist relativ simpel: Sie bestehen aus
einem Buchstaben oder Unterstrich, gefolgt von einer beliebigen
Anzahl von Buchstaben, Ziffern oder Unterstrichen. Der Grund,
warum Bezeichner nicht mit Ziffern beginnen d"urfen ist der, da"s
dadurch eine Eindeutigkeit der Sprache nicht mehr gew"ahrleistet
w"are. So k"onnte z.~B. \fu{12to} einerseits als Zahl \fu{12}
gefolgt vom Schl"usselwort \rw{to}, andererseits aber auch als
Bezeichner gedeutet werden. Dieser Konflikt lie"se sich zwar
durch Verwendung von Leerzeichen vermeiden, w"urde aber die
lexikalische Analyse ungleich schwerer machen. Warum ich nicht
auch andere Symbole in Bezeichnern erlaubt habe, ist einfach: Je
weniger Symbole in Bezeichnern verwendet werden d"urfen, desto mehr
hat man f"ur andere Zwecke frei.
                            
\section{Weiterf\protect{"uhrende} Projekte}
\label{langext}

Keine Programmiersprache ist perfekt, und schon gar nicht SEL,
f"ur dessen jetzige Form ich bewu"st die Versionsnummer 0.3
gew"ahlt habe. Leider war es mir aus Zeit- und Platzgr"unden
nicht m"oglich, auch nur einen Teil der von mir geplanten
Erweiterungen in diese Arbeit einzubringen, soda"s ich sie in
diesem Kapitel kurz zusammenfassen m"ochte, einerseits, um zu
zeigen, was SEL noch alles fehlt und andererseits, um den Leser
anzuregen, sich mit der Materie eingehender auseinanderzusetzen.

\begin{description}
\item[Zeiger auf Funktionen] stellen ein sehr praktisches
Sprachkonzept dar. Sie erm"oglichen es z.~B., in einer
graphischen Benutzeroberfl"ache, wie etwa Windows NT oder meinen
Schani-Tools, jedem Bildschirmfenster eine Funktion zuzuweisen,
die die Verwaltung desselben "ubernimmt. Sprachlich gesehen liegt
die einzige Schwierigkeit bei der Typangabe eines solchen
Funktionszeigers, denn der Aufruf einer Funktion durch einen
Zeiger kann genauso wie bei einer "`richtigen"' Funktion durch
den Namen der Zeigervariablen, gefolgt von einer in Klammern
gesetzten Parameterliste erfolgen, und die Adresse einer Funktion
kann mittels deren Namen, ohne nachfolgende Parameterliste in
Klammern, ermittelt werden. Doch auch der Typ f"ur einen
Funktionszeiger kann relativ leicht ohne Zweideutigkeiten in die
Grammatik eingebaut werden. So k"onnte etwa \fu{char(int,float)}
f"ur einen Zeiger auf eine Funktion, die einen \tint\ und einen
\tfloat\ als Parameter erwartet und einen \tchar\ zur"uckliefert,
stehen.\footnote{Tats"achlich f"uhrt eine Erweiterung wie oben
beschrieben zu einem kleinen Konflikt in der Grammatik, der
jedoch leicht beseitigt werden kann.}

\item[Variante Strukturen] sind Strukturen, die zwar mehrere
Elemente haben, diese jedoch denselben Speicherbereich
einnehmen. So etwas ist immer dann von Vorteil, wenn eine
Struktur ben"otigt wird, die Objekte beschreibt, die zwar in
gewisser Hinsicht "ahnlich sind, jedoch andere Parameter haben.
Ein Beispiel hierf"ur ist die in Kapitel \ref{intexprs}
besprochene Datenstruktur f"ur SEL-Ausdr"ucke.

\item[Makros mit Parametern] w"urden, wenn sie so flexibel
aufgebaut sind, da"s neue Befehle durch sie
"`vorget"auscht"' werden k"onnen, perfekt in eine Makrosprache
passen. Ein gutes Beispiel f"ur solche Makros sind die sog. User
Defines Commands (UDC's) von Clipper.

\item[Objektorientierte Programmierung] ist zwar ein edles Ziel,
jedoch nicht leicht zu verwirklichen.
\end{description}

\chapter{Der Interpreter}

Das Ziel dieses zweiten Kapitels ist es, f"ur die im ersten
Kapitel erkl"arte Programmiersprache SEL einen Interpreter zu
entwerfen. Der Titel dieser Arbeit deutet darauf hin, da"s dieser
Entwurf theoretisch ist, also in dieser Arbeit keine
Implementation eines Interpreters zu finden ist. Die Gr"unde
hierf"ur sind verschieden. Schon lange bevor ich "uberhaupt die
Idee zu dieser Fachbereichsarbeit hatte, wollte ich einen
Interpreter schreiben. Als ich mit Herrn Prof. Freiler "uber das
Thema dieser Arbeit diskutierte, legte er mir die Beschr"ankung
auf, da"s, wenn das Thema dieser Arbeit ein Programm w"are,
dieses in \pascal\ geschrieben sein m"u"ste. Da ich aber den
Interpreter in erster Linie f"ur mich und nicht f"ur diese Arbeit
programmieren wollte, hatte ich vor, ihn in C zu schreiben. Aus
diesem Grund einigte ich mich mit Herrn Prof. Freiler darauf,
da"s ich den Interpreter nur theoretisch, also unabh"angig von
irgendeiner Implementation, behandeln w"urde. Im Nachhinein
betrachtet bin ich froh "uber diese Entscheidung, da es kaum
m"oglich gewesen w"are, in einer im Umfang derma"sen beschr"ankten
Arbeit einen ganzen Interpreter und eine angemessene Erkl"arung
dazu unterzubringen.

Doch auch mit der Beschr"ankung auf eine theoretische Abhandlung
ist der Umfang dieser Arbeit zu gering, um ausf"uhrlich auf alle
Details einzugehen. Aus diesem Grund m"ochte ich im Gegensatz 
zum ersten Kapitel, in dem ich alle meine Entscheidungen 
begr"undet und evtl. auch andere L"osungsans"atze beschrieben 
habe, in diesem nur den Entwurf pr"asentieren und lediglich in
einigen seltenen F"allen andere m"ogliche L"osungen diskutieren.

Dieses Kapitel beschreibt alle wichtigen Algorithmen und
Datenstrukturen des Interpreters, geht jedoch nicht n"aher auf
grundlegende Datenstrukturen wie etwa dynamische Listen ein.
Sollte der Leser mit diesen Strukturen nicht vertraut sein,
empfehle ich ihm die Lekt"ure von \cite{sedge}, Kapitel 3--5 und
16.

\section{Der Begriff des Interpreters}
\label{whatisit}

Um ein Programm einer h"oheren Programmiersprache auf einem
Computer auszuf"uhren, kann man sich zweier verschiedener
Werkzeuge bedienen: Interpreter und Compiler. Der Begriff
"`Compiler"' wird vom Duden Informatik wie folgt definiert:

\begin{quotation}
Programm, das Programme aus einer Programmiersprache A in eine
Programmiersprache B "ubersetzt.\footnote{siehe \cite{duden},
Seite 745}
\end{quotation}

Das bedeutet, da"s ein Compiler ein Programm gar nicht direkt
ausf"uhrt, sondern es lediglich in eine andere Programmiersprache
"ubersetzt. In der Regel ist diese Sprache eine Assemler- oder
Maschinensprache, die direkt vom Computer "`verstanden"' wird. Im
Gegensatz dazu ist ein Interpreter

\begin{quotation}
Ein Programm, welches ein Programm einer anderen
Programmiersprache nach den notwendigen syntaktischen
"Uberpr"ufungen sofort ausf"uhrt. Im Gegensatz zu einem
"Ubersetzer\footnote{andere Bezeichnung f"ur Compiler} mu"s das
Quellprogramm nicht erst in eine andere Programmiersprache
"ubersetzt werden, sondern der Interpreter analysiert
nacheinander jede Anweisung und Deklaration des Quellprogramms
und f"uhrt diese unmittelbar aus.\footnote{siehe \cite{duden},
Seite 322}
\end{quotation}

In der Praxis jedoch wird kaum ein Interpreter ein Programm
ausf"uhren, ohne es vorher in eine geeignetere Form umgewandelt
zu haben. Man kann sich leicht vorstellen, da"s ein Interpreter,
der etwa bei einem Funktionsaufruf alle Module nach der Deklaration
dieser Funktion durchsuchen mu"s, zu langsam ist, um sinnvoll
eingesetzt werden zu k"onnen. Auch das wiederholte Lesen und
Erkennen der Schl"usselw"orter, die im Quelltext verwendet
werden, ist zu zeitaufwendig, als da"s es w"ahrend der
Ausf"uhrung des Programms durchgef"uhrt werden k"onnte. In den
folgenden Kapiteln werden wir deshalb besprechen, wie sich ein
Interpreter die Arbeit so einteilen kann, da"s eine m"oglichst
hohe Ausf"uhrungsgeschwindigkeit erreicht wird.

Ich m"ochte anmerken, da"s das in diesem Kapitel besprochene
Design f"ur einen Interpreter zwar theoretisch, jedoch keineswegs
realit"atsfern ist. Ich habe selbst einen SEL-Interpreter auf
diese Weise implementiert.

\section{Einsatzgebiete eines Interpreters}

Wenn man bedenkt, da"s die Ausf"uhrung eines Programms durch
einen Interpreter oft mehrere hundert Male l"anger dauert als die
eines mit einem Compiler erzeugten Maschinenspracheprogramms, fragt
man sich, welchen Zweck ein Interpreter eigentlich erf"ullen
kann. Ein Haupteinsatzgebiet f"ur Interpreter sind Computer, die
f"ur die sinnvolle Verwendung eines Compilers zu langsam sind
und/oder zu wenig Speicher bieten. Beispiele hierf"ur sind alte
Home-Computer wie z.~B. der Commodore~C-64 aber auch
Taschenrechner wie etwa der Texas Instruments~TI-85 oder der
Hewlett Packard~HP-48. Oft werden Interpreter auch zum
Entwickeln von Programmen eingesetzt, da sie den Vorteil haben,
Programme ohne Zeitverz"ogerung ausf"uhren zu k"onnen. Ein
anderes wichtiges Einsatzgebiet f"ur Interpreter ist die
Steuerung von Me"sger"aten und Robotern. Dort stellt die langsamere
Ausf"uhrung von Programmen kaum einen Nachteil dar, andererseits
ist es ein gro"ser Vorteil, da"s Programme ge"andert werden
k"onnen und sofort wieder einsatzbereit sind.

\section{Bestandteile eines Interpreters}

Bevor wir uns daranmachen, uns dar"uber zu unterhalten, wie der
Interpreter aussehen und welche Leistung er bieten soll, m"ussen
wir "uber den groben Aufbau eines Interpreters Bescheid wissen.
Wie bereits oben erw"ahnt wird in der Praxis das Quellprogramm
vor dem Interpretieren in eine f"ur den Interpreter geeignete
Form umgewandelt ({\em Tokenizing}). Dieser Vorgang wird in der
Regel von der lexikalischen Analyse erledigt. Dabei wird der
Quellcode sequentiell abgearbeitet und Schl"usselw"orter und
Bezeichner codiert. Zus"atzlich werden die Positionen von wichtigen
Schl"usselw"ortern wie z.~B. \rw{function}, \rw{global} und
\rw{local} gespeichert. Im Fall von SEL wird gleichzeitig
zur lexikalischen Analyse der Pr"aprozessor ausgef"uhrt. Er hat die
Aufgabe, Include-Dateien einzubinden und Makros in den
Programmtext einzuf"ugen.

Im Anschlu"s an die lexikalische Analyse und das Tokenizing kann
mit der Ausf"uhrung des Programms begonnen werden. Bei einer
Programmiersprache wie SEL m"ussen zuerst nat"urlich spezielle
Initialisierungen vorgenommen werden, n"amlich das Einrichten der
modul- und programmglobalen Variablen und das Analysieren der
Funktionsk"opfe. Zus"atlich speichert der Interpreter die
Stellen, an denen die jeweiligen Funktionen beginnen, soda"s
nicht das gesamte Programm durchsucht werden mu"s. Jetzt wird das
Programm durch einen Aufruf der \fu{main}-Funktion gestartet. Die
einzelnen Anweisungen und Ausdr"ucke werden vom Parser analysiert
und sofort ausgef"uhrt. Dabei mu"s nat"urlich auf spezielle
Routinen, die z.~B. f"ur die Verwaltung der Variablen und
"ahnliches verantwortlich sind, zur"uckgegriffen werden.

Zusammenfassend l"a"st sich also sagen, da"s ein Interpreter aus
drei wichtigen Teilen besteht: Der lexikalischen Analyse, die das
Tokenizing vornimmt, dem Parser, der das Programm syntaktisch
analysiert und den Verwaltungsroutinen, die die Hintergrundarbeit
leisten.

\section{Zielsetzung f\protect{"ur} das Design}
\label{intgoals}

Wie schon im ersten Kapitel erw"ahnt, ist SEL mit dem Ziel, sie
als Makrosprache zu verwenden, entworfen worden. Es w"are also
angebracht, den Interpreter von Beginn an so zu entwerfen, da"s
er nicht ein eigenst"andiges Programm darstellt, sondern eher
eine Art Modul, das man in andere Programme aufnehmen kann. Das
setzt nat"urlich voraus, da"s er nicht eine Unmenge an Speicher
verbraucht und auch, da"s er den Speicher, den er verbraucht,
nach der Ausf"uhrung des SEL-Programms wieder freigibt, soda"s er
wieder dem Hauptprogramm zur Verf"ugung steht.

Der Interpreter sollte nat"urlich, um ohne Bedenken in einem
Programm eingesetzt werden zu k"onnen, Programmierfehler so
schnell wie m"oglich erkennen k"onnen, den Ausf"uhrungsproze"s
angemessen stoppen, allen Speicher, den das SEL-Programm
angefordert hat, und auch den, den der Interpreter im Zuge der
Ausf"uhrung des Programms benutzt hat, freigeben und den
Fehler dem Hauptprogramm melden.

Wie auch die Programmiersprache SEL sollte auch der Interpreter
erweiterungsf"ahig sein, d.~h., er sollte dem Hauptprogramm
erlauben, bestimmte Funktionen auch SEL-Programmen zug"anglich zu
machen, wie beispielsweise f"ur das Zeichnen einer Linie in einem
Zeichenprogramm.

Da ich mich in letzter Zeit mit dem Multitasking-Betriebssystem
Windows NT auseinandergesetzt habe, wollte ich auch meinen
Interpreter multitaskingf"ahig machen, soda"s es in einem
Programm m"oglich ist, mehrere SEL-Programme gleichzeitig
auszuf"uhren. Um mich nicht falsch zu verstehen: Das Multitasking
wird in keiner Weise vom Interpreter in die Hand genommen,
sondern obliegt einzig und allein dem Betriebssystem. Der
Interpreter allerdings mu"s so entworfen werden, da"s er das
Multitasking nicht erschwert oder gar unm"oglich macht, was, wie
sich zeigen wird, nicht sehr schwer ist.

Im Laufe dieses Kapitels werde ich auf einige, aber nicht auf alle
Forderungen dieses Unterkapitels eingehen, weil sie zu trivial
sind. So ist z.~B. die Freigabe des benutzten Speichers nur eine
Formsache und mu"s nicht extra von mir erl"autert werden.

\section{Die lexikalische Analyse}
\label{scanner}

Die lexikalische Analyse hat die Aufgabe, den Zeichenstrom des
Quelltextes in einen Tokenstrom zu verwandeln. Als Token
bezeichnet man die kleinste Einheit, die vom Parser bearbeitet
wird. Token sind z.~B. die Schl"usselw"orter, Bezeichner,
Konstanten u.~s.~w. Der Grund, warum man die Erkennung von
Schl"usselw"ortern etc. nicht auch dem Parser "uberl"a"st, was
durchaus m"oglich w"are, ist der, da"s ein Parser diese Aufgabe
bei weitem nicht so schnell ausf"uhren kann wie ein eigens daf"ur
angefertigter {\em Scanner}.\footnote{Begriff f"ur ein Programm
oder einen Programmteil, der die lexikalische Analyse
durchf"uhrt} Au"serdem wird der Bau eines Parsers wesentlich
vereinfacht, wenn man einen Scanner verwendet.

Um die Arbeitsweise eines Scanners zu beschreiben, m"ochte ich
zun"achst ein Beispiel heranziehen. Nehmen wir an, der Scanner
h"atte den Zeichenstrom \fu{for (x=1 to 10)} zu bearbeiten. Er
untersucht nun das erste Zeichen und merkt, da es sich um einen
Buchstaben handelt, da"s er es entweder mit einem Bezeichner oder
mit einem Schl"usselwort zu tun hat. Die einfachste Methode,
dieses Token zu analysieren ist, so lange Buchstaben, Ziffern und
Unterstriche zu lesen, bis er an einem Zeichen ankommt, das
nicht mehr aus einer der genannten Kategorien stammt. Der
Scanner w"urde also zuerst das \fu{f}, dann das \fu{o} und
anschlie"send das \fu{r} lesen. Da das n"achste Zeichen weder
Buchstabe, Ziffer noch Unterstrich ist, stoppt er hier. Die
einfachste Methode, herauszufinden, ob man es hier mit einem
Schl"usselwort oder einem Bezeichner zu tun hat, ist, eine
Tabelle anzulegen, die alle Schl"usselw"orter der Sprache
umfa"st. Findet der Scanner das Wort in der Tabelle, ist es ein
Schl"usselwort, andernfalls ein Bezeichner. Das n"achste Zeichen
ist ein Leerzeichen. In einer formlosen Programmiersprache wie
SEL werden Leerzeichen und Tabulatoren "ubersprungen, soda"s der
Scanner nichts anderes zu tun hat, als zum n"achsten Zeichen
vorzur"ucken. Zwar sind Zeilenumbr"uche f"ur den Parser nicht
relevant, doch f"ur den Pr"aprozessor sind sie wichtig, soda"s
sie gemeldet werden sollten.

Da es nur ein Symbol in SEL gibt,
das mit einer "offnenden, runden Klammer beginnt, n"amlich dieses
selbst, wei"s der Scanner sofort, mit welchen Token er es zu tun
hat und kann vorr"ucken. Interessant wird es nun bei der Analyse
des Symbols \fu{=}, denn hier kann der Scanner nicht sofort davon
ausgehen, da"s es sich um den Zuweisungsoperator handelt, da es
noch ein anderes Symbol, n"amlich \fu{==} gibt, das mit einem
\fu{=} anf"angt. So mu"s der Scanner auch hier das n"achste
Zeichen betrachten um entscheiden zu k"onnen, welches Symbol er
vor sich hat.

Wie dieses Beispiel zeigt, ist es f"ur einen
Scanner vital zu wissen, welches Zeichen als n"achstes vorliegt,
wobei er sich dieses aber auch f"ur das n"achste Token aufheben
k"onnen mu"s. Die L"osung ist einfach: Man arbeitet mit einem
Zeichen im voraus, d.~h. es wird immer das kommende Zeichen
betrachtet, nicht das aktuelle. Wird es akzeptiert, liest man das
n"achste Zeichen ein, wenn nicht, kehrt man mit dem erkannten
Token zur aufrufenden Funktion (in Falle von SEL der
Pr"aprozessor) zur"uck. Beim n"achsten Aufruf des Scanners wird
mit dem vorher zur"uckgestellten Zeichen weitergearbeitet.

Um den Proze"s der lexikalischen Analyse zus"atzlich zu
beschleunigen, sollte die Eingabe gepuffert werden. Dabei wird
ein Teil der Eingabe in einen Zwischenspeicher (Puffer) gelesen,
um von dort aus vom Scanner abgearbeitet zu werden. Wird das Ende
des Puffers erreicht, mu"s abermals ein Teil der Eingabe gelesen
werden, so lange, bis das Ende der Datei erreicht ist. In diesem
Fall sollte der Scanner der "ubergeordneten Funktion Bescheid
geben.

Bei bestimmten Tokentypen gen"ugt es nicht, zu melden, da"s
dieses Token erkannt wurde, sondern es mu"s auch ein Attribut
zur"uckgeliefert werden, im Falle eines Bezeichners sein Name, im
Fall einer \tint -Konstante ihr Wert u.~s.~w., es sollte also
f"ur ein Token eine geeignete Datenstruktur verwendet werden, in
der alle n"otigen Attribute Platz finden. F"ur unseren
Interpreter k"onnte diese Struktur etwa so aussehen:

\newpage
\begin{verbatim}
struct Token
  iType  : int;
  pcName : *char;
  cChar  : char;
  iInt   : int;
  fFloat : float;
end
\end{verbatim}

Ist das Token ein Bezeichner, enth"alt \fu{pcName}
dessen Namen, im Falle einer \tchar -Konstanten enth"alt
\fu{cChar} ihren Wert u.~s.~w.

\subsection{Der Pr\protect{"aprozessor}}

Der Pr"aprozessor ist in gewisser Weise eine Zwischenstation von
Scanner und Parser. In SEL hat er die Aufgabe, den Inhalt von
Include-Dateien in den Programmtext einzuf"ugen, Kommentare zu
entfernen und Makros zu bearbeiten.

Das Entfernen von Kommentaren ist im Grunde kein Problem. Wird
vom Scanner das Symbol \fu{//} erkannt, dann werden so lange
Token gelesen, bis ein Zeilenumbruch gemeldet wird, wonach der
Kommentar beendet ist. Zus"atzlich zu den Kommentaren sollten
auch alle Zeilenumbr"uche herausgefiltert werden, da sie den
Parser st"oren w"urden.

Etwas komplizierter ist das Einf"ugen von Include-Dateien. Da es
m"oglich sein sollte, Include-Dateien zu verschachteln, gibt es
zwei M"oglichkeiten zur Implementation. Entweder, der
Pr"aprozessor arbeitet rekursiv, soda"s er sich selbst nochmals
aufruft, wenn eine Include-Datei zu bearbeiten ist, oder aber er
benutzt einen Stack zur Zwischenspeicherung der einzelnen
Include-Level, was am Ende eigentlich auf dasselbe hinausl"auft.
Ein kleines Problem ergibt sich mit dem Puffer f"ur den
Scanner: Wenn der Pr"aprozessor nur einen einzigen Puffer f"ur alle
Include-Level zusammen benutzt, mu"s er beim Eintritt in eine
Include-Datei geleert und mit
der Eingabe aus der Include-Datei gef"ullt werden. Ist diese zu Ende,
mu"s der Puffer abermals geleert werden und mit dem
"ubergeordneten Level gef"ullt werden. Das Problem dabei ist,
da"s man in der Datei dabei ein wenig zur"ucksetzen mu"s, weil
ein Teil des Puffers gel"oscht wurde. Ist das aus
irgendwelchen Gr"unden nicht m"oglich oder zu ineffizient,
mu"s man f"ur jeden einzelnen Include-Level einen Puffer
benutzen, was nat"urlich mehr Speicher verbraucht.

Zum Bearbeiten von Makros mu"s eine Datenstruktur
angelegt werden, in der alle Makros aufgezeichnet und auch schnell wieder
gefunden werden k"onnen. Ein bin"arer Baum und eine Hash-Tabelle
bieten sich hierzu an. F"ur jeden einzelnen Eintrag in dieser
Struktur mu"s zus"atzlich noch gespeichert werden, durch welchen
Text er zu ersetzen ist. Da der Pr"aprozessor auf den Scanner
aufbaut, ist es einfacher und effizienter, nicht den Zeichenstrom
zu speichern, sondern den Tokenstrom, der vom Scanner geliefert
wird. Dazu verwendet man am einfachsten eine einfach verkettete
Liste aus Token:

\begin{verbatim}
struct TokenList
  tokToken : Token;
  Next     : *TokenList;
end
\end{verbatim}

Trifft der Pr"aprozessor nun auf einen Bezeichner, mu"s er zun"achst
"uberpr"ufen, ob nicht vielleicht ein Makro unter diesem Namen
exisitiert. Dazu sucht er
nach ihm in der Datenstruktur. Ist er dort vorhanden,
setzt der Pr"aprozessor anstatt seiner die Tokenliste ein,
andernfalls wird er wie ein normaler Bezeichner behandelt.

Nachdem es m"oglich sein soll, innerhalb eines Makros ein anderes
zu benutzen, m"ussen wir noch etwas einbauen. Die einfachste
L"osung ist, da"s wir zum Einf"ugen eines Tokens in unseren
Zwischencode eine eigene Funktion verwenden, die, falls sie auf
einen Bezeichner trifft, der als Makro definiert ist, sich f"ur
jedes der Token aus der Liste selbst wieder aufruft. Dadurch ist
garantiert, da"s, falls in dieser Liste ein Makro auftaucht,
auch dieses entsprechend behandelt wird.
                              
\subsection{Der Zwischencode}
\label{tokencode}

Wie schon oben erw"ahnt, verwenden wir eine Art Zwischencode, um
dem Interpreter eine schnelle Ausf"uhrung zu erm"oglichen. Dabei
werden Schl"usselw"orter und Symbole mit einer eindeutigen Nummer
codiert, soda"s der Tokentyp sofort ersichtlich ist. Bei
bestimmten Tokentypen m"ussen zus"atzliche Informationen
gespeichert werden, wie z.~B. bei Konstanten oder Bezeichnern.

Um die Gr"o"se des Zwischencodes nicht zu sehr aufzubl"ahen,
werden auch die Namen der Bezeichner mit einer Nummer
codiert. Dazu wird eine Tabelle aller im Modul vorkommenden
Bezeichner angelegt und jedem Eintrag eine eindeutige Nummer
zugewiesen. Diese Nummer wird anschlie"send im Zwischencode
verwendet, um Bezeichner zu charakterisieren. Zu beachten ist,
da"s nat"urlich auch der Parser "uber diese Tabelle verf"ugen
mu"s, da er sonst nur die Nummern der Bezeichner wei"s. Wird
allerdings in allen Modulen des Programms dieselbe Tabelle
verwendet, so kann der Interpreter das Programm auch ohne das
Wissen um die Namen der Bezeichner ausf"uhren.

Eine besondere Hilfeleistung erstattet der Zwischencode dem
Interpreter bei Schleifen und Kontrollanweisungen: F"ur jedes \rw{if}
wird zus"atzlich die Position im Zwischencode abgespeichert, an der das
zugeh"orige \rw{end} bzw. \rw{else} steht. Das hat den Vorteil,
da"s der Interpreter, wenn er z~B. auf eine \rw{if}-Anweisung
trifft, deren Bedingung logisch falsch ist, ohne lange zu suchen
zum zugeh"origen \rw{end} bzw. \rw{else} springen kann, was
verst"andlicherweise die Ausf"uhrungsgeschwindigkeit erheblich
steigert, besonders bei Schleifen.

Wie bereits erw"ahnt, werden dem Interpreter zus"atzlich die
Position aller Vorkommnisse der Schl"usselw"orter \rw{global},
\rw{local} und \rw{function} mitgeteilt, soda"s er ohne zu suchen die
Initialisierung der programm- und modulglobalen Variablen und der
Funktionsk"opfe vornehmen kann.

Neben der erheblichen Geschwindigkeitssteigerung, die der
Zwischencode erm"oglicht, hat er noch den Vorteil, da"s er
weniger Speicherplatz als der Quellcode ben"otigt, bei etwas
gr"osseren Programmen etwa die H"alfte, was besonders unter
MS-DOS entgegenkommend ist.

Ich m"ochte hier kurz noch darauf eingehen, wie die oben
angesprochenen Verweise von Anweisungen auf die zugeh"origen
\rw{end}s bzw. \rw{else}s erzeugt werden k"onnen. Diese auf den
ersten Blick umst"andliche Aufgabe l"a"st sich relativ leicht
l"osen, indem man einen Stack zur Speicherung der
Verschachtelungen der Anweisungen verwendet. Trifft der Scanner z.~B.
auf ein \rw{if}, dann pusht er die Position, an der er sich
gerade im Zwischencode befindet, auf den Stack und l"a"st Platz
f"ur den sp"ater folgenden Verweis frei. Als n"achstes trifft
er z.~B. auf ein \rw{for} und tut dasselbe. Irgendwann mu"s
er jedoch zu dem \rw{end} gelangen, da"s die jeweils letzte
Anweisung schlie"st. In diesem Fall speichert er das \rw{end} im
Zwischencode, poppt vom Stack die Position der letzten
Anweisung und f"ugt an dieser Stelle die Position ein, an der
er soeben das \rw{end} gespeichert hat. Diese Technik wird
"ubrigens {\em Backpatching} genannt und oft im Compilerbau
eingesetzt.\footnote{Tats"achlich ist unser Scanner, der
einen SEL-Quelltext in Zwischencode umwandelt, streng gesehen
auch ein Compiler.}

\bigskip

Eine effiziente und einfache Darstellung des Zwischencodes ist
die Verwendung jeweils eines Bytes f"ur ein Token, wobei bei
bestimmten Tokentypen, wie z.~B. dem \rw{if}-Token, ein
zus"atzlicher Verweis auf ein anderes Token eingebaut wird. Wie
viele Bytes f"ur diesen Verweis verwendet werden, h"angt von den
Anforderungen an den Interpreter ab. Ich verwende hierzu in
meinem Interpreter 32 Bits, soda"s der Zwischencode theoretisch
bis zu 4 GByte gro"s sein kann, was sicherlich f"ur alle Anwendungen
mehr als ausreichend ist. F"ur die meisten
Anwendungen w"urden auch 16 Bits reichen, da Module in der Regel
sowieso klein gehalten werden. Bei anderen Token wie etwa einer
String-Konstante mu"s au"ser dem Tokentyp noch ein Wert
gespeichert werden, wie in diesem Fall der String.

\section{Der Parser}
\label{parser}

Der wohl komplizierteste Teil des Interpreters ist der
Parser. Er hat zur Aufgabe, aus dem Tokenstrom des Zwischencodes
die einzelnen Anweisungen zu erkennen. Sehr wichtig f"ur den Bau
eines Parsers ist die Grammatik, denn sie gibt an, was der Parser
eigentlich erkennen soll. Aus verschiedenen Gr"unden werden wir
Parser schreiben, die Teile der Sprache SEL erkennen k"onnen,
n"amlich Typen, Ausdr"ucke und Folgen von Anweisungen. Die Parser
f"ur Funktionsk"opfe und \rw{global}- bzw. \rw{local}-Bl"ocke
sind trivial, soda"s ich hier nicht speziell auf sie eingehen
werde. Ein Grund f"ur diese Trennung ist, da"s der
SEL-Interpreter vor allem f"ur die Verwendung in anderen
Programmen gedacht ist und dort ist es in vielen Situationen von
Vorteil, wenn der Benutzer Ausdr"ucke eingeben kann, z.~B. als
Filterbedingung in einer Datenbank.
                              
\subsection{Der Typen-Parser}

Beim Design eines Parsers ist es zun"achst wichtig, "uber die
Grammatik Bescheid zu wissen, die der Parser erkennen soll. In
diesem Unterkapitel besch"aftigen wir uns mit einem sehr
einfachen Teil der Grammatik von SEL, n"amlich den Typen. Das
gibt uns die M"oglichkeit, uns "uber die verschiedenen Arten von
Parsern und "uber die Zusammenh"ange zwischen Grammatik und
Parser zu informieren. Der "Ubersichtlichkeit halber sei an
dieser Stelle die Grammatik dargestellt:

\begin{verbatim}
type:          `char'
             | `int'
             | `float'
             | `void'
             | `[' INT_CONSTANT `]' type
             | `*' type
             | IDENTIFIER
             ;
\end{verbatim}

Diese Grammatik hat sieben Produktionen, das Startsymbol
ist \fu{type}.

\subsubsection{Top-Down und Bottom-Up Parser}
\label{topdbotu}

Im Interpreter- und Compilerbau werden haupts"achlich zwei
verschiedene Arten von Parsern verwendet: Top-Down und Bottom-Up
Parser. W"ahrend Top-Down Parser beim Startsymbol beginnen und
immer weiter hinabsteigen, bis sie auf die Token treffen, fangen
Bottom-Up Parser bei den Token an und arbeiten sich bis zum
Startsymbol hinauf. Das l"a"st sich am besten
verdeutlichen, indem man versucht, einen Parse-Tree zu
konstruieren, also einen Baum, der f"ur einen g"ultigen
Ausdruck der Sprache, in diesem Fall etwa \fu{*[10]Test}, den
Weg vom Startsymbol bis hinunter zu den einzelnen Token,
verdeutlicht. Der Parse-Tree f"ur den Beispieltyp ist in
Abbildung \ref{parsetree} dargestellt.

\begin{figure}

\centerline{\includegraphics[height=4cm]{partree1.eps}}

\caption{Der Parse-Tree f"ur \fu{*[10]Test}}
\label{parsetree}
\end{figure}

Die Wurzel des Baumes ist immer das Startsymbol, in diesem Fall
also \fu{type}. Dieser \fu{type} kann auf der ersten Ebene
zerlegt werden in das Sternsymbol und einen zweiten \fu{type},
gem"a"s der Produktion \fu{`*' type}. Dieser zweite \fu{type}
kann zelegt werden in eine "offnende und eine schlie"sende
Klammer und eine \fu{INT\_CONSTANT} dazwischen, gefolgt von noch
einem \fu{type}. Dieser wiederum ist ein \fu{IDENTIFIER}, n"amlich
\fu{Test}.

Nun zu den Vorgehensweisen der beiden Parsertypen. Ein Bottom-Up
Parser arbeitet sich wie gesagt von den einzelnen Token zum
Startsymbol hinauf. Dazu benutzt er einen Stapel, auf dem er zwei
verschiedene Aktionen durchf"uhren kann: {\em Shift} und
{\em Reduce}. Aus diesem Grund werden Bottom-Up Parser auch oft
als Shift-Reduce Parser bezeichnet. Die Analyse des
Beispieltyps f"uhrt der Parser so durch: Zun"achst liest er den
Stern, den er shiftet. Auch die beiden Klammern und die
\fu{INT\_CONSTANT} werden von ihm geshiftet. Nun liest er einen
Bezeichner und entscheidet, diesen zu einem Typen zu reducen,
gem"a"s der Produktion \fu{IDENTIFIER}. M"u"ste er einen
Parse-Tree erstellen, h"atte er an dieser Stelle einen Baum
mit der Wurzel 3\footnote{siehe Abbildung \ref{parsetree}}
erzeugt. Auf
dem Stack des Parsers liegen nun die Token \fu{`*'}, \fu{`['},
\fu{INT\_CONSTANT}, \fu{`]'} und das Nonterminal \fu{type}, in der
angegebenen Reihenfolge. An dieser Stelle entscheidet der Parser,
die letzten vier Symbole zu reducen. Um einen Parse-Tree zu
erhalten, m"u"ste er nun Baum 3 verwenden, um den gr"o"seren
Baum 2 zu konstruieren. Zuletzt w"urde der Parser nochmals
reducen und den Baum 1 erstellen. Der Parser hat nun
seine Analyse beendet, da das einzige, was auf seinem Stack
liegt, das Startsymbol ist. Die Hauptaufgabe bei der Erstellung
eines Bottom-Up Parsers ist, zu entscheiden, wann er zu shiften
und wann zu reducen hat. Dies ist eine schwere Aufgabe und wird
in der Praxis von einem Parser-Generator (auch Compiler-Compiler
genannt) wie etwa YACC erledigt. Da wir jedoch unseren Parser
selbst programmieren wollen, beschr"anken wir uns auf die
wesentlich einfacheren Top-Down Parser.

Ein Top-Down Parser geht bei der Analyse von \fu{*[10]Test}
folgenderma"sen vor: Er erkennt das Sternsymbol und wei"s nun,
da"s er nach der Produktion \fu{`*' type} vorzugehen hat. An
dieser Stelle w"urde er den Baum 1 erzeugen. Er liest nun den
Stern ein und mu"s nun einen weiteren Typ analysieren. Hier wird
klar, da"s ein Top-Down Parser in der Regel rekursiv
implementiert wird. Zwar ist es m"oglich, diese Rekursion mit
einem Stack zu beseitigen, doch ein solcher Parser ist schwerer
zu programmieren und zu warten, soda"s ich eine rekursive
Implementation gew"ahlt habe.

Der Parser ruft sich nun selbst auf und erkennt anhand der
"offnenden eckigen Klammer, da"s er die Produktion
\fu{`[' INT\_CONSTANT `]' type} w"ahlen mu"s und erzeugt
gleichzeitig Baum 2. Nacheinander liest er nun die drei folgenden
Token und ruft sich abermals selbst auf. Nun wendet er die
Produktion \fu{IDENTIFIER} an und erstellt Baum 3. Damit ist die
Analyse beendet.

\subsubsection{Look-Aheads}

Wie sich in diesem Beispiel gezeigt hat, ist es f"ur einen Parser wichtig, zu
wissen, welches Token er als n"achstes zu bearbeiten hat, denn
danach hat er zu entscheiden, welche Produktion er ausw"ahlt. Zu
diesem Zweck verwendet man sog. {\em Look-Aheads}. F"ur einen
SEL-Parser und f"ur fast alle Programmiersprachen ist ein Look-Ahead
v"ollig ausreichend, es gen"ugt also, ein Token
vorauszusehen. Die Implementation eines Look-Ahead ist simpel:
Anstatt nur des aktuellen Tokens verwendet der Parser zus"atzlich das
Look-Ahead. Soll nun das aktuelle Token "ubersprungen werden,
weist er einfach den Wert des Look-Aheads dem aktuellen Token
zu und liest einen neuen Look-Ahead aus dem Zwischencode ein.

\subsubsection{Die Konstruktion des Parsers}

Mit diesen Informationen ist es nun ein Leichtes, einen
rekursiven Top-Down Parser f"ur Typen zu bauen. Er besteht aus
einer einzigen Funktion, die durch Betrachten des Look-Ahead
entscheidet, was passieren soll. Ist der Look-Ahead eines der
Token \fu{`char'}, \fu{`int'}, \fu{`float'}, \fu{`void'} oder
\fu{IDENTIFIER}, mu"s der Parser dieses Token nur lesen und die Analyse
ist beendet. Im Fall einer "offnenden eckigen Klammer liest er
drei Token und ruft sich anschlie"send selbst auf.\footnote{Es
ist leicht m"oglich, eine am Ende stehende Rekursion durch einen
Sprung an den Beginn der Funktion, z.~B. durch eine
\rw{repeat}-Schleife, evtl. durch einen Stack unterst"utzt, zu
ersetzen, was im Normalfall die bessere L"osung ist.} Analog dazu
ist die Produktion \fu{`*' type} zu implementieren. Sollte der
Look-Ahead keinem der Token entsprechen, hat es der Parser
offensichtlich mit einem Fehler zu tun, den er entsprechend
behandeln mu"s. Ein Interpreter wird in diesem Fall eine
Fehlermeldung ausgeben und die Ausf"uhrung beenden.

Nat"urlich ist klar, da"s unser Parser nicht nur die Aufgabe hat,
einen Typ zu erkennen, sondern auch, mitzuteilen, um welchen Typ
es sich dabei handelt. Das jedoch f"allt unter die semantische
Analyse und wird in Kapitel \ref{inttypes} behandelt.

\subsection{Der Ausdruck-Parser}
\label{exprparser}

Nachdem wir nun die grundlegenden Funktionen eines Parsers
kennen, wollen wir einen etwas komplexeren Parser konstruieren:
Den Parser f"ur SEL-Ausdr"ucke. Die Grammatik hierzu kann im
Anhang \ref{grammar} oder in \cite{selkr}, Seite 14 nachgelesen
werden. Ausdr"ucke in SEL werden in der Grammatik durch die
beiden Nonterminals \fu{expr} und \fu{primary} repr"asentiert. In
den Produktionen f"ur \fu{primary} taucht noch das dritte
Nonterminal \fu{type} auf, doch dieses stellt hier kein
Problem mehr dar, da der entsprechende Parser bereits im letzten
Unterkapitel skizziert wurde.

Besonders interessant sind die Produktionen f"ur \fu{expr}, denn
sie bestehen gr"o"stenteils aus dem Nonterminal \fu{expr} selbst.
Wie sich leicht ersehen l"a"st, ist diese Grammatik auf keinen
Fall eindeutig, was die beiden Parse-Trees f"ur den Ausdruck
\fu{a=b+1} in Abbildung \ref{twopartrees}
verdeutlichen.\footnote{In einer eindeutigen Grammatik darf es
f"ur einen korrekten Satz nur einen einzigen Parse-Tree geben.}

\begin{figure}

\centerline{\includegraphics[height=3cm]{partree2.eps}}

\caption{Zwei Parse-Trees f"ur \fu{a=b+1}}
\label{twopartrees}
\end{figure}

Alle diese Konflikte jedoch lassen sich durch die in
\cite{selkr},~Seite~21 angegebene Priorit"atstabelle l"osen. Die
Frage, die sich nun stellt, ist, wie man diese Rangfolge in den
Parser einbringen kann. Betrachten wir der Einfachheit halber
nur die folgenden drei Produktionen von \fu{expr}:

\begin{verbatim}
expr:   expr `=' expr
      | expr `+' expr
      | expr `*' expr
      | primary
      ;
\end{verbatim}

Die Priorit"atstabelle sagt, da"s die h"ochste Priorit"at der
Operator \fu{*} hat, danach folgen \fu{+} und~\fu{=}. Der
Konflikt kann nun gel"ost werden, indem man f"ur jede Priorit"atsstufe ein
eigenes Nonterminal einf"uhrt, wobei die Produktionen f"ur die
Operatoren niedriger Priorit"at auf die Nonterminals der
Operatoren h"oherer Priorit"at zur"uckgreifen, wie hier
dargestellt:

\begin{verbatim}
expr3:   expr2 `=' expr2
       | expr2
       ;
expr2:   expr1 `+' expr1
       | expr1
       ;
expr1:   primary `*' primary
       | primary
       ;
\end{verbatim}

Diese Grammatik kann, angewendet auf den Ausdruck \fu{a=b+1},
nur einen Parse-Tree erzeugen, n"amlich den, der der
Klammerung des Ausdrucks \fu{a=(b+1)} entspricht.
Ungl"ucklicherweise geht mit dieser Grammatik die F"ahigkeit
verloren, einen Ausdruck wie \fu{a+b+c} darzustellen. Um auch
diesen Mangel zu beseitigen, ersetzt man in jeder Produktion mit
einem Operator die linke bzw. rechte Seite mit dem Nonterminal
des jeweiligen Levels, je nachdem, ob der Operator links- oder
rechtsassoziativ ist. Daraus resultiert folgende Grammatik:

\begin{verbatim}
expr3:   expr2 `=' expr3
       | expr2
       ;
expr2:   expr2 `+' expr1
       | expr1
       ;
expr1:   expr1 `*' primary
       | primary
       ;
\end{verbatim}

Diese Grammatik ist eindeutig und erzeugt sowohl f"ur den
Ausdruck \fu{a=b+1} als auch f"ur \fu{a+b+c} die richtigen
Parse-Trees, man kann also mit dieser Methode die "`gro"se"'
Grammatik f"ur den Bau des Anweisungs-Parsers umwandeln. Ich m"ochte die
sich daraus ergebende Grammatik in dieser Arbeit nicht
auff"uhren, da sie sehr leicht hergeleitet werden kann und zu
viel Platz beanspruchen w"urde.

Mit den bisher besprochenen Techniken kann man bereits einen
Parser f"ur einen Teil der Grammatik konstruieren, sp"atestens
bei den Linksrekursionen, wie z.~B. in den
Produktionen f"ur \fu{expr2} in der obigen Grammatik, wird man
allerdings Schwierigkeiten bekommen. Eine Rechtsrekursion wie etwa bei
\fu{expr3} ist keine Schwierigkeit: Man kann sie entweder
durch eine echte Rekursion implementieren oder durch eine
Schleife, unterst"utzt durch einen Stack. F"ur die erste Produktion
von \fu{expr2} m"u"ste man jedoch eine Funktion programmieren,
die als erstes sich selbst aufruft, was nat"urlich in einer
endlosen Rekursion m"unden w"urde.

Um dieses Problem zu l"osen, versuchen wir zu ergr"unden, welche
Menge von Ausdr"ucken eigentlich die beiden Produktionen von
\fu{expr2} beschreiben. Der einfachste Ausdruck, den die
Grammatik beschreiben kann, ist offensichtlich \fu{expr1}. Er
l"a"st sich mit Hilfe der beiden Produktionen nicht weiter
expandieren. Benutzt man als Basis nun die erste Produktion,
kann man den ersten \fu{expr2} expandieren, entweder
durch die erste oder die zweite Produktion. Benutzt man die
zweite Produktion, bekommt man als Ergebnis den Ausdruck
\fu{expr1 `+' expr1}, der sich abermals nicht weiter expandieren
l"a"st. F"uhrt man diesen Proze"s fort, wird man bald merken,
da"s diese beiden Produktionen eine Abfolge von beliebig vielen
\fu{expr1}, verbunden jeweils durch ein \fu{`+'}, beschreiben.
Zwar beschreiben die Produktionen f"ur \fu{expr3} im Grunde
dieselbe Menge, doch ist der Operator \fu{=} rechts-assoziativ,
w"ahrend \fu{+} links-assoziativ ist, was sich z.~B. beim
Erzeugen eines Parse-Trees zeigt.

Mit diesem Wissen k"onnen wir nun einen Parser konstruieren, der
einen \fu{expr2} akzeptiert, m"ussen aber wie gesagt die
Links-Assoziativit"at beachten. Um nicht unn"otig viele Worte
machen zu m"ussen, folgt ein Beispiel einer Implementation, die
einen (primitiven) Parse-Tree erzeugt. F"ur den Parse-Tree wird
folgende Struktur verwendet, die keiner Erkl"arung bedarf:

\begin{verbatim}
struct ParseTree
  Operator : int;
  Value    : int;
  Left     : *ParseTree;
  Right    : *ParseTree;
end
\end{verbatim}

Die Funktion \fu{expr2} erzeugt einen Parse-Tree f"ur einen
\fu{expr2} und greift dabei auf die Funktion \fu{expr1} zur"uck,
die dasselbe f"ur einen \fu{expr1} tut. Zus"atzlich werden die
Funktionen \fu{LookAhead} und \fu{Accept} verwendet, die ich
nicht n"aher erl"autern mu"s:

\begin{verbatim}
function expr2 () : *ParseTree

  Tree : *ParseTree;
  Temp : *ParseTree;

begin
  Tree = expr1();
  while (LookAhead() == PLUS)
    Accept(PLUS);
    Temp = new ParseTree;
    Temp->Operator = PLUS;
    Temp->Left = Tree;
    Temp->Right = expr1();
    Tree = Temp;
  end
  return Tree;
end
\end{verbatim}

Auf diese Art und Weise k"onnen Funktionen f"ur jeden der
Priorit"atsgrade erstellt werden. Im wesentlichen mit derselben
Methode werden auch die Links-Rekursionen in den Produktionen von
\fu{primary} beseitigt. Eine Kleinigkeit in der Implementation
von \fu{primary} m"ochte ich noch besprechen: Die Produktion
\fu{IDENTIFIER `(' [expr \{`,' expr\}] `)'} enth"alt ein
optionales Element und darin eingeschlossen ein Element, da"s
sich beliebig oft wiederholen kann. Es ist zwar m"oglich, solche
Konstruktionen durch Umformung der Grammatik mit den oben
besprochenen Methoden zu analysieren, doch die direkte
Implementation ist nicht schwer zu realisieren und wird deshalb
von mir vorgezogen.

Die beiden eckigen Klammern in der Produktion bedeuten, da"s auf
die "offnende runde Klammer, die nach dem \fu{IDENTIFIER} zu
stehen hat, entweder eine schlie"sende runde Klammer oder ein
\fu{expr} folgt. Welche der beiden M"oglichkeiten zutrifft,
kann der Parser dadurch erkennen, da"s er den Look-Ahead
betrachtet.
Da ein \fu{expr} niemals mit einer schlie"senden Klammer beginnen
kann, ist klar, da"s, wenn der Look-Ahead die schlie"sende
Klammer ist, die erste Variante zutrifft, andernfalls die zweite.
In letzterem Fall mu"s die Funktion \fu{expr} aufgerufen werden, um
den ersten \fu{expr} zu erkennen. Nun steht der Parser vor einem
"ahnlichen Problem. Entweder, es folgt ein Beistrich und noch ein
\fu{expr}, dann vielleicht noch ein Beistrich u.~s.~w., oder es
folgt die schlie"sende Klammer. Wie oben kann auch dieses
Problem mit dem Look-Ahead gel"ost werden. Entspricht er einem Beistrich,
mu"s noch ein \fu{expr} analysiert werden, wonach der Parser wieder
auf einen Beistrich testen mu"s. Ist der Look-Ahead die
schlie"sende Klammer, mu"s er diese nur noch lesen und ist
mit der Analyse des \fu{primary} fertig. Wenn jedoch der
Look-Ahead keinem der beiden Symbole entspricht, hat er
offensichtlich einen Fehler gefunden, der dementsprechend zu
behandeln ist.

\subsection{Der Anweisungsparser}

Der Parser f"ur Anweisungen ist relativ einfach zu realisieren,
da er weder Linksrekursionen enth"alt, noch irgendwelche
Priorit"atsregeln zu befolgen sind. Die Anweisungen von SEL sind
so konzipiert, da"s der Parser schon anhand des ersten Tokens
entscheiden kann, um welche Anweisung es sich handelt. Er
unterscheidet sich jedoch insofern von den bis jetzt besprochenen
Parsern, als da"s bei bestimmten Anweisungen nur Teile dieser
Anweisung geparst werden m"ussen, wie z.~B. bei \rw{if}, wo der Wert der
Bedingung bestimmt, ob der \rw{if}- oder der \rw{else}-Zweig,
wenn vorhanden, ausgef"uhrt werden soll.

Betrachten wir als Beispiel die zitierte \rw{if}-Anweisung. Der
Parser erkennt anhand des Look-Ahead, da"s es sich um eine
\rw{if}-Anweisung handelt. Er "uberliest also die beiden Token
\rw{if} und \fu{(} und ruft anschlie"send den im vorigen
Unterkapitel besprochenen Ausdruck-Parser auf, um den \fu{expr}
zu lesen. Wie sich in Kapitel \ref{semantics} zeigen wird,
berechnet dieser Parser den Wert des Ausdrucks, der nun
benutzt werden kann, um zu entscheiden, ob die dem \rw{if} folgenden
Anweisungen ausgef"uhrt werden sollen. Ist der Wert ungleich \fu{0}, ist
die Sache leicht. In diesem Fall mu"s sich der Anweisungs-Parser
nur selbst aufrufen, um die folgenden Anweisungen bis zum
\rw{end} bzw. \rw{else} zu parsen. Andernfalls benutzt er die
zus"atzliche Information, die ihm der in Kapitel \ref{tokencode}
besprochene Zwischencode zum \rw{if}-Token liefert. Sie gibt die
Position an, an der das zugeh"orige \rw{end} bzw. \rw{else}
steht. So m"ussen nicht langwierig alle folgenden Anweisungen bis
zum jeweiligen Token durchsucht werden, sondern es kann direkt an
die entsprechende Position gesprungen werden. Steht an dieser
Stelle ein \rw{end}, ist die Analyse der \rw{if}-Anweisung
beendet, andernfalls mu"s der \rw{else}-Zweig ausgef"uhrt werden.

Noch etwas komplizierter verh"alt es sich mit der
\rw{switch}-Anweisung. Dort mu"s das \rw{switch}-Token die
Position des ersten zugeh"origen \rw{case} enthalten, dieses
wiederum die des n"achsten \rw{case} bzw. \rw{default}
u.~s.~w. Auf diese Weise ist es m"oglich, einen Zweig nach
dem anderen auf "Ubereinstimmung zu untersuchen und, falls eine
"ubereinstimmende Konstante gefunden wurde, diesen Zweig
auszuf"uhren, indem die Funktion sich selbst aufruft. Damit die
\rw{switch}-Anweisung auch wie vorgeschrieben ausgef"uhrt wird,
mu"s der Anweisungs-Parser beim normalen Parsen ein
\rw{case}-Schl"usselwort einfach "uberspringen und bei einem
\rw{break} zur "ubergeordneten Funktion zur"uckkehren.

Ein Problemfall ist die \rw{return}-Anweisung. Wenn der
Parser auf sie trifft, mu"s er auch zur "ubergeordneten Funktion
zur"uckkehren, doch er mu"s das im Gegensatz zur
\rw{break}-Anweisung so lange tun, bis der Anweisungs-Parser
komplett verlassen wurde. Wir m"ussen uns also "uberlegen, wie
wir es anstellen k"onnten, da"s der Parser wei"s, wann die
untergeordnete Instanz ein \rw{return} entdeckt hat. Die
einfachste L"osung ist, das mit dem R"uckgabewert zu tun. Im
n"achsten Kapitel wird sich zeigen, da"s SEL-Werte vom
Interpreter in Strukturen gehalten werden. Am sinnvollsten ist
es, als R"uckgabetyp f"ur den Parser einen Zeiger auf eine solche
Struktur zu benutzen. Trifft er nun auf ein \rw{break}, gibt er
einen Nullzeiger zur"uck und der "ubergeordnete Parser wei"s, das
alles in Ordunung ist. Wurde aber eine \rw{return}-Anweisung
entdeckt, wird ein Zeiger auf den zugeh"origen Wert
zur"uckgegeben. Die "ubergeordnete Funktion wei"s nun, da"s auch
sie beendet werden mu"s und gibt ihrerseits den Wert weiter. Wird
eine Funktion geparst, die als R"uckgabetyp \tvoid\ hat, mu"s
ein Zeiger auf einen Dummy-Wert zur"uckgegeben werden.

Der Anweisungs-Parser mu"s auch mit einem Null-Zeiger verlassen
werden, wenn er auf ein \rw{end} oder \rw{else} trifft, die zu
keiner von ihm behandelten Anweisungen geh"ort, denn dann f"allt
es in den Zust"andigkeitsbereich der "ubergeordneten
Rekursionsebene.

Hier wird auch klar, warum ich auf eine
\rw{goto}-Anweisung\footnote{siehe Kapitel \ref{goto}} verzichtet
habe: In einem rekursiven Parser w"are es f"ur einen Sprung
n"otig, zun"achst zu ermitteln, auf welcher Rekursionsebene das
Label liegt, und ob es vielleicht innerhalb einer anderen Anweisung
liegt, was einen besonders schlimmen Fall darstellt. Betrachten wir
dazu folgendes (theoretische) Programmfragment:

\begin{verbatim}
if (x == 5)
  goto label;
end
if (y == 3)
  do_something();
  if (z == 4)
    label:
    do_something_else();
  end
end
\end{verbatim}

In diesem Fall m"u"ste der Interpreter, wenn er die
\rw{goto}-Anweisung ausf"uhrt, zun"achst eine Rekursionsebene
verlassen und dann in zwei andere einsteigen, was nat"urlich sehr
schwer zu realisieren ist. "Uberdies m"u"ste der Parser sehr
schnell herausfinden k"onnen, wieviel Ebenen zu verlassen und in
welche er einzusteigen hat. Zu diesem Problem ist mir bis jetzt
noch keine zufriedenstellende L"osung eingefallen.

\section{Die Verwaltung}
\label{semantics}

In den letzten Unterkapiteln wurde behandelt, wie der Interpreter
den Quelltext (genauer gesagt, den Zwischencode) lexikalisch
analysiert, in Zwischencode verwandelt und diesen syntaktisch
analysiert, jedoch wurde noch nicht gesagt, wie der Interpreter
die durch diese Analyse gewonnenen Informationen dazu verwenden
kann, das Programm auszuf"uhren. Das ist die Aufgabe dieses
Unterkapitels. Hier wird besprochen, was der Interpreter bei der
Ausf"uhrung des Programms tut, und wie dies implementiert werden
kann.

Die meisten der in diesem Unterkapitel besprochenen Algorithmen
werden w"ahrend der Syntaxanalyse, also des im letzten
Unterkapitel besprochenen Parsing, ausgef"uhrt, was auch der
Definition des Interpreters\footnote{siehe Kapitel
\ref{whatisit}} entspricht.
        
\subsection{Module}
\label{modules}

Da SEL es erlaubt, modulglobale Variablen und Funktionen zu
benutzen, ist es vonn"oten, ein SEL-Programm auch im Speicher in
Module aufzuteilen. Dazu m"ussen f"ur jedes Modul mindestens der
Zwischencode und die modulglobalen Variablen und Funktionen
gespeichert werden. Im n"achsten Unterkapitel werden sich dazu
noch einige zus"atzliche Informationen gesellen, die bei der
Fehlerbehandlung helfen.

Die programmglobalen Funktionen und Variablen werden au"serhalb
der Module f"ur das gesamte Programm gemeinsam verwaltet, was
wesentlich effizienter und auch einsichtiger ist (und es nebenbei
einfacher und schneller macht, Namenskonflike zu bemerken).
Weiters gibt es zwei Variablen, die f"ur das gesamte Programm
gelten. Sie geben die Position an, an der sich der Parser
gerade mit der Ausf"uhrung aufh"alt. Die erste dieser Variablen
nenne ich \fu{MP} (Module Pointer). Sie enth"alt einen Wert, der
das Modul, in dem sich die gerade abgearbeiteten Anweisungen
befinden, eindeutig identifiziert. Am effizientesten ist es, ein
Array aus Modulbeschreibungen anzulegen, das durch diese
Variable indiziert wird. Die zweite Variable, \fu{IP}
(Instruction Pointer), gibt die Position innerhalb des
Zwischencodes des aktuellen Moduls an.\footnote{Genauerweise
m"u"ste diese Variable den Namen \fu{TP} f"ur Token Pointer
tragen, da sie sich nicht auf die aktuelle Anweisung, sondern auf
das aktuelle Token bezieht.}

\bigskip

In Kapitel \ref{intgoals} wurde als Ziel f"ur den Interpreter
festgelegt, da"s er Multitasking unterst"utzt. Um die dazu
notwendigen Schritte anzugeben, mu"s erst gekl"art werden, wie
Multitasking funktioniert. Zun"achst m"ochte ich feststellen, da"s
ich mit Multitasking nicht die F"ahigkeit eines Betriebssystems
meine, mehrere Programme gleichzeitig zu bearbeiten, sondern die
F"ahigkeit, innerhalb eines Programms mehrere Funktionen zur
gleichen Zeit auszuf"uhren. Unter den Betriebssystemen OS/2 und
Windows NT wird daf"ur der Begriff "`Multithreading"' gepr"agt.
Im Fall eines SEL-Interpreters l"auft dieses Multitasking darauf
hinaus, da"s durch diesen Interpreter (der ja ein Programm
darstellt) mehrere SEL-Programme gleichzeitig ausgef"uhrt werden
k"onnen.

Zwei Funktionen, die gleichzeitig ausgef"uhrt werden, k"onnen
zwar auf dieselben globalen Variablen zugreifen, soda"s eine
gemeinsame Benutzung eines Datenpools, etwa einer
Funktionssammlung f"ur ein Mathematikprogramm, gew"ahrleistet
ist, jedoch besitzt jede der beiden Funktionen einen eigenen
Stack, was garantiert, da"s sie ungest"ort von den anderen
Funktionen ihre eigenen lokalen Variablen verwenden kann.

Wenn nun der Interpreter zwei verschiedene SEL-Programme
gleichzeitig ausf"uhren soll, tut er das, indem er zwei Tasks
startet, die jeweils ein Programm bearbeiten. Um zu
gew"ahrleisten, da"s sich diese beiden Funktionen dabei nicht in
die Quere kommen, d"urfen die Daten, die das aktuelle Programm
betreffen, nat"urlich nicht in globalen Variablen gespeichert
werden, sondern m"ussen der Funktion, die f"ur die Ausf"uhrung
zust"andig ist, als Parameter mitgegeben werden. F"ur unseren
rekursiven Parser bedeutet das, da"s diese Daten sehr oft als
Parameter "ubergeben werden, was nicht nur den Stack belastet,
sondern auch die Ausf"uhrungsgeschwindigkeit. Um diese Verluste
gering zu halten, empfielt es sich, nicht direkt die Daten,
sondern Zeiger auf sie zu verwenden, was aber in den meisten
Programmiersprachen ohnedies die einzige M"oglichkeit ist.

\subsection{Typen}
\label{inttypes}

Eine f"ur die Ausf"uhrungsgeschwindigkeit sehr wichtige
Entscheidung bei einem Interpreter f"ur eine Programmiersprache,
die ein im Vergleich zu anderen Interpretersprachen so komplexes
Typkonzept verwendet, wie SEL, ist, wie diese Typen
repr"asentiert werden sollen. In Compilern werden hierf"ur
h"aufig verkettete Listen verwendet, die jedoch f"ur einen
Interpreter nicht sehr brauchbar sind, da sie wegen ihres
Aufbaus, der auf dynamischer Speicherverwaltung basiert, zu
langsam sind und au"serdem zu viel Speicher verbrauchen.

Das von mir verwendete Konzept basiert darauf, da"s Typen durch
Byte-Arrays mit begrenzter L"ange repr"asentiert werden. Da es
in SEL nur eine geringe Zahl an Grundtypen und auch nicht viele
M"oglichkeiten, diese zu verkn"upfen, gibt, kann ein Byte
jeweils eine Zusammensetzungsstufe oder einen Grundtypen
enthalten. Ein Beispiel: Der Typ \type{\poi \arr{10}int} wird
durch ein Array mit der Gr"o"se von 5 Bytes repr"asentiert,
dessen erstes Byte angibt, da"s es sich um einen Zeiger handelt.
Das zweite Element gibt an, da"s ein Array im Spiel ist. F"ur
dieses Array werden auch die n"achsten zwei Bytes ben"otigt, die
die Gr"o"se angeben. Das f"unfte Element schlie"slich steht f"ur
den Grundtypen \tint\ und ist gleichzeitig die Ende-Markierung,
da es sich dabei um einen Grundtypen handelt. In gewisser Weise
"ahnlich wie \tint\ werden Strukturen behandelt, die ja auch
immer am Ende des Typs stehen. Bei ihnen mu"s, "ahnlich dem
Array, ein Wert angegeben werden, der den Typ der Struktur
bestimmt. Zu diesem Zweck vergibt der Interpreter f"ur die Strukturen
Nummern, wobei die Werte von 0 bis 9999 f"ur programmglobale
Strukturen reserviert sind. Die weiteren Zuordungen richten sich
danach, wie viel Strukturen in den einzelnen Modulen definiert
sind.

Wie schon oben erw"ahnt, besitzen die Typenarrays eine begrenzte
L"ange, was es erm"oglicht, sie schon vor ihrer Verwendung
bereitzustellen, soda"s sie nicht w"ahrend der Ausf"uhrung des
SEL-Programms dynamisch allokiert werden m"ussen. In meiner
Implementation habe ich mich f"ur eine L"ange von 20 Bytes
entschieden. Das mag nicht viel erscheinen, vor allem in
Anbetracht der Tatsache, da"s es in SEL m"oglich ist, Typen
beliebig zu verschachteln, doch ist es m"oglich, mit diesen 20
Bytes einen Zeiger auf einen Zeiger auf einen Zeiger auf ein Array
aus Arrays aus Strukturen zu repr"asentieren, wobei noch 8 Byte
frei bleiben. Dieses Beispiel soll zeigen, da"s 20 Bytes f"ur die
Praxis eine absolut ausreichende Gr"o"se ist.

Beim Analysieren von Ausdr"ucken ist es in einem SEL-Interpreter
sehr oft n"otig, Typen bereitzustellen, von denen aber viele
schon im n"achsten Moment wieder freigegeben werden. Ist es nicht
m"oglich, diese Operationen effizient auszuf"uhren, leidet
darunter die Geschwindigkeit des Interpreters sehr. Ich habe zu
diesem Zweck einen Algorithmus entwickelt, der es erm"oglicht, in
konstanter Zeit, also unabh"angig davon, wieviele von den
bereitgestellten Typenarrays bereits vergeben sind, Typenarrays
zu allokieren und freizugeben:

\bigskip

Sei $n$ die Anzahl der maximal bereitzustellenden Typen und $s$
die Gr"o"se eines Typenarrays, dann wird vor der Ausf"uhrung des
SEL-Programms ein Byte-Array mit der Gr"o"se $n \cdot s$ erzeugt,
wobei $t$ die Speicheradresse sei, an der das Array beginnt.
Weiters sei $r$ die Speicheradresse, an der ein Array der Gr"o"se
$n$ beginnt, dessen Elemente vern"unftigerweise Maschinenworte
sind. Dieses Array wird initialisiert, soda"s $r[x] = x$.
Zus"atzlich sei $i$ eine numerische Variable, die mit
dem Wert $0$ initialisiert werden.

Soll ein Typenarray allokiert werden, wird zun"achst "uberpr"uft,
ob $i = n$. In diesem Fall sind bereits alle Arrays vergeben und
eine entsprechende Fehlermeldung wird ausgegeben. Andernfalls ist
$t + r[i] \cdot s$ die Adresse, bei der das angeforderte Array
beginnt. Zus"atzlich ist $i = i + 1$ zu setzen.

Soll das Typenarray mit der Speicheradresse $f$ freigegeben
werden, so ist zuerst $i = i - 1$ und dann $r[i] = (f - t) / s$
zu setzen.

Der angegebene Algorithmus verl"a"st sich darauf, da"s nicht
Arrays freigegeben werden sollen, die nicht allokiert sind, und
da"s jeweils maximal $s$ Bytes eines Arrays genutzt werden.

\bigskip

Dieser Algorithmus verwendet zwei Datenstrukturen. Das Array $t$
ist sozusagen ein Zusammenschlu"s aller Arrays, die zur
Verf"ugung gestellt werden k"onnen, was sich nat"urlich auch
mittels eines zweidimensionalen Arrays realisieren lie"se. Das
Array $r$ ist ein Puffer, der die Nummern (0 bis $n - 1$)
aller noch freien Typenarrays enth"alt. Die Variable $i$ gibt den
Index des n"achsten zu vergebenden Typenarrays an. Wird dieses
angefordert, wird $i$ um 1 erh"oht. Umgekehrt wird $i$ um 1
vermindert, wenn ein Array freizugeben ist. Dessen Nummer wird
anschlie"send in $r[i]$ gespeichert, soda"s es bei der n"achsten
Anforderung sofort wieder vergeben wird. Wie schon gesagt, nimmt
der Algorithmus an, da"s keine Fehler vom Programmierer begangen
werden, wie z~B. die Freigabe eines bereits freigegebenen Arrays.

In meiner Implementation von SEL habe ich f"ur $n$ den Wert 400
benutzt, f"ur $s$ wie gesagt 20. Diese Werte sind f"ur die
meisten Anwendungen ausreichend, und wenn n"otig, k"onnten sie
vom Benutzer entsprechend ge"andert werden. Die
Speicheranforderungen dieses Algorithmus sind durchaus
vertretbar. Bei den angegebenen Werten erw"achst f"ur die beiden
Arrays unter Windows NT ein Speicherbedarf von genau 9600 Byte,
was bei den heute "ublichen Speicherkapazit"aten kaum ein Problem
darstellt.

\subsection{Ausdr"ucke}
\label{intexprs}

"Ahnlich wichtig wie die Typen sind in SEL die Ausdr"ucke, die
nat"urlich auch verwaltet werden m"ussen. Dabei werden wir jedoch
nicht ganze Ausdr"ucke verwenden, sondern uns lediglich auf die
Werte der Ausdr"ucke beschr"anken, die im Laufe des
Parsevorganges errechnet werden.

Um zu zeigen, wie wichtig die effiziente Verwaltung von
Ausdr"ucken ist, hier ein konkretes Beispiel: Ein
QuickSort-Programm f"ur 500 zuf"allig angeordnete Zahlen hat auf
meinem Computer eine Laufzeit von etwa 25 Sekunden, wobei w"ahrend
dieser Zeit ungef"ahr 60000 Ausdr"ucke allokiert und wieder
freigegeben werden m"ussen.

Zum Allokieren und Freigeben von Ausdr"ucken verwende ich aus
diesem Grund denselben Algorithmus wie f"ur Typen, wobei ich von
Haus aus maximal 200 Ausdr"ucke zur Verf"ugung stelle, was in den
meisten F"allen ausreichend ist.

Sehr wichtig f"ur eine effiziente Verarbeitung ist nat"urlich
auch der Aufbau eines Ausdrucks. Ich verwende hierzu eine
Struktur. Ein obligater Bestandteil eines Ausdrucks ist dessen
Typ, soda"s ein Typenarray Teil dieser Struktur ist. Weiters gibt
es je eine Variable f"ur die Basistypen \tchar , \tint\ und
\tfloat . Hierzu wird es aus Effizienzgr"unden am besten sein,
passende Typen der Implementationssprache zu verwenden. F"ur den
Fall, da"s der Ausdruck ein Array, ein Zeiger oder eine Struktur
ist, habe ich zus"atzlich einen typenlosen Zeiger aufgenommen,
der, falls die Implementationssprache keine typenlosen Zeiger
zul"a"st, durch einen Zeiger auf ein Byte ersetzt werden kann.
Im Falle eines Funktionsaufrufs ist es n"otig, mehrere Ausdr"ucke
zu verketten, um die Parameter der Funktion anzugeben, soda"s
zus"atzlich ein Zeiger auf einen Ausdruck in die Struktur
eingebaut werden sollte, soda"s sie sich zu einer verketteten
Liste erweitern l"a"st.

Soviel zur Verwaltung der Ausdr"ucke im Speicher. Eine weit
wichtigere Frage jedoch ist, was der Interpreter mit diesen
Ausdr"ucken tut. Nun, zun"achst einmal mu"s jede der in
Kapitel \ref{exprparser} besprochenen Produktionsfunktionen so
konzipiert werden, da"s sie der "ubergeordneten Funktion einen Ausdruck
mit dem Wert des von ihr geparsten Ausdrucks zur"uckgibt. Diese
verwendet ihn ihrerseits, um den Wert ihres Ausdrucks zu
berechnen u.~s.~w. In der Kette der Funktionsaufrufe gibt es dann
irgendwann eine Funktion, die den Wert des Ausdrucks benutzt, um
damit etwas zu entscheiden, wie z.~B., ob der Zweig einer
\rw{if}-Anweisung ausgef"uhrt werden soll.

Solange man lediglich Berechnungen mit den Ausdr"ucken
anstellt, ergeben sich mit der oben skizzierten Methode
keinerlei Probleme, wie ist es jedoch m"oglich, den
Zuweisungsoperator zu implementieren? Bauen wir die Funktion zum
Parsen eines \fu{primary} so auf, da"s sie f"ur eine Variable
lediglich deren Wert zur"uckgibt, werden wir Probleme haben, denn
dann wei"s der Interpreter nicht, wo diese Variable im Speicher liegt und
kann ihr so auch keinen Wert zuweisen. Um dieses Problem zu
"uberwinden, wird intern ein zus"atzlicher Typ,
n"amlich der Referenztyp, der einem Zeiger sehr "ahnlich ist,
verwendet.
Trifft die Funktion zum Parsen eines \fu{primary} nun im Zwischencode auf eine Variable des Typs
\tint , dann gibt sie nicht deren Wert zur"uck, sondern ihre
Adresse und vermerkt, da"s es sich dabei um eine Referenz auf
einen \tint\ handelt, der Einfachheit halber von mir mit
\type{\reference int} bezeichnet.
Mu"s nun mit diesem Ausdruck gerechnet werden, gen"ugt es, ihn zu
dereferenzieren, soda"s ein Ausdruck mit dem Typ \tint\ entsteht.
Wird ihm ein Wert zugewiesen, mu"s dieser lediglich an die
Speicheradresse geschrieben werden, die er angibt. Soll der
un"are Operator \fu{\&} auf ihn angewendet werden, dann ist es
lediglich n"otig, das erste Byte im Typenarray so zu "andern,
da"s daraus ein Zeiger wird.
                
\subsection{Variablen}
\label{maintvars}

Variablen sind von den oben besprochenen Ausdr"ucken und Typen
insofern verschieden, als da"s sie nicht intern vom Interpreter
verwendet werden, sondern direkt mit dem Quellprogramm in
Verbindung stehen. Aus dem Modulkonzept und der Tatsache, da"s
Funktionsvariablen nur in der jeweiligen Funktion verf"ugbar
sind, folgt, da"s die Verwaltung von Variablen auf mehrern Stufen
erfolgen mu"s. Die umfassendste Stufe, allerdings auch diejenige
mit der niedrigsten Priorit"at, ist die der programmglobalen
Variablen. Auf sie folgt die Stufe der modulglobalen Variablen.
Die h"ochste Priorit"at besitzt die Stufe der Funktionsvariablen,
die allerdings in der Regel auch die wenigsten Variablen
aufnimmt.

Wird nun im Programm eine Variable verwendet, mu"s zun"achst
gepr"uft werden, ob es eine Funktionsvariable ihres Namens gibt.
Ist das nicht der Fall, kann es sich nur mehr um eine globale
Variable handeln. Gibt es eine modulglobale Variable, die den
gesuchten Namen hat, wird sie verwendet, andernfalls mu"s es eine
programmglobale Variable geben. Ist auch das nicht zutreffend,
hat es der Interpreter mit einem Programmfehler zu tun (siehe Kapitel
\ref{errortrapping}).

Die komplizierteste Verwaltungsstufe ist die der
Funktionsvariablen (es ist am einfachsten, bei der Verwaltung der
Variablen nicht zwischen Parametern und direkten Funktionsvariablen
zu unterscheiden), was daran liegt, da"s nicht alle
Funktionsvariablen in einen "`Topf"' geworfen werden k"onnen,
sondern da"s f"ur jede Funktion eine eigene Variablenliste
angelegt werden mu"s. Die einfachste L"osung f"ur dieses Problem
ist, einen Stack zu verwenden, auf dem f"ur jeden
Funktionsaufruf eine Liste aller Funktionsvariablen angelegt
wird, "ahnlich dem Code, den ein Compiler erzeugen w"urde. Analog
dazu werden bei der R"uckkehr aus einer Funktion die vorher
erzeugten Variablen wieder entfernt. Um nun
nach einer Funktionsvariablen zu suchen, mu"s nur der oberste
Stacklevel in Betracht gezogen werden, da er alle Variablen
beinhaltet, die zu der gerade auszuf"uhrenden Funktion geh"oren.

Wie schon in Kapitel \ref{modules} angesprochen, stellt der
Interpreter f"ur jedes Modul eine eigene Datenstruktur zur
Verf"ugung, die unter anderem alle modulglobalen Variablen dieses
Moduls verwaltet. Ist also nach einer modulglobalen Variablen zu
suchen, kann man die Suche auf das durch \fu{MP} angegebene Modul
beschr"anken.

Der einfachste Fall liegt bei programmglobalen Variablen vor.
F"ur ihre Verwaltung gen"ugt es, eine einzige Liste
bereitzustellen, die die programmglobalen Variablen aller Module
bereitstellt, ohne Unterscheidung, in welchem Modul diese
definiert sind.

\bigskip

Bis jetzt habe ich, wenn es um die Details der
Variablenverwaltung ging, lediglich von "`Listen"' gesprochen,
wobei ich damit weniger dynamische Listen, als vielmehr
geeignete Datenstrukturen f"ur das schnelle Auffinden von
Variablen gemeint habe. Diese Datenstrukturen m"ochte ich jetzt
genauer besprechen.

Das Problem, auf das wir hier treffen ist, eine
Datenstruktur (und nat"urlich auch die dazugeh"origen
Algorithmen) aussuchen zu m"ussen, die es erm"oglicht, einen
Datensatz in einer Menge von Datens"atzen m"oglichst
schnell zu finden. Dieses Problem ist ein gutes Beispiel f"ur das
Grundproblem eines Programmdesigners: F"ur eine bestimmte Aufgabe
den besten Algorithmus bzw. die beste Datenstruktur (was in den
meisten F"allen Hand in Hand geht) auszuw"ahlen. G"abe es eine
absolut beste L"osung, dann h"atte der Designer seine Aufgabe
schnell erledigt, doch leider gibt es eine solche nur in den
seltensten F"allen.

Das wichtigste Entscheidungskriterium f"ur unsere Datenstruktur
ist, wieviel Datens"atze zu verwalten sind. F"ur die
Variablenverwaltung eines Interpreters bieten sich grunds"atzlich
zwei Strukturen an: Bin"are Suchb"aume und Hash-Tabellen. Um
entscheiden zu k"onnen, welche der beiden Strukturen f"ur unsere
Anwendungen die bessere ist, ist es n"otig, ein wenig zu rechnen.
Der Einfachheit halber gehen wir vom g"unstigsten Fall aus (die
Ergebnisse sind von denen des durchschnittlichen Falls nicht sehr
verschieden). Dieser ist, da"s der Benutzer die Variablennamen so
gew"ahlt hat, da"s sich eine gleichm"a"sige Anordung in der
Hash-Tabelle ergibt bzw. der Suchbaum ausgeglichen ist.

In einer Hash-Tabelle, bei der die Eintr"age f"ur gleiche
Hash-Werte linear verkettet sind, werden f"ur eine Suche nach
einem Datensatz im Mittel $\max (N / 2M + 1/2, 1)$ Vergleiche
ben"otigt, wohingegen f"ur die Suche in einem Baum etwa
$\log_2N - 1$ Vergleiche notwendig sind, wobei $N$ die Anzahl der
Datens"atze und $M$ die Gr"o"se der Hash-Tabelle angibt. Auf den
ersten Blick scheint der Baum die bessere L"osung zu sein, da die
Anzahl der ben"otigten Vergleiche f"ur gro"se $N$ in etwa
proportional zum bin"aren Logarithmus der Anzahl der Datens"atze
ist, bei einer Hash-Tabelle allerdings die Anzahl der Vergleiche
direkt proportional zur Anzahl der Datens"atze ist, was f"ur
gro"se $N$ nat"urlich bedeutend schlechter ist, als die
Leistungsdaten des Baums. Nehmen wir allerdings an, wir benutzen
eine Hash-Tabelle mit 37 Eintr"agen, was in Bezug auf den
Speicherplatz "uberhaupt kein Problem darstellt. In diesem Fall
sind f"ur eine Suche in einer Hash-Tabelle mit 100 Datens"atzen
im Durchschnitt weniger als zwei Vergleiche n"otig, wohingegen in
einem Baum mehr als sechs Vergleiche ben"otigt werden. Erst wenn
mehr als 550 Datens"atze vorhanden sind, ist die Suche in einem
Baum g"unstiger als in einer Hash-Tabelle.

Es ist kaum anzunehmen, da"s jemand f"ur einen Interpreter ein
Programm schreibt, welches mehr als 300 globale Variablen verwendet,
soda"s es f"ur unsere Zwecke die bei weitem bessere L"osung ist,
f"ur die Verwaltung programm- und modulglobaler Variablen
Hash-Tabellen zu verwenden, wobei eine Gr"o"se von 37 durchaus
praktikabel ist, sowohl in Bezug auf Schnelligkeit als auch auf
die vern"unftige Ausnutzung des Speicherplatzes. F"ur
Funktionsvariablen sind 37 Eintr"age wohl "uberdimensioniert.
Hierf"ur bieten sich die Primzahlen 7 und 11 an.

Zuletzt stellt sich noch die Frage nach der Hash-Funktion. Ich
verwende sowohl der Einfachheit als auch der Effizienz halber die
Modulo-Funktion, d. h., ich betrachte eine Zeichenkette als lange
Ganzzahl, deren Rest bez"uglich des Moduls 37 den Index in der
Hash-Tabelle angibt. Die folgende Implementation geht davon aus, 
da"s ein \tchar\ 256 verschiedene Werte annehmen kann.

\begin{verbatim}
function hash (s : *char, m : int) : int

  h : int = 0;

begin
  while (*s)
    h = (h * 256 + *s) % m;
    s = s + 1;
  end
  return h;
end
\end{verbatim}

\subsection{Funktionen}

Die Funktionsverwaltung hat in erster Linie die Aufgabe, eine
Funktion eines bestimmten Namens zu finden und deren Adresse im
Zwischencode zu ermitteln. Weiters ist es wichtig, Namen und
Typen der Funktionsparameter zu erfahren, um erstens "uberpr"ufen
zu k"onnen, ob die Funktion mit der richtigen Anzahl und den
richtigen Typen der Parameter aufgerufen worden ist, und zweitens,
um die Parameter auf dem Funktionsstack anlegen zu
k"onnen.\footnote{siehe Kapitel \ref{maintvars}}

Zus"atzlich sollte es m"oglich sein, nicht nur Funktionen, die in SEL
implementiert sind, aufzurufen, sondern auch solche, die in der
Implementationssprache des Interpreters geschrieben sind,
da eine Implementation der Grundfunktionen von SEL in
SEL selbst f"ur einen Interpreter viel zu ineffizient w"are.
Deshalb schreibt man Funktionen wie etwa zum Kopieren von
Strings in der Implementationssprache des Interpreters (in
meinem Falle C) und ruft diese auf. Um das nahtlos in die
Verwaltung einbinden zu k"onnen, ist es notwendig, da"s die
Implementationssprache Zeiger auf Funktionen unterst"utzt. Ist
diese M"oglichkeit nicht vorhanden (wie etwa im Falle von
\pascal ), dann mu"s jede einzelne Funktion im Programmcode extra
behandelt werden, was sowohl umst"andlicher als auch langsamer
ist.\footnote{In Turbo Pascal ab der Version 5.5 ergibt sich die
M"oglichkeit, statt eines Zeigers auf eine Funktion einen Zeiger
auf ein Objekt zu verwenden, welches eine Methode enth"alt, die
die gew"unschte Funktion ersetzt.}

Um eine Funktion eines bestimmten Namens zu finden, werden
Hash-Tabellen, wie im letzten Unterkapitel beschrieben,
verwendet: Eine f"ur programmglobale Funktionen und eine
pro Modul f"ur modulglobale Funktionen, der Einfachheit halber
jeweils mit 37 Eintr"agen, was bei einem Programm mit 10 Modulen
unter Windows NT einen Speicherplatz von 1628 Bytes beansprucht
(was ohne schlechtes Gewissen als "`unbedenklich"' eingestuft
werden kann). Eintr"age f"ur denselben Hash-Wert werden
unsortiert linear verkettet.\footnote{Eine Sortierung h"atte bei
einer linearen Verkettung auch nicht den geringsten Sinn, da
ohnedies alle Eintr"age bis zum gew"unschten durchlaufen werden
m"ussen, unabh"angig davon, in welcher Reihenfolge die Eintr"age
gespeichert sind.}

Die Datenstruktur, die eine einzelne Funktion beschreibt, mu"s
neben deren Namen und R"uckgabetyp auch eine Angabe dar"uber
enthalten, wo der Code der Funktion zu finden ist. Hierf"ur wird
zun"achst eine Variable ben"otigt, die Auskunft dar"uber gibt, ob
die Funktion in SEL geschrieben ist oder in der
Implementationssprache des Interpreters. Im ersten Fall mu"s die
Datenstruktur Auskunft dar"uber geben k"onnen, in welchem Modul
und an welcher Stelle der Code beginnt, andernfalls wird ein
Zeiger auf die jeweilige Funktion ben"otigt. Zus"atzlich m"ussen
Informationen "uber Namen und Typen der Parameter gespeichert
sein. Ich verwende hierzu eine einfach verkettete Liste.

\bigskip

Interessant ist die Realisierung von Funktionsaufrufen. Dazu ist
zun"achst zu bedenken, da"s der Parser rekursiv implementiert ist
und der Funktionsaufruf ein \fu{primary} ist, sich also auf der
untersten Grammatikebene der Ausdr"ucke befindet. Aus diesem
Grund wird man schwer umhinkommen, auch Funktionsaufrufe rekursiv
zu behandeln, was nat"urlich bei einem rekursiven SEL-Programm zu
einer starken Belastung des Prozessorstacks f"uhren kann, weshalb
dieser m"oglichst gro"s dimensioniert werden sollte. Mu"s nun
innerhalb eines \fu{primary} eine Funktion aufgerufen werden,
wird in den oben besprochenen Hash-Tabellen die Funktion gesucht
und, wenn vorhanden, wird an die angegebene Stelle im
Zwischencode gesprungen. Dort werden zun"achst der Funktionskopf
und die Deklaration der direkten Funktionsvariablen interpretiert
und anschlie"send der Code der Funktion mittels des
Anweisungs-Parsers ausgef"uhrt. Ist die Ausf"uhrung der Funktion
beendet, mu"s zur aufrufenden Funktion zur"uckgekehrt, also die
Werte der Variablen \fu{MP} und \fu{IP} restauriert werden.
Hierzu m"ussen klarerweise die alten Werte irgendwo
zwischengespeichert werden. Weiters m"ussen die lokalen Variablen
der Funktion abgelegt werden, und zwar so, da"s auch mehrere
Instanzen dieser Funktion existieren k"onnen und jede ihre
eigenen Variablen zur Verf"ugung hat. Wie schon im letzten
Kapitel angedeutet, wird hierzu ein Stack verwendet, in Anlehnung
an den Prozessorstack, der bei compilierten Programmen dieselben
Aufgaben wahrnimmt. Relevant ist jeweils nur der oberste Level
des Stacks, der die R"ucksprungadresse und die lokalen Variablen
der jeweils aktiven Funktion enth"alt. Wird aus dieser eine
andere Funktion aufgerufen, kommt ein Level zum Stack hinzu,
andernfalls wird ein Level entfernt. Um sich nicht auf die
relativ langsame dynamische Allokation verlassen zu m"ussen, ist
es am besten, diesen Stack als Array zu
konzipieren.\footnote{Eine Konzeption als dynamische Liste hat
auf den ersten Blick den Vorteil, da"s die Rekursionstiefe
lediglich vom vorhandenen Speicherplatz abh"angig ist, was jedoch
bei einem rekursiven Interpreter nicht zutrifft: dort ist der
Prozessorstack das Ma"s aller Dinge.}

\section{Fehlerbehandlung}
\label{errortrapping}

Da nicht davon auszugehen ist, da"s das Programm, das
interpretiert werden soll, fehlerfrei ist, mu"s der Interpreter
auch die M"oglichkeit haben, Fehler zu entdecken und sie
entsprechend zu behandeln. Im Gegensatz zu Compilern, die wenn
m"oglich nach dem Erkennen eines Fehlers mit dem "Ubersetzen
fortfahren sollten, um evtl. weitere Fehler zu entdecken, ist das
bei einem Interpreter nicht n"otig und auch wenig sinnvoll, da
das auszuf"uhrende Programm in seiner Stabilit"at gest"ort und
vermutlich abst"urzen w"urde. Eine sogenannte "`{\em Error
Recovery}"' bleibt uns also erspart. Die Ziele f"ur die
Fehlerbehandlung sind also das Erkennen von Fehlern und die
Ausgabe einer f"ur den Benutzer m"oglichst hilfreichen
Fehlermeldung.

Grunds"atzlich mu"s ein Interpreter zwei Arten von Fehlern
behandeln: syntaktische und semantische Fehler. Syntaktische
Fehler treten dann auf, wenn die im Anhang \ref{grammar}
angegebene Grammatik von SEL durch das Programm verletzt wird.
Diese Art von Fehler ist in der Regel leicht zu erkennen und kann
dem Benutzer relativ verst"andlich erkl"art werden. Ein Beispiel:
Der Interpreter ist gerade besch"aftigt, die Anweisung
\fu{while [x != 10)} auszuf"uhren. Der Parser erkennt das
Schl"usselwort \rw{while}, dem laut Grammatik immer
eine "offnende runde Klammer folgen mu"s. Offenbar hat sich der Benutzer
jedoch vertippt und statt der runden eine "offnende eckige
Klammer eingegeben. Es ist leicht, den Parser so zu erweitern,
da"s er an dieser Stelle den Fehler erkennt. Die entsprechende
Fehlermeldung sollte lauten: \fu{'(' expected}.

Semantische Fehler treten auf, wenn das Programm versucht,
unerlaubte Operationen auszuf"uhren, wie etwa das Casten eines
\tfloat\ in einen Zeiger, die Verwendung einer Variablen,
die gar nicht existiert, oder die Anwendung eines Operators auf
einen nicht erlaubten Typen. Tritt ein solcher Fall auf,
l"a"st sich das in der Regel ohnedies leicht feststellen und auch
eine passende Fehlermeldung ist nicht schwer zu schreiben.

F"ur den Benutzer des SEL-Interpreters ist es klarerweise nicht
nur wichtig, da"s ein Fehler aufgetreten und welcher Art
dieser ist, sondern auch, in welchem Modul in welcher Zeile der
Fehler steckt. Da der Interpreter weder syntaktische noch
semantische Fehler erkennen kann, w"ahrend er den Quellcode
untersucht, und der Zwischencode aus Effizienzgr"unden keine
Zeilenumbr"uche enth"alt, m"ussen die Zeilennummern gesondert
gespeichert werden. Ich verwende hierzu ein durch die
Zeilennummer indiziertes Array, das f"ur jede
Zeile eines Moduls die Anzahl der Bytes angibt, die diese Zeile
im Zwischencode beansprucht. Da der Zwischencode in der Regel
k"urzer als der Quellcode und kaum eine Programmzeile l"anger
als 100 Zeichen ist, gen"ugt es, pro Arrayelement ein Byte zur
Verf"ugung zu stellen. Tritt nun an der Stelle \fu{MP==3} und
\fu{IP==1364} ein Fehler auf, dann mu"s im Modul Nr. 3 das
Zeilennummernarray so lange durchlaufen werden, bis die Summe der
Elemente gr"o"ser als 1364 ist. Der folgenden Funktion, die die
gesuchte Zeilennummer zur"uckgibt, ist
das Zeilennummernarray als Zeiger und der Wert von \fu{IP} zu
"ubergeben:

\begin{verbatim}
function line (a : *char, ip : int) : int

  l : int = 0;
  c : int = 0;
  
begin
  repeat
    c = c + a[l];
    l = l + 1;
  until (c > ip)
  return l;
end
\end{verbatim}

In einem rekursiven Interpreter ist neben dem Erkennen und
Anzeigen eines Fehlers noch ein weiteres Problem zu "uberwinden:
das au"serordentliche Beenden des Interpreters. Es ist davon
auszugehen, da"s der Interpreter schon eine gewisse
Rekursionsebene erreicht hat, wenn der Fehler auftritt, soda"s es
nicht so einfach ist, ihn zu beenden und evtl. zu einer (den
Interpreter) aufrufenden Funktion zur"uckzukehren. Die
ANSI-Bibliothek von C bietet hierzu das praktische Konzept der
"`{\em Long-Jumps}"' an, das ungef"ahr einem \rw{goto}
gleichkommt, das nicht nur Spr"unge innerhalb einer Funktion
erlaubt, sondern auch aus einer Funktion in eine andere,
"ubergeordnete Funktion hinein. In C++ z.~B. ist es nicht einmal
mehr n"otig, sich auf Bibliotheksfunktionen zu st"utzen, da ein
"ahnliches Konzept (das sog. "`{\em Exception-Handling}"') in der
Programmiersprache selbst eingebaut ist. Leider besitzen nicht
alle Sprachen einen derartigen Luxus. In \pascal\ etwa mu"s man
sich einer anderen Methode bedienen. Eine einfache und
naheliegende L"osung ist, eine Variable zu definieren, die
angibt, ob bei der Ausf"uhrung des Programms bisher ein Fehler
aufgetreten ist. Eine Funktion, die einen Fehler entdeckt, hat
diesen zu melden und die Variable auf einen Wert zu setzen, der
anzeigt, da"s ein Fehler gefunden wurde. Kehrt diese Funktion zu
der aufrufenden Funktion zur"uck, mu"s diese "uberpr"ufen, ob die
Variable gesetzt ist und in diesem Fall zu ihrer aufrufenden
Funktion zur"uckkehren, so lange, bis man an der Funktion
angelangt ist, die die Ausf"uhrung des Programms initiiert hat.
Zwar ist diese Methode simpel, doch hat sie einen gewissen
Overhead zur Folge, der nicht nur den Code gr"o"ser, sondern auch
un"ubersichtlicher macht.
                                   
\newpage
\section{Weiterf\protect{"uhrende} Projekte}

Ebenso wie im ersten Kapitel m"ochte ich auch hier einige meiner
Ideen darlegen, die leider zu umfangreich sind, als da"s ich sie
ausf"uhrlich in dieser Arbeit h"atte behandeln k"onnen.

\begin{description}
\item[Verwendung einer Baumstruktur] anstatt des Zwischencodes
k"onnte die Geschwindigkeit des Interpreters um einiges erh"ohen.
Diese Baumstruktur w"urde in gewisser Weise einem
Parse-Tree\footnote{siehe Kapitel \ref{topdbotu}} "ahneln, bei
dem jedoch die unn"otigen Bl"atter entfernt sind. Zus"atzlich
k"onnte vor der Ausf"uhrung des Programms der Baum optimiert
werden, z.~B. durch Vereinfachung von konstanten Ausdr"ucken
(\fu{2*3+4} kann ohne Ver"anderung der Laufzeiteigenschaften zu
\fu{10} umgeformt werden). Eine Darstellung als Baum hat f"ur den
Interpreter den gro"sen Vorteil, da"s die lexikalische Analyse
und der Parsevorgang nur einmal ausgef"uhrt werden m"ussen. Ist die
Baumstruktur gut konzipiert, k"onnte der Parser sogar als
Front-End f"ur einen Compiler dienen.

\item[Ein Compiler f"ur SEL] ist mein n"achstes Ziel. Leider sind
jedoch die meisten Maschinensprachen entweder zu komplex (wie
etwa die des Intel 80386) oder zu umst"andlich (wie die des Intel
8086), als da"s es Spa"s machen k"onnte, f"ur diese Prozessoren
einen Compiler zu schreiben. Aus diesem Grund erw"age ich zwei
M"oglichkeiten, um diese Schwierigkeiten zu umgehen.
Die einfachere der beiden ist, ein SEL-Programm anstatt in
Maschinencode in eine andere h"ohere Programmiersprache,
vorzugsweise C, zu "ubersetzen. Besonders interessant w"are es,
diesen Compiler selbst in SEL zu schreiben und ihn dann durch sich
selbst "ubersetzen zu lassen (mit Hilfe des Interpreters).

Die zweite M"oglichkeit ist,
anstatt f"ur einen real existierenden Prozessor Maschinencode
f"ur einen einfachen, auf die Programmiersprache zugeschnittenen,
hypothetischen Prozessor zu erzeugen, wobei dieser Prozessor
nat"urlich zun"achst entworfen und auf einem existierenden
Prozessor emuliert werden m"u"ste. Obwohl auch hier ein Interpreter
im Spiel ist, ist die Ausf"uhrungsgeschwindigkeit eines
Pseudo-Code weit h"oher als die eines Interpreters, der den
Quellcode (oder auch einen Zwischencode) der Programmiersprache
interpretiert. Der Vorteil dieser beiden M"oglichkeiten ist, da"s
sie gr"o"stenteils prozessorunabh"angig sind. So existiert heute
beinahe f"ur jeden Prozessor eine ausreichende Menge an
C-Compilern und f"ur den P-Code m"u"ste lediglich der Interpreter
neu geschrieben, oder, falls auch er in einer h"oheren
Programmiersprache geschrieben ist, lediglich neu compiliert
werden. Aus Effizienzgr"unden ist es jedoch ratsam, einen
Interpreter f"ur P-Code in Assembler zu implementieren.

\item[Ein Pretty-Printer] ist ein Programm, das einen Quelltext
einer bestimmten Programmiersprache in eine ansprechende Form
bringt, also Einr"uckungen vornimmt, wenn n"otig, und Kommentare
einheitlich gestaltet. F"ur C existiert eine Reihe solcher
Programme, das bekannteste ist wohl C-Indent. Es ist eine
interessante, wenn auch nicht "uberm"a"sig schwere Aufgabe, einen
solchen Pretty-Printer f"ur SEL zu schreiben.
\end{description}

\begin{appendix}

\chapter{Die Grammatik von SEL}
\label{grammar}

\begin{verbatim}
module:        {block}
             ;
block:         function
             | local
             | global
             ;
function:      `function' IDENTIFIER `(' [formal_params] `)' `:' type
                 {func_var}
               `begin'
                 {statment}
               `end'
             ;
formal_params: IDENTIFIER `:' [`&']  type [`,' formal_params]
             ;
func_var:      IDENTIFIER `:' type [`=' expr] `;'
             ;
func_proto:    `function' IDENTIFIER
               `(' [[`&' type {`,' [`&'] type}] `)' `:' type
             ;
local:         `local' {local_decl} `end'
             ;
local_decl:    struct
             | func_proto
             | var
             ;
global:        `global' {global_decl} `end'
             ;
global_decl:   struct
             | func_proto
             | var
             | extern
             ;
\end{verbatim}
\newpage
\begin{verbatim}
statement:     [expr] `;'
             | if
             | while
             | for
             | repeat
             | switch
             | prefix
             | `break' `;'
             | `return' [expr] `;'
             ;
if:            `if' `(' expr `)' {statement} [`else' {statement}] `end'
             ;
switch:        `switch' `(' expr `)' {statement} `end'
             ;
prefix:        `case' constant `:'
             | `default' `:'
             ;
while:         `while' `(' expr `)' {statement} `end'
             ;
repeat:        `repeat' {statement} `until' `(' expr `)'
             ;
for:           `for' `(' primary `=' expr (`to' | `downto') expr
                         [`step' expr] `)'
                 {statement}
               `end'
             ;
expr:          expr (`=' | `&&' | `||' | `^^' | `==' | `<' | `>'
                   | `<=' | `>=' | `!=' | `&' | `|' | `^' | `<<'
                   | `>>' | `+' | `-' | `*' | `/' | `%') expr
             | (`!' | `~' | `-' | `*' | `&') expr
             | `delete' expr
             | primary
             ;
primary:       IDENTIFIER
             | `(' expr `)'
             | CHAR_CONSTANT
             | INT_CONSTANT
             | FLOAT_CONSTANT
             | STRING_CONSTANT
             | primary `->' IDENTIFIER
             | primary `:' type
             | primary `[' expr `]'
             | IDENTIFIER `(' [expr {`,' expr}] `)'
             | `sizeof' type
             | `new' type
             ;
var:           IDENTIFIER `:' type [`=' constant] `;'
             ;
extern:        `extern' IDENTIFIER `:' type `;'
             ;
\end{verbatim}
\newpage
\begin{verbatim}
struct:        `struct' IDENTIFIER {IDENTIFIER `:' type `;'} `end'
             ;
type:          `char'
             | `int'
             | `float'
             | `void'
             | `[' INT_CONSTANT `]' type
             | `*' type
             | IDENTIFIER
             ;
constant:      CHAR_CONSTANT
             | [`-'] INT_CONSTANT
             | [`-'] FLOAT_CONSTANT
             | STRING_CONSTANT
             ;
INT_CONSTANT:    [0-9]+
CHAR_CONTANT:    '(.|\.)'
FLOAT_CONSTANT:  [0-9]*(`.'[0-9]*)?((e|E)(+|-)?[0-9]+)?
STRING_CONSTANT: "([^"\n]|\.)*"
IDENTIFIER:      [A-Za-z_][A-Za-z0-9_]*
\end{verbatim}

\end{appendix}

\begin{thebibliography}{9}
\addcontentsline{toc}{chapter}{Literaturverzeichnis}
\bibitem{selkr} Mark Probst. {\em SEL-Kurzreferenz}, Version~0.3.
Beilage.
\bibitem{kandr} Brian W.~Kernighan, Dennis M.~Ritchie. {\em
Programmieren in C}, deutsche Ausgabe bearbeitet
von A.~T.~Schreiner, Ernst Janich. Carl Hanser Verlag, Wien,~1983.
ISBN~3-446-13878-1.
\bibitem{sedge} Robert Sedgewick. {\em Algorithmen}, deutsche
Ausgabe. Addison-Wesley~GmbH, Deutschland, 1.~Auflage~1991.
ISBN~3-89319-301-4.
\bibitem{duden} {\em Duden Informatik}, 2.~Auflage. Dudenverlag,
Mannheim,~1993. ISBN~3-411-05232-5.
\end{thebibliography}

% \newpage
% \addcontentsline{toc}{chapter}{Index}
% \printindex

\end{document}

\begin{document}

\title{Design einer Programmiersprache\\
       und theoretische Besprechung\\
       eines Interpreters anhand der\\
       Programmiersprache SEL}
\author{Mark Probst}

\maketitle

\tableofcontents

\chapter{Die Programmiersprache}

Dieses Kapitel hat das Ziel, das Design der Programmiersprache
SEL, wie in \cite{selkr} beschrieben, zu erkl"aren bzw. zu
begr�nden, und zwar in Hinblick auf den Bau eines Interpreters
und eines Compilers. Aus diesem Grund setze ich eine gewisse
Kenntnis sowohl einer strukturierten Programmiersprache,
insbesondere nat"urlich der Sprache SEL, als auch der
grundlegenden Funktionsweise eines Interpreters bzw. eines
Compilers beim Leser voraus. Der Leser sollte also zumindest
\cite{selkr} gelesen haben, da"s, obwohl es als Kurzreferenz
konzipiert ist, durchaus auch geeignet ist, einen "Uberblick
"uber die Programmiersprache SEL zu vermitteln.

\section{Der Begriff der Programmiersprache}

Grunds"atzlich kann eine Programmiersprache definiert werden als

\begin{quotation}
Sprache zur Formulierung von Rechenvorschriften, d. h. von
Datenstrukturen und Algorithmen, die von einem Computer
ausgef"uhrt werden k"onnen.\footnote{siehe \cite{duden},
Seite~544}
\end{quotation}

Die in dieser Arbeit behandelte Programmiersprache SEL geh"ort
zur Gruppe der imperativen Programmiersprachen, deren wichtigste
Untergruppe wohl die der strukturierten Programmiersprachen ist.
Imperative Programmiersprachen werden im Duden Informatik
wie folgt definiert:

\begin{quotation}
Bei diesen Sprachen besteht ein Programm aus einer Folge von
Befehlen an den Computer, wie z.~B. "`Schreibe in die Variable a
den Wert 3"', "`Springe an die Stelle x im
Programm"'.\footnote{siehe \cite{duden}, Seite~545}
\end{quotation}

Strukturierte Programmiersprachen stellen f"ur den Programmierer
zus"atzlich Hilfen in Form von Kontrollanweisungen, Schleifen und
benutzerdefinierten Funktionen zur Verf"ugung. Diese machen
sowohl das Schreiben der Programme einfacher und schneller wie
auch die Programmtexte leichter lesbar. Zur Zeit sind sie die am
meisten verwendete Art der Programmiersprachen.

\section{Ziele f\protect{"ur} das Design von SEL}

Bevor man an die Arbeit an einem so komplexen Projekt wie das Design
einer Programmiersprache herangeht, sollte man sich im Klaren
dar"uber sein, welche Ziele man eigentlich verfolgt. So habe ich
auch an meine zuk"unftige Programmiersprache einige Anforderungen
gestellt, von denen ich die meisten in der gegenw"artigen Version
der Sprache verwirklicht sehe.

Die wichtigste Entscheidung bei einem Projekt ist die nach dem
Zweck desselben. Im Fall von SEL wollte ich weniger eine
Programmiersprache f"ur das Schreiben von kommerzieller Software
schaffen, als vielmehr eine kleine, feine Sprache, die ich in
vielen Situationen benutzen k"onnte. Haupts"achlich schwebte mir
dabei die Verwendung als Makrosprache in diversen Programmen vor.
Eine solche Makrosprache sollte vor allem leicht erlernbar sein
und keine gro"se Umstellung von Seiten des Benutzers verlangen,
sollte aber andererseits nicht so primitiv aufgebaut sein, da"s
sie eine vern"unftige Verwendung nicht mehr m"oglich macht, was
die meisten Makrosprachen leider sind.

Um eine geringe Umstellungszeit m"oglich zu machen, mu"s die
Programmiersprache an eine bereits bestehende, weit verbreitete
Sprache angelehnt sein. Im Fall von SEL habe ich mich
haupts"achlich an C und \pascal\ gehalten, weil mir diese am
gel"aufigsten und vermutlich die am weitest verbreiteten
Sprachen sind.

Da ich ein gro"ser Freund von C und ihrer pr"agnanten und kurzen
Ausdrucksweise bin, habe ich mir zum Ziel gesetzt, auch SEL zu
einer effizienten, kurzen Ausdrucksweise zu verhelfen. Eine
Konsequenz aus dieser Entscheidung ist es, anstatt von
Schl"usselw"ortern, so wie in \pascal , Symbole wie in C zu
verwenden, soferne dies nicht die "Ubersichtlichkeit und
Konsistenz innerhalb der Sprache zu sehr gef"ahrdet. Die zweite
Konsequenz ist es, redundante Sprachteile zu entfernen, solange
auch das nicht die Lesbarkeit der Programmtexte verletzt.

Ein weiteres wichtiges Kriterium bei einer Makrosprache ist die
Erweiterungsf"ahigkeit, da Makrosprachen in so verschiedenen
Anwendungen wie z.~B. Textverarbeitungen und Zeichenprogrammen
eingesetzt werden, da"s es schlichtweg unm"oglich ist, eine
Sprache zu konstruieren, die f"ur alle diese Anwendungsgebiete
passend ist. Diese Erweiterungsf"ahigkeit sollte einerseits durch
die Erweiterung mit Hilfe von Funktionen wie z.~B.
f"ur das Zeichnen von Linien oder Kreisen und andererseits durch
die M"oglichkeit der "Anderung in der Grammatik, wie etwa
f"ur das flexiblere Handling von Vektoren und Matrizen, gegeben
sein. Einige Vorschl"age f"ur "Anderungen der letzteren Art sind
in Kapitel \ref{langext} gegeben.

Zuletzt soll die Programmiersprache SEL so aufgebaut sein, da"s
es nicht nur m"oglich ist, einen Interpreter, sondern auch mit
vertretbarem Aufwand einen Compiler zu implementieren. Nicht nur,
weil das bei einer Programmiersprache immer eine gute Idee ist,
sondern auch in Hinblick darauf, da"s ich selbst einmal einen
Compiler schreiben m"ochte.

Als Abschlu"s dieses Unterkapitels seien nochmals alle Forderungen
kurz zusammengefa"st, um in sp"ateren Abschnitten darauf
zur"uckgreifen zu k"onnen.

\begin{itemize}
\item Geringer Sprachumfang, jedoch auch f"ur mittelgro"se
Programme geeignet
\item Anlehnung an C und \pascal\
\item Verwendung von Symbolen anstatt von Schl"usselw"ortern
(wenn m"oglich)
\item Vermeidung von Redundanzen
\item Konsistenz innerhalb der Sprache
\item Erweiterungsf"ahigkeit
\item Sowohl f"ur Interpreter als auch f"ur Compiler geeignet
\end{itemize}

\section{Das Programm}

Das gr"o"ste in sich abgeschlossene Objekt, da"s ein Compiler bzw.
Interpreter zu behandeln hat, ist das Programm. Uns stellt sich nun die
Frage, wie dieses Programm aufgebaut sein soll.

Zun"achst ist zu
entscheiden, ob das Programm in Module eingeteilt werden soll oder
nicht. Eine solche Einteilung hat den Vorteil, da"s der Programmierer
Implementationsdetails eines Programmteils oder einer Funktionssammlung
vom Rest des Programms abschirmen kann und den Anwender dieser
Funktionssammlung gar nicht erst in Versuchung bringt, in die internen
Abl"aufe einzugreifen. Nat"urlich ist das nicht der einzige Vorteil
dieser Technik. So kann es z.~B. bei einer Modularisierung des Programms
auch kaum zu Namens"uberschneidungen bei Funktionen und Variablen
kommen. Last but not least wird auch das Programm leichter
"uberschaubar, weil man sich jeweils nur auf ein Modul konzentrieren
mu"s. Ich habe mich f"ur eine Modularisierung entschieden. Erstens, weil
es die M"oglichkeiten des Programmierers um einen nicht zu
untersch"atzenden Faktor erweitert und zweitens, weil es nicht schwer
ist, ein solches Modulkonzept in einen Interpreter oder einen Compiler
einzubauen---jedenfalls nicht das hier verwendete.

In SEL wird ein Modul in mehrere Bl"ocke eingeteilt, die in beliebiger
Reihenfolge auftreten d"urfen, wobei es drei Arten
von Bl"ocken zu unterscheiden gibt:

\begin{itemize}
\item Die \rw{global}-Bl"ocke werden verwendet, um programmglobale Variablen,
Funktionen oder Datentypen
zu deklarieren, also solche, die im ganzen Programm bekannt
sind.
\item Die \rw{local}-Bl"ocke werden analog zu \rw{global}-Bl"ocken
verwendet, um
modulglobale Variablen, Funktionen oder Datentypen zu deklarieren.
\item Funktionen werden "ahnlich wie in C und \pascal\ definiert.
\end{itemize}

Der Hauptgrund, warum ich mich f"ur diese Art der Modulstruktur
entschieden habe ist der, da"s es f"ur einen Interpreter weit
praktikabler ist, zur Initialisierung der Variablen vor dem Start eines
Programms einfach nach den Schl"usselw"ortern \rw{global} und
\rw{local} zu suchen, als zwischen den Funktionsdefinitionen. Dieses
Problem h"atte nat"urlich auch eine Programmstruktur wie die in \pascal ,
wo Variablen- und Typendeklarationen ihren festen Platz haben, gel"ost,
doch erscheint mir dieser Ansatz als zu starr, vermutlich, weil ich die
Freiheiten von C gewohnt bin.

Noch ein kleines Wort zur allgemeinen Struktur eines Moduls: Wie
bereits oben erw"ahnt, besteht ein Modul aus einer beliebigen
Abfolge von Bl"ocken, wobei es durchaus m"oglich ist, da"s
mehrere Bl"ocke ein und desselben Typs vorkommen. F"ur das
gesamte Modul gilt jedoch eine Regel: Ein Objekt ist erst ab der
Zeile bekannt, in der es deklariert wurde, das hei"st, da"s es
z.~B. nicht m"oglich ist, eine Variable in Zeile 154 zu
verwenden, wenn sie erst in Zeile 186 deklariert wird.
Diese Regelung macht es einem Compiler m"oglich, ein Programm in
einem Durchgang zu "ubersetzen ({\em One-Pass Compiler}). Eine
Ausnahme bilden die Strukturen (siehe Kapitel \ref{structs}).

\subsection{global-Bl\protect{"ocke}}

Ein \rw{global}-Block\footnote{siehe Anhang \ref{grammar} oder
\cite{selkr}, Seite 9} wird durch das Schl"usselwort \rw{global}
eingeleitet und durch \rw{end} beendet. Dazwischen befindet sich
eine beliebige Abfolge von Variablen-, Funktions- und
Strukturdeklarationen. Alle in diesem Abschnitt deklarierten
Objekte sind dazu bestimmt, auch in anderen Modulen des Programms
verwendet zu werden, was jedoch nicht unbedingt bedeutet, da"s
sie in diesen auch ohne weiteres verf"ugbar sind. Soll ein
programmglobales Objekt auch in einem anderen Modul verwendet
werden, so ist es n"otig, es auch dort zu deklarieren und---im
Falle von Variablen---als \rw{extern} zu bezeichnen. Der Grund
f"ur diese Vorgangsweise ist, da"s sich ein Compiler, der die
Deklarationen aller anderen Module verwenden kann, nicht
konstruieren l"a"st, ohne nicht ein eigenes Objektformat zu
benutzen, wie das z.~B. bei Turbo \pascal\ der Fall ist. Das
Standard-Objektformat von MS-DOS und Windows NT ist dazu nicht in
der Lage. Au"serdem w"are es erheblich schwerer, Module
zu schreiben, die gegenseitig voneinander abh"angig sind, soda"s
ich von dieser L"osung Abstand genommen habe.

Auffallend an den drei Arten der Deklarationen, die innerhalb
eines \rw{global}-Blocks vorkommen d"urfen, ist, da"s die
Variablendeklaration als einzige mit einem Semikolon
abgeschlossen wird, und das, obwohl es ohne Gef"ahrdung der
Eindeutigkeit m"oglich w"are, dieses wegzulassen. Der Grund
hierf"ur ist einerseits die Einhaltung der Konsistenz der
Sprache\footnote{siehe auch Kapitel \ref{funcs}} und andererseits
die Offenhaltung f"ur Erweiterungen der Sprache. So ist es im
Moment zwar lediglich m"oglich, eine Variable mit einer \tchar -,
\tint -, \tfloat - oder Stringkonstante zu initialisieren, jedoch
ist es durchaus denkbar, dies in zuk"unftigen Versionen auf
Arrays, Strukturen und Zeiger auszuweiten, wo die Eindeutigkeit
nicht mehr unbedingt gegeben sein mu"s.

Der Grund, warum Variablen, die in einem anderen Modul definiert
worden sind und in einem bestimmten Modul verwendet werden, als
\rw{extern} zu deklarieren sind und dies bei Strukturen und
Funktionen nicht n"otig ist, ist einfach: Bei einer
Strukturdeklaration wird dem Compiler bzw. Interpreter lediglich
"`mitgeteilt"', wie diese aufgebaut ist. Es wird in diesem Fall
kein vom Programm verwendetes Objekt im Speicher erzeugt. Bei
einer Variablendefinition hingegen wird dem Compiler sowohl der
Typ der Variable mitgeteilt als auch der Befehl, sie im Speicher
zu erzeugen. Ohne Unterscheidung zwischen Variablendefinition und
\rw{extern}-Deklaration w"u"ste der Compiler nicht, ob er nun
eine Variable im Speicher zu erzeugen h"atte oder sich lediglich
vormerken sollte, da"s diese Variable in einem anderen Modul
existiert und in diesem verwendet wird. Bei Funktionen ist es
kaum anders: Eine Deklaration innerhalb eines \rw{local}- oder
\rw{global}-Blocks sagt dem Compiler lediglich, da"s es
"`irgendwo anders"' eine solche Funktion gibt und da"s sie hier
verwendet wird. Eine Funktionsdeklaration innerhalb eines
\rw{global}-Blocks dient weiterhin dazu, zu kennzeichnen, da"s
diese Funktion, falls sie in diesem Modul implementiert ist, f"ur
die Verwendung in anderen Modulen bestimmt ist. Eine zweite
M"oglichkeit, dies zu erzielen, w"are, dem Funktionskopf bei der
Implementation noch ein spezielles Schl"usselwort hinzuzuf"ugen,
doch erscheint mir die in SEL verwendete L"osung als die bessere
und auf jeden Fall als die "ubersichtlichere.

Variablen d"urfen in \rw{local} bzw. \rw{global}-Bl"ocken
lediglich mit Konstanten, in
Funktionen hingegen auch mit Ausdr"ucken initialisiert werden.
Dies ist einerseits der Einfachheit halber so, andererseits, weil
es immer eine riskante Sache ist, Werte zu berechnen, noch bevor
die einzelnen Teile des Programms initialisiert sind. Dar"uber hinaus w"are
es notwendig, eine Regelung zu erstellen, die angibt, welche
Variablen als erste initialisiert werden und welches Modul einem
anderen vorangeht. Dieser Aufwand erscheint mir allerdings etwas
"ubertrieben, nicht zuletzt deshalb, weil eine solche
Initialisierung leicht in der \fu{main}-Funktion\footnote{siehe
Kapitel \ref{main}} nachgeholt werden kann.

\subsection{local-Bl\protect{"ocke}}

Ein \rw{local}-Block\footnote{siehe Anhang \ref{grammar} oder
\cite{selkr}, Seite 9} unterscheidet sich von einem
\rw{global}-Block, abgesehen davon, da"s er durch ein anderes
Schl"usselwort eingeleitet wird, dadurch, da"s eine
\rw{extern}-Deklaration nicht erlaubt ist. Dies ist auch ganz
verst"andlich, denn es ist ein Widerspruch in sich, eine
modulglobale Variable, die in einem anderen Modul definiert
ist, zu verwenden.

Die Tatsache, da"s Funktionen, die nicht explizit in einem
\rw{global}-Block aufscheinen, automatisch modulglobal sind,
verleitet zum Schlu"s, da"s Funktionsdeklarationen in
\rw{local}-Bl"ocken redundant sind. Es k"onnte jedoch ohne
weiteres die Notwendigkeit entstehen, da"s zwei Funktionen sich
gegenseitig aufrufen. In diesem Fall w"are jedoch ein
Aufruf illegal, da Funktionen erst ab dem Punkt ihrer
Deklaration im Modul bekannt sind. Dieser Konflikt ist sehr
einfach zu l"osen, indem man eine der Funktionen (oder der
"Ubersichtlichkeit halber beide) in einem \rw{local}-Block
deklariert und diese erst nach der anderen implementiert. Ein
Beispiel: Die Funktion \fu{A} und \fu{B} rufen sich gegenseitig
auf. Nun wird vor der Implementation die Funktion \fu{A} als
\rw{local} deklariert, soda"s sie im Rest des Moduls bekannt ist.
Nun wird die Funktion \fu{B} implementiert. Der Aufruf von \fu{A}
stellt kein Problem dar, der oben geschilderten Deklaration
wegen. Die Implementation von \fu{A} ist nun auch kein
Problem mehr, da an diesem Punkt schon beide Funktionen bekannt
sind.

\subsection{Funktionen}
\label{funcs}

Im Gegensatz zu \pascal\ wird in SEL, so wie in C, nicht zwischen
Funktionen und Prozeduren unterschieden. Dies hat---anders als es
zu erwarten w"are---fast nur Vorteile. Auf der einen Seite ist es
m"oglich, Prozeduren durch Funktionen mit dem R"uckgabetyp
\tvoid \footnote{siehe Kapitel \ref{elemtypes}} nachzubilden,
andererseits vereinfacht es den Bau eines
Interpreters bzw. eines Compilers. Der einzige Nachteil ist ein
geringer Verlust an Lesbarkeit.

Eine Funktion\footnote{siehe Kapitel \ref{exprfuncs} und
Anhang \ref{grammar} oder
\cite{selkr}, Seiten 7--9 und 17} gliedert sich augenscheinlich in
drei gro"se Teile:
Den Funktionskopf, die Definitionen der direkten
Funktionsvariablen\footnote{siehe Kapitel \ref{variables}}
und den Funktionsk"orper.

Der Funktionskopf besteht aus dem Schl"usselwort \rw{function},
dem Namen der Funktion, einer in Klammern eingeschlossenen Liste
der Parameter inklusive deren Typen und dem R"uckgabetyp. Der
Doppelpunkt vor dem Typ soll andeuten, da"s eine Funktion in
gewisser Weise einer Variablen eines bestimmten Typs "ahnelt.
Es sei angemerkt, da"s dieser Doppelpunkt ohne weiteres
ausgelassen werden k"onnte, ohne die Eindeutigkeit der Sprache zu
gef"ahrden.\footnote{Tats"achlich k"onnte sogar die schlie"sende
Klammer ausgelassen werden, doch diente das sicherlich nicht der
Lesbarkeit.}

Bei genauerer Betrachtung des Funktionskopfes und der
Funktionsdeklaration in einem \rw{global} bzw. \rw{local}-Block
f"allt auf, da"s in einem Funktionskopf die Namen der
Parameter aufscheinen, in den beiden letztgenannten Bl"ocken dies
jedoch nicht der Fall ist. Betrachten wir zur Begr"undung dieser
Entscheidung einen Interpreter und einen Compiler. Der
Interpreter kann alle Module des Programms auf einmal betrachten,
er kennt also die Namen der Parameter auch dann, wenn sie in
einem anderen Modul erw"ahnt sind als in dem, von dem aus die
Funktion aufgerufen wird. Ein Compiler jedoch besitzt diesen
Vorteil nicht: Er beschr"ankt sich jeweils auf ein einzelnes
Modul. Bei einem Funktionsaufruf ist aber das Einzige, was der
Compiler zu tun hat, zu "uberpr"ufen, ob die Parameter die
richtigen Typen haben, um sie dann in der richtigen Reihenfolge
auf dem Stack abzulegen. Wirkliche Bedeutung haben die Namen der
Parameter nur f"ur die eigentliche Funktion und insofern stellt
die Beschr"ankung auf die Typen keine Einschr"ankung dar, es ist
jedoch durchaus denkbar, da"s es der "Ubersichtlichkeit wegen in
einer zuk"unftigen Version von SEL m"oglich sein wird, auch die
Namen anzugeben, wobei diese Informationen aber vom Interpreter
bzw. Compiler ignoriert werden.

Dem Funktionskopf folgen die Definitionen der direkten
Funktionsvariablen.
Diese werden bei jedem Funktionsaufruf neu auf dem Stack angelegt
und belegen daher bei einer rekursiven Funktion in den
verschiedenen Instanzen nicht denselben Speicherbereich---erst
diese Tatsache macht Rekursionen so effektiv. Es ist m"oglich,
die Variablen mit Werten zu initialisieren, die sich aus Ausdr"ucken ergeben.
Dies stellt weder einen Compiler noch einen Interpreter vor
Probleme, soda"s ich hier nicht n"aher darauf eingehen m"ochte.
Das Semikolon am Ende jeder Variablendefinition w"are nicht
unbedingt notwendig, ergibt jedoch sowohl in Hinsicht auf die
"Ubersichtlichkeit als auch auf die Erweiterungsf"ahigkeit der
Sprache durchaus Sinn.

Der Funktionsk"orper wird durch das Schl"usselwort \rw{begin}
eingeleitet und durch \rw{end} beendet. Dazwischen befindet sich
eine beliebige Anzahl von Anweisungen, die bei Aufruf der
Funktion ausgef"uhrt werden.

\subsubsection{Die main-Funktion}
\label{main}

Jedes Programm mu"s einen Punkt haben, an dem es beginnt, soda"s
nicht unkontrolliert irgendwelche Funktionen ausgef"uhrt werden.
Hierf"ur gibt es grunds"atzlich zwei M"oglichkeiten. Entweder,
man f"uhrt bestimmte Konstrukte ein, um diesen Hauptteil des
Programms zu kennzeichnen, so wie es z.~B. in \pascal\
praktiziert wird, oder man verwendet eine bestimmte Funktion
dazu, wie in C. Ich habe mich f"ur letzteren
L"osungsansatz entschieden, weil er es erlaubt, das gesamte
Programm als eine Art Funktion zu sehen, die vom Anwender
aufgerufen werden kann. Es wird auch m"oglich, die
Kommandozeilenparameter als einfache Funktionsparameter zu
"ubergeben---wenn dies auch in der aktuellen Version von SEL
noch nicht eingef"uhrt ist. Weiters kann durch dieses Prinzip dem
Betriebssystem ein Wert zur"uckgegeben werden---der sogenannte
{\em Exit-Code}. In \pascal\ sind hierf"ur bestimmte Funktionen
bzw. Prozeduren notwendig.

Der Einfachheit halber und wegen der Verwandschaft mit C wird die
Hauptfunktion des Programms \fu{main} genannt. Sie nimmt keine
Parameter an und hat entweder den R"uckgabetyp \tvoid\ oder
\tint . Ist ersteres der Fall, wird immer der Exit-Code
\fu{0} zur"uckgegeben, andernfalls wird dieser durch die
\rw{return}-Anweisung bestimmt.

\section{Anweisungen}
\label{stmts}

Anweisungen kontrollieren das Geschehen innerhalb eines
Programms. Sie befehlen\footnote{daher die Bezeichnung
"`imperative Programmiersprache"'} dem Computer, was er tun soll.
In SEL wird grunds"atzlich zwischen drei Arten von Anweisungen
unterschieden: Kontrollanweisungen, Schleifen und Ausdr"ucken.
W"ahrend die ersten beiden in diesem Unterkapitel besprochen
werden, widmet sich das gesamte n"achste Unterkapitel den
Ausdr"ucken. Der Korrektheit wegen sei hier gesagt, da"s es
eigentlich noch eine vierte Art von Anweisungen gibt: Die {\em
Nullanweisung}, die jedoch, da sie keinen bestimmten Zweck
erf"ullt, au"ser vielleicht die Lesbarkeit des Programms zu
erh"ohen, f"ur unsere Betrachtungen nicht weiter wichtig ist.

Es ist auffallend, da"s SEL keine Anweisungen zur Verf"ugung
stellt, die in irgendeiner Weise mit der Hardware des Computers
kommunizieren k"onnen, wie z.~B. Anweisungen zum Einlesen von der
Tastatur, oder zur Ausgabe von Zeichen auf den Bildschirm oder in
eine Datei. Die Gr"unde hierf"ur sind vielf"altig. Erstens haben
wir es uns zum Ziel gesetzt, den Sprachumfang von SEL m"oglichst
gering zu halten und Ein/Ausgabeanweisungen alleine
w"urden---sollten sie einigerma"sen vern"unftig sein---die
Grammatik von SEL betr"achtlich aufbl"ahen. Weiters habe ich SEL
mit dem Ziel vor Augen entworfen, sie einmal als Makrosprache
einzusetzen. Makrosprachen allerdings m"ussen in vielen
verschiedenen Umgebungen sinnvoll einzusetzen sein, wie z.~B. in
einer Textverarbeitung oder in einem Zeichenprogramm, das unter
einer graphischen Benutzeroberfl"ache l"auft. Nun sind aber die
M"oglichkeiten zur Textausgabe in einer graphischen
Benutzeroberfl"ache grundverschieden zu denen von z.~B. DOS.
Ausgabeanweisungen zu entwerfen, die f"ur alle diese
M"oglichkeiten halbwegs geeignet sind, erscheint mir fast
unm"oglich. Deshalb ist es unumg"anglich, jede Implementation von
SEL mit verschiedenen Grundfunktionen auszustatten, die auf die
jeweilige Umgebung zugeschnitten sind. Eine Standardisierung
dieser Grundfunktionen ist immer eine gute Idee, doch sicher sind
diese Standardfunktionen nicht "uberall anwendbar.\footnote{Ein
Beispiel hierf"ur w"are der ANSI-Standard von C, dessen
Ein/Ausgabefunktionen auf eine textorientierte Oberfl"ache
zugeschnitten sind. Die \fu{printf}-Funktion z.~B. ist unter der
graphischen Benutzeroberfl"ache Windows nutzlos.}

\bigskip

An dieser Stelle scheint es mir wichtig, ein Konzept aus anderen
strukturierten Programmiersprachen zu erw"ahnen. Das
Blockkonzept\footnote{nicht zu verwechseln mit den Bl"ocken in
SEL} erm"oglicht es, eine Anweisung entweder als eben nur eine
Anweisung oder als eine von zwei Schl"usselw"ortern oder Symbolen
begrenzte, beliebige Abfolge von Anweisungen zu sehen. So verlangt
eine \rw{if}-Anweisung in C oder \pascal\ nur eine Anweisung als
Zweig, es ist jedoch kein Problem, mit Hilfe der Blockstruktur
auch mehrere Anweisungen ausf"uhren zu lassen. Dieses
Blockkonzept ist in manchen Situationen sehr praktisch, aber es
bl"aht die Grammatik einer Programmiersprache auch etwas auf und
macht es so nat"urlich schwerer, einen Interpreter oder Compiler
zu schreiben. Aus diesem Grund habe ich mich entschieden, es in
SEL nicht einzubauen. Die Quelltexte sind so zwar etwas l"anger,
doch es entsteht keine wirkliche Beeintr"achtigung.

\subsection{Kontrollanweisungen}

Kontrollanweisungen haben die Aufgabe, zu kontrollieren, ob und
was ausgef"uhrt wird. In SEL gibt es drei Anweisungen dieser Art.
Die \rw{if}-Anweisung "uberpr"uft eine Aussage und f"uhrt, wenn
sie einen positiven Wahrheitswert ergibt, bestimmte Anweisungen
aus. Auf Wunsch werden auch Anweisungen ausgef"uhrt, wenn das
Ergebnis der "Uberpr"ufung negativ ist. Die \rw{switch}-Anweisung
f"uhrt f"ur bestimmte Werte eines Ausdrucks verschiedene
Anweisungen aus. Die \rw{return}-Anweisung springt zur
aufrufenden Funktion zur"uck.

\subsubsection{Die if-Anweisung}

Die \rw{if}-Anweisung\footnote{siehe Anhang \ref{grammar} oder
\cite{selkr}, Seite 10} wird durch das Schl"usselwort \rw{if}
eingeleitet, auf welches ein in Klammern eingeschlossener Ausdruck
folgt. Die Alternative zu dieser L"osung w"are gewesen, den
Ausdruck sofort auf das \rw{if} folgen zu lassen und durch ein
bestimmtes Schl"usselwort (z.~B. \rw{then}) oder ein Symbol zu
beenden. Der Doppelpunkt jedoch w"urde Mehrdeutigkeiten
hervorrufen und ein Semikolon erscheint mir als unpassend, und da
wir uns entschienden haben, Symbole zu verwenden, wo nicht
unbedingt Schl"usselw"orter erforderlich sind, f"allt auch das
\rw{then} weg. Die einfachste und unauff"alligste L"osung in
diesem Fall ist die Klammerung.

Auf die schlie"sende Klammer folgt eine beliebige Anzahl von
Anweisungen, beendet entweder durch \rw{end} oder \rw{else}. Da
auch \rw{if} eine Anweisung ist, ist es m"oglich, \rw{if}s zu
verschachteln. Im zweiten Kapitel wird sich zeigen, da"s es f"ur
einen Interpreter bzw. Compiler kein Problem ist, mit einer
solchen Schachtelung fertigzuwerden. Folgt auf die Anweisungen ein
\rw{end}, so ist die \rw{if}-Anweisung beendet. Es w"are
m"oglich, auch das \rw{end} durch ein Symbol zu ersetzen,
doch alle sinnvollen Symbole, bis auf die Klammern, sind bereits
f"ur andere Zwecke vergeben. Nat"urlich k"onnte man auch eine
schlie"sende Klammer benutzen, doch um ein solches Konstrukt
sinnvoll zu machen, ben"otigt man auch eine "offnende Klammer.
Eine runde Klammer w"urde den falschen Eindruck erwecken, eckige
Klammern sind bereits vergeben, es bleiben die geschwungenen
Klammern, die auch in C zu diesem Zweck verwendet werden. Da ich
aber auch in allen anderen F"allen einem \rw{end} bzw. einem
\rw{begin}-\rw{end} Paar den Vorrang gegeben habe, f"uhrte eine
Verwendung hier h"ochstens zu einer Verwirrung des Benutzers.

Folgt auf die Anweisungen ein \rw{else}, dann folgen auf
dieses weitere Anweisungen, durch ein \rw{end} begrenzt.
Nat"urlich k"onnte man auch anstelle des \rw{else} ein Symbol
verwenden, doch das w"urde den Zweck der Konstruktion zu sehr
verschleiern:

\begin{verbatim}
if x == 1
{
  do_what_you_want_to();
::
  do_something_else();
}
\end{verbatim}

Ich glaube, bei Betrachtung dieses Beispiels zeigt sich, warum
ich die Schl"usselw"orter den Symbolen vorgezogen habe.

Manche Programmiersprachen, wie z.~B. \modula , stellen
zus"atzlich zum normalen \rw{else} ein sogenannte \rw{elseif}
oder auch \rw{elsif} bereit. Es erlaubt, sollte die
\rw{if}-Bedingung negativ sein, noch eine zweite Bedingung zu
"uberpr"ufen und, sollte auch diese logisch falsch sein, eine
dritte u. s. w. Ich habe mich entschieden, ein solches
\rw{elseif} nicht in SEL einzubauen, aus zwei Gr"unden.
Erstens l"a"st sich diese \rw{if}-\rw{elseif} Kombination in
vielen F"allen durch ein besser passendes
\rw{switch}\footnote{siehe Kapitel \ref{switch}} ersetzen und
zweitens ist es m"oglich, den \rw{elseif}-Effekt durch ein
verschachteltes \rw{if} zu erzeugen---\rw{elseif} selbst ist
nichts anderes als eine Vereinfachung dieser Konstruktion.

\subsubsection{Die switch-Anweisung}
\label{switch}

Die \rw{switch}-Anweisung\footnote{siehe Anhang \ref{grammar} oder
\cite{selkr}, Seite 11} erlaubt es, abh"angig vom Wert eines
Ausdrucks verschiedene Anweisungen auszuf"uhren. Dazu wird
zun"achst der Ausdruck evaluiert und ein \rw{case} gesucht,
dessen Wert mit dem des Ausdrucks "aquivalent ist. Ist kein
passendes \rw{case} vorhanden, wird zu einem evtl. vorhandenen
\rw{default} gesprungen. Dort wird mit der Ausf"uhrung
fortgefahren und erst bei einem \rw{break} oder dem zum
\rw{switch} geh"orenden \rw{end} beendet, also nach dem \rw{end}
fortgesetzt. Ist erst einmal ein bestimmter \rw{case}- oder
\rw{default}-Zweig eingeschlagen, werden alle weiteren
\rw{case}'s und \rw{default}s ignoriert. Das erm"oglicht es, f"ur
mehrere Werte ein und dieselben Anweisungen auff"uhren zu lassen,
indem man die einzelnen \rw{case}s hintereinander anreiht.

Die \rw{case}- und \rw{default}-Schl"usselw"orter werden in SEL,
ebenso wie in C, als Pr"afixe f"ur Anweisungen behandelt: bei
normaler Ausf"uhrung werden sie einfach ignoriert. Sie sind also
im Grunde nichts anderes als eine Art Sprunglabel, so wie sie in
vielen Programmiersprachen f"ur die
\rw{goto}-Anweisung\footnote{siehe Kapitel \ref{goto}}
verwendet werden. Eine andere---auf den ersten Blick
nat"urlichere---L"osung w"are, auf die \rw{switch}-Anweisung eine
Liste von \rw{case}'s oder \rw{default}s folgen zu lassen, die
ihrerseits jeweils von Anweisungen gefolgt sind. Um es m"oglich
zu machen, da"s eine Anweisungsfolge f"ur mehrere verschiedene
Werte ausgef"uhrt wird, m"u"ste man statt einem \rw{case} jeweils
mehrere \rw{case}'s erwarten. Welche der beiden L"osungen man
bevorzugt, ist Geschmackssache. Ich habe mich f"ur die
Pr"afix-L"osung entschieden, weil ich sie erstens von C gewohnt
bin und sie zweitens ein wenig leistungsf"ahiger ist.

Das Schl"usselwort \rw{case} vor dem jeweiligen Wert ist
notwendig, um Zweideutigkeiten in der Sprache zu vermeiden. Um
dieses Schl"usselwort auslassen zu k"onnen, m"u"sten wir die oben
besprochene zweite L"osung w"ahlen, jedoch m"u"sten die einzelnen
Anweisungsfolgen wiederum durch ein anderes Schl"usselwort, z.~B.
\rw{end} beendet werden, soda"s sich keine wirkliche Ersparnis
erg"abe.

Auch daf"ur, da"s in \rw{case}-Zweigen nur Konstanten erlaubt
sind, gibt es Gr"unde. Erstens h"atten wir durch den Doppelpunkt
Zweideutigkeiten---er w"urde als Typecast\footnote{siehe Kapitel
\ref{casts}} interpretiert werden---doch dieses Problem ist durch
Wahl eines anderen Symbols ohne weiteres zu "ubergehen. Der weit
gewichtigere Grund ist, da"s die f"ur die \rw{switch}-Anweisung
typische Effizienz verloren ginge: Bis der Interpreter zum
richtigen Zweig gelangt, m"u"ste er eine Vielzahl von Ausdr"ucken
auswerten, was sehr viel mehr Zeit in Anspruch nimmt, als der
simple Vergleich mit einer Konstante. Au"serdem w"urde diese
Entscheidung das Schreiben eines Compilers etwas verkomplizieren.
In F"allen, in denen eine solche \rw{switch}-Anweisung n"otig
w"are, sollte man auf \rw{if}-\rw{else}-Verschachtelungen umsteigen.

\subsubsection{Die return-Anweisung}

Die \rw{return}-Anweisung hat die Aufgabe, zur aufrufenden
Funktion zur"uckzukehren und evtl. einen Wert zu "ubermitteln,
den sog. {\em Return-Code}. Eine andere M"oglichkeit w"are, den
return-code durch eine Zuweisung an eine Variable, die den Namen
der Funktion hat, zu bestimmen und mittels einer anderen
Anweisung bei Bedarf zur aufrufenden Funktion zur"uckzuspringen.
Da man im Normalfall beides an der gleichen Stelle tut, habe ich
mich f"ur die \rw{return}-L"osung entschieden.

Die einfachste Methode w"are jedoch, "uberhaupt keine Anweisung
bereitzustellen, mit der in die aufrufende Funktion
zur"uckgekehrt werden kann und den R"uckgabewert am Ende der
Funktion festzulegen. F"ur einen rekursiven Interpreter ist dies
die beste L"osung, doch wie wir im n"achsten Kapitel sehen
werden, stellt uns auch der \rw{return}-Ansatz vor keine gro"sen
Probleme.

\subsection{Schleifen}

Schleifen kontrollieren nicht, was, sondern wie oft etwas
ausgef"uhrt wird. Die drei Arten der Schleifen, die in SEL
erlaubt sind, entsprechen den in den meisten strukturierten
Programmiersprachen vorhandenen Schleifen.

\subsubsection{Die repeat-Schleife}

Die \rw{repeat}-Schleife in SEL entspricht der
\rw{repeat}-Schleife in \pascal , bis auf das Detail, da"s die
\rw{until}-Bedingung in Klammern zu setzen ist und kein Semikolon
danach erforderlich ist. Diese Konzeption deckt sich mit den
restlichen Anweisungen von SEL und wurde daher von mir bevorzugt.

In manchen Programmiersprachen, wie z.~B. C, existiert die
\rw{repeat}-Schleife in einer abgewandelten Form, n"amlich der
der \rw{do-while}-Schleife. Der einzige Unterschied
besteht darin, da"s die \rw{repeat}-Schleife eine
Abbruchbedingung erwartet, die \rw{do-while}-Schleife jedoch eine
Fortf"uhrungsbedingung, das hei"st, eine Schleife kann durch die
andere ersetzt werden, wenn die Bedingung negiert wird. Welche
Schleife die bessere ist, kommt auf die Situation an, wirkliche
Vor- oder Nachteile gibt es nicht, die Auswahl ist also reine
Geschmackssache.

\subsubsection{Die while-Schleife}

Die \rw{while}-Schleife ist eng verwandt mit der
\rw{repeat}-Schleife. Der einzige Unterschied besteht darin, da"s
die \rw{repeat}-Schleife mindestens einmal durchlaufen wird, die
\rw{while}-Schleife jedoch nicht unbedingt. Syntaktisch gesehen
ist die \rw{while}-Schleife sehr "ahnlich der von C, die
Unterschiede sind in der in Kapitel \ref{stmts} angesprochenen
Blockstruktur von C begr"undet. So folgt auf die schlie"sende
Klammer der Schleifenbedingung in C in der Regel eine "offnende
geschwungene Klammer, in SEL jedoch ist diese ausgespart, da
nicht notwendig. Die schlie"sende geschwungene Klammer, die in C
die Schleife beendet, ist mit dem Schl"usselwort \rw{end} in SEL
zu vergleichen. Die Gr"unde, warum ich die Konstruktion mittels
Schl"usselw"ortern vorgezogen habe, ist bereits in vorhergehenden
Kapiteln er"ortert, ich gehe deshalb hier nicht n"aher darauf
ein.

\subsubsection{Die for-Schleife}

Die \rw{for}-Schleife in SEL erlaubt es, Anweisungen in einer
Anzahl, die schon im voraus bekannt ist, zu wiederholen.
Damit unterscheidet sie sich signifikant von den \rw{while} und
der \rw{repeat}-Schleifen, bei denen die Anzahl der Wiederholungen
nicht unbedingt bekannt sind. Das verhilft der
\rw{for}-Schleife---zumindest auf einem Interpreter---zu einer
sehr schnellen Ausf"uhrung, da die Bedingungen f"ur die Schleife
nur einmal berechnet werden m"ussen, w"ahrenddessen dies bei den
beiden anderen Schleifenarten bei jedem Durchlauf geschehen mu"s.

Dieses eine Merkmal der \rw{for}-Schleife ist vielen
Programmiersprachen gemeinsam, doch eine bekannte Sprache macht
eine Ausnahme: C. In C werden nicht eine Z"ahlervariable und
deren untere und obere Grenzen und evtl. noch die Schrittweite
angegeben, sondern es werden drei Ausdr"ucke verwendet: Der erste
wird lediglich beim Start der Schleife evaluiert, der dritte nach
jedem Durchlauf. Der zweite Ausdruck schlie"slich gibt die
Bedingung an, bei der die Schleife nochmals ausgef"uhrt werden
soll---ist sie logisch falsch, wird die Schleife abgebrochen. In
C ist also eine \rw{for}-Schleife nichts anderes als eine
erweiterte \rw{while}-Schleife und hat den Vorteil, da"s mit
ihrer Hilfe auch verkettete Listen und "ahnliches leicht
durchlaufen werden k"onnen. Auch ein "`klassische"'
\rw{for}-Schleife l"a"st sich so leicht nachbilden, doch diese
Art der Schleifenkonstruktion hat f"ur einen Interpreter den
gro"sen Nachteil, da"s bei jedem Schleifendurchlauf gleich zwei
Ausdr"ucke evaluiert werden m"ussen.

Ungl"ucklicherweise habe ich diesen Nachteil nicht sofort bemerkt und
f"uhrte auch in SEL die von C gewohnte \rw{for}-Konstruktion ein. Erst
sp"ater merkte ich, um wie viel effizienter diese Schleife f"ur einen
Interpreter w"are, w"urde ich sie "ahnlich wie in \pascal\ aufbauen. Die
"Anderungen am Interpreter und an der Dokumentation lie"sen sich zum
Gl"uck in weniger als einer Stunde durchf"uhren, soda"s Verz"ogerungen
f"ur den Zeitplan dieser Arbeit ausblieben und ich zur Erkenntnis
gelangte, da"s man aus Fehlern doch am besten lernt.

Noch ein kleines Wort zur Syntax: Die Bedingungen der Schleife sind in
Klammern gesetzt, weil das in SEL auch bei anderen Anweisungen der Fall
ist und dadurch die Konsistenz gewahrt wird, und die Grammatik ohne
Begrenzung nicht eindeutig w"are. Die Z"ahlervariable wird als
\fu{primary} angegeben, weil es auch m"oglich sein sollte, Array- und
Strukturelemente als Z"ahler zu benutzen, und nicht als \fu{expr}, weil
dann der folgende Zuweisungsoperator als ein Teil dieses Ausdrucks
interpretiert w"urde. Sollte es wirklich notwendig sein, einen \fu{expr}
zu benutzen, so kann dieser in Klammer gesetzt werden, und entspricht
damit einem \fu{primary}.\footnote{siehe Anhang \ref{grammar}} Anstatt
der Schl"usselw"orter \rw{to} bzw. \rw{downto} Symbole zu verwenden,
w"are erstens un"ubersichtlich und w"urde zweitens---au"ser bei
Verwendung neuer Symbole---die Eindeutigkeit zunichte machen, weshalb
ich die oben angesprochenen Schl"usselw"orter verwendete. Gleiches gilt
f"ur \rw{step}.

\subsection{Die goto-Anweisung}
\label{goto}

Die "Uberschrift dieses Absatzes mag darauf deuten, da"s es in SEL eine
\rw{goto}-Anweisung gibt, doch der Schein tr"ugt: Aus Gr"unden, die in
Kapitel \ref{parser} aufgef"uhrt sind und aus dem Grund, da"s diese
Anweisung in dem meisten F"allen ohnedies nicht notwendig ist und nur zu
einem schlechten Programmierstil verf"uhrt, habe ich mich gegen eine
Einbindung in SEL entschieden.    

\section{Ausdr\protect{"ucke}}

Ausdr"ucke sind ein Konzept, das in fast allen strukturierten
Programmiersprachen verwendet wird und das nicht ohne Grund: Es
vereinfacht die Grammatik der Sprache und macht damit das Schreiben
eines Interpreters bzw. eines Compilers einfacher. Au"serdem erleichtert
es dem Anwender der Sprache das Schreiben von Programmen sehr und
unterst"utzt die Lesbarkeit. So k"onnen Ausdr"ucke in vielen
verschiedenen Situationen verwendet werden: Als Bedingungen in
Schleifen, als Funktionsaufrufe oder als die rechten Seiten von
Zuweisungen.\footnote{Tats"achlich ist in SEL sogar die Zuweisung ein
Ausdruck, im Gegensatz zu \pascal\ oder \modula , doch dazu sp"ater
mehr.}

Ein Ausdruck in SEL kann "ahnlich einem Ausdruck im
mathematischen Sinn aufgefa"st werden, er ist eine Verkn"upfung
von Operanden durch Operatoren. Jeder Operand f"ur sich ist
selbst ein Ausdruck\footnote{bis auf die Operanden der Operatoren
\fu{new} und \fu{sizeof}, die Typen sind.} und das Ergebnis der
Verkn"upfung mit einem Operator ist wieder ein Ausdruck. So ist
es in SEL m"oglich, Ausdr"ucke beliebig tief zu schachteln. Diese
Eigenschaften der Ausdr"ucke legen es nahe, einen Ausdruck als
Baum\footnote{siehe \cite{sedge}, Seiten 57--72} darzustellen.
Dabei stellt jeder innere Knoten sowie die Wurzel einen Operator
dar und jedes Blatt einen Wert. Betrachtet man nun jeden
einzelnen Knoten (und dessen Nachfolger), so gelangt man zum
Ergebnis, da"s dieser auch wieder ein Ausdruck f"ur sich ist. Ein
Beispiel f"ur diese Darstellung gibt Sedgewick in \cite{sedge},
Seite 63, Abbildung 4.4.

Jeder Ausdruck in SEL hat nicht nur einen Wert, sondern auch
einen Typ\footnote{siehe auch Kapitel \ref{types}}, mit einer Ausnahme:
Ein Ausdruck des Typs \tvoid\ hat keinen Wert, was auch gar keinen Sinn
h"atte.

\subsection{Operatoren}

Grunds"atzlich gibt es vier Arten, in SEL Operatoren zu
klassifizieren\footnote{genauer gesagt vier "`vern"unftige"' Arten}:
nach der Anzahl ihrer Argumente, nach ihrer funktionellen
Zugeh"origkeit, nach ihrer Priorit"at und nach ihrem Ergebnistyp.

\paragraph{Klassifikation nach der Anzahl der Argumente}

Auf den ersten Blick gesehen gibt es in SEL nach dem Gesichtspunkt der
Zahl der Argumente zwei Arten von Operatoren: un"are und bin"are. Zu den
un"aren z"ahlen z.~B. die Operatoren \fu{+} und \fu{<<}, zu den bin"aren
z"ahlen \fu{!} und \fu{-}. Der letztgenannte Operator existiert auch in
einer bin"aren Version, man k"onnte also sagen, da"s es Operatoren gibt,
die sowohl un"ar als auch bin"ar sind. Ich pers"onlich betrachte diese
beiden "`Versionen"' als zwei verschiedene Operatoren, da auch ihre
Funktionen verschieden sind. Betrachtet man Ausdr"ucke in SEL genauer,
wird man bemerken, da"s ein
Funktionsaufruf durchaus mehrere Argumente erwarten kann, jedoch stellt
sich hier die Frage, ob man einen Funktionsaufruf (besser: eine Funktion) als
Operator auffassen sollte. Ich bin der Meinung, eine Funktion ist
weniger ein Operator als ein Wert, und habe ihr deshalb ein eigenes
Unterkapitel gewidmet.\footnote{siehe Kapitel \ref{exprfuncs}}

\paragraph{Klassifikation nach der funktionellen Zugeh"origkeit}

Auf diese Weise kann man die Operatoren von SEL in mehrere
Gruppen einteilen, wobei die Grenzen eher flie"send sind. Mir
erscheint folgende Einteilung am vern"unftigsten: arithmetische
Operatoren wie \fu{+} oder \fu{<<}, logische Operatoren wie
\fu{\&\&} oder \fu{!}, Vergleichsoperatoren wie \fu{>=} oder
\fu{==}, Operatoren zur Speicherverwaltung wie \fu{new} oder
\fu{[]} und als zwei eigene Gruppen der Zuweisungsoperator \fu{=}
und der Typecasting-Operator \fu{:}. Wie gesagt sind die
"Uberg"ange zwischen diesen Gruppen sehr flie"send. So k"onnte
man beispielsweise meinen, der Operator \fu{[]} geh"ore zur
Gruppe der arithmetischen Operatoren, da er ja in gewisser Weise
Vektorarithmetik m"oglich macht. Auch w"are es m"oglich, den
Typecasting-Operator zur Speicherverwaltung z"ahlen, da er auch
den Speicher modifiziert.

\paragraph{Klassifikation nach der Priorit"at}

Operatoren in SEL besitzen verschiedene
Priorit"atsstufen\footnote{siehe \cite{selkr}, Seite 21}, um eine
"uberm"a"sige Klammerung zu vermeiden. Die Priorit"aten in SEL sind
stark an die von C angelehnt\footnote{siehe \cite{kandr}, Seite 54},
weichen jedoch teilweise von diesen ab. Meistens im selben Atemzug mit
der Priorit"at der Operatoren wird die Richtung der Auswertung genannt.
Diese gibt an, ob der Operator links- oder rechts-assoziativ ist. Ein
Beispiel f"ur einen rechts-assoziativen Operator ist \fu{=}: Der
Ausdruck \fu{a=b=0} ist "aquivalent zu \fu{a=(b=0)}. So ist es ohne
Klammerung m"oglich, mehreren Variablen in einer Anweisung denselben
Wert zuzuweisen. Ein
links-assoziativer Operator ist das bin"are \fu{-}. W"are er das nicht,
w"urde der Ausdruck \fu{5-3-1} nicht zu \fu{(5-3)-1}, sondern zu
\fu{5-(3-1)} "aquivalent sein und so das falsche Ergebnis, n"amlich
\fu{3} anstatt \fu{1}, liefern.

\paragraph{Klassifikation nach dem Ergebnistyp}

Diese Einteilung ist wohl die komplizierteste und es w"are wohl besser,
man w"urde als Kriterium nicht den Ergebnistyp selbst, sondern das, 
was den Ergebnistyp bestimmt, heranziehen. 
So sind dies etwa beim Operator \fu{+} die
Typen seiner Operanden, beim Operator \rw{new} hingegen ist es der Typ,
mit dem er verbunden wird. Weiters gibt es Operatoren, die immer denselben
Ergebnistyp haben, wie z.~B. die logischen Operatoren oder der Operator
\rw{sizeof}. Es gibt sogar einen Operator, dessen Ergebnistyp immer
\tvoid\ ist, n"amlich \rw{delete}.

\bigskip

In der folgenden Aufstellung der Operatoren verwende ich die
Klassifikation nach der funktionellen Zugeh"origkeit, weil sie
f"ur diesen Zweck am geeignetsten ist.

\subsubsection{Arithmetische Operatoren}

SEL bietet neben den numerischen Operatoren, zu denen die vier
Grundrechnungsoperatoren f"ur Addition, Subtraktion, Multiplikation und
Division z"ahlen, zus"atzlich der Modulo-Operator, der den Rest bei einer
ganzzahligen Division ergibt, und der Negationsoperator, der
einer Umkehr des Vorzeichens gleichkommt, auch
Operatoren zur Bitmanipulation, n"amlich bitweise Links- und
Rechtsverschiebung und bitweise Und-, Oder- und
XOR-Verkn"upfungen. Die bin"aren numerischen Operatoren, also
\fu{+}, \fu{-}, \fu{*}, \fu{/} und \fu{\%}, besitzen eine
niedrigere Priorit"at als alle un"aren Operatoren, jedoch eine
h"ohere Priorit"at als alle anderen bin"aren Operatoren.
Untereinander sind die Additions- und Subtraktionsoperatoren auf
derselben Priorit"atsstufe, die um einen Grad niedriger steht,
als die der anderen numerischen bin"aren Operatoren, in Anlehnung
an die in der Mathematik "ublichen Notation. Die un"aren
Operatoren haben deshalb eine h"ohere Priorit"at, weil das
erstens unseren Erwartungen entspricht---man betrachtet intuitiv einen
un"aren Operator einem Wert als "`n"aherstehender"', als einen
bin"aren---und es zweitens in fast allen F"allen praktischer ist.
So ist z.~B. das Ergebnis des Ausdrucks \fu{-x*y} dem von
\fu{-(x*y)} "aquivalent. Bes"a"sen aber die bin"aren Operatoren
eine h"ohere Priorit"at, m"u"ste man anstatt \fu{-x-y} den
Ausdruck \fu{(-x)-y} schreiben, was etwas unnat"urlich aussieht.
Hingegen k"onnte man den Ausdruck \fu{-(x+y)} zwar bei anderer
Operatorenrangfolge ohne Klammern schreiben, doch selbst ohne die
Rangfolge zu ver"andern, k"onnte man durch "Anderung des \fu{+} in
ein \fu{-} die Klammerung vermeiden: \fu{-x-y}. Die Rangfolge
erweist sich also als nat"urlich und praktisch. Die
Vergleichsoperatoren haben deshalb niedrigere Priorit"at als die
numerischen Operatoren, weil es kaum vorkommt, da"s man mit
Vergleichsergebnissen rechnet. F"ur die logischen Operatoren gilt
dieselbe Begr"undung. Die Festsetzung der Priorit"at der
Bitmanipulations-Operatoren relativ zu den numerischen Operatoren
ist eher willk"urlich. Es w"urde kaum ein Vor- oder Nachteil
entstehen, h"atten die Operatoren zur Bitmanipulation h"ohere
Priorit"at. Der Grund, warum ich diese Operatoren nicht auf
dieselbe Priorit"atsstufe gesetzt habe ist der, da"s sie zu
verschieden sind, als da"s dies dienlich w"are.

Die numerischen Operatoren haben eine besondere Eigenschaft: Sie
stellen sich auf die Typen ihrer Operanden ein und ver"andern
diese sogar, wenn n"otig. So stellt sich zum Beispiel die Frage,
welchen Typ eine Addition eines \tint\ und eines \tfloat\ haben
sollte. Vern"unftigerweise wandelt der Additionsoperator in
diesem Fall den \tint\ in einen \tfloat\ um, und f"uhrt erst
anschlie"send die Addition aus. Auf diese Weise gehen keine
Nachkommastellen des \tfloat\ verloren. Gleiches passiert mit
einem \tchar\ und einem \tint : Der \tchar\ wird in den Typ
\tint\ umgewandelt, da nur so gew"ahrleistet ist, da"s mit dem
richtigen Wert des \tint\ gearbeitet wird. F"ur die Typen bei
einer numerischen Operation gilt in SEL also die Regel \hbox{$
\tchar < \tint < \tfloat $}.\footnote{siehe \cite{selkr}, Seite
18} Zus"atzlich zu diesen Typen k"onnen die Operatoren \fu{+} und
\fu{-} auch mit Zeigern operieren, doch dazu mehr in Kapitel
\ref{memmgrops}.

Die Operatoren zur Bitmanipulation arbeiten in SEL lediglich mit
dem Typ \tint\, und das aus guten Gr"unden: Erstens ist es f"ur
einen Computer zwar leicht m"oglich, mit den Bits eines \tfloat\
zu operieren, doch es ist f"ur einen Menschen ungleich
schwieriger, diese Bits auch richtig zu interpretieren---auf den
Wert der Zahl bezogen. Zweitens sind die "ublichen 8 Bit eines
\tchar\ im Regelfall zu wenig f"ur die Anwendungen der
Bitmanipulation und drittens vereinfacht es nicht nur den Bau
eines Interpreters bzw. Compilers, es macht diese auch schneller.

\subsubsection{Logische Operatoren}

Neben den Operatoren zur logischen Und- und Oder-Verkn"upfung
bietet SEL auch einen Operator zur logischen XOR-Verkn"upfung an.
Ich habe mich dazu entschieden, weil es leicht in einen
Interpreter bzw. Compiler einzubauen ist und andererseits der
Operator in manchen Situationen sehr praktisch ist und logische
Ausdr"ucke verk"urzen kann. Die logischen Operatoren sind auf der
zweitniedrigsten Priorit"atsstufe angesiedelt---unter ihnen
befinden sich lediglich die Operatoren \fu{=} und \fu{delete}.
Diese Stufe habe ich gew"ahlt, weil logische Operatoren
haupts"achlich in Schleifen- und \rw{if}-Bedingungen eingesetzt
werden und es dort am praktischsten ist, wenn sie die niedrigste
Priorit"at haben. Nur selten ist es notwendig, da"s einer
Variablen ein logischer Wert zugewiesen wird, doch auch dieser
Fall ist ohne Klammerung zu bew"altigen, da ja der
Zuweisungsoperator eine noch niedrigere Priorit"at hat. Der
einzige Fall, in dem sich diese Rangfolge als unpraktisch
herausstellen w"urde ist, da"s in einer Schleifen- oder
\rw{if}-Bedingung ein logischer Wert einer Variablen zugewiesen
wird und dieser Wahrheitswert mit einem anderen mittels eines
logischen Operators verkn"upft werden mu"s. Ein Beispiel hierf"ur
w"are \fu{if (x \&\& (y = z))}. Tats"achlich tritt dieser Fall so
selten auf---in meiner bislang 4-j"ahrigen C-Laufbahn bin ich
noch nie in eine solche Situation gelangt---da"s man ihn getrost
au"ser Acht lassen kann.

Der logische Negationsoperator \fu{!} liegt auf derselben
Priorit"atsstufe wie die Operatoren zur bitweisen und
arithmetischen Negation und die beiden Operatoren \fu{*} und
\fu{\&}, die im Kapitel \ref{memmgrops} besprochen werden, welche
um eine Stufe tiefer liegt als die der Operatoren \fu{->}, \fu{:}
und \fu{[]}. Um das zu begr"unden, betrachten wir den Ausdruck
\fu{!x->y}. Offenbar wird hier versucht, aus der Struktur \fu{x}
auf das Element \fu{y} zuzugreifen und dieses zu negieren. H"atte
der Operator \fu{!} eine h"ohere Priorit"at als \fu{->}, so
w"urde sich die Negation auf die Struktur beziehen, und es ist
nicht m"oglich, eine Struktur zu negieren.

Logische Operatoren arbeiten in SEL nur mit dem Typ \tint .
Zwar w"urde es auch gen"ugen, den Typ \tchar\ zu diesem Zweck
heranzuziehen, doch da in kaum einer Anwendung Wahrheitswerte
gespeichert werden und \tint\ normalerweise der Gr"o"se des
Maschinenwortes des jeweiligen Prozessors entspricht, und daher
schneller zu bearbeiten ist, als ein \tchar , ist er die bessere
Wahl. Sollte es jedoch n"otig sein, Wahrheitswerte abzuspeichern,
kann das ohne weiteres mit Hilfe von \tchar s geschehen, die durch
ein Typecasting entstanden sind, oder aber durch Bitfelder, was
noch speicherplatzsparender ist.

\subsubsection{Vergleichsoperatoren}

SEL bietet insgesamt sechs Vergleichsoperatoren an. Diese sind um
eine Priorit"atsstufe h"oher als die logischen Operatoren
angesiedelt, da sie meistens in Schleifen- und
\rw{if}-Bedingungen eingesetzt werden und dort diese Rangfolge
von Vorteil ist. So wird es wahrscheinlich oft vorkommen, da"s
"uberpr"uft werden soll, ob zwei Werte zwei anderen Werten
entsprechen, aber eher selten wird man "uberpr"ufen wollen, ob
die Oder-Verkn"upfungen jeweils zweier Werte dieselben Ergebnisse
haben.

Vergleichsoperatoren k"onnen in SEL lediglich Werte der Typen
\tchar , \tint\ und \tfloat\ vergleichen, wobei das im vorletzten
Unterkapitel vorgestellte "`Gr"o"senverh"altnis"' gilt. Es ist
hingegen nicht m"oglich, z.~B. Strings miteinander zu
vergleichen, weil das, um es vern"unftig zu machen, eine
Ver"anderung im Typkonzept von SEL mit sich z"oge und au"serdem
sehr gut von einer Funktion erledigt werden kann---SEL stellt
alle n"otigen Konstrukte hierf"ur bereit.

\subsubsection{Speicherverwaltungsoperatoren}
\label{memmgrops}

Unter den Begriff "`Speicherverwaltungsoperatoren"' fallen all
jene Operatoren, die mit Zeigern, Arrays und Strukturen arbeiten.
Dazu z"ahlen nicht nur die Operatoren zur Ermittlung der
Speicheradresse einer Variablen und zur Dereferenzierung eines
Operators, \fu{\&} und \fu{*}, zum Zugriff auf ein Array- bzw.
Strukturelement, \fu{[]} und \fu{->}, sondern auch die Operatoren
zum Allokieren und Freigeben von Speicher, \fu{new} und
\fu{delete}, und der Operator zum Ermitteln der Gr"o"se eines
Typen, \fu{sizeof}. Au"serdem fallen unter diesen Begriff die
beiden bin"aren Operatoren \fu{+} und \fu{-}, mit denen
zus"atzlich zur numerischen Arithmetik auch Zeigerarithmetik
betrieben werden kann.

Die h"ochste Priorit"atsstufe aller Operatoren haben \fu{new} und
\fu{sizeof}, auf der n"achsten Stufe liegen \fu{->} und \fu{[]}.
Es h"atte keinen Sinn, diese Rangfolge zu "andern, denn die
Operanden von \fu{new} und \fu{sizeof} sind keine Ausdr"ucke,
sondern Typen, und es ist nicht m"oglich, den Operator \fu{->}
oder \fu{[]} auf einen Typ anzuwenden, soda"s dies schon von Haus
aus unterbunden werden sollte. Die beiden un"aren Operatoren
\fu{*} und \fu{\&} liegen um eine Priorit"atsstufe tiefer,
gemeinsam mit allen anderen un"aren Operatoren. Eine andere Stufe
f"ur diese Operatoren w"are nicht angebracht, aus Gr"unden, die
ich schon genannt habe.

Ein in gewisser Hinsicht au"sergew"ohnlicher Operator---gemessen
an seinen Gegenst"ucken aus C und \pascal ---ist der Operator
\fu{->}. Er wird benutzt, um auf Elemente einer Struktur
zuzugreifen. Derselbe Operator wird in C und \pascal\ durch den
Punkt (\fu{.}) symbolisiert. In C jedoch gibt es zus"atzlich den
Operator~\fu{->}. Er wird benutzt, um auf Elemente aus einem
Zeiger auf eine Struktur zuzugreifen---in \pascal\ mu"s hierf"ur
zuerst der Zeiger dereferenziert werden. Der SEL-Operator ist nun
insofern au"senstehend, als er die Funktionalit"at des
Punktoperators und des Pfeiloperators aus C vereinigt. So
erm"oglicht er es, angewendet auf eine Struktur, auf ein Element
daraus zuzugreifen. Angewendet auf einen Zeiger auf eine Struktur
dereferenziert er ersteren und greift dann erst auf das jeweilige
Element zu. Diese Vereinigung macht insofern Sinn, als da"s
keine Funktionalit"at verlorengeht, denn der Punktoperator
hat, angewendet auf einen Zeiger, keine Funktion. Auch das Design
eines Interpreters bzw. eines Compilers wird dadurch nicht
erschwert, denn diese wissen, wenn sie auf den Pfeiloperator
treffen, mit welchem Typ gearbeitet wird. Ein Compiler m"u"ste
also nur, tr"afe er auf einen Pfeiloperator, dessen linker Operand
ein Zeiger ist, den Code zur Dereferenzierung einbauen.

\subsubsection{Der Zuweisungsoperator}

Der Zuweisungsoperator von SEL ist ein Sonderfall unter den
Operatoren, aus mehreren Gr"unden. Der erste Grund ist, da"s es
in manchen Programmiersprachen, darunter auch \pascal , gar
keinen Zuweisungsoperator in diesem Sinne gibt. In \pascal\ ist
eine Zuweisung eine Anweisung und kein Ausdruck. Das macht es in
\pascal\ auf der einen Seite schwerer, den Zuweisungsoperator
vom Compiler bzw. Interpreter unentdeckt mit dem
Vergleichsoperator zu verwechseln, denn der Zuweisungsoperator
ist in Ausdr"ucken nicht erlaubt und Ausdr"ucke sind in \pascal\
keine g"ultigen Anweisungen. Andererseits jedoch ist es nicht
m"oglich, mehreren Variablen in einer Anweisung denselben Wert
zuzuweisen, wie man es von C oder auch aus der Mathematik gewohnt
ist (\fu{x=y=0}). Au"serdem ist es nicht m"oglich, z.~B. in einer
\rw{if}-Anweisung gleichzeitig eine Zuweisung vorzunehmen und den
Wert der zuzuweisenden Variablen zu "uberpr"ufen (z.~B.
\hbox{\fu{if ((x=y) == 3)}}). Manche legen das sogar als Vorteil
aus, weil es eine unsch"one Programmierung erschwert, und
insofern gebe ich ihnen Recht, doch ich pers"onlich finde solche
Konstrukte in manchen Situationen doch sehr praktisch, und da der
Aufwand f"ur einen Zuweisungsoperator etwa dem einer
Zuweisungsanweisung gleichkommt, oder sogar noch geringer ist,
habe ich ihn in SEL eingebaut.

Der zweite Grund f"ur die Besonderheit des \fu{=}-Operators ist,
da"s er mit allen Typen zusammenarbeitet. Anders als der
Vergleichsoperator, der auf die Grundtypen beschr"ankt ist, kann
der Zuweisunsoperator eine Variable jeden beliebigen Typs einer
anderen zuweisen---Typgleichheit vorausgesetzt. Ich habe diese
"`Allm"achtigkeit"' des Operators deswegen in SEL erm"oglicht, weil
es leicht ist, das in einen Interpreter einzubauen---im Fall der
Konzeption, die wir im zweiten Kapitel besprechen werden, sogar
noch leichter, als sich auf die Grundtypen zu beschr"anken. Auch
f"ur einen Compiler ist das nicht weiter schwer: Den Code f"ur die
Grundtypen kann der Compiler direkt einf"ugen, f"ur Strukturen
oder Arrays k"onnte eine Funktion aus der Laufzeitumgebung
aufgerufen werden. Hier dr"angt sich nat"urlich die Frage auf,
warum ich dieses Konzept nicht auf den Vergleichsoperator
\fu{==} ausgeweitet habe. Die Antwort ist einfach: H"atte ich das
getan, w"are es unausweichlich gewesen, das auch f"ur die
Operatoren \fu{<}, \fu{>}, \fu{<=}, \fu{>=} und \fu{!=} zu tun.
Mit letzterem h"atte ich keine Probleme gehabt, doch wann ist
eine Struktur, die die Adresse einer Person enth"alt, gr"o"ser
als eine andere desselben Typs? Byte-f"ur-Byte-Vergleich w"are
hier sicher nicht angebracht.\footnote{In C++ kann dieses Problem
gel"ost werden, weil es dort m"oglich ist, Operatoren zu
"uberladen, d. h., da"s der Benutzer imstande ist, eine Funktion
zu schreiben, die zwei Strukturen (beispielsweise des Typs
\fu{Adresse}) vergleicht, die dann anstatt des
Vergleichsoperators aufgerufen wird.}

\subsubsection{Der Typecasting-Operator}

Der Typecasting-Operator \fu{:} wird benutzt, um einen Wert
explizit in einen anderen Typ umzuwandeln. Er wird genauer in
Kapitel \ref{casts} beschrieben. Seine Priorit"at ist mit der der
beiden Operatoren \fu{->} und \fu{[]} "aquivalent, erstens, weil
er ein un"arer Operator ist und zweitens, weil es eher selten
vorkommt, da"s gr"o"sere Ausdr"ucke gecastet werden m"ussen. Wenn
n"otig, wird dies sowieso von SEL "ubernommen.

\subsection{Konstanten}
\label{const}

Grunds"atzlich gibt es in SEL vier verschiedene Arten von
Konstanten: \tchar -Konstanten, \tint -Konstanten,
\tfloat -Konstanten und String-Konstanten. Eine \tchar -Konstante
ist ein unter einfache Anf"uhrungsstriche gesetztes Zeichen bzw.
Steuerzeichen und hat---es ist nicht anders zu erwarten---den Typ
\tchar . Eine \tint -Konstante ist eine simple Abfolge von
Ziffern. Zahlenkonstanten im lexikalischen Sinn sind in SEL immer
positiv---negative Zahlen macht erst der un"are Operator \fu{-}
m"oglich, und insofern ergibt sich daraus kein Nachteil. Der
Grund, warum ich Konstanten nicht auch auf negative Zahlen
ausgeweitet habe ist der, da"s sich die lexikalische
Analyse\footnote{siehe Kapitel \ref{scanner}}
dadurch sehr erschweren w"urde. Nehmen wir als Beispiel den
Ausdruck \fu{7-4}. Der Scanner erkennt als erstes die
\tint -Konstante \fu{7}. Danach m"u"ste ein Operator folgen oder
der Ausdruck zu Ende sein, was aber der Scanner nicht wissen
kann, da er vom Parser in dieser Hinsicht abgekapselt ist. Aus
diesem Grund w"urde er anstatt den Operator \fu{-} zu erkennen
eine weitere \tint -Konstante, n"amlich \fu{-4} lesen, was
nat"urlich falsch ist. Das Gesagte gilt auch f"ur
\tfloat -Konstanten.

String-Konstanten sind in gewisser Hinsicht etwas Besonderes,
n"amlich in Anbetracht der Tatsache, da"s es einen Stringtyp in
SEL gar nicht gibt. Normalerweise werden Strings als Arrays von
\tchar s behandelt, die durch ein Byte mit dem Wert \fu{0}
abgeschlossen sind, es w"are allerdings nicht schwer, Funktionen
zu programmieren, die \pascal -Strings verarbeiten. Da jedoch die
erstere Stringart die verbreitetere ist, wollen wir bei dieser
bleiben. Wie wir wissen, kann ein Array in SEL in vielen
Situationen wie ein Zeiger behandelt werden. Was also w"are
falsch daran, einen String als Pointer-Konstante zu sehen?
Tats"achlich wird das nicht nur in SEL, sondern auch in C so
praktiziert, von wo ich dieses Konzept "`importiert"' habe. Dieser
Pointer zeigt auf einen Speicherbereich, in dem der eigentliche
String gespeichert ist. Nat"urlich mu"s dieser Speicherbereich
w"ahrend der gesamten Programmausf"uhrung an derselben
Speicheradresse bleiben, und nicht irgendeiner L"oschaktion zum
Opfer fallen, soda"s die Pointerkonstante im ganzen
Programm zu jeder Zeit "`aktuell"' ist.\footnote{In Wirklichkeit
mu"s der String nicht w"ahrend der gesamten Programmausf"uhrung
im Speicher sein, sondern erst nach seinem ersten Auftreten im
Code, doch der Einfachheit halber wird er die ganze Zeit im
Speicher gehalten.}

\subsection{Variablen}
\label{variables}

In SEL k"onnen zwei Arten von Variablen unterschieden werden, die
ihrerseits in jeweils zwei Untertypen unterteilt werden k"onnen.
Lokale Variablen sind nur innerhalb einer bestimmten Funktion
bekannt und existieren f"ur jede Instanz dieser Funktion
getrennt.\footnote{Mehrere Instanzen einer Funktion sind dann
m"oglich, wenn sich diese direkt oder indirekt selbst aufruft.}
Die Gruppe der lokalen Variablen werden weiter unterteilt in
Parameter und direkte Funktionsvariablen.\footnote{Eine bessere
Bezeichnung als "`direkte Funktionsvariablen"' habe ich nicht
gefunden.} Die Werte der Parameter werden von der aufrufenden
Funktion bestimmt, wohingegen die direkten Funktionsvariablen der
Funktion alleine unterstehen. Lokale Variablen werden auf dem
Stack gespeichert. Wird eine Funktion aufgerufen, werden die
Parameter von der aufrufenden Funktion auf den Stack gelegt und
anschlie"send zum Code der Funktion
gesprungen.\footnote{Nat"urlich wird auch die R"ucksprungadresse
auf dem Stack gespeichert} Die Funktion ihrerseits legt die
direkten Funktionsvariablen auf den Stack und f"uhrt dann ihre
Anweisungen aus. Wird sie beendet, entfernt sie sowohl
direkte Funktionsvariablen als auch Parameter vom Stack
und springt dann zur aufrufenden Funktion zur"uck. Das Konzept,
den Stack als Speicher f"ur lokale Variablen zu benutzen macht
es m"oglich, da"s bei einer rekursiven Funktion jede Instanz ihre
eigenen Variablen zur Verf"ugung hat.

Globale Variablen werden in SEL unterteilt in modulglobale und
programmglobale Variablen. Erstere sind nur in der Quelldatei
(Modul) bekannt, in der sie definiert sind, wohingegen letztere
auch in anderen Modulen benutzt werden k"onnen. Sowohl modul- als
auch programmglobale Variablen werden im statischen Speicher
aufbewahrt. Daf"ur, da"s modulglobale Variablen nur in ihrem
jeweiligen Modul referenziert werden, haben der Interpreter bzw.
der Compiler oder der Linker zu sorgen.

Ich glaube, "uber die Verwendung von Variablen in Ausdr"ucken
mu"s ich hier nicht viele Worte verlieren: das geschieht ganz
einfach---wie in fast allen Programmiersprachen "ublich---durch
die Benutzung des Variablennamens.

\subsection{Funktionen}
\label{exprfuncs}

Die Aufrufkonvention f"ur Funktionen ist einfach: Der
Funktionsname gefolgt von in Klammern gesetzten Parametern,
jeweils durch einen Beistrich getrennt, bilden einen Funktionsaufruf.
Hat die Funktion keine Parameter, sind trotzdem "offnende und
schlie"sende Klammern anzugeben, aus mehreren Gr"unden. Erstens
kann der Interpreter bzw. Compiler daran ohne nachzusehen
erkennen, da"s es sich um einen Funktionsaufruf handelt und nicht
um die Verwendung einer Variablen. Zweitens ist dies in Hinblick
auf die evtl. Erweiterung der Sprache durch
Funktionszeiger\footnote{siehe Kapitel \ref{langext}}, also
Zeigervariablen, die die Adressen von Funktionen aufnehmen,
n"otig. In diesem Fall signalisiert das Klammerpaar, da"s wir
vorhaben, diese Funktion aufzurufen und nicht, sie als Zeiger zu
verwenden.

Die Anzahl und die Typen der Funktionsparameter sind nat"urlich
einzuhalten. Die Regelung bei letzteren ist nicht so streng:
Stimmen die Typen nicht "uberein, versucht SEL mittels impliziten
Typecastings\footnote{siehe Kapitel \ref{casts}}, diese in die
richtigen umzuwandeln, was jedoch nicht unbedingt gelingen mu"s,
da sich der Interpreter bzw. der Compiler an bestimmte Regeln zur
Typumwandlung halten mu"s.

\section{Typen}
\label{types}

In diesem Abschnitt m"ochte ich das Typkonzept von SEL erkl"aren.
Zun"achst stellt sich beim Design einer Programmiersprache die
Frage, welches Typkonzept "uberhaupt zur Anwendung kommen soll.
So gibt es z.~B. in \basic\ einige wenige Typen, aus denen der
Benutzer ausw"ahlen kann, es ist dort jedoch nicht m"oglich,
selber Typen zu definieren. Zwar k"onnen Arrays definiert werden,
doch ist ein Array in \basic\ nichts anderes als eine Menge von
Variablen, deren Namen sich jeweils durch eine Nummer
unterscheiden---ein Array als Ganzes zu behandeln, ist in \basic\
nicht m"oglich.

Dann gibt es Programmiersprachen---\dbase\ ist
ein Beispiel hierf"ur---in denen Variablen zwar einen bestimmten
Typ haben, dieser sich aber jederzeit "andern kann. Dort ist es
z.~B. m"oglich, einer Variablen einen numerischen Wert zuzuweisen
und mit ihr Berechnungen anzustellen. Sp"ater kann man ihr bei
Bedarf einen String zuweisen und diesen am Bildschirm ausgeben.
Zwar bietet dieses Konzept auf den ersten Blick f"ur den Benutzer
Vorteile\footnote{Auch diese Vorteile sind nur scheinbar, denn
sie verf"uhren zu einem sehr schmutzigen Programmierstil} und es
ist f"ur eine reine Interpretersprache eine gute Wahl, doch einem
Compiler macht ein solches Konzept gro"se Probleme: In
Programmiersprachen wie C oder \pascal\ haben Variablen einen
vorgegebenen, konstanten Typ. Vom Compiler werden diese Variablen
an eine bestimmte Speicheradresse gelegt, die sich im Laufe des
Programms nicht ver"andert, was ohne weiteres m"oglich ist, da
sich ja die Speicheranforderungen der Variablen auch nicht
ver"andern. K"onnen sich aber die Typen der Variablen "andern,
"andern sich auch die Gr"o"sen der Variablen, d. h., da"s es
nicht m"oglich ist, sie an eine konstante Stelle im Speicher zu
setzen. W"ahrend ein C-Compiler, wenn er die Variable zu
referenzieren hat, ganz einfach ihre Speicheradresse einf"ugen
kann, mu"s ein compiliertes \dbase -Programm eine Tabelle
verwalten, in der alle Variablen und deren momentane Adressen und
Typen verzeichnet sind. Wird nun der Ausdruck \fu{x+1}
abgearbeitet, mu"s zun"achst in der Tabelle nachgesehen werden,
an welcher Adresse sich die Variable \fu{x} befindet. Zus"atzlich
mu"s auch gepr"uft werden, welchen Typ diese hat, da z.~B. eine
ganzzahlige Addition einen anderen Code verlangt, als eine
Addition mit Gleitkommazahlen. Erst jetzt kann die Addition
ausgef"uhrt werden. Es ist also nur verst"andlich, da"s Compiler
f"ur eine solche Sprache einen Code erzeugen, der zehnmal
langsamer\footnote{Die Zahl 10 mag wie eine "Ubertreibung
klingen, tats"achlich ist es nicht selten, da"s solche Codes noch
um einiges langsamer sind. So ben"otigt z.~B. ein Programm, da"s
einen QuickSort-Algorithmus ausf"uhrt in Clipper (einem
\dbase -Compiler) sage und schreibe 120 mal l"anger als ein
vergleichbares Programm in C} ist als der eines C- oder
\pascal -Compilers.

Eine weitere wichtige Kategorie ist die der typenlosen
Programmiersprachen---AWK f"allt teilweise darunter. Dort ist es
m"oglich, einer Variablen einen String zuzuweisen, der eine Zahl
enth"alt und in der n"achsten Zeile mit dieser Variable zu
rechnen als w"are sie numerischen Typs. Dieses Konzept ist
insofern unpraktisch, als man f"ur verschiedene Typen (obwohl
es diese eigentlich nicht gibt) auch verschiedene Operanden
braucht. Bei einem Ausdruck \fu{x+y} w"are also nicht klar, ob
die beiden Strings zu konkatenieren sind oder ob deren numerische
Werte addiert werden sollen. Auch in Bezug auf den Bau eines
Compilers ist dieses Konzept von Nachteil, aus "ahnlichen
wie im vorigen Absatz dargelegten Gr"unden.

Manche alte Programmiersprachen stellen als einzigen Typ das
Maschinenwort zur Verf"ugung, was einerseits die Gefahr birgt,
da"s in einer solchen Sprache geschriebene Programme sehr vom
Prozessor abh"angig sind und andererseits den Nachteil hat, da"s
viele Aufgaben, die der Interpreter bzw. Compiler ohne 
Effizienzverluste "ubernehmen k"onnte, vom Programmierer 
bestritten werden m"ussen.

Wie sich aus den vorigen Abs"atzen leicht herauslesen l"a"st,
habe ich mich f"ur das Typkonzept entschieden, das auch in C
oder \pascal\ verwendet wird. Dort mu"s jede Variable vor ihrer
Verwendung definiert werden, wobei gleichzeitig ihr Typ
festgelegt wird, der sich im Laufe des Programms nicht "andert.
Au"serdem ist es dem Benutzer m"oglich, aus den Basistypen der
Programmiersprache eigene Typen herzuleiten, wenn auch nicht
unbeschr"ankt.

\subsection{Elementare Typen}
\label{elemtypes}

Wie bereits gesagt, ben"otigt man zur Konstruktion von
komplizierteren Typen eine Menge an Basistypen, mit denen man
mittels der Konstrukte, die SEL bereitstellt, arbeiten kann.
Programmiersprachen wie C oder \pascal\ bieten eine relativ
gro"se Zahl solcher Basistypen an. So gibt es in C zwei
verschiedene Typen, die ein einzelnes Zeichen repr"asentieren,
drei Grundtypen f"ur ganze Zahlen, wobei von jedem dieser Typen
ein vorzeichenloser sowie ein vorzeichenbehafteter Typ
exisitiert, und mehrere Gleitkomma-Typen. Die Vielfalt in
\pascal\ ist "ahnlich gro"s, wobei zus"atzlich noch ein Typ f"ur
logische Werte exisitiert (\type{boolean}). Ich habe mich
entschieden, in SEL nur vier Elementartypen aufzunehmen:
\tchar , \tint , \tfloat\ und \tvoid .

Der Typ \tchar\ repr"asentiert ein einzelnes Zeichen und nimmt
dementsprechend nur ein Byte ein. Arrays aus \tchar s k"onnen
z.~B. als Strings interpretiert werden.

Ein \tint\ repr"asentiert eine ganze Zahl mit Vorzeichen und hat
in der Regel die Gr"o"se eines Maschinenwortes.\footnote{Eine
Ausnahme bildet MS-DOS, wo ein Maschinenwort nur 16 Bits lang
ist, d.~h. der Wertbereich w"urde von etwa -32000 bis 32000
reichen, was zwar f"ur viele Anwendungen reicht, aber ungen"ugend
ist, da es ja nur einen Typ f"ur ganze Zahlen gibt. Auf MS-DOS
Maschinen ist deshalb die Gr"o"se eines \tint\ zwei
Maschinenworte, also 32 Bits.} Zus"atzlich hat \tint\ die
Aufgabe, logische Werte zu repr"asentieren, wobei der Wert \fu{0}
einem "`logisch falsch"' und ein Wert ungleich \fu{0} "`logisch
wahr"' entspricht. Unter anderem hat diese Konzeption den
Vorteil, da"s ein Ausdruck des Typs \tint\ nicht nur den Wert des
\tint\ wiederspiegelt, sondern auch die Aussage, ob dieser
ungleich \fu{0} ist oder nicht. So ist es z.~B. m"oglich, mit
einer sehr kurzen \rw{if}-Anweisung zu pr"ufen, ob ein Zeiger ein
Nullzeiger ist. W"ahrend hierf"ur in \pascal\ die Anweisung
\fu{if p <> nil then} benutzt wird, kann in SEL einfach
\fu{if (p)} geschrieben werden. Zieht man es aus Gr"unden der
"Ubersichtlichkeit vor, trotzdem einen boolee'schen Typ zu haben,
kann man diesen in SEL problemlos mit den Direktiven

\begin{verbatim}
#define boolean        int
#define false          0
#define true           1
\end{verbatim}

\noindent simulieren.\footnote{Die Betonung liegt auf "`simulieren"',
da wir uns mit diesen Direktiven keinen eigenen Typ schaffen sondern
lediglich \fu{boolean} zu einem Pseudonym f"ur \rw{int} machen.}

Der Typ \tfloat\ repr"asentiert eine vorzeichenbehaftete
Gleitkommazahl. Ihre Gr"o"se im Speicher h"angt vom jeweiligen
Prozessor ab.

Einen Spezialfall unter den Typen stellt \tvoid\ dar: Er
repr"asentiert gar nichts! F"ur sein Dasein sprechen zwei
Gr"unde. Erstens macht er es m"oglich, Funktionen ohne
R"uckgabewert zu konstruieren, also Konstrukte, die mit den
Prozeduren in \pascal\ oder \modula\ vergleichbar sind. Au"serdem
ist er der Grundtyp f"ur typenlose Zeiger, also Zeiger, die auf
etwas nicht Bestimmtes zeigen. Verwendung finden solche Zeiger
z.~B. in Funktionen zur Speicherallokation. Diese geben einen
Zeiger auf einen Speicherbereich bestimmter Gr"o"se zur"uck,
wobei aber nicht bekannt ist, f"ur welchen Typ dieser Bereich
reserviert wird.

\subsection{Arrays}

Ein Array ist ein Zusammenschlu"s mehrerer Variablen desselben
Typs, die im Speicher sinnvollerweise aufeinander folgen. Ich
habe diesen Typ in SEL aufgenommen, weil er erstens nicht schwer
zu implementieren und zweitens in sehr vielen Situationen
unverzichtbar ist. Auffallend an der Schreibweise des Arraytyps
ist, da"s die Dimension des Arrays dem Grundtyp nicht nach- sondern
vorgestellt sind. Gem"a"s unserem Grundsatz, Symbole anstatt
Schl"usselw"ortern zu verwenden, wo es nicht st"ort, habe ich hier
auf Ausschm"uckungen wie das in \pascal\ verwendete \fu{array
... of ...} verzichtet und statt dessen eine in eckige Klammern
geschlossene Array-Dimension benutzt. Auf ein Angeben der
Grenzen, wie in \pascal\ "ublich, habe ich verzichtet, das erste
Element besitzt den Index \fu{0}. Diese Vereinfachung hat zwei
Gr"unde: Einerseits ist es nicht schwer, auf speziell indizierte
Arrays zu verzichten, andererseits w"urde sich das auf die
Konsistenz in Bezug auf die Verwendung von Zeigern anstatt Arrays
auswirken: Angenommen, wir h"atten ein Array \fu{a : [5..10]int}.
In SEL ist es nun m"oglich, dieses Array als Zeigervariable zu
betrachten und so zu behandeln, d. h., wir k"onnten mit ihm
Zeigerarithmetik betreiben. Per definitionem ist ein Zeiger
\fu{p} gleich \fu{p + 0}, was bedeutet, da"s wir mit \fu{*(a + 0)} auf das
erste Element des Arrays zugreifen k"onnen. Wir k"onnen auf
dieses Element auch direkt "uber das Array zugreifen, also mit
\fu{a[5]}. Ich denke, hier zeigt sich schon, worauf ich
hinauswill: Die Zahlen, die zur Indizierung verwendet werden,
stimmen nicht "uberein. Um das zu umgehen, m"u"ste man entweder
auf die Gleichsetzung von Array und Zeiger verzichten, was ich
nicht wollte, oder auch in die Zeiger eine komplizierte
Nummerierung einbauen, was mir und sicher auch dem Benutzer der
Programmiersprache zu umst"andlich ist.

Zur"uck zu der Dimension des Arrays, die, wie oben erw"ahnt, dem
Grundtyp vorangestellt ist. Der Grund hierf"ur ist in erster
Linie beim Typecasting\footnote{siehe Kapitel \ref{casts}}
von SEL zu suchen. Nehmen wir als Beispiel den Ausdruck
\fu{a : *int[10]}. Hier w"are es unklar, ob die eckige Klammer noch
zum Typ geh"ort, oder ob versucht wird, auf das zehnte Element
des Zeigers (der als Array behandelt wird) zuzugreifen.
Nat"urlich kann auch dieses Problem durch Klammerung umgangen
werden, doch warum etwas komplizierter machen, als es ist? Die
vorangestellte Arraydimension in SEL hat n"amlich noch einen
Vorteil: Es ist ohne eine Rangordnung oder Klammerung m"oglich,
einen Typ eindeutig zu bestimmen. Als Beispiel soll der C-Typ
\fu{*int[10]} dienen. Ohne den beiden Operatoren eine
unterschiedliche Priorit"at zuzuweisen, ist es unm"oglich zu
sagen, ob der Typ nun \fu{*(int[10])} oder \fu{(*int)[10]} ist.
In SEL ist die Sache eindeutig: \fu{[10]*int} kann nur als Array
aus zehn Zeigern auf \tint\ gedeutet werden.

\subsection{Zeiger}

Zeiger stellen ein wichtiges Konstrukt dar, mit dessen Hilfe es
m"oglich ist, sehr komplexe, verkettete und verzweigte
Datenstrukturen zu bilden, die ohne sie nur sehr schwer zu
verwirklichen w"aren. Der Bau eines Interpreters oder gar eines
Compilers z.~B. w"are ohne die Verwendung von Zeigern kaum zu
bewerkstelligen. Aus diesem Grund habe ich Zeiger in SEL
eingebaut, was f"ur eine Interpretersprache eigentlich sehr
ungew"ohnlich ist.

Sehr eng in Verbindung mit Zeigern steht in SEL der Operator
\fu{*}, der einerseits in der Typdefinition eines Zeigers
verwendet wird und andererseits benutzt wird, einen Zeiger zu
dereferenzieren. Letzteres geschieht, indem dem Zeiger der
Sternoperator vorangestellt wird, anders als in \pascal , wo der
Dereferenzierungsoperator als Suffix verwendet wird. Der Grund,
warum in SEL der Operator vorangestellt wird, ist die Tatsache,
da"s er dasselbe Symbol hat wie der Multiplikationsoperator, und
der Parser\footnote{siehe Kapitel \ref{parser}} nicht w"u"ste,
wie etwa der Ausdruck \fu{p * - x} gemeint ist. Einerseits k"onnte es
die Differenz zwischen der Variablen, auf die \fu{p} zeigt und
\fu{x} sein, andererseits aber auch \fu{p} multipliziert mit
\fu{-x}. Nat"urlich h"atte ich auch ein anderes Symbol f"ur die
Dereferenzierung verwenden k"onnen, doch das wollte ich nicht,
weil ich den Operator \fu{*} von C gewohnt bin, au"serdem fand
ich es schon immer ein wenig verwirrend, da"s bei der
Typdeklaration der Operator vorangestellt, bei der
Dereferenzierung aber nachgestellt sein mu"s.

Der zweite wichtige Operator in Verbindung mit Zeigern ist
\fu{\&}. Wird er einer Variablen vorangestellt, dann ergibt er
einen Zeiger auf diese Variable. Mit seiner Hilfe ist es also
m"oglich, einen Zeiger auf eine bestimmte Variable zeigen zu
lassen. Verst"andlicherweise ist es nicht m"oglich, \fu{\&} in
Verbindung mit einer Konstante zu benutzen, denn eine Konstante
hat keinen bestimmten Platz im Speicher, sie wird vom Interpreter
direkt in einen Wert umgewandelt und von einem Compiler direkt in
den Code eingef"ugt. Zu verwechseln ist dies jedoch nicht mit den
String-Konstanten\footnote{siehe Kapitel \ref{const}}, die
eigentlich nicht konstante Strings sind, sondern konstante Zeiger
auf einen String, der an einer bestimmten Stelle im Speicher
liegt. Doch zur"uck zum Operator \fu{\&}: Sein zweiter
Verwendungszweck ist die Kennzeichnung eines Parameters einer
Funktion, der {\em Call-By-Reference} "ubergeben werden soll. In
\pascal\ z.~B. "ubernimmt diese Aufgabe das Schl"usselwort
\rw{var}, in C gibt es etwas Derartiges nicht, erst in C++ wurden
Referenzparameter eingef"uhrt. Auch dort werden sie mit \fu{\&}
gekennzeichnet. In einer Programmiersprache, in der Zeiger zur
Verf"ugung stehen, sind Referenzparameter nicht zwingend
notwendig: Es k"onnten Zeiger auf die betreffenden Variablen
"ubergeben werden, wie es in C auch praktiziert wird. Ich habe
mich aus zwei Gr"unden entschlossen, Referenzparameter doch in
SEL aufzunehmen. Erstens k"onnen sie in manchen Situationen
Programme um sehr viel "ubersichtlicher und leichter wartbar
machen. Zweitens, und das war f"ur mich wohl der gewichtigere
Grund, arbeitet der im n"achsten Kapitel theoretisch beschriebene
Interpreter intern ohnedies mit Referenzen, soda"s es keine
Schwierigkeit war, Referenzparameter einzubauen.

\subsubsection{Dynamische Variablen}

Eng verbunden mit den Zeigern ist das Konzept der dynamischen
Variablen. Dynamisch deshalb, weil sie erzeugt werden k"onnen,
wenn man sie braucht und aus dem Speicher entfernt, wenn sie
"uberfl"ussig werden. Statische Variablen belegen
Speicherbereiche, die vom Compiler bzw. Interpreter im voraus
bestimmt werden, wohingegen der Speicherbereich f"ur eine
dynamische Variable erst w"ahrend der Programmausf"uhrung
bestimmt werden kann. Das impliziert nat"urlich, da"s dem
Benutzer geeignete Mittel in die Hand gegeben werden m"ussen, um
dynamische Variablen zu erzeugen. Grunds"atzlich gibt es hierf"ur
zwei M"oglichkeiten. Entweder, man baut das Konzept der
dynamischen Speicherverwaltung nicht fest in die
Programmiersprache ein, sondern "uberl"a"st es dem Benutzer bzw.
dem Designer der Laufzeitbibliothek, geeignete Funktionen zu
implementieren, wie es in C der Fall ist, oder man
ber"ucksichtigt dynamische Speicherverwaltung schon beim Design
der Programmiersprache und stellt Konstrukte zur Verf"ugung, um
dynamische Variablen zu erzeugen, wie z.~B. in C++.

Ich habe mich f"ur den zweiten L"osungsansatz entschieden, weil
ich finde, da"s das Konzept der dynamischen Speicherverwaltung zu
wichtig ist, um es beim Design einer Programmiersprache nicht
miteinzubeziehen. Ein daraus resultierender Vorteil ist, da"s
Programme leichter lesbar werden und es leicher ist, sauberer zu
programmieren, als z.~B. in C.\footnote{Was nicht hei"sen will,
da"s es in C nicht m"oglich ist, sauber zu programmieren, es ist
nur nicht so leicht wie etwa in C++} Angelehnt an das Vorbild C++
werden in SEL zur dynamischen Speicherverwaltung die beiden
un"aren Operatoren \rw{new} und \rw{delete}\footnote{siehe
Kapitel \ref{memmgrops}} verwendet. Der Operand von \rw{new} ist ein
Typ, also im Gegensatz zu denen der meisten anderen Operatoren
kein Ausdruck. Der Ergebnistyp ist ein Zeiger auf den
"ubergebenen Typ, der auf den neu allokierten Speicherbereich
zeigt. Dieser Bereich kann nun so lange wie n"otig benutzt
werden. Wird er nicht mehr verwendet, sollte er mit dem Operator
\rw{delete} freigegeben werden, um sp"ater anderen dynamischen
Variablen zur Verf"ugung zu stehen. Der Grund, warum ich aus
\rw{new} und \rw{delete} keine Anweisungen gemacht
habe---besonders bei \rw{delete} w"urde sich das anbieten---ist,
da"s alle anderen Anweisungen in SEL zur Flu"skontrolle im
Programm verwendet werden, soda"s dies gewisserma"sen ein Bruch
in der Konsistenz der Sprache w"are.

\subsection{Strukturen}
\label{structs}

Eine Struktur\footnote{Der Name "`Struktur"' stammt von C, in
\pascal\ wird f"ur dasselbe der Name "`Record"' verwendet} ist
der Zusammenschlu"s mehrerer Variablen, auch verschiedenen Typs,
unter einem Namen. Sie bilden zusammen mit Zeigern und
dynamischen Variablen die Basis f"ur weit komplexere, verzweigte
und verkettete Datenstrukturen. "Uber ihre Definition m"ochte ich
nicht viele Worte verlieren, da sie einfach und zweckm"a"sig ist
und es au"serdem nicht viele andere M"oglichkeiten gibt,
Strukturen zu definieren.  Was jedoch viel interessanter ist, ist
die Benutzung von Zeigern auf eine Struktur in der Struktur
selbst. Ein Beispiel:

\begin{verbatim}
struct List
  i    : int;
  Next : *List;
end
\end{verbatim}

Diese Struktur ist eine m"ogliche Definition einer einfach
verketteten Liste. Neben einem oder mehreren Werten (in diesem
Fall ein \tint ) enth"alt sie einen Zeiger auf eine Struktur
ihres eigenen Typs. Einen Interpreter oder Compiler, der mit nur
einem Pass arbeitet und diese Struktur akzeptiert, zu schreiben,
ist noch leicht: Man k"onnte sich schon beim Lesen des auf das
Schl"usselwort \rw{struct} folgenden Bezeichners \fu{List}
notieren, da"s es diesen Typ von Struktur gibt. Beim Zeiger
\fu{Next} kann diese Notiz benutzt werden, um seinen Typ
festzulegen. Beim Lesen des Schl"usselwortes \rw{end} wird die
Notitz anschlie"send zu einer kompletten Definition erg"anzt. Was
tun wir jedoch, wenn es zwei Strukturen gibt, die jeweils einen
Zeiger auf die andere Struktur enthalten, wie es nicht selten
vorkommt? Eine der beiden Strukturen mu"s zuerst gelesen werden,
und dort treffen wir auf einen unbekannten Typen. In C wird dieses
Problem dadurch gel"ost, da"s man dem Compiler vor der Definition
der beiden Strukturen einen Hinweis darauf gibt, da"s diese
Strukturen irgendwo einmal definiert werden, soda"s er schon vor
der Definition seine Notiz machen kann. In SEL jedoch k"onnen wir
dieses Problem weit einfacher l"osen: Da es nicht m"oglich ist,
einem Typ in SEL einen anderen Namen zu geben, steht, wenn ein
unbekannter Typ gelesen wird, zweifelsfrei fest, da"s es sich
dabei um eine Struktur handeln mu"s, soda"s unser Compiler bzw.
Interpreter bereits hier die Notiz vornehmen kann.

\subsection{Typecasting}
\label{casts}

Bei einer Gegen"uberstellung der Typkonzepte von C und \pascal\
f"allt auf, da"s \pascal\ in der Typ"uberpr"ufung weit strenger
ist als C. So ist es in C ohne weiteres m"oglich, einem \tchar\
einen numerischen Wert zuzuweisen, wohingegen in \pascal\ zu
diesem Zweck ein eigener Typ existiert. In der Sprache \pascal\
ist eine Umwandlung eines Typs in einen anderen gar nicht
vorgesehen,\footnote{au"ser in F"allen, wo die beiden Typen zur
gleichen Gruppe geh"oren, wie z.~B. \tint\ und \type{longint}} in
C wird sie in vielen F"allen sogar vom Compiler "ubernommen, an
anderen Stellen kann sie der Benutzer mittels Typecasting
durchf"uhren. Der Grund, warum ich in SEL ein flexibles
Typkonzept eingebaut und nicht die strenge "Uberpr"ufung von
\pascal\ "ubernommen habe, ist, da"s es bei vier Basistypen sehr
schwer ist, ohne Typumwandlung auszukommen. Au"serdem k"onnen
Programme durch implizites Typecasting um einiges k"urzer und
"ubersichtlicher werden, man denke nur an die Umwandlung zwischen
\tchar\ und \tint , wof"ur in \pascal\ eigene Funktionen
notwendig sind.

Um aber nicht in einem vollst"andigen Chaos zu enden, darf nicht
wahllos jeder Typ implizit in jeden anderen Typen umgewandelt
werden, nur damit im Programm ohne Fehlermeldung fortgefahren
werden kann. Die Regeln, die in SEL f"ur implizites sowie
explizites Typecasting gelten, sind in \cite{selkr}, Seite 24
zusammengefa"st.

Der Grund, warum ich nicht das in C "ubliche Typecasting, einem
Ausdruck den neuen Typ in Klammern voranzustellen, "ubernommen
habe, ist der, da"s ich einfach etwas Neues ausprobieren wollte,
und das Typecasting mittels des \fu{:}-Operators exzellent zu den
"ubrigen Konstrukten in SEL (Variablen- und Strukturdefinition,
Parameterdefinition in Funktionen u.~s.~w.) pa"st. Es zeigt sich,
da"s Programme bei richtiger Verwendung des Typecasting-Operators
weit "ubersichtlicher sind als in C, meine Wahl hat sich also in
keiner Weise als Fehlgriff erwiesen.

\subsection{Mengen, Listen, B\protect{"aume} etc.}

Komplexere Typen, wie z.~B. Mengen, Listen oder B"aume sind in
SEL nicht vorhanden, aus guten Gr"unden. Mengen einerseits
werden relativ selten verwendet und lassen sich leicht
vermeiden. Au"serdem sind sie sehr leicht mit den Mitteln von
SEL zu implementieren, soda"s das Einbeziehen in SEL nur
"au"serst wenig Sinn h"atte. Listen und B"aume hingegen werden
zwar oft verwendet, sind allerdings so komplex und existieren in
einer derartig gro"sen Anzahl von Variationen, da"s es mit
einfachen Mitteln kaum m"oglich ist, sie vern"unftig zu
unterst"utzen.

\section{Bezeichner und Schl\protect{"usselw\protect{"orter}}}

Wie in vielen Programmiersprachen "ublich, sind auch in SEL
Schl"usselw"orter reserviert, d.~h., sie k"onnen nicht als
Bezeichner verwendet werden, aus zwei Gr"unden. Erstens wird
dadurch die Programmierung eines Interpreters bzw. Compilers
erheblich erschwert, und zweitens w"aren dadurch sehr
verwirrende Konstruktionen wie z.~B.

\begin{verbatim}
for (for = to to step step end)
  ...
end
\end{verbatim}

\noindent m"oglich, was nicht im Sinn des Designers sein kann.

Schl"usselw"orter und Bezeichner sind in SEL case-sensitiv,
d.~h., da"s zwei W"orter, die sich nur dadurch unterscheiden,
da"s ein oder mehrere Buchstaben statt gro"s klein geschrieben
sind oder umgekehrt, als unterschiedlich betrachtet werden. Das
hat drei Vorteile: Erstens wird dadurch die Entwicklung der
lexikalischen Analyse\footnote{siehe Kapitel \ref{scanner}}
ein wenig einfacher, zweitens k"onnen dadurch verschiedene Arten
von Bezeichnern (Bezeichner f"ur Funktionen, Variablen, Makros
u.~s.~w.) durch verschiedene Schreibweise gekennzeichnet werden.
Drittens wird durch eine vorgegebene Schreibweise bei
unerfahrenen oder schlampigen Programmierern die
"Ubersichtlichkeit der Quelltexte gesteigert.

Die Syntax der Bezeichner ist relativ simpel: Sie bestehen aus
einem Buchstaben oder Unterstrich, gefolgt von einer beliebigen
Anzahl von Buchstaben, Ziffern oder Unterstrichen. Der Grund,
warum Bezeichner nicht mit Ziffern beginnen d"urfen ist der, da"s
dadurch eine Eindeutigkeit der Sprache nicht mehr gew"ahrleistet
w"are. So k"onnte z.~B. \fu{12to} einerseits als Zahl \fu{12}
gefolgt vom Schl"usselwort \rw{to}, andererseits aber auch als
Bezeichner gedeutet werden. Dieser Konflikt lie"se sich zwar
durch Verwendung von Leerzeichen vermeiden, w"urde aber die
lexikalische Analyse ungleich schwerer machen. Warum ich nicht
auch andere Symbole in Bezeichnern erlaubt habe, ist einfach: Je
weniger Symbole in Bezeichnern verwendet werden d"urfen, desto mehr
hat man f"ur andere Zwecke frei.
                            
\section{Weiterf\protect{"uhrende} Projekte}
\label{langext}

Keine Programmiersprache ist perfekt, und schon gar nicht SEL,
f"ur dessen jetzige Form ich bewu"st die Versionsnummer 0.3
gew"ahlt habe. Leider war es mir aus Zeit- und Platzgr"unden
nicht m"oglich, auch nur einen Teil der von mir geplanten
Erweiterungen in diese Arbeit einzubringen, soda"s ich sie in
diesem Kapitel kurz zusammenfassen m"ochte, einerseits, um zu
zeigen, was SEL noch alles fehlt und andererseits, um den Leser
anzuregen, sich mit der Materie eingehender auseinanderzusetzen.

\begin{description}
\item[Zeiger auf Funktionen] stellen ein sehr praktisches
Sprachkonzept dar. Sie erm"oglichen es z.~B., in einer
graphischen Benutzeroberfl"ache, wie etwa Windows NT oder meinen
Schani-Tools, jedem Bildschirmfenster eine Funktion zuzuweisen,
die die Verwaltung desselben "ubernimmt. Sprachlich gesehen liegt
die einzige Schwierigkeit bei der Typangabe eines solchen
Funktionszeigers, denn der Aufruf einer Funktion durch einen
Zeiger kann genauso wie bei einer "`richtigen"' Funktion durch
den Namen der Zeigervariablen, gefolgt von einer in Klammern
gesetzten Parameterliste erfolgen, und die Adresse einer Funktion
kann mittels deren Namen, ohne nachfolgende Parameterliste in
Klammern, ermittelt werden. Doch auch der Typ f"ur einen
Funktionszeiger kann relativ leicht ohne Zweideutigkeiten in die
Grammatik eingebaut werden. So k"onnte etwa \fu{char(int,float)}
f"ur einen Zeiger auf eine Funktion, die einen \tint\ und einen
\tfloat\ als Parameter erwartet und einen \tchar\ zur"uckliefert,
stehen.\footnote{Tats"achlich f"uhrt eine Erweiterung wie oben
beschrieben zu einem kleinen Konflikt in der Grammatik, der
jedoch leicht beseitigt werden kann.}

\item[Variante Strukturen] sind Strukturen, die zwar mehrere
Elemente haben, diese jedoch denselben Speicherbereich
einnehmen. So etwas ist immer dann von Vorteil, wenn eine
Struktur ben"otigt wird, die Objekte beschreibt, die zwar in
gewisser Hinsicht "ahnlich sind, jedoch andere Parameter haben.
Ein Beispiel hierf"ur ist die in Kapitel \ref{intexprs}
besprochene Datenstruktur f"ur SEL-Ausdr"ucke.

\item[Makros mit Parametern] w"urden, wenn sie so flexibel
aufgebaut sind, da"s neue Befehle durch sie
"`vorget"auscht"' werden k"onnen, perfekt in eine Makrosprache
passen. Ein gutes Beispiel f"ur solche Makros sind die sog. User
Defines Commands (UDC's) von Clipper.

\item[Objektorientierte Programmierung] ist zwar ein edles Ziel,
jedoch nicht leicht zu verwirklichen.
\end{description}

\chapter{Der Interpreter}

Das Ziel dieses zweiten Kapitels ist es, f"ur die im ersten
Kapitel erkl"arte Programmiersprache SEL einen Interpreter zu
entwerfen. Der Titel dieser Arbeit deutet darauf hin, da"s dieser
Entwurf theoretisch ist, also in dieser Arbeit keine
Implementation eines Interpreters zu finden ist. Die Gr"unde
hierf"ur sind verschieden. Schon lange bevor ich "uberhaupt die
Idee zu dieser Fachbereichsarbeit hatte, wollte ich einen
Interpreter schreiben. Als ich mit Herrn Prof. Freiler "uber das
Thema dieser Arbeit diskutierte, legte er mir die Beschr"ankung
auf, da"s, wenn das Thema dieser Arbeit ein Programm w"are,
dieses in \pascal\ geschrieben sein m"u"ste. Da ich aber den
Interpreter in erster Linie f"ur mich und nicht f"ur diese Arbeit
programmieren wollte, hatte ich vor, ihn in C zu schreiben. Aus
diesem Grund einigte ich mich mit Herrn Prof. Freiler darauf,
da"s ich den Interpreter nur theoretisch, also unabh"angig von
irgendeiner Implementation, behandeln w"urde. Im Nachhinein
betrachtet bin ich froh "uber diese Entscheidung, da es kaum
m"oglich gewesen w"are, in einer im Umfang derma"sen beschr"ankten
Arbeit einen ganzen Interpreter und eine angemessene Erkl"arung
dazu unterzubringen.

Doch auch mit der Beschr"ankung auf eine theoretische Abhandlung
ist der Umfang dieser Arbeit zu gering, um ausf"uhrlich auf alle
Details einzugehen. Aus diesem Grund m"ochte ich im Gegensatz 
zum ersten Kapitel, in dem ich alle meine Entscheidungen 
begr"undet und evtl. auch andere L"osungsans"atze beschrieben 
habe, in diesem nur den Entwurf pr"asentieren und lediglich in
einigen seltenen F"allen andere m"ogliche L"osungen diskutieren.

Dieses Kapitel beschreibt alle wichtigen Algorithmen und
Datenstrukturen des Interpreters, geht jedoch nicht n"aher auf
grundlegende Datenstrukturen wie etwa dynamische Listen ein.
Sollte der Leser mit diesen Strukturen nicht vertraut sein,
empfehle ich ihm die Lekt"ure von \cite{sedge}, Kapitel 3--5 und
16.

\section{Der Begriff des Interpreters}
\label{whatisit}

Um ein Programm einer h"oheren Programmiersprache auf einem
Computer auszuf"uhren, kann man sich zweier verschiedener
Werkzeuge bedienen: Interpreter und Compiler. Der Begriff
"`Compiler"' wird vom Duden Informatik wie folgt definiert:

\begin{quotation}
Programm, das Programme aus einer Programmiersprache A in eine
Programmiersprache B "ubersetzt.\footnote{siehe \cite{duden},
Seite 745}
\end{quotation}

Das bedeutet, da"s ein Compiler ein Programm gar nicht direkt
ausf"uhrt, sondern es lediglich in eine andere Programmiersprache
"ubersetzt. In der Regel ist diese Sprache eine Assemler- oder
Maschinensprache, die direkt vom Computer "`verstanden"' wird. Im
Gegensatz dazu ist ein Interpreter

\begin{quotation}
Ein Programm, welches ein Programm einer anderen
Programmiersprache nach den notwendigen syntaktischen
"Uberpr"ufungen sofort ausf"uhrt. Im Gegensatz zu einem
"Ubersetzer\footnote{andere Bezeichnung f"ur Compiler} mu"s das
Quellprogramm nicht erst in eine andere Programmiersprache
"ubersetzt werden, sondern der Interpreter analysiert
nacheinander jede Anweisung und Deklaration des Quellprogramms
und f"uhrt diese unmittelbar aus.\footnote{siehe \cite{duden},
Seite 322}
\end{quotation}

In der Praxis jedoch wird kaum ein Interpreter ein Programm
ausf"uhren, ohne es vorher in eine geeignetere Form umgewandelt
zu haben. Man kann sich leicht vorstellen, da"s ein Interpreter,
der etwa bei einem Funktionsaufruf alle Module nach der Deklaration
dieser Funktion durchsuchen mu"s, zu langsam ist, um sinnvoll
eingesetzt werden zu k"onnen. Auch das wiederholte Lesen und
Erkennen der Schl"usselw"orter, die im Quelltext verwendet
werden, ist zu zeitaufwendig, als da"s es w"ahrend der
Ausf"uhrung des Programms durchgef"uhrt werden k"onnte. In den
folgenden Kapiteln werden wir deshalb besprechen, wie sich ein
Interpreter die Arbeit so einteilen kann, da"s eine m"oglichst
hohe Ausf"uhrungsgeschwindigkeit erreicht wird.

Ich m"ochte anmerken, da"s das in diesem Kapitel besprochene
Design f"ur einen Interpreter zwar theoretisch, jedoch keineswegs
realit"atsfern ist. Ich habe selbst einen SEL-Interpreter auf
diese Weise implementiert.

\section{Einsatzgebiete eines Interpreters}

Wenn man bedenkt, da"s die Ausf"uhrung eines Programms durch
einen Interpreter oft mehrere hundert Male l"anger dauert als die
eines mit einem Compiler erzeugten Maschinenspracheprogramms, fragt
man sich, welchen Zweck ein Interpreter eigentlich erf"ullen
kann. Ein Haupteinsatzgebiet f"ur Interpreter sind Computer, die
f"ur die sinnvolle Verwendung eines Compilers zu langsam sind
und/oder zu wenig Speicher bieten. Beispiele hierf"ur sind alte
Home-Computer wie z.~B. der Commodore~C-64 aber auch
Taschenrechner wie etwa der Texas Instruments~TI-85 oder der
Hewlett Packard~HP-48. Oft werden Interpreter auch zum
Entwickeln von Programmen eingesetzt, da sie den Vorteil haben,
Programme ohne Zeitverz"ogerung ausf"uhren zu k"onnen. Ein
anderes wichtiges Einsatzgebiet f"ur Interpreter ist die
Steuerung von Me"sger"aten und Robotern. Dort stellt die langsamere
Ausf"uhrung von Programmen kaum einen Nachteil dar, andererseits
ist es ein gro"ser Vorteil, da"s Programme ge"andert werden
k"onnen und sofort wieder einsatzbereit sind.

\section{Bestandteile eines Interpreters}

Bevor wir uns daranmachen, uns dar"uber zu unterhalten, wie der
Interpreter aussehen und welche Leistung er bieten soll, m"ussen
wir "uber den groben Aufbau eines Interpreters Bescheid wissen.
Wie bereits oben erw"ahnt wird in der Praxis das Quellprogramm
vor dem Interpretieren in eine f"ur den Interpreter geeignete
Form umgewandelt ({\em Tokenizing}). Dieser Vorgang wird in der
Regel von der lexikalischen Analyse erledigt. Dabei wird der
Quellcode sequentiell abgearbeitet und Schl"usselw"orter und
Bezeichner codiert. Zus"atzlich werden die Positionen von wichtigen
Schl"usselw"ortern wie z.~B. \rw{function}, \rw{global} und
\rw{local} gespeichert. Im Fall von SEL wird gleichzeitig
zur lexikalischen Analyse der Pr"aprozessor ausgef"uhrt. Er hat die
Aufgabe, Include-Dateien einzubinden und Makros in den
Programmtext einzuf"ugen.

Im Anschlu"s an die lexikalische Analyse und das Tokenizing kann
mit der Ausf"uhrung des Programms begonnen werden. Bei einer
Programmiersprache wie SEL m"ussen zuerst nat"urlich spezielle
Initialisierungen vorgenommen werden, n"amlich das Einrichten der
modul- und programmglobalen Variablen und das Analysieren der
Funktionsk"opfe. Zus"atlich speichert der Interpreter die
Stellen, an denen die jeweiligen Funktionen beginnen, soda"s
nicht das gesamte Programm durchsucht werden mu"s. Jetzt wird das
Programm durch einen Aufruf der \fu{main}-Funktion gestartet. Die
einzelnen Anweisungen und Ausdr"ucke werden vom Parser analysiert
und sofort ausgef"uhrt. Dabei mu"s nat"urlich auf spezielle
Routinen, die z.~B. f"ur die Verwaltung der Variablen und
"ahnliches verantwortlich sind, zur"uckgegriffen werden.

Zusammenfassend l"a"st sich also sagen, da"s ein Interpreter aus
drei wichtigen Teilen besteht: Der lexikalischen Analyse, die das
Tokenizing vornimmt, dem Parser, der das Programm syntaktisch
analysiert und den Verwaltungsroutinen, die die Hintergrundarbeit
leisten.

\section{Zielsetzung f\protect{"ur} das Design}
\label{intgoals}

Wie schon im ersten Kapitel erw"ahnt, ist SEL mit dem Ziel, sie
als Makrosprache zu verwenden, entworfen worden. Es w"are also
angebracht, den Interpreter von Beginn an so zu entwerfen, da"s
er nicht ein eigenst"andiges Programm darstellt, sondern eher
eine Art Modul, das man in andere Programme aufnehmen kann. Das
setzt nat"urlich voraus, da"s er nicht eine Unmenge an Speicher
verbraucht und auch, da"s er den Speicher, den er verbraucht,
nach der Ausf"uhrung des SEL-Programms wieder freigibt, soda"s er
wieder dem Hauptprogramm zur Verf"ugung steht.

Der Interpreter sollte nat"urlich, um ohne Bedenken in einem
Programm eingesetzt werden zu k"onnen, Programmierfehler so
schnell wie m"oglich erkennen k"onnen, den Ausf"uhrungsproze"s
angemessen stoppen, allen Speicher, den das SEL-Programm
angefordert hat, und auch den, den der Interpreter im Zuge der
Ausf"uhrung des Programms benutzt hat, freigeben und den
Fehler dem Hauptprogramm melden.

Wie auch die Programmiersprache SEL sollte auch der Interpreter
erweiterungsf"ahig sein, d.~h., er sollte dem Hauptprogramm
erlauben, bestimmte Funktionen auch SEL-Programmen zug"anglich zu
machen, wie beispielsweise f"ur das Zeichnen einer Linie in einem
Zeichenprogramm.

Da ich mich in letzter Zeit mit dem Multitasking-Betriebssystem
Windows NT auseinandergesetzt habe, wollte ich auch meinen
Interpreter multitaskingf"ahig machen, soda"s es in einem
Programm m"oglich ist, mehrere SEL-Programme gleichzeitig
auszuf"uhren. Um mich nicht falsch zu verstehen: Das Multitasking
wird in keiner Weise vom Interpreter in die Hand genommen,
sondern obliegt einzig und allein dem Betriebssystem. Der
Interpreter allerdings mu"s so entworfen werden, da"s er das
Multitasking nicht erschwert oder gar unm"oglich macht, was, wie
sich zeigen wird, nicht sehr schwer ist.

Im Laufe dieses Kapitels werde ich auf einige, aber nicht auf alle
Forderungen dieses Unterkapitels eingehen, weil sie zu trivial
sind. So ist z.~B. die Freigabe des benutzten Speichers nur eine
Formsache und mu"s nicht extra von mir erl"autert werden.

\section{Die lexikalische Analyse}
\label{scanner}

Die lexikalische Analyse hat die Aufgabe, den Zeichenstrom des
Quelltextes in einen Tokenstrom zu verwandeln. Als Token
bezeichnet man die kleinste Einheit, die vom Parser bearbeitet
wird. Token sind z.~B. die Schl"usselw"orter, Bezeichner,
Konstanten u.~s.~w. Der Grund, warum man die Erkennung von
Schl"usselw"ortern etc. nicht auch dem Parser "uberl"a"st, was
durchaus m"oglich w"are, ist der, da"s ein Parser diese Aufgabe
bei weitem nicht so schnell ausf"uhren kann wie ein eigens daf"ur
angefertigter {\em Scanner}.\footnote{Begriff f"ur ein Programm
oder einen Programmteil, der die lexikalische Analyse
durchf"uhrt} Au"serdem wird der Bau eines Parsers wesentlich
vereinfacht, wenn man einen Scanner verwendet.

Um die Arbeitsweise eines Scanners zu beschreiben, m"ochte ich
zun"achst ein Beispiel heranziehen. Nehmen wir an, der Scanner
h"atte den Zeichenstrom \fu{for (x=1 to 10)} zu bearbeiten. Er
untersucht nun das erste Zeichen und merkt, da es sich um einen
Buchstaben handelt, da"s er es entweder mit einem Bezeichner oder
mit einem Schl"usselwort zu tun hat. Die einfachste Methode,
dieses Token zu analysieren ist, so lange Buchstaben, Ziffern und
Unterstriche zu lesen, bis er an einem Zeichen ankommt, das
nicht mehr aus einer der genannten Kategorien stammt. Der
Scanner w"urde also zuerst das \fu{f}, dann das \fu{o} und
anschlie"send das \fu{r} lesen. Da das n"achste Zeichen weder
Buchstabe, Ziffer noch Unterstrich ist, stoppt er hier. Die
einfachste Methode, herauszufinden, ob man es hier mit einem
Schl"usselwort oder einem Bezeichner zu tun hat, ist, eine
Tabelle anzulegen, die alle Schl"usselw"orter der Sprache
umfa"st. Findet der Scanner das Wort in der Tabelle, ist es ein
Schl"usselwort, andernfalls ein Bezeichner. Das n"achste Zeichen
ist ein Leerzeichen. In einer formlosen Programmiersprache wie
SEL werden Leerzeichen und Tabulatoren "ubersprungen, soda"s der
Scanner nichts anderes zu tun hat, als zum n"achsten Zeichen
vorzur"ucken. Zwar sind Zeilenumbr"uche f"ur den Parser nicht
relevant, doch f"ur den Pr"aprozessor sind sie wichtig, soda"s
sie gemeldet werden sollten.

Da es nur ein Symbol in SEL gibt,
das mit einer "offnenden, runden Klammer beginnt, n"amlich dieses
selbst, wei"s der Scanner sofort, mit welchen Token er es zu tun
hat und kann vorr"ucken. Interessant wird es nun bei der Analyse
des Symbols \fu{=}, denn hier kann der Scanner nicht sofort davon
ausgehen, da"s es sich um den Zuweisungsoperator handelt, da es
noch ein anderes Symbol, n"amlich \fu{==} gibt, das mit einem
\fu{=} anf"angt. So mu"s der Scanner auch hier das n"achste
Zeichen betrachten um entscheiden zu k"onnen, welches Symbol er
vor sich hat.

Wie dieses Beispiel zeigt, ist es f"ur einen
Scanner vital zu wissen, welches Zeichen als n"achstes vorliegt,
wobei er sich dieses aber auch f"ur das n"achste Token aufheben
k"onnen mu"s. Die L"osung ist einfach: Man arbeitet mit einem
Zeichen im voraus, d.~h. es wird immer das kommende Zeichen
betrachtet, nicht das aktuelle. Wird es akzeptiert, liest man das
n"achste Zeichen ein, wenn nicht, kehrt man mit dem erkannten
Token zur aufrufenden Funktion (in Falle von SEL der
Pr"aprozessor) zur"uck. Beim n"achsten Aufruf des Scanners wird
mit dem vorher zur"uckgestellten Zeichen weitergearbeitet.

Um den Proze"s der lexikalischen Analyse zus"atzlich zu
beschleunigen, sollte die Eingabe gepuffert werden. Dabei wird
ein Teil der Eingabe in einen Zwischenspeicher (Puffer) gelesen,
um von dort aus vom Scanner abgearbeitet zu werden. Wird das Ende
des Puffers erreicht, mu"s abermals ein Teil der Eingabe gelesen
werden, so lange, bis das Ende der Datei erreicht ist. In diesem
Fall sollte der Scanner der "ubergeordneten Funktion Bescheid
geben.

Bei bestimmten Tokentypen gen"ugt es nicht, zu melden, da"s
dieses Token erkannt wurde, sondern es mu"s auch ein Attribut
zur"uckgeliefert werden, im Falle eines Bezeichners sein Name, im
Fall einer \tint -Konstante ihr Wert u.~s.~w., es sollte also
f"ur ein Token eine geeignete Datenstruktur verwendet werden, in
der alle n"otigen Attribute Platz finden. F"ur unseren
Interpreter k"onnte diese Struktur etwa so aussehen:

\newpage
\begin{verbatim}
struct Token
  iType  : int;
  pcName : *char;
  cChar  : char;
  iInt   : int;
  fFloat : float;
end
\end{verbatim}

Ist das Token ein Bezeichner, enth"alt \fu{pcName}
dessen Namen, im Falle einer \tchar -Konstanten enth"alt
\fu{cChar} ihren Wert u.~s.~w.

\subsection{Der Pr\protect{"aprozessor}}

Der Pr"aprozessor ist in gewisser Weise eine Zwischenstation von
Scanner und Parser. In SEL hat er die Aufgabe, den Inhalt von
Include-Dateien in den Programmtext einzuf"ugen, Kommentare zu
entfernen und Makros zu bearbeiten.

Das Entfernen von Kommentaren ist im Grunde kein Problem. Wird
vom Scanner das Symbol \fu{//} erkannt, dann werden so lange
Token gelesen, bis ein Zeilenumbruch gemeldet wird, wonach der
Kommentar beendet ist. Zus"atzlich zu den Kommentaren sollten
auch alle Zeilenumbr"uche herausgefiltert werden, da sie den
Parser st"oren w"urden.

Etwas komplizierter ist das Einf"ugen von Include-Dateien. Da es
m"oglich sein sollte, Include-Dateien zu verschachteln, gibt es
zwei M"oglichkeiten zur Implementation. Entweder, der
Pr"aprozessor arbeitet rekursiv, soda"s er sich selbst nochmals
aufruft, wenn eine Include-Datei zu bearbeiten ist, oder aber er
benutzt einen Stack zur Zwischenspeicherung der einzelnen
Include-Level, was am Ende eigentlich auf dasselbe hinausl"auft.
Ein kleines Problem ergibt sich mit dem Puffer f"ur den
Scanner: Wenn der Pr"aprozessor nur einen einzigen Puffer f"ur alle
Include-Level zusammen benutzt, mu"s er beim Eintritt in eine
Include-Datei geleert und mit
der Eingabe aus der Include-Datei gef"ullt werden. Ist diese zu Ende,
mu"s der Puffer abermals geleert werden und mit dem
"ubergeordneten Level gef"ullt werden. Das Problem dabei ist,
da"s man in der Datei dabei ein wenig zur"ucksetzen mu"s, weil
ein Teil des Puffers gel"oscht wurde. Ist das aus
irgendwelchen Gr"unden nicht m"oglich oder zu ineffizient,
mu"s man f"ur jeden einzelnen Include-Level einen Puffer
benutzen, was nat"urlich mehr Speicher verbraucht.

Zum Bearbeiten von Makros mu"s eine Datenstruktur
angelegt werden, in der alle Makros aufgezeichnet und auch schnell wieder
gefunden werden k"onnen. Ein bin"arer Baum und eine Hash-Tabelle
bieten sich hierzu an. F"ur jeden einzelnen Eintrag in dieser
Struktur mu"s zus"atzlich noch gespeichert werden, durch welchen
Text er zu ersetzen ist. Da der Pr"aprozessor auf den Scanner
aufbaut, ist es einfacher und effizienter, nicht den Zeichenstrom
zu speichern, sondern den Tokenstrom, der vom Scanner geliefert
wird. Dazu verwendet man am einfachsten eine einfach verkettete
Liste aus Token:

\begin{verbatim}
struct TokenList
  tokToken : Token;
  Next     : *TokenList;
end
\end{verbatim}

Trifft der Pr"aprozessor nun auf einen Bezeichner, mu"s er zun"achst
"uberpr"ufen, ob nicht vielleicht ein Makro unter diesem Namen
exisitiert. Dazu sucht er
nach ihm in der Datenstruktur. Ist er dort vorhanden,
setzt der Pr"aprozessor anstatt seiner die Tokenliste ein,
andernfalls wird er wie ein normaler Bezeichner behandelt.

Nachdem es m"oglich sein soll, innerhalb eines Makros ein anderes
zu benutzen, m"ussen wir noch etwas einbauen. Die einfachste
L"osung ist, da"s wir zum Einf"ugen eines Tokens in unseren
Zwischencode eine eigene Funktion verwenden, die, falls sie auf
einen Bezeichner trifft, der als Makro definiert ist, sich f"ur
jedes der Token aus der Liste selbst wieder aufruft. Dadurch ist
garantiert, da"s, falls in dieser Liste ein Makro auftaucht,
auch dieses entsprechend behandelt wird.
                              
\subsection{Der Zwischencode}
\label{tokencode}

Wie schon oben erw"ahnt, verwenden wir eine Art Zwischencode, um
dem Interpreter eine schnelle Ausf"uhrung zu erm"oglichen. Dabei
werden Schl"usselw"orter und Symbole mit einer eindeutigen Nummer
codiert, soda"s der Tokentyp sofort ersichtlich ist. Bei
bestimmten Tokentypen m"ussen zus"atzliche Informationen
gespeichert werden, wie z.~B. bei Konstanten oder Bezeichnern.

Um die Gr"o"se des Zwischencodes nicht zu sehr aufzubl"ahen,
werden auch die Namen der Bezeichner mit einer Nummer
codiert. Dazu wird eine Tabelle aller im Modul vorkommenden
Bezeichner angelegt und jedem Eintrag eine eindeutige Nummer
zugewiesen. Diese Nummer wird anschlie"send im Zwischencode
verwendet, um Bezeichner zu charakterisieren. Zu beachten ist,
da"s nat"urlich auch der Parser "uber diese Tabelle verf"ugen
mu"s, da er sonst nur die Nummern der Bezeichner wei"s. Wird
allerdings in allen Modulen des Programms dieselbe Tabelle
verwendet, so kann der Interpreter das Programm auch ohne das
Wissen um die Namen der Bezeichner ausf"uhren.

Eine besondere Hilfeleistung erstattet der Zwischencode dem
Interpreter bei Schleifen und Kontrollanweisungen: F"ur jedes \rw{if}
wird zus"atzlich die Position im Zwischencode abgespeichert, an der das
zugeh"orige \rw{end} bzw. \rw{else} steht. Das hat den Vorteil,
da"s der Interpreter, wenn er z~B. auf eine \rw{if}-Anweisung
trifft, deren Bedingung logisch falsch ist, ohne lange zu suchen
zum zugeh"origen \rw{end} bzw. \rw{else} springen kann, was
verst"andlicherweise die Ausf"uhrungsgeschwindigkeit erheblich
steigert, besonders bei Schleifen.

Wie bereits erw"ahnt, werden dem Interpreter zus"atzlich die
Position aller Vorkommnisse der Schl"usselw"orter \rw{global},
\rw{local} und \rw{function} mitgeteilt, soda"s er ohne zu suchen die
Initialisierung der programm- und modulglobalen Variablen und der
Funktionsk"opfe vornehmen kann.

Neben der erheblichen Geschwindigkeitssteigerung, die der
Zwischencode erm"oglicht, hat er noch den Vorteil, da"s er
weniger Speicherplatz als der Quellcode ben"otigt, bei etwas
gr"osseren Programmen etwa die H"alfte, was besonders unter
MS-DOS entgegenkommend ist.

Ich m"ochte hier kurz noch darauf eingehen, wie die oben
angesprochenen Verweise von Anweisungen auf die zugeh"origen
\rw{end}s bzw. \rw{else}s erzeugt werden k"onnen. Diese auf den
ersten Blick umst"andliche Aufgabe l"a"st sich relativ leicht
l"osen, indem man einen Stack zur Speicherung der
Verschachtelungen der Anweisungen verwendet. Trifft der Scanner z.~B.
auf ein \rw{if}, dann pusht er die Position, an der er sich
gerade im Zwischencode befindet, auf den Stack und l"a"st Platz
f"ur den sp"ater folgenden Verweis frei. Als n"achstes trifft
er z.~B. auf ein \rw{for} und tut dasselbe. Irgendwann mu"s
er jedoch zu dem \rw{end} gelangen, da"s die jeweils letzte
Anweisung schlie"st. In diesem Fall speichert er das \rw{end} im
Zwischencode, poppt vom Stack die Position der letzten
Anweisung und f"ugt an dieser Stelle die Position ein, an der
er soeben das \rw{end} gespeichert hat. Diese Technik wird
"ubrigens {\em Backpatching} genannt und oft im Compilerbau
eingesetzt.\footnote{Tats"achlich ist unser Scanner, der
einen SEL-Quelltext in Zwischencode umwandelt, streng gesehen
auch ein Compiler.}

\bigskip

Eine effiziente und einfache Darstellung des Zwischencodes ist
die Verwendung jeweils eines Bytes f"ur ein Token, wobei bei
bestimmten Tokentypen, wie z.~B. dem \rw{if}-Token, ein
zus"atzlicher Verweis auf ein anderes Token eingebaut wird. Wie
viele Bytes f"ur diesen Verweis verwendet werden, h"angt von den
Anforderungen an den Interpreter ab. Ich verwende hierzu in
meinem Interpreter 32 Bits, soda"s der Zwischencode theoretisch
bis zu 4 GByte gro"s sein kann, was sicherlich f"ur alle Anwendungen
mehr als ausreichend ist. F"ur die meisten
Anwendungen w"urden auch 16 Bits reichen, da Module in der Regel
sowieso klein gehalten werden. Bei anderen Token wie etwa einer
String-Konstante mu"s au"ser dem Tokentyp noch ein Wert
gespeichert werden, wie in diesem Fall der String.

\section{Der Parser}
\label{parser}

Der wohl komplizierteste Teil des Interpreters ist der
Parser. Er hat zur Aufgabe, aus dem Tokenstrom des Zwischencodes
die einzelnen Anweisungen zu erkennen. Sehr wichtig f"ur den Bau
eines Parsers ist die Grammatik, denn sie gibt an, was der Parser
eigentlich erkennen soll. Aus verschiedenen Gr"unden werden wir
Parser schreiben, die Teile der Sprache SEL erkennen k"onnen,
n"amlich Typen, Ausdr"ucke und Folgen von Anweisungen. Die Parser
f"ur Funktionsk"opfe und \rw{global}- bzw. \rw{local}-Bl"ocke
sind trivial, soda"s ich hier nicht speziell auf sie eingehen
werde. Ein Grund f"ur diese Trennung ist, da"s der
SEL-Interpreter vor allem f"ur die Verwendung in anderen
Programmen gedacht ist und dort ist es in vielen Situationen von
Vorteil, wenn der Benutzer Ausdr"ucke eingeben kann, z.~B. als
Filterbedingung in einer Datenbank.
                              
\subsection{Der Typen-Parser}

Beim Design eines Parsers ist es zun"achst wichtig, "uber die
Grammatik Bescheid zu wissen, die der Parser erkennen soll. In
diesem Unterkapitel besch"aftigen wir uns mit einem sehr
einfachen Teil der Grammatik von SEL, n"amlich den Typen. Das
gibt uns die M"oglichkeit, uns "uber die verschiedenen Arten von
Parsern und "uber die Zusammenh"ange zwischen Grammatik und
Parser zu informieren. Der "Ubersichtlichkeit halber sei an
dieser Stelle die Grammatik dargestellt:

\begin{verbatim}
type:          `char'
             | `int'
             | `float'
             | `void'
             | `[' INT_CONSTANT `]' type
             | `*' type
             | IDENTIFIER
             ;
\end{verbatim}

Diese Grammatik hat sieben Produktionen, das Startsymbol
ist \fu{type}.

\subsubsection{Top-Down und Bottom-Up Parser}
\label{topdbotu}

Im Interpreter- und Compilerbau werden haupts"achlich zwei
verschiedene Arten von Parsern verwendet: Top-Down und Bottom-Up
Parser. W"ahrend Top-Down Parser beim Startsymbol beginnen und
immer weiter hinabsteigen, bis sie auf die Token treffen, fangen
Bottom-Up Parser bei den Token an und arbeiten sich bis zum
Startsymbol hinauf. Das l"a"st sich am besten
verdeutlichen, indem man versucht, einen Parse-Tree zu
konstruieren, also einen Baum, der f"ur einen g"ultigen
Ausdruck der Sprache, in diesem Fall etwa \fu{*[10]Test}, den
Weg vom Startsymbol bis hinunter zu den einzelnen Token,
verdeutlicht. Der Parse-Tree f"ur den Beispieltyp ist in
Abbildung \ref{parsetree} dargestellt.

\begin{figure}

\centerline{\includegraphics[height=4cm]{partree1.eps}}

\caption{Der Parse-Tree f"ur \fu{*[10]Test}}
\label{parsetree}
\end{figure}

Die Wurzel des Baumes ist immer das Startsymbol, in diesem Fall
also \fu{type}. Dieser \fu{type} kann auf der ersten Ebene
zerlegt werden in das Sternsymbol und einen zweiten \fu{type},
gem"a"s der Produktion \fu{`*' type}. Dieser zweite \fu{type}
kann zelegt werden in eine "offnende und eine schlie"sende
Klammer und eine \fu{INT\_CONSTANT} dazwischen, gefolgt von noch
einem \fu{type}. Dieser wiederum ist ein \fu{IDENTIFIER}, n"amlich
\fu{Test}.

Nun zu den Vorgehensweisen der beiden Parsertypen. Ein Bottom-Up
Parser arbeitet sich wie gesagt von den einzelnen Token zum
Startsymbol hinauf. Dazu benutzt er einen Stapel, auf dem er zwei
verschiedene Aktionen durchf"uhren kann: {\em Shift} und
{\em Reduce}. Aus diesem Grund werden Bottom-Up Parser auch oft
als Shift-Reduce Parser bezeichnet. Die Analyse des
Beispieltyps f"uhrt der Parser so durch: Zun"achst liest er den
Stern, den er shiftet. Auch die beiden Klammern und die
\fu{INT\_CONSTANT} werden von ihm geshiftet. Nun liest er einen
Bezeichner und entscheidet, diesen zu einem Typen zu reducen,
gem"a"s der Produktion \fu{IDENTIFIER}. M"u"ste er einen
Parse-Tree erstellen, h"atte er an dieser Stelle einen Baum
mit der Wurzel 3\footnote{siehe Abbildung \ref{parsetree}}
erzeugt. Auf
dem Stack des Parsers liegen nun die Token \fu{`*'}, \fu{`['},
\fu{INT\_CONSTANT}, \fu{`]'} und das Nonterminal \fu{type}, in der
angegebenen Reihenfolge. An dieser Stelle entscheidet der Parser,
die letzten vier Symbole zu reducen. Um einen Parse-Tree zu
erhalten, m"u"ste er nun Baum 3 verwenden, um den gr"o"seren
Baum 2 zu konstruieren. Zuletzt w"urde der Parser nochmals
reducen und den Baum 1 erstellen. Der Parser hat nun
seine Analyse beendet, da das einzige, was auf seinem Stack
liegt, das Startsymbol ist. Die Hauptaufgabe bei der Erstellung
eines Bottom-Up Parsers ist, zu entscheiden, wann er zu shiften
und wann zu reducen hat. Dies ist eine schwere Aufgabe und wird
in der Praxis von einem Parser-Generator (auch Compiler-Compiler
genannt) wie etwa YACC erledigt. Da wir jedoch unseren Parser
selbst programmieren wollen, beschr"anken wir uns auf die
wesentlich einfacheren Top-Down Parser.

Ein Top-Down Parser geht bei der Analyse von \fu{*[10]Test}
folgenderma"sen vor: Er erkennt das Sternsymbol und wei"s nun,
da"s er nach der Produktion \fu{`*' type} vorzugehen hat. An
dieser Stelle w"urde er den Baum 1 erzeugen. Er liest nun den
Stern ein und mu"s nun einen weiteren Typ analysieren. Hier wird
klar, da"s ein Top-Down Parser in der Regel rekursiv
implementiert wird. Zwar ist es m"oglich, diese Rekursion mit
einem Stack zu beseitigen, doch ein solcher Parser ist schwerer
zu programmieren und zu warten, soda"s ich eine rekursive
Implementation gew"ahlt habe.

Der Parser ruft sich nun selbst auf und erkennt anhand der
"offnenden eckigen Klammer, da"s er die Produktion
\fu{`[' INT\_CONSTANT `]' type} w"ahlen mu"s und erzeugt
gleichzeitig Baum 2. Nacheinander liest er nun die drei folgenden
Token und ruft sich abermals selbst auf. Nun wendet er die
Produktion \fu{IDENTIFIER} an und erstellt Baum 3. Damit ist die
Analyse beendet.

\subsubsection{Look-Aheads}

Wie sich in diesem Beispiel gezeigt hat, ist es f"ur einen Parser wichtig, zu
wissen, welches Token er als n"achstes zu bearbeiten hat, denn
danach hat er zu entscheiden, welche Produktion er ausw"ahlt. Zu
diesem Zweck verwendet man sog. {\em Look-Aheads}. F"ur einen
SEL-Parser und f"ur fast alle Programmiersprachen ist ein Look-Ahead
v"ollig ausreichend, es gen"ugt also, ein Token
vorauszusehen. Die Implementation eines Look-Ahead ist simpel:
Anstatt nur des aktuellen Tokens verwendet der Parser zus"atzlich das
Look-Ahead. Soll nun das aktuelle Token "ubersprungen werden,
weist er einfach den Wert des Look-Aheads dem aktuellen Token
zu und liest einen neuen Look-Ahead aus dem Zwischencode ein.

\subsubsection{Die Konstruktion des Parsers}

Mit diesen Informationen ist es nun ein Leichtes, einen
rekursiven Top-Down Parser f"ur Typen zu bauen. Er besteht aus
einer einzigen Funktion, die durch Betrachten des Look-Ahead
entscheidet, was passieren soll. Ist der Look-Ahead eines der
Token \fu{`char'}, \fu{`int'}, \fu{`float'}, \fu{`void'} oder
\fu{IDENTIFIER}, mu"s der Parser dieses Token nur lesen und die Analyse
ist beendet. Im Fall einer "offnenden eckigen Klammer liest er
drei Token und ruft sich anschlie"send selbst auf.\footnote{Es
ist leicht m"oglich, eine am Ende stehende Rekursion durch einen
Sprung an den Beginn der Funktion, z.~B. durch eine
\rw{repeat}-Schleife, evtl. durch einen Stack unterst"utzt, zu
ersetzen, was im Normalfall die bessere L"osung ist.} Analog dazu
ist die Produktion \fu{`*' type} zu implementieren. Sollte der
Look-Ahead keinem der Token entsprechen, hat es der Parser
offensichtlich mit einem Fehler zu tun, den er entsprechend
behandeln mu"s. Ein Interpreter wird in diesem Fall eine
Fehlermeldung ausgeben und die Ausf"uhrung beenden.

Nat"urlich ist klar, da"s unser Parser nicht nur die Aufgabe hat,
einen Typ zu erkennen, sondern auch, mitzuteilen, um welchen Typ
es sich dabei handelt. Das jedoch f"allt unter die semantische
Analyse und wird in Kapitel \ref{inttypes} behandelt.

\subsection{Der Ausdruck-Parser}
\label{exprparser}

Nachdem wir nun die grundlegenden Funktionen eines Parsers
kennen, wollen wir einen etwas komplexeren Parser konstruieren:
Den Parser f"ur SEL-Ausdr"ucke. Die Grammatik hierzu kann im
Anhang \ref{grammar} oder in \cite{selkr}, Seite 14 nachgelesen
werden. Ausdr"ucke in SEL werden in der Grammatik durch die
beiden Nonterminals \fu{expr} und \fu{primary} repr"asentiert. In
den Produktionen f"ur \fu{primary} taucht noch das dritte
Nonterminal \fu{type} auf, doch dieses stellt hier kein
Problem mehr dar, da der entsprechende Parser bereits im letzten
Unterkapitel skizziert wurde.

Besonders interessant sind die Produktionen f"ur \fu{expr}, denn
sie bestehen gr"o"stenteils aus dem Nonterminal \fu{expr} selbst.
Wie sich leicht ersehen l"a"st, ist diese Grammatik auf keinen
Fall eindeutig, was die beiden Parse-Trees f"ur den Ausdruck
\fu{a=b+1} in Abbildung \ref{twopartrees}
verdeutlichen.\footnote{In einer eindeutigen Grammatik darf es
f"ur einen korrekten Satz nur einen einzigen Parse-Tree geben.}

\begin{figure}

\centerline{\includegraphics[height=3cm]{partree2.eps}}

\caption{Zwei Parse-Trees f"ur \fu{a=b+1}}
\label{twopartrees}
\end{figure}

Alle diese Konflikte jedoch lassen sich durch die in
\cite{selkr},~Seite~21 angegebene Priorit"atstabelle l"osen. Die
Frage, die sich nun stellt, ist, wie man diese Rangfolge in den
Parser einbringen kann. Betrachten wir der Einfachheit halber
nur die folgenden drei Produktionen von \fu{expr}:

\begin{verbatim}
expr:   expr `=' expr
      | expr `+' expr
      | expr `*' expr
      | primary
      ;
\end{verbatim}

Die Priorit"atstabelle sagt, da"s die h"ochste Priorit"at der
Operator \fu{*} hat, danach folgen \fu{+} und~\fu{=}. Der
Konflikt kann nun gel"ost werden, indem man f"ur jede Priorit"atsstufe ein
eigenes Nonterminal einf"uhrt, wobei die Produktionen f"ur die
Operatoren niedriger Priorit"at auf die Nonterminals der
Operatoren h"oherer Priorit"at zur"uckgreifen, wie hier
dargestellt:

\begin{verbatim}
expr3:   expr2 `=' expr2
       | expr2
       ;
expr2:   expr1 `+' expr1
       | expr1
       ;
expr1:   primary `*' primary
       | primary
       ;
\end{verbatim}

Diese Grammatik kann, angewendet auf den Ausdruck \fu{a=b+1},
nur einen Parse-Tree erzeugen, n"amlich den, der der
Klammerung des Ausdrucks \fu{a=(b+1)} entspricht.
Ungl"ucklicherweise geht mit dieser Grammatik die F"ahigkeit
verloren, einen Ausdruck wie \fu{a+b+c} darzustellen. Um auch
diesen Mangel zu beseitigen, ersetzt man in jeder Produktion mit
einem Operator die linke bzw. rechte Seite mit dem Nonterminal
des jeweiligen Levels, je nachdem, ob der Operator links- oder
rechtsassoziativ ist. Daraus resultiert folgende Grammatik:

\begin{verbatim}
expr3:   expr2 `=' expr3
       | expr2
       ;
expr2:   expr2 `+' expr1
       | expr1
       ;
expr1:   expr1 `*' primary
       | primary
       ;
\end{verbatim}

Diese Grammatik ist eindeutig und erzeugt sowohl f"ur den
Ausdruck \fu{a=b+1} als auch f"ur \fu{a+b+c} die richtigen
Parse-Trees, man kann also mit dieser Methode die "`gro"se"'
Grammatik f"ur den Bau des Anweisungs-Parsers umwandeln. Ich m"ochte die
sich daraus ergebende Grammatik in dieser Arbeit nicht
auff"uhren, da sie sehr leicht hergeleitet werden kann und zu
viel Platz beanspruchen w"urde.

Mit den bisher besprochenen Techniken kann man bereits einen
Parser f"ur einen Teil der Grammatik konstruieren, sp"atestens
bei den Linksrekursionen, wie z.~B. in den
Produktionen f"ur \fu{expr2} in der obigen Grammatik, wird man
allerdings Schwierigkeiten bekommen. Eine Rechtsrekursion wie etwa bei
\fu{expr3} ist keine Schwierigkeit: Man kann sie entweder
durch eine echte Rekursion implementieren oder durch eine
Schleife, unterst"utzt durch einen Stack. F"ur die erste Produktion
von \fu{expr2} m"u"ste man jedoch eine Funktion programmieren,
die als erstes sich selbst aufruft, was nat"urlich in einer
endlosen Rekursion m"unden w"urde.

Um dieses Problem zu l"osen, versuchen wir zu ergr"unden, welche
Menge von Ausdr"ucken eigentlich die beiden Produktionen von
\fu{expr2} beschreiben. Der einfachste Ausdruck, den die
Grammatik beschreiben kann, ist offensichtlich \fu{expr1}. Er
l"a"st sich mit Hilfe der beiden Produktionen nicht weiter
expandieren. Benutzt man als Basis nun die erste Produktion,
kann man den ersten \fu{expr2} expandieren, entweder
durch die erste oder die zweite Produktion. Benutzt man die
zweite Produktion, bekommt man als Ergebnis den Ausdruck
\fu{expr1 `+' expr1}, der sich abermals nicht weiter expandieren
l"a"st. F"uhrt man diesen Proze"s fort, wird man bald merken,
da"s diese beiden Produktionen eine Abfolge von beliebig vielen
\fu{expr1}, verbunden jeweils durch ein \fu{`+'}, beschreiben.
Zwar beschreiben die Produktionen f"ur \fu{expr3} im Grunde
dieselbe Menge, doch ist der Operator \fu{=} rechts-assoziativ,
w"ahrend \fu{+} links-assoziativ ist, was sich z.~B. beim
Erzeugen eines Parse-Trees zeigt.

Mit diesem Wissen k"onnen wir nun einen Parser konstruieren, der
einen \fu{expr2} akzeptiert, m"ussen aber wie gesagt die
Links-Assoziativit"at beachten. Um nicht unn"otig viele Worte
machen zu m"ussen, folgt ein Beispiel einer Implementation, die
einen (primitiven) Parse-Tree erzeugt. F"ur den Parse-Tree wird
folgende Struktur verwendet, die keiner Erkl"arung bedarf:

\begin{verbatim}
struct ParseTree
  Operator : int;
  Value    : int;
  Left     : *ParseTree;
  Right    : *ParseTree;
end
\end{verbatim}

Die Funktion \fu{expr2} erzeugt einen Parse-Tree f"ur einen
\fu{expr2} und greift dabei auf die Funktion \fu{expr1} zur"uck,
die dasselbe f"ur einen \fu{expr1} tut. Zus"atzlich werden die
Funktionen \fu{LookAhead} und \fu{Accept} verwendet, die ich
nicht n"aher erl"autern mu"s:

\begin{verbatim}
function expr2 () : *ParseTree

  Tree : *ParseTree;
  Temp : *ParseTree;

begin
  Tree = expr1();
  while (LookAhead() == PLUS)
    Accept(PLUS);
    Temp = new ParseTree;
    Temp->Operator = PLUS;
    Temp->Left = Tree;
    Temp->Right = expr1();
    Tree = Temp;
  end
  return Tree;
end
\end{verbatim}

Auf diese Art und Weise k"onnen Funktionen f"ur jeden der
Priorit"atsgrade erstellt werden. Im wesentlichen mit derselben
Methode werden auch die Links-Rekursionen in den Produktionen von
\fu{primary} beseitigt. Eine Kleinigkeit in der Implementation
von \fu{primary} m"ochte ich noch besprechen: Die Produktion
\fu{IDENTIFIER `(' [expr \{`,' expr\}] `)'} enth"alt ein
optionales Element und darin eingeschlossen ein Element, da"s
sich beliebig oft wiederholen kann. Es ist zwar m"oglich, solche
Konstruktionen durch Umformung der Grammatik mit den oben
besprochenen Methoden zu analysieren, doch die direkte
Implementation ist nicht schwer zu realisieren und wird deshalb
von mir vorgezogen.

Die beiden eckigen Klammern in der Produktion bedeuten, da"s auf
die "offnende runde Klammer, die nach dem \fu{IDENTIFIER} zu
stehen hat, entweder eine schlie"sende runde Klammer oder ein
\fu{expr} folgt. Welche der beiden M"oglichkeiten zutrifft,
kann der Parser dadurch erkennen, da"s er den Look-Ahead
betrachtet.
Da ein \fu{expr} niemals mit einer schlie"senden Klammer beginnen
kann, ist klar, da"s, wenn der Look-Ahead die schlie"sende
Klammer ist, die erste Variante zutrifft, andernfalls die zweite.
In letzterem Fall mu"s die Funktion \fu{expr} aufgerufen werden, um
den ersten \fu{expr} zu erkennen. Nun steht der Parser vor einem
"ahnlichen Problem. Entweder, es folgt ein Beistrich und noch ein
\fu{expr}, dann vielleicht noch ein Beistrich u.~s.~w., oder es
folgt die schlie"sende Klammer. Wie oben kann auch dieses
Problem mit dem Look-Ahead gel"ost werden. Entspricht er einem Beistrich,
mu"s noch ein \fu{expr} analysiert werden, wonach der Parser wieder
auf einen Beistrich testen mu"s. Ist der Look-Ahead die
schlie"sende Klammer, mu"s er diese nur noch lesen und ist
mit der Analyse des \fu{primary} fertig. Wenn jedoch der
Look-Ahead keinem der beiden Symbole entspricht, hat er
offensichtlich einen Fehler gefunden, der dementsprechend zu
behandeln ist.

\subsection{Der Anweisungsparser}

Der Parser f"ur Anweisungen ist relativ einfach zu realisieren,
da er weder Linksrekursionen enth"alt, noch irgendwelche
Priorit"atsregeln zu befolgen sind. Die Anweisungen von SEL sind
so konzipiert, da"s der Parser schon anhand des ersten Tokens
entscheiden kann, um welche Anweisung es sich handelt. Er
unterscheidet sich jedoch insofern von den bis jetzt besprochenen
Parsern, als da"s bei bestimmten Anweisungen nur Teile dieser
Anweisung geparst werden m"ussen, wie z.~B. bei \rw{if}, wo der Wert der
Bedingung bestimmt, ob der \rw{if}- oder der \rw{else}-Zweig,
wenn vorhanden, ausgef"uhrt werden soll.

Betrachten wir als Beispiel die zitierte \rw{if}-Anweisung. Der
Parser erkennt anhand des Look-Ahead, da"s es sich um eine
\rw{if}-Anweisung handelt. Er "uberliest also die beiden Token
\rw{if} und \fu{(} und ruft anschlie"send den im vorigen
Unterkapitel besprochenen Ausdruck-Parser auf, um den \fu{expr}
zu lesen. Wie sich in Kapitel \ref{semantics} zeigen wird,
berechnet dieser Parser den Wert des Ausdrucks, der nun
benutzt werden kann, um zu entscheiden, ob die dem \rw{if} folgenden
Anweisungen ausgef"uhrt werden sollen. Ist der Wert ungleich \fu{0}, ist
die Sache leicht. In diesem Fall mu"s sich der Anweisungs-Parser
nur selbst aufrufen, um die folgenden Anweisungen bis zum
\rw{end} bzw. \rw{else} zu parsen. Andernfalls benutzt er die
zus"atzliche Information, die ihm der in Kapitel \ref{tokencode}
besprochene Zwischencode zum \rw{if}-Token liefert. Sie gibt die
Position an, an der das zugeh"orige \rw{end} bzw. \rw{else}
steht. So m"ussen nicht langwierig alle folgenden Anweisungen bis
zum jeweiligen Token durchsucht werden, sondern es kann direkt an
die entsprechende Position gesprungen werden. Steht an dieser
Stelle ein \rw{end}, ist die Analyse der \rw{if}-Anweisung
beendet, andernfalls mu"s der \rw{else}-Zweig ausgef"uhrt werden.

Noch etwas komplizierter verh"alt es sich mit der
\rw{switch}-Anweisung. Dort mu"s das \rw{switch}-Token die
Position des ersten zugeh"origen \rw{case} enthalten, dieses
wiederum die des n"achsten \rw{case} bzw. \rw{default}
u.~s.~w. Auf diese Weise ist es m"oglich, einen Zweig nach
dem anderen auf "Ubereinstimmung zu untersuchen und, falls eine
"ubereinstimmende Konstante gefunden wurde, diesen Zweig
auszuf"uhren, indem die Funktion sich selbst aufruft. Damit die
\rw{switch}-Anweisung auch wie vorgeschrieben ausgef"uhrt wird,
mu"s der Anweisungs-Parser beim normalen Parsen ein
\rw{case}-Schl"usselwort einfach "uberspringen und bei einem
\rw{break} zur "ubergeordneten Funktion zur"uckkehren.

Ein Problemfall ist die \rw{return}-Anweisung. Wenn der
Parser auf sie trifft, mu"s er auch zur "ubergeordneten Funktion
zur"uckkehren, doch er mu"s das im Gegensatz zur
\rw{break}-Anweisung so lange tun, bis der Anweisungs-Parser
komplett verlassen wurde. Wir m"ussen uns also "uberlegen, wie
wir es anstellen k"onnten, da"s der Parser wei"s, wann die
untergeordnete Instanz ein \rw{return} entdeckt hat. Die
einfachste L"osung ist, das mit dem R"uckgabewert zu tun. Im
n"achsten Kapitel wird sich zeigen, da"s SEL-Werte vom
Interpreter in Strukturen gehalten werden. Am sinnvollsten ist
es, als R"uckgabetyp f"ur den Parser einen Zeiger auf eine solche
Struktur zu benutzen. Trifft er nun auf ein \rw{break}, gibt er
einen Nullzeiger zur"uck und der "ubergeordnete Parser wei"s, das
alles in Ordunung ist. Wurde aber eine \rw{return}-Anweisung
entdeckt, wird ein Zeiger auf den zugeh"origen Wert
zur"uckgegeben. Die "ubergeordnete Funktion wei"s nun, da"s auch
sie beendet werden mu"s und gibt ihrerseits den Wert weiter. Wird
eine Funktion geparst, die als R"uckgabetyp \tvoid\ hat, mu"s
ein Zeiger auf einen Dummy-Wert zur"uckgegeben werden.

Der Anweisungs-Parser mu"s auch mit einem Null-Zeiger verlassen
werden, wenn er auf ein \rw{end} oder \rw{else} trifft, die zu
keiner von ihm behandelten Anweisungen geh"ort, denn dann f"allt
es in den Zust"andigkeitsbereich der "ubergeordneten
Rekursionsebene.

Hier wird auch klar, warum ich auf eine
\rw{goto}-Anweisung\footnote{siehe Kapitel \ref{goto}} verzichtet
habe: In einem rekursiven Parser w"are es f"ur einen Sprung
n"otig, zun"achst zu ermitteln, auf welcher Rekursionsebene das
Label liegt, und ob es vielleicht innerhalb einer anderen Anweisung
liegt, was einen besonders schlimmen Fall darstellt. Betrachten wir
dazu folgendes (theoretische) Programmfragment:

\begin{verbatim}
if (x == 5)
  goto label;
end
if (y == 3)
  do_something();
  if (z == 4)
    label:
    do_something_else();
  end
end
\end{verbatim}

In diesem Fall m"u"ste der Interpreter, wenn er die
\rw{goto}-Anweisung ausf"uhrt, zun"achst eine Rekursionsebene
verlassen und dann in zwei andere einsteigen, was nat"urlich sehr
schwer zu realisieren ist. "Uberdies m"u"ste der Parser sehr
schnell herausfinden k"onnen, wieviel Ebenen zu verlassen und in
welche er einzusteigen hat. Zu diesem Problem ist mir bis jetzt
noch keine zufriedenstellende L"osung eingefallen.

\section{Die Verwaltung}
\label{semantics}

In den letzten Unterkapiteln wurde behandelt, wie der Interpreter
den Quelltext (genauer gesagt, den Zwischencode) lexikalisch
analysiert, in Zwischencode verwandelt und diesen syntaktisch
analysiert, jedoch wurde noch nicht gesagt, wie der Interpreter
die durch diese Analyse gewonnenen Informationen dazu verwenden
kann, das Programm auszuf"uhren. Das ist die Aufgabe dieses
Unterkapitels. Hier wird besprochen, was der Interpreter bei der
Ausf"uhrung des Programms tut, und wie dies implementiert werden
kann.

Die meisten der in diesem Unterkapitel besprochenen Algorithmen
werden w"ahrend der Syntaxanalyse, also des im letzten
Unterkapitel besprochenen Parsing, ausgef"uhrt, was auch der
Definition des Interpreters\footnote{siehe Kapitel
\ref{whatisit}} entspricht.
        
\subsection{Module}
\label{modules}

Da SEL es erlaubt, modulglobale Variablen und Funktionen zu
benutzen, ist es vonn"oten, ein SEL-Programm auch im Speicher in
Module aufzuteilen. Dazu m"ussen f"ur jedes Modul mindestens der
Zwischencode und die modulglobalen Variablen und Funktionen
gespeichert werden. Im n"achsten Unterkapitel werden sich dazu
noch einige zus"atzliche Informationen gesellen, die bei der
Fehlerbehandlung helfen.

Die programmglobalen Funktionen und Variablen werden au"serhalb
der Module f"ur das gesamte Programm gemeinsam verwaltet, was
wesentlich effizienter und auch einsichtiger ist (und es nebenbei
einfacher und schneller macht, Namenskonflike zu bemerken).
Weiters gibt es zwei Variablen, die f"ur das gesamte Programm
gelten. Sie geben die Position an, an der sich der Parser
gerade mit der Ausf"uhrung aufh"alt. Die erste dieser Variablen
nenne ich \fu{MP} (Module Pointer). Sie enth"alt einen Wert, der
das Modul, in dem sich die gerade abgearbeiteten Anweisungen
befinden, eindeutig identifiziert. Am effizientesten ist es, ein
Array aus Modulbeschreibungen anzulegen, das durch diese
Variable indiziert wird. Die zweite Variable, \fu{IP}
(Instruction Pointer), gibt die Position innerhalb des
Zwischencodes des aktuellen Moduls an.\footnote{Genauerweise
m"u"ste diese Variable den Namen \fu{TP} f"ur Token Pointer
tragen, da sie sich nicht auf die aktuelle Anweisung, sondern auf
das aktuelle Token bezieht.}

\bigskip

In Kapitel \ref{intgoals} wurde als Ziel f"ur den Interpreter
festgelegt, da"s er Multitasking unterst"utzt. Um die dazu
notwendigen Schritte anzugeben, mu"s erst gekl"art werden, wie
Multitasking funktioniert. Zun"achst m"ochte ich feststellen, da"s
ich mit Multitasking nicht die F"ahigkeit eines Betriebssystems
meine, mehrere Programme gleichzeitig zu bearbeiten, sondern die
F"ahigkeit, innerhalb eines Programms mehrere Funktionen zur
gleichen Zeit auszuf"uhren. Unter den Betriebssystemen OS/2 und
Windows NT wird daf"ur der Begriff "`Multithreading"' gepr"agt.
Im Fall eines SEL-Interpreters l"auft dieses Multitasking darauf
hinaus, da"s durch diesen Interpreter (der ja ein Programm
darstellt) mehrere SEL-Programme gleichzeitig ausgef"uhrt werden
k"onnen.

Zwei Funktionen, die gleichzeitig ausgef"uhrt werden, k"onnen
zwar auf dieselben globalen Variablen zugreifen, soda"s eine
gemeinsame Benutzung eines Datenpools, etwa einer
Funktionssammlung f"ur ein Mathematikprogramm, gew"ahrleistet
ist, jedoch besitzt jede der beiden Funktionen einen eigenen
Stack, was garantiert, da"s sie ungest"ort von den anderen
Funktionen ihre eigenen lokalen Variablen verwenden kann.

Wenn nun der Interpreter zwei verschiedene SEL-Programme
gleichzeitig ausf"uhren soll, tut er das, indem er zwei Tasks
startet, die jeweils ein Programm bearbeiten. Um zu
gew"ahrleisten, da"s sich diese beiden Funktionen dabei nicht in
die Quere kommen, d"urfen die Daten, die das aktuelle Programm
betreffen, nat"urlich nicht in globalen Variablen gespeichert
werden, sondern m"ussen der Funktion, die f"ur die Ausf"uhrung
zust"andig ist, als Parameter mitgegeben werden. F"ur unseren
rekursiven Parser bedeutet das, da"s diese Daten sehr oft als
Parameter "ubergeben werden, was nicht nur den Stack belastet,
sondern auch die Ausf"uhrungsgeschwindigkeit. Um diese Verluste
gering zu halten, empfielt es sich, nicht direkt die Daten,
sondern Zeiger auf sie zu verwenden, was aber in den meisten
Programmiersprachen ohnedies die einzige M"oglichkeit ist.

\subsection{Typen}
\label{inttypes}

Eine f"ur die Ausf"uhrungsgeschwindigkeit sehr wichtige
Entscheidung bei einem Interpreter f"ur eine Programmiersprache,
die ein im Vergleich zu anderen Interpretersprachen so komplexes
Typkonzept verwendet, wie SEL, ist, wie diese Typen
repr"asentiert werden sollen. In Compilern werden hierf"ur
h"aufig verkettete Listen verwendet, die jedoch f"ur einen
Interpreter nicht sehr brauchbar sind, da sie wegen ihres
Aufbaus, der auf dynamischer Speicherverwaltung basiert, zu
langsam sind und au"serdem zu viel Speicher verbrauchen.

Das von mir verwendete Konzept basiert darauf, da"s Typen durch
Byte-Arrays mit begrenzter L"ange repr"asentiert werden. Da es
in SEL nur eine geringe Zahl an Grundtypen und auch nicht viele
M"oglichkeiten, diese zu verkn"upfen, gibt, kann ein Byte
jeweils eine Zusammensetzungsstufe oder einen Grundtypen
enthalten. Ein Beispiel: Der Typ \type{\poi \arr{10}int} wird
durch ein Array mit der Gr"o"se von 5 Bytes repr"asentiert,
dessen erstes Byte angibt, da"s es sich um einen Zeiger handelt.
Das zweite Element gibt an, da"s ein Array im Spiel ist. F"ur
dieses Array werden auch die n"achsten zwei Bytes ben"otigt, die
die Gr"o"se angeben. Das f"unfte Element schlie"slich steht f"ur
den Grundtypen \tint\ und ist gleichzeitig die Ende-Markierung,
da es sich dabei um einen Grundtypen handelt. In gewisser Weise
"ahnlich wie \tint\ werden Strukturen behandelt, die ja auch
immer am Ende des Typs stehen. Bei ihnen mu"s, "ahnlich dem
Array, ein Wert angegeben werden, der den Typ der Struktur
bestimmt. Zu diesem Zweck vergibt der Interpreter f"ur die Strukturen
Nummern, wobei die Werte von 0 bis 9999 f"ur programmglobale
Strukturen reserviert sind. Die weiteren Zuordungen richten sich
danach, wie viel Strukturen in den einzelnen Modulen definiert
sind.

Wie schon oben erw"ahnt, besitzen die Typenarrays eine begrenzte
L"ange, was es erm"oglicht, sie schon vor ihrer Verwendung
bereitzustellen, soda"s sie nicht w"ahrend der Ausf"uhrung des
SEL-Programms dynamisch allokiert werden m"ussen. In meiner
Implementation habe ich mich f"ur eine L"ange von 20 Bytes
entschieden. Das mag nicht viel erscheinen, vor allem in
Anbetracht der Tatsache, da"s es in SEL m"oglich ist, Typen
beliebig zu verschachteln, doch ist es m"oglich, mit diesen 20
Bytes einen Zeiger auf einen Zeiger auf einen Zeiger auf ein Array
aus Arrays aus Strukturen zu repr"asentieren, wobei noch 8 Byte
frei bleiben. Dieses Beispiel soll zeigen, da"s 20 Bytes f"ur die
Praxis eine absolut ausreichende Gr"o"se ist.

Beim Analysieren von Ausdr"ucken ist es in einem SEL-Interpreter
sehr oft n"otig, Typen bereitzustellen, von denen aber viele
schon im n"achsten Moment wieder freigegeben werden. Ist es nicht
m"oglich, diese Operationen effizient auszuf"uhren, leidet
darunter die Geschwindigkeit des Interpreters sehr. Ich habe zu
diesem Zweck einen Algorithmus entwickelt, der es erm"oglicht, in
konstanter Zeit, also unabh"angig davon, wieviele von den
bereitgestellten Typenarrays bereits vergeben sind, Typenarrays
zu allokieren und freizugeben:

\bigskip

Sei $n$ die Anzahl der maximal bereitzustellenden Typen und $s$
die Gr"o"se eines Typenarrays, dann wird vor der Ausf"uhrung des
SEL-Programms ein Byte-Array mit der Gr"o"se $n \cdot s$ erzeugt,
wobei $t$ die Speicheradresse sei, an der das Array beginnt.
Weiters sei $r$ die Speicheradresse, an der ein Array der Gr"o"se
$n$ beginnt, dessen Elemente vern"unftigerweise Maschinenworte
sind. Dieses Array wird initialisiert, soda"s $r[x] = x$.
Zus"atzlich sei $i$ eine numerische Variable, die mit
dem Wert $0$ initialisiert werden.

Soll ein Typenarray allokiert werden, wird zun"achst "uberpr"uft,
ob $i = n$. In diesem Fall sind bereits alle Arrays vergeben und
eine entsprechende Fehlermeldung wird ausgegeben. Andernfalls ist
$t + r[i] \cdot s$ die Adresse, bei der das angeforderte Array
beginnt. Zus"atzlich ist $i = i + 1$ zu setzen.

Soll das Typenarray mit der Speicheradresse $f$ freigegeben
werden, so ist zuerst $i = i - 1$ und dann $r[i] = (f - t) / s$
zu setzen.

Der angegebene Algorithmus verl"a"st sich darauf, da"s nicht
Arrays freigegeben werden sollen, die nicht allokiert sind, und
da"s jeweils maximal $s$ Bytes eines Arrays genutzt werden.

\bigskip

Dieser Algorithmus verwendet zwei Datenstrukturen. Das Array $t$
ist sozusagen ein Zusammenschlu"s aller Arrays, die zur
Verf"ugung gestellt werden k"onnen, was sich nat"urlich auch
mittels eines zweidimensionalen Arrays realisieren lie"se. Das
Array $r$ ist ein Puffer, der die Nummern (0 bis $n - 1$)
aller noch freien Typenarrays enth"alt. Die Variable $i$ gibt den
Index des n"achsten zu vergebenden Typenarrays an. Wird dieses
angefordert, wird $i$ um 1 erh"oht. Umgekehrt wird $i$ um 1
vermindert, wenn ein Array freizugeben ist. Dessen Nummer wird
anschlie"send in $r[i]$ gespeichert, soda"s es bei der n"achsten
Anforderung sofort wieder vergeben wird. Wie schon gesagt, nimmt
der Algorithmus an, da"s keine Fehler vom Programmierer begangen
werden, wie z~B. die Freigabe eines bereits freigegebenen Arrays.

In meiner Implementation von SEL habe ich f"ur $n$ den Wert 400
benutzt, f"ur $s$ wie gesagt 20. Diese Werte sind f"ur die
meisten Anwendungen ausreichend, und wenn n"otig, k"onnten sie
vom Benutzer entsprechend ge"andert werden. Die
Speicheranforderungen dieses Algorithmus sind durchaus
vertretbar. Bei den angegebenen Werten erw"achst f"ur die beiden
Arrays unter Windows NT ein Speicherbedarf von genau 9600 Byte,
was bei den heute "ublichen Speicherkapazit"aten kaum ein Problem
darstellt.

\subsection{Ausdr"ucke}
\label{intexprs}

"Ahnlich wichtig wie die Typen sind in SEL die Ausdr"ucke, die
nat"urlich auch verwaltet werden m"ussen. Dabei werden wir jedoch
nicht ganze Ausdr"ucke verwenden, sondern uns lediglich auf die
Werte der Ausdr"ucke beschr"anken, die im Laufe des
Parsevorganges errechnet werden.

Um zu zeigen, wie wichtig die effiziente Verwaltung von
Ausdr"ucken ist, hier ein konkretes Beispiel: Ein
QuickSort-Programm f"ur 500 zuf"allig angeordnete Zahlen hat auf
meinem Computer eine Laufzeit von etwa 25 Sekunden, wobei w"ahrend
dieser Zeit ungef"ahr 60000 Ausdr"ucke allokiert und wieder
freigegeben werden m"ussen.

Zum Allokieren und Freigeben von Ausdr"ucken verwende ich aus
diesem Grund denselben Algorithmus wie f"ur Typen, wobei ich von
Haus aus maximal 200 Ausdr"ucke zur Verf"ugung stelle, was in den
meisten F"allen ausreichend ist.

Sehr wichtig f"ur eine effiziente Verarbeitung ist nat"urlich
auch der Aufbau eines Ausdrucks. Ich verwende hierzu eine
Struktur. Ein obligater Bestandteil eines Ausdrucks ist dessen
Typ, soda"s ein Typenarray Teil dieser Struktur ist. Weiters gibt
es je eine Variable f"ur die Basistypen \tchar , \tint\ und
\tfloat . Hierzu wird es aus Effizienzgr"unden am besten sein,
passende Typen der Implementationssprache zu verwenden. F"ur den
Fall, da"s der Ausdruck ein Array, ein Zeiger oder eine Struktur
ist, habe ich zus"atzlich einen typenlosen Zeiger aufgenommen,
der, falls die Implementationssprache keine typenlosen Zeiger
zul"a"st, durch einen Zeiger auf ein Byte ersetzt werden kann.
Im Falle eines Funktionsaufrufs ist es n"otig, mehrere Ausdr"ucke
zu verketten, um die Parameter der Funktion anzugeben, soda"s
zus"atzlich ein Zeiger auf einen Ausdruck in die Struktur
eingebaut werden sollte, soda"s sie sich zu einer verketteten
Liste erweitern l"a"st.

Soviel zur Verwaltung der Ausdr"ucke im Speicher. Eine weit
wichtigere Frage jedoch ist, was der Interpreter mit diesen
Ausdr"ucken tut. Nun, zun"achst einmal mu"s jede der in
Kapitel \ref{exprparser} besprochenen Produktionsfunktionen so
konzipiert werden, da"s sie der "ubergeordneten Funktion einen Ausdruck
mit dem Wert des von ihr geparsten Ausdrucks zur"uckgibt. Diese
verwendet ihn ihrerseits, um den Wert ihres Ausdrucks zu
berechnen u.~s.~w. In der Kette der Funktionsaufrufe gibt es dann
irgendwann eine Funktion, die den Wert des Ausdrucks benutzt, um
damit etwas zu entscheiden, wie z.~B., ob der Zweig einer
\rw{if}-Anweisung ausgef"uhrt werden soll.

Solange man lediglich Berechnungen mit den Ausdr"ucken
anstellt, ergeben sich mit der oben skizzierten Methode
keinerlei Probleme, wie ist es jedoch m"oglich, den
Zuweisungsoperator zu implementieren? Bauen wir die Funktion zum
Parsen eines \fu{primary} so auf, da"s sie f"ur eine Variable
lediglich deren Wert zur"uckgibt, werden wir Probleme haben, denn
dann wei"s der Interpreter nicht, wo diese Variable im Speicher liegt und
kann ihr so auch keinen Wert zuweisen. Um dieses Problem zu
"uberwinden, wird intern ein zus"atzlicher Typ,
n"amlich der Referenztyp, der einem Zeiger sehr "ahnlich ist,
verwendet.
Trifft die Funktion zum Parsen eines \fu{primary} nun im Zwischencode auf eine Variable des Typs
\tint , dann gibt sie nicht deren Wert zur"uck, sondern ihre
Adresse und vermerkt, da"s es sich dabei um eine Referenz auf
einen \tint\ handelt, der Einfachheit halber von mir mit
\type{\reference int} bezeichnet.
Mu"s nun mit diesem Ausdruck gerechnet werden, gen"ugt es, ihn zu
dereferenzieren, soda"s ein Ausdruck mit dem Typ \tint\ entsteht.
Wird ihm ein Wert zugewiesen, mu"s dieser lediglich an die
Speicheradresse geschrieben werden, die er angibt. Soll der
un"are Operator \fu{\&} auf ihn angewendet werden, dann ist es
lediglich n"otig, das erste Byte im Typenarray so zu "andern,
da"s daraus ein Zeiger wird.
                
\subsection{Variablen}
\label{maintvars}

Variablen sind von den oben besprochenen Ausdr"ucken und Typen
insofern verschieden, als da"s sie nicht intern vom Interpreter
verwendet werden, sondern direkt mit dem Quellprogramm in
Verbindung stehen. Aus dem Modulkonzept und der Tatsache, da"s
Funktionsvariablen nur in der jeweiligen Funktion verf"ugbar
sind, folgt, da"s die Verwaltung von Variablen auf mehrern Stufen
erfolgen mu"s. Die umfassendste Stufe, allerdings auch diejenige
mit der niedrigsten Priorit"at, ist die der programmglobalen
Variablen. Auf sie folgt die Stufe der modulglobalen Variablen.
Die h"ochste Priorit"at besitzt die Stufe der Funktionsvariablen,
die allerdings in der Regel auch die wenigsten Variablen
aufnimmt.

Wird nun im Programm eine Variable verwendet, mu"s zun"achst
gepr"uft werden, ob es eine Funktionsvariable ihres Namens gibt.
Ist das nicht der Fall, kann es sich nur mehr um eine globale
Variable handeln. Gibt es eine modulglobale Variable, die den
gesuchten Namen hat, wird sie verwendet, andernfalls mu"s es eine
programmglobale Variable geben. Ist auch das nicht zutreffend,
hat es der Interpreter mit einem Programmfehler zu tun (siehe Kapitel
\ref{errortrapping}).

Die komplizierteste Verwaltungsstufe ist die der
Funktionsvariablen (es ist am einfachsten, bei der Verwaltung der
Variablen nicht zwischen Parametern und direkten Funktionsvariablen
zu unterscheiden), was daran liegt, da"s nicht alle
Funktionsvariablen in einen "`Topf"' geworfen werden k"onnen,
sondern da"s f"ur jede Funktion eine eigene Variablenliste
angelegt werden mu"s. Die einfachste L"osung f"ur dieses Problem
ist, einen Stack zu verwenden, auf dem f"ur jeden
Funktionsaufruf eine Liste aller Funktionsvariablen angelegt
wird, "ahnlich dem Code, den ein Compiler erzeugen w"urde. Analog
dazu werden bei der R"uckkehr aus einer Funktion die vorher
erzeugten Variablen wieder entfernt. Um nun
nach einer Funktionsvariablen zu suchen, mu"s nur der oberste
Stacklevel in Betracht gezogen werden, da er alle Variablen
beinhaltet, die zu der gerade auszuf"uhrenden Funktion geh"oren.

Wie schon in Kapitel \ref{modules} angesprochen, stellt der
Interpreter f"ur jedes Modul eine eigene Datenstruktur zur
Verf"ugung, die unter anderem alle modulglobalen Variablen dieses
Moduls verwaltet. Ist also nach einer modulglobalen Variablen zu
suchen, kann man die Suche auf das durch \fu{MP} angegebene Modul
beschr"anken.

Der einfachste Fall liegt bei programmglobalen Variablen vor.
F"ur ihre Verwaltung gen"ugt es, eine einzige Liste
bereitzustellen, die die programmglobalen Variablen aller Module
bereitstellt, ohne Unterscheidung, in welchem Modul diese
definiert sind.

\bigskip

Bis jetzt habe ich, wenn es um die Details der
Variablenverwaltung ging, lediglich von "`Listen"' gesprochen,
wobei ich damit weniger dynamische Listen, als vielmehr
geeignete Datenstrukturen f"ur das schnelle Auffinden von
Variablen gemeint habe. Diese Datenstrukturen m"ochte ich jetzt
genauer besprechen.

Das Problem, auf das wir hier treffen ist, eine
Datenstruktur (und nat"urlich auch die dazugeh"origen
Algorithmen) aussuchen zu m"ussen, die es erm"oglicht, einen
Datensatz in einer Menge von Datens"atzen m"oglichst
schnell zu finden. Dieses Problem ist ein gutes Beispiel f"ur das
Grundproblem eines Programmdesigners: F"ur eine bestimmte Aufgabe
den besten Algorithmus bzw. die beste Datenstruktur (was in den
meisten F"allen Hand in Hand geht) auszuw"ahlen. G"abe es eine
absolut beste L"osung, dann h"atte der Designer seine Aufgabe
schnell erledigt, doch leider gibt es eine solche nur in den
seltensten F"allen.

Das wichtigste Entscheidungskriterium f"ur unsere Datenstruktur
ist, wieviel Datens"atze zu verwalten sind. F"ur die
Variablenverwaltung eines Interpreters bieten sich grunds"atzlich
zwei Strukturen an: Bin"are Suchb"aume und Hash-Tabellen. Um
entscheiden zu k"onnen, welche der beiden Strukturen f"ur unsere
Anwendungen die bessere ist, ist es n"otig, ein wenig zu rechnen.
Der Einfachheit halber gehen wir vom g"unstigsten Fall aus (die
Ergebnisse sind von denen des durchschnittlichen Falls nicht sehr
verschieden). Dieser ist, da"s der Benutzer die Variablennamen so
gew"ahlt hat, da"s sich eine gleichm"a"sige Anordung in der
Hash-Tabelle ergibt bzw. der Suchbaum ausgeglichen ist.

In einer Hash-Tabelle, bei der die Eintr"age f"ur gleiche
Hash-Werte linear verkettet sind, werden f"ur eine Suche nach
einem Datensatz im Mittel $\max (N / 2M + 1/2, 1)$ Vergleiche
ben"otigt, wohingegen f"ur die Suche in einem Baum etwa
$\log_2N - 1$ Vergleiche notwendig sind, wobei $N$ die Anzahl der
Datens"atze und $M$ die Gr"o"se der Hash-Tabelle angibt. Auf den
ersten Blick scheint der Baum die bessere L"osung zu sein, da die
Anzahl der ben"otigten Vergleiche f"ur gro"se $N$ in etwa
proportional zum bin"aren Logarithmus der Anzahl der Datens"atze
ist, bei einer Hash-Tabelle allerdings die Anzahl der Vergleiche
direkt proportional zur Anzahl der Datens"atze ist, was f"ur
gro"se $N$ nat"urlich bedeutend schlechter ist, als die
Leistungsdaten des Baums. Nehmen wir allerdings an, wir benutzen
eine Hash-Tabelle mit 37 Eintr"agen, was in Bezug auf den
Speicherplatz "uberhaupt kein Problem darstellt. In diesem Fall
sind f"ur eine Suche in einer Hash-Tabelle mit 100 Datens"atzen
im Durchschnitt weniger als zwei Vergleiche n"otig, wohingegen in
einem Baum mehr als sechs Vergleiche ben"otigt werden. Erst wenn
mehr als 550 Datens"atze vorhanden sind, ist die Suche in einem
Baum g"unstiger als in einer Hash-Tabelle.

Es ist kaum anzunehmen, da"s jemand f"ur einen Interpreter ein
Programm schreibt, welches mehr als 300 globale Variablen verwendet,
soda"s es f"ur unsere Zwecke die bei weitem bessere L"osung ist,
f"ur die Verwaltung programm- und modulglobaler Variablen
Hash-Tabellen zu verwenden, wobei eine Gr"o"se von 37 durchaus
praktikabel ist, sowohl in Bezug auf Schnelligkeit als auch auf
die vern"unftige Ausnutzung des Speicherplatzes. F"ur
Funktionsvariablen sind 37 Eintr"age wohl "uberdimensioniert.
Hierf"ur bieten sich die Primzahlen 7 und 11 an.

Zuletzt stellt sich noch die Frage nach der Hash-Funktion. Ich
verwende sowohl der Einfachheit als auch der Effizienz halber die
Modulo-Funktion, d. h., ich betrachte eine Zeichenkette als lange
Ganzzahl, deren Rest bez"uglich des Moduls 37 den Index in der
Hash-Tabelle angibt. Die folgende Implementation geht davon aus, 
da"s ein \tchar\ 256 verschiedene Werte annehmen kann.

\begin{verbatim}
function hash (s : *char, m : int) : int

  h : int = 0;

begin
  while (*s)
    h = (h * 256 + *s) % m;
    s = s + 1;
  end
  return h;
end
\end{verbatim}

\subsection{Funktionen}

Die Funktionsverwaltung hat in erster Linie die Aufgabe, eine
Funktion eines bestimmten Namens zu finden und deren Adresse im
Zwischencode zu ermitteln. Weiters ist es wichtig, Namen und
Typen der Funktionsparameter zu erfahren, um erstens "uberpr"ufen
zu k"onnen, ob die Funktion mit der richtigen Anzahl und den
richtigen Typen der Parameter aufgerufen worden ist, und zweitens,
um die Parameter auf dem Funktionsstack anlegen zu
k"onnen.\footnote{siehe Kapitel \ref{maintvars}}

Zus"atzlich sollte es m"oglich sein, nicht nur Funktionen, die in SEL
implementiert sind, aufzurufen, sondern auch solche, die in der
Implementationssprache des Interpreters geschrieben sind,
da eine Implementation der Grundfunktionen von SEL in
SEL selbst f"ur einen Interpreter viel zu ineffizient w"are.
Deshalb schreibt man Funktionen wie etwa zum Kopieren von
Strings in der Implementationssprache des Interpreters (in
meinem Falle C) und ruft diese auf. Um das nahtlos in die
Verwaltung einbinden zu k"onnen, ist es notwendig, da"s die
Implementationssprache Zeiger auf Funktionen unterst"utzt. Ist
diese M"oglichkeit nicht vorhanden (wie etwa im Falle von
\pascal ), dann mu"s jede einzelne Funktion im Programmcode extra
behandelt werden, was sowohl umst"andlicher als auch langsamer
ist.\footnote{In Turbo Pascal ab der Version 5.5 ergibt sich die
M"oglichkeit, statt eines Zeigers auf eine Funktion einen Zeiger
auf ein Objekt zu verwenden, welches eine Methode enth"alt, die
die gew"unschte Funktion ersetzt.}

Um eine Funktion eines bestimmten Namens zu finden, werden
Hash-Tabellen, wie im letzten Unterkapitel beschrieben,
verwendet: Eine f"ur programmglobale Funktionen und eine
pro Modul f"ur modulglobale Funktionen, der Einfachheit halber
jeweils mit 37 Eintr"agen, was bei einem Programm mit 10 Modulen
unter Windows NT einen Speicherplatz von 1628 Bytes beansprucht
(was ohne schlechtes Gewissen als "`unbedenklich"' eingestuft
werden kann). Eintr"age f"ur denselben Hash-Wert werden
unsortiert linear verkettet.\footnote{Eine Sortierung h"atte bei
einer linearen Verkettung auch nicht den geringsten Sinn, da
ohnedies alle Eintr"age bis zum gew"unschten durchlaufen werden
m"ussen, unabh"angig davon, in welcher Reihenfolge die Eintr"age
gespeichert sind.}

Die Datenstruktur, die eine einzelne Funktion beschreibt, mu"s
neben deren Namen und R"uckgabetyp auch eine Angabe dar"uber
enthalten, wo der Code der Funktion zu finden ist. Hierf"ur wird
zun"achst eine Variable ben"otigt, die Auskunft dar"uber gibt, ob
die Funktion in SEL geschrieben ist oder in der
Implementationssprache des Interpreters. Im ersten Fall mu"s die
Datenstruktur Auskunft dar"uber geben k"onnen, in welchem Modul
und an welcher Stelle der Code beginnt, andernfalls wird ein
Zeiger auf die jeweilige Funktion ben"otigt. Zus"atzlich m"ussen
Informationen "uber Namen und Typen der Parameter gespeichert
sein. Ich verwende hierzu eine einfach verkettete Liste.

\bigskip

Interessant ist die Realisierung von Funktionsaufrufen. Dazu ist
zun"achst zu bedenken, da"s der Parser rekursiv implementiert ist
und der Funktionsaufruf ein \fu{primary} ist, sich also auf der
untersten Grammatikebene der Ausdr"ucke befindet. Aus diesem
Grund wird man schwer umhinkommen, auch Funktionsaufrufe rekursiv
zu behandeln, was nat"urlich bei einem rekursiven SEL-Programm zu
einer starken Belastung des Prozessorstacks f"uhren kann, weshalb
dieser m"oglichst gro"s dimensioniert werden sollte. Mu"s nun
innerhalb eines \fu{primary} eine Funktion aufgerufen werden,
wird in den oben besprochenen Hash-Tabellen die Funktion gesucht
und, wenn vorhanden, wird an die angegebene Stelle im
Zwischencode gesprungen. Dort werden zun"achst der Funktionskopf
und die Deklaration der direkten Funktionsvariablen interpretiert
und anschlie"send der Code der Funktion mittels des
Anweisungs-Parsers ausgef"uhrt. Ist die Ausf"uhrung der Funktion
beendet, mu"s zur aufrufenden Funktion zur"uckgekehrt, also die
Werte der Variablen \fu{MP} und \fu{IP} restauriert werden.
Hierzu m"ussen klarerweise die alten Werte irgendwo
zwischengespeichert werden. Weiters m"ussen die lokalen Variablen
der Funktion abgelegt werden, und zwar so, da"s auch mehrere
Instanzen dieser Funktion existieren k"onnen und jede ihre
eigenen Variablen zur Verf"ugung hat. Wie schon im letzten
Kapitel angedeutet, wird hierzu ein Stack verwendet, in Anlehnung
an den Prozessorstack, der bei compilierten Programmen dieselben
Aufgaben wahrnimmt. Relevant ist jeweils nur der oberste Level
des Stacks, der die R"ucksprungadresse und die lokalen Variablen
der jeweils aktiven Funktion enth"alt. Wird aus dieser eine
andere Funktion aufgerufen, kommt ein Level zum Stack hinzu,
andernfalls wird ein Level entfernt. Um sich nicht auf die
relativ langsame dynamische Allokation verlassen zu m"ussen, ist
es am besten, diesen Stack als Array zu
konzipieren.\footnote{Eine Konzeption als dynamische Liste hat
auf den ersten Blick den Vorteil, da"s die Rekursionstiefe
lediglich vom vorhandenen Speicherplatz abh"angig ist, was jedoch
bei einem rekursiven Interpreter nicht zutrifft: dort ist der
Prozessorstack das Ma"s aller Dinge.}

\section{Fehlerbehandlung}
\label{errortrapping}

Da nicht davon auszugehen ist, da"s das Programm, das
interpretiert werden soll, fehlerfrei ist, mu"s der Interpreter
auch die M"oglichkeit haben, Fehler zu entdecken und sie
entsprechend zu behandeln. Im Gegensatz zu Compilern, die wenn
m"oglich nach dem Erkennen eines Fehlers mit dem "Ubersetzen
fortfahren sollten, um evtl. weitere Fehler zu entdecken, ist das
bei einem Interpreter nicht n"otig und auch wenig sinnvoll, da
das auszuf"uhrende Programm in seiner Stabilit"at gest"ort und
vermutlich abst"urzen w"urde. Eine sogenannte "`{\em Error
Recovery}"' bleibt uns also erspart. Die Ziele f"ur die
Fehlerbehandlung sind also das Erkennen von Fehlern und die
Ausgabe einer f"ur den Benutzer m"oglichst hilfreichen
Fehlermeldung.

Grunds"atzlich mu"s ein Interpreter zwei Arten von Fehlern
behandeln: syntaktische und semantische Fehler. Syntaktische
Fehler treten dann auf, wenn die im Anhang \ref{grammar}
angegebene Grammatik von SEL durch das Programm verletzt wird.
Diese Art von Fehler ist in der Regel leicht zu erkennen und kann
dem Benutzer relativ verst"andlich erkl"art werden. Ein Beispiel:
Der Interpreter ist gerade besch"aftigt, die Anweisung
\fu{while [x != 10)} auszuf"uhren. Der Parser erkennt das
Schl"usselwort \rw{while}, dem laut Grammatik immer
eine "offnende runde Klammer folgen mu"s. Offenbar hat sich der Benutzer
jedoch vertippt und statt der runden eine "offnende eckige
Klammer eingegeben. Es ist leicht, den Parser so zu erweitern,
da"s er an dieser Stelle den Fehler erkennt. Die entsprechende
Fehlermeldung sollte lauten: \fu{'(' expected}.

Semantische Fehler treten auf, wenn das Programm versucht,
unerlaubte Operationen auszuf"uhren, wie etwa das Casten eines
\tfloat\ in einen Zeiger, die Verwendung einer Variablen,
die gar nicht existiert, oder die Anwendung eines Operators auf
einen nicht erlaubten Typen. Tritt ein solcher Fall auf,
l"a"st sich das in der Regel ohnedies leicht feststellen und auch
eine passende Fehlermeldung ist nicht schwer zu schreiben.

F"ur den Benutzer des SEL-Interpreters ist es klarerweise nicht
nur wichtig, da"s ein Fehler aufgetreten und welcher Art
dieser ist, sondern auch, in welchem Modul in welcher Zeile der
Fehler steckt. Da der Interpreter weder syntaktische noch
semantische Fehler erkennen kann, w"ahrend er den Quellcode
untersucht, und der Zwischencode aus Effizienzgr"unden keine
Zeilenumbr"uche enth"alt, m"ussen die Zeilennummern gesondert
gespeichert werden. Ich verwende hierzu ein durch die
Zeilennummer indiziertes Array, das f"ur jede
Zeile eines Moduls die Anzahl der Bytes angibt, die diese Zeile
im Zwischencode beansprucht. Da der Zwischencode in der Regel
k"urzer als der Quellcode und kaum eine Programmzeile l"anger
als 100 Zeichen ist, gen"ugt es, pro Arrayelement ein Byte zur
Verf"ugung zu stellen. Tritt nun an der Stelle \fu{MP==3} und
\fu{IP==1364} ein Fehler auf, dann mu"s im Modul Nr. 3 das
Zeilennummernarray so lange durchlaufen werden, bis die Summe der
Elemente gr"o"ser als 1364 ist. Der folgenden Funktion, die die
gesuchte Zeilennummer zur"uckgibt, ist
das Zeilennummernarray als Zeiger und der Wert von \fu{IP} zu
"ubergeben:

\begin{verbatim}
function line (a : *char, ip : int) : int

  l : int = 0;
  c : int = 0;
  
begin
  repeat
    c = c + a[l];
    l = l + 1;
  until (c > ip)
  return l;
end
\end{verbatim}

In einem rekursiven Interpreter ist neben dem Erkennen und
Anzeigen eines Fehlers noch ein weiteres Problem zu "uberwinden:
das au"serordentliche Beenden des Interpreters. Es ist davon
auszugehen, da"s der Interpreter schon eine gewisse
Rekursionsebene erreicht hat, wenn der Fehler auftritt, soda"s es
nicht so einfach ist, ihn zu beenden und evtl. zu einer (den
Interpreter) aufrufenden Funktion zur"uckzukehren. Die
ANSI-Bibliothek von C bietet hierzu das praktische Konzept der
"`{\em Long-Jumps}"' an, das ungef"ahr einem \rw{goto}
gleichkommt, das nicht nur Spr"unge innerhalb einer Funktion
erlaubt, sondern auch aus einer Funktion in eine andere,
"ubergeordnete Funktion hinein. In C++ z.~B. ist es nicht einmal
mehr n"otig, sich auf Bibliotheksfunktionen zu st"utzen, da ein
"ahnliches Konzept (das sog. "`{\em Exception-Handling}"') in der
Programmiersprache selbst eingebaut ist. Leider besitzen nicht
alle Sprachen einen derartigen Luxus. In \pascal\ etwa mu"s man
sich einer anderen Methode bedienen. Eine einfache und
naheliegende L"osung ist, eine Variable zu definieren, die
angibt, ob bei der Ausf"uhrung des Programms bisher ein Fehler
aufgetreten ist. Eine Funktion, die einen Fehler entdeckt, hat
diesen zu melden und die Variable auf einen Wert zu setzen, der
anzeigt, da"s ein Fehler gefunden wurde. Kehrt diese Funktion zu
der aufrufenden Funktion zur"uck, mu"s diese "uberpr"ufen, ob die
Variable gesetzt ist und in diesem Fall zu ihrer aufrufenden
Funktion zur"uckkehren, so lange, bis man an der Funktion
angelangt ist, die die Ausf"uhrung des Programms initiiert hat.
Zwar ist diese Methode simpel, doch hat sie einen gewissen
Overhead zur Folge, der nicht nur den Code gr"o"ser, sondern auch
un"ubersichtlicher macht.
                                   
\newpage
\section{Weiterf\protect{"uhrende} Projekte}

Ebenso wie im ersten Kapitel m"ochte ich auch hier einige meiner
Ideen darlegen, die leider zu umfangreich sind, als da"s ich sie
ausf"uhrlich in dieser Arbeit h"atte behandeln k"onnen.

\begin{description}
\item[Verwendung einer Baumstruktur] anstatt des Zwischencodes
k"onnte die Geschwindigkeit des Interpreters um einiges erh"ohen.
Diese Baumstruktur w"urde in gewisser Weise einem
Parse-Tree\footnote{siehe Kapitel \ref{topdbotu}} "ahneln, bei
dem jedoch die unn"otigen Bl"atter entfernt sind. Zus"atzlich
k"onnte vor der Ausf"uhrung des Programms der Baum optimiert
werden, z.~B. durch Vereinfachung von konstanten Ausdr"ucken
(\fu{2*3+4} kann ohne Ver"anderung der Laufzeiteigenschaften zu
\fu{10} umgeformt werden). Eine Darstellung als Baum hat f"ur den
Interpreter den gro"sen Vorteil, da"s die lexikalische Analyse
und der Parsevorgang nur einmal ausgef"uhrt werden m"ussen. Ist die
Baumstruktur gut konzipiert, k"onnte der Parser sogar als
Front-End f"ur einen Compiler dienen.

\item[Ein Compiler f"ur SEL] ist mein n"achstes Ziel. Leider sind
jedoch die meisten Maschinensprachen entweder zu komplex (wie
etwa die des Intel 80386) oder zu umst"andlich (wie die des Intel
8086), als da"s es Spa"s machen k"onnte, f"ur diese Prozessoren
einen Compiler zu schreiben. Aus diesem Grund erw"age ich zwei
M"oglichkeiten, um diese Schwierigkeiten zu umgehen.
Die einfachere der beiden ist, ein SEL-Programm anstatt in
Maschinencode in eine andere h"ohere Programmiersprache,
vorzugsweise C, zu "ubersetzen. Besonders interessant w"are es,
diesen Compiler selbst in SEL zu schreiben und ihn dann durch sich
selbst "ubersetzen zu lassen (mit Hilfe des Interpreters).

Die zweite M"oglichkeit ist,
anstatt f"ur einen real existierenden Prozessor Maschinencode
f"ur einen einfachen, auf die Programmiersprache zugeschnittenen,
hypothetischen Prozessor zu erzeugen, wobei dieser Prozessor
nat"urlich zun"achst entworfen und auf einem existierenden
Prozessor emuliert werden m"u"ste. Obwohl auch hier ein Interpreter
im Spiel ist, ist die Ausf"uhrungsgeschwindigkeit eines
Pseudo-Code weit h"oher als die eines Interpreters, der den
Quellcode (oder auch einen Zwischencode) der Programmiersprache
interpretiert. Der Vorteil dieser beiden M"oglichkeiten ist, da"s
sie gr"o"stenteils prozessorunabh"angig sind. So existiert heute
beinahe f"ur jeden Prozessor eine ausreichende Menge an
C-Compilern und f"ur den P-Code m"u"ste lediglich der Interpreter
neu geschrieben, oder, falls auch er in einer h"oheren
Programmiersprache geschrieben ist, lediglich neu compiliert
werden. Aus Effizienzgr"unden ist es jedoch ratsam, einen
Interpreter f"ur P-Code in Assembler zu implementieren.

\item[Ein Pretty-Printer] ist ein Programm, das einen Quelltext
einer bestimmten Programmiersprache in eine ansprechende Form
bringt, also Einr"uckungen vornimmt, wenn n"otig, und Kommentare
einheitlich gestaltet. F"ur C existiert eine Reihe solcher
Programme, das bekannteste ist wohl C-Indent. Es ist eine
interessante, wenn auch nicht "uberm"a"sig schwere Aufgabe, einen
solchen Pretty-Printer f"ur SEL zu schreiben.
\end{description}

\begin{appendix}

\chapter{Die Grammatik von SEL}
\label{grammar}

\begin{verbatim}
module:        {block}
             ;
block:         function
             | local
             | global
             ;
function:      `function' IDENTIFIER `(' [formal_params] `)' `:' type
                 {func_var}
               `begin'
                 {statment}
               `end'
             ;
formal_params: IDENTIFIER `:' [`&']  type [`,' formal_params]
             ;
func_var:      IDENTIFIER `:' type [`=' expr] `;'
             ;
func_proto:    `function' IDENTIFIER
               `(' [[`&' type {`,' [`&'] type}] `)' `:' type
             ;
local:         `local' {local_decl} `end'
             ;
local_decl:    struct
             | func_proto
             | var
             ;
global:        `global' {global_decl} `end'
             ;
global_decl:   struct
             | func_proto
             | var
             | extern
             ;
\end{verbatim}
\newpage
\begin{verbatim}
statement:     [expr] `;'
             | if
             | while
             | for
             | repeat
             | switch
             | prefix
             | `break' `;'
             | `return' [expr] `;'
             ;
if:            `if' `(' expr `)' {statement} [`else' {statement}] `end'
             ;
switch:        `switch' `(' expr `)' {statement} `end'
             ;
prefix:        `case' constant `:'
             | `default' `:'
             ;
while:         `while' `(' expr `)' {statement} `end'
             ;
repeat:        `repeat' {statement} `until' `(' expr `)'
             ;
for:           `for' `(' primary `=' expr (`to' | `downto') expr
                         [`step' expr] `)'
                 {statement}
               `end'
             ;
expr:          expr (`=' | `&&' | `||' | `^^' | `==' | `<' | `>'
                   | `<=' | `>=' | `!=' | `&' | `|' | `^' | `<<'
                   | `>>' | `+' | `-' | `*' | `/' | `%') expr
             | (`!' | `~' | `-' | `*' | `&') expr
             | `delete' expr
             | primary
             ;
primary:       IDENTIFIER
             | `(' expr `)'
             | CHAR_CONSTANT
             | INT_CONSTANT
             | FLOAT_CONSTANT
             | STRING_CONSTANT
             | primary `->' IDENTIFIER
             | primary `:' type
             | primary `[' expr `]'
             | IDENTIFIER `(' [expr {`,' expr}] `)'
             | `sizeof' type
             | `new' type
             ;
var:           IDENTIFIER `:' type [`=' constant] `;'
             ;
extern:        `extern' IDENTIFIER `:' type `;'
             ;
\end{verbatim}
\newpage
\begin{verbatim}
struct:        `struct' IDENTIFIER {IDENTIFIER `:' type `;'} `end'
             ;
type:          `char'
             | `int'
             | `float'
             | `void'
             | `[' INT_CONSTANT `]' type
             | `*' type
             | IDENTIFIER
             ;
constant:      CHAR_CONSTANT
             | [`-'] INT_CONSTANT
             | [`-'] FLOAT_CONSTANT
             | STRING_CONSTANT
             ;
INT_CONSTANT:    [0-9]+
CHAR_CONTANT:    '(.|\.)'
FLOAT_CONSTANT:  [0-9]*(`.'[0-9]*)?((e|E)(+|-)?[0-9]+)?
STRING_CONSTANT: "([^"\n]|\.)*"
IDENTIFIER:      [A-Za-z_][A-Za-z0-9_]*
\end{verbatim}

\end{appendix}

\begin{thebibliography}{9}
\addcontentsline{toc}{chapter}{Literaturverzeichnis}
\bibitem{selkr} Mark Probst. {\em SEL-Kurzreferenz}, Version~0.3.
Beilage.
\bibitem{kandr} Brian W.~Kernighan, Dennis M.~Ritchie. {\em
Programmieren in C}, deutsche Ausgabe bearbeitet
von A.~T.~Schreiner, Ernst Janich. Carl Hanser Verlag, Wien,~1983.
ISBN~3-446-13878-1.
\bibitem{sedge} Robert Sedgewick. {\em Algorithmen}, deutsche
Ausgabe. Addison-Wesley~GmbH, Deutschland, 1.~Auflage~1991.
ISBN~3-89319-301-4.
\bibitem{duden} {\em Duden Informatik}, 2.~Auflage. Dudenverlag,
Mannheim,~1993. ISBN~3-411-05232-5.
\end{thebibliography}

% \newpage
% \addcontentsline{toc}{chapter}{Index}
% \printindex

\end{document}

\begin{document}

\title{Design einer Programmiersprache\\
       und theoretische Besprechung\\
       eines Interpreters anhand der\\
       Programmiersprache SEL}
\author{Mark Probst}

\maketitle

\tableofcontents

\chapter{Die Programmiersprache}

Dieses Kapitel hat das Ziel, das Design der Programmiersprache
SEL, wie in \cite{selkr} beschrieben, zu erkl"aren bzw. zu
begr�nden, und zwar in Hinblick auf den Bau eines Interpreters
und eines Compilers. Aus diesem Grund setze ich eine gewisse
Kenntnis sowohl einer strukturierten Programmiersprache,
insbesondere nat"urlich der Sprache SEL, als auch der
grundlegenden Funktionsweise eines Interpreters bzw. eines
Compilers beim Leser voraus. Der Leser sollte also zumindest
\cite{selkr} gelesen haben, da"s, obwohl es als Kurzreferenz
konzipiert ist, durchaus auch geeignet ist, einen "Uberblick
"uber die Programmiersprache SEL zu vermitteln.

\section{Der Begriff der Programmiersprache}

Grunds"atzlich kann eine Programmiersprache definiert werden als

\begin{quotation}
Sprache zur Formulierung von Rechenvorschriften, d. h. von
Datenstrukturen und Algorithmen, die von einem Computer
ausgef"uhrt werden k"onnen.\footnote{siehe \cite{duden},
Seite~544}
\end{quotation}

Die in dieser Arbeit behandelte Programmiersprache SEL geh"ort
zur Gruppe der imperativen Programmiersprachen, deren wichtigste
Untergruppe wohl die der strukturierten Programmiersprachen ist.
Imperative Programmiersprachen werden im Duden Informatik
wie folgt definiert:

\begin{quotation}
Bei diesen Sprachen besteht ein Programm aus einer Folge von
Befehlen an den Computer, wie z.~B. "`Schreibe in die Variable a
den Wert 3"', "`Springe an die Stelle x im
Programm"'.\footnote{siehe \cite{duden}, Seite~545}
\end{quotation}

Strukturierte Programmiersprachen stellen f"ur den Programmierer
zus"atzlich Hilfen in Form von Kontrollanweisungen, Schleifen und
benutzerdefinierten Funktionen zur Verf"ugung. Diese machen
sowohl das Schreiben der Programme einfacher und schneller wie
auch die Programmtexte leichter lesbar. Zur Zeit sind sie die am
meisten verwendete Art der Programmiersprachen.

\section{Ziele f\protect{"ur} das Design von SEL}

Bevor man an die Arbeit an einem so komplexen Projekt wie das Design
einer Programmiersprache herangeht, sollte man sich im Klaren
dar"uber sein, welche Ziele man eigentlich verfolgt. So habe ich
auch an meine zuk"unftige Programmiersprache einige Anforderungen
gestellt, von denen ich die meisten in der gegenw"artigen Version
der Sprache verwirklicht sehe.

Die wichtigste Entscheidung bei einem Projekt ist die nach dem
Zweck desselben. Im Fall von SEL wollte ich weniger eine
Programmiersprache f"ur das Schreiben von kommerzieller Software
schaffen, als vielmehr eine kleine, feine Sprache, die ich in
vielen Situationen benutzen k"onnte. Haupts"achlich schwebte mir
dabei die Verwendung als Makrosprache in diversen Programmen vor.
Eine solche Makrosprache sollte vor allem leicht erlernbar sein
und keine gro"se Umstellung von Seiten des Benutzers verlangen,
sollte aber andererseits nicht so primitiv aufgebaut sein, da"s
sie eine vern"unftige Verwendung nicht mehr m"oglich macht, was
die meisten Makrosprachen leider sind.

Um eine geringe Umstellungszeit m"oglich zu machen, mu"s die
Programmiersprache an eine bereits bestehende, weit verbreitete
Sprache angelehnt sein. Im Fall von SEL habe ich mich
haupts"achlich an C und \pascal\ gehalten, weil mir diese am
gel"aufigsten und vermutlich die am weitest verbreiteten
Sprachen sind.

Da ich ein gro"ser Freund von C und ihrer pr"agnanten und kurzen
Ausdrucksweise bin, habe ich mir zum Ziel gesetzt, auch SEL zu
einer effizienten, kurzen Ausdrucksweise zu verhelfen. Eine
Konsequenz aus dieser Entscheidung ist es, anstatt von
Schl"usselw"ortern, so wie in \pascal , Symbole wie in C zu
verwenden, soferne dies nicht die "Ubersichtlichkeit und
Konsistenz innerhalb der Sprache zu sehr gef"ahrdet. Die zweite
Konsequenz ist es, redundante Sprachteile zu entfernen, solange
auch das nicht die Lesbarkeit der Programmtexte verletzt.

Ein weiteres wichtiges Kriterium bei einer Makrosprache ist die
Erweiterungsf"ahigkeit, da Makrosprachen in so verschiedenen
Anwendungen wie z.~B. Textverarbeitungen und Zeichenprogrammen
eingesetzt werden, da"s es schlichtweg unm"oglich ist, eine
Sprache zu konstruieren, die f"ur alle diese Anwendungsgebiete
passend ist. Diese Erweiterungsf"ahigkeit sollte einerseits durch
die Erweiterung mit Hilfe von Funktionen wie z.~B.
f"ur das Zeichnen von Linien oder Kreisen und andererseits durch
die M"oglichkeit der "Anderung in der Grammatik, wie etwa
f"ur das flexiblere Handling von Vektoren und Matrizen, gegeben
sein. Einige Vorschl"age f"ur "Anderungen der letzteren Art sind
in Kapitel \ref{langext} gegeben.

Zuletzt soll die Programmiersprache SEL so aufgebaut sein, da"s
es nicht nur m"oglich ist, einen Interpreter, sondern auch mit
vertretbarem Aufwand einen Compiler zu implementieren. Nicht nur,
weil das bei einer Programmiersprache immer eine gute Idee ist,
sondern auch in Hinblick darauf, da"s ich selbst einmal einen
Compiler schreiben m"ochte.

Als Abschlu"s dieses Unterkapitels seien nochmals alle Forderungen
kurz zusammengefa"st, um in sp"ateren Abschnitten darauf
zur"uckgreifen zu k"onnen.

\begin{itemize}
\item Geringer Sprachumfang, jedoch auch f"ur mittelgro"se
Programme geeignet
\item Anlehnung an C und \pascal\
\item Verwendung von Symbolen anstatt von Schl"usselw"ortern
(wenn m"oglich)
\item Vermeidung von Redundanzen
\item Konsistenz innerhalb der Sprache
\item Erweiterungsf"ahigkeit
\item Sowohl f"ur Interpreter als auch f"ur Compiler geeignet
\end{itemize}

\section{Das Programm}

Das gr"o"ste in sich abgeschlossene Objekt, da"s ein Compiler bzw.
Interpreter zu behandeln hat, ist das Programm. Uns stellt sich nun die
Frage, wie dieses Programm aufgebaut sein soll.

Zun"achst ist zu
entscheiden, ob das Programm in Module eingeteilt werden soll oder
nicht. Eine solche Einteilung hat den Vorteil, da"s der Programmierer
Implementationsdetails eines Programmteils oder einer Funktionssammlung
vom Rest des Programms abschirmen kann und den Anwender dieser
Funktionssammlung gar nicht erst in Versuchung bringt, in die internen
Abl"aufe einzugreifen. Nat"urlich ist das nicht der einzige Vorteil
dieser Technik. So kann es z.~B. bei einer Modularisierung des Programms
auch kaum zu Namens"uberschneidungen bei Funktionen und Variablen
kommen. Last but not least wird auch das Programm leichter
"uberschaubar, weil man sich jeweils nur auf ein Modul konzentrieren
mu"s. Ich habe mich f"ur eine Modularisierung entschieden. Erstens, weil
es die M"oglichkeiten des Programmierers um einen nicht zu
untersch"atzenden Faktor erweitert und zweitens, weil es nicht schwer
ist, ein solches Modulkonzept in einen Interpreter oder einen Compiler
einzubauen---jedenfalls nicht das hier verwendete.

In SEL wird ein Modul in mehrere Bl"ocke eingeteilt, die in beliebiger
Reihenfolge auftreten d"urfen, wobei es drei Arten
von Bl"ocken zu unterscheiden gibt:

\begin{itemize}
\item Die \rw{global}-Bl"ocke werden verwendet, um programmglobale Variablen,
Funktionen oder Datentypen
zu deklarieren, also solche, die im ganzen Programm bekannt
sind.
\item Die \rw{local}-Bl"ocke werden analog zu \rw{global}-Bl"ocken
verwendet, um
modulglobale Variablen, Funktionen oder Datentypen zu deklarieren.
\item Funktionen werden "ahnlich wie in C und \pascal\ definiert.
\end{itemize}

Der Hauptgrund, warum ich mich f"ur diese Art der Modulstruktur
entschieden habe ist der, da"s es f"ur einen Interpreter weit
praktikabler ist, zur Initialisierung der Variablen vor dem Start eines
Programms einfach nach den Schl"usselw"ortern \rw{global} und
\rw{local} zu suchen, als zwischen den Funktionsdefinitionen. Dieses
Problem h"atte nat"urlich auch eine Programmstruktur wie die in \pascal ,
wo Variablen- und Typendeklarationen ihren festen Platz haben, gel"ost,
doch erscheint mir dieser Ansatz als zu starr, vermutlich, weil ich die
Freiheiten von C gewohnt bin.

Noch ein kleines Wort zur allgemeinen Struktur eines Moduls: Wie
bereits oben erw"ahnt, besteht ein Modul aus einer beliebigen
Abfolge von Bl"ocken, wobei es durchaus m"oglich ist, da"s
mehrere Bl"ocke ein und desselben Typs vorkommen. F"ur das
gesamte Modul gilt jedoch eine Regel: Ein Objekt ist erst ab der
Zeile bekannt, in der es deklariert wurde, das hei"st, da"s es
z.~B. nicht m"oglich ist, eine Variable in Zeile 154 zu
verwenden, wenn sie erst in Zeile 186 deklariert wird.
Diese Regelung macht es einem Compiler m"oglich, ein Programm in
einem Durchgang zu "ubersetzen ({\em One-Pass Compiler}). Eine
Ausnahme bilden die Strukturen (siehe Kapitel \ref{structs}).

\subsection{global-Bl\protect{"ocke}}

Ein \rw{global}-Block\footnote{siehe Anhang \ref{grammar} oder
\cite{selkr}, Seite 9} wird durch das Schl"usselwort \rw{global}
eingeleitet und durch \rw{end} beendet. Dazwischen befindet sich
eine beliebige Abfolge von Variablen-, Funktions- und
Strukturdeklarationen. Alle in diesem Abschnitt deklarierten
Objekte sind dazu bestimmt, auch in anderen Modulen des Programms
verwendet zu werden, was jedoch nicht unbedingt bedeutet, da"s
sie in diesen auch ohne weiteres verf"ugbar sind. Soll ein
programmglobales Objekt auch in einem anderen Modul verwendet
werden, so ist es n"otig, es auch dort zu deklarieren und---im
Falle von Variablen---als \rw{extern} zu bezeichnen. Der Grund
f"ur diese Vorgangsweise ist, da"s sich ein Compiler, der die
Deklarationen aller anderen Module verwenden kann, nicht
konstruieren l"a"st, ohne nicht ein eigenes Objektformat zu
benutzen, wie das z.~B. bei Turbo \pascal\ der Fall ist. Das
Standard-Objektformat von MS-DOS und Windows NT ist dazu nicht in
der Lage. Au"serdem w"are es erheblich schwerer, Module
zu schreiben, die gegenseitig voneinander abh"angig sind, soda"s
ich von dieser L"osung Abstand genommen habe.

Auffallend an den drei Arten der Deklarationen, die innerhalb
eines \rw{global}-Blocks vorkommen d"urfen, ist, da"s die
Variablendeklaration als einzige mit einem Semikolon
abgeschlossen wird, und das, obwohl es ohne Gef"ahrdung der
Eindeutigkeit m"oglich w"are, dieses wegzulassen. Der Grund
hierf"ur ist einerseits die Einhaltung der Konsistenz der
Sprache\footnote{siehe auch Kapitel \ref{funcs}} und andererseits
die Offenhaltung f"ur Erweiterungen der Sprache. So ist es im
Moment zwar lediglich m"oglich, eine Variable mit einer \tchar -,
\tint -, \tfloat - oder Stringkonstante zu initialisieren, jedoch
ist es durchaus denkbar, dies in zuk"unftigen Versionen auf
Arrays, Strukturen und Zeiger auszuweiten, wo die Eindeutigkeit
nicht mehr unbedingt gegeben sein mu"s.

Der Grund, warum Variablen, die in einem anderen Modul definiert
worden sind und in einem bestimmten Modul verwendet werden, als
\rw{extern} zu deklarieren sind und dies bei Strukturen und
Funktionen nicht n"otig ist, ist einfach: Bei einer
Strukturdeklaration wird dem Compiler bzw. Interpreter lediglich
"`mitgeteilt"', wie diese aufgebaut ist. Es wird in diesem Fall
kein vom Programm verwendetes Objekt im Speicher erzeugt. Bei
einer Variablendefinition hingegen wird dem Compiler sowohl der
Typ der Variable mitgeteilt als auch der Befehl, sie im Speicher
zu erzeugen. Ohne Unterscheidung zwischen Variablendefinition und
\rw{extern}-Deklaration w"u"ste der Compiler nicht, ob er nun
eine Variable im Speicher zu erzeugen h"atte oder sich lediglich
vormerken sollte, da"s diese Variable in einem anderen Modul
existiert und in diesem verwendet wird. Bei Funktionen ist es
kaum anders: Eine Deklaration innerhalb eines \rw{local}- oder
\rw{global}-Blocks sagt dem Compiler lediglich, da"s es
"`irgendwo anders"' eine solche Funktion gibt und da"s sie hier
verwendet wird. Eine Funktionsdeklaration innerhalb eines
\rw{global}-Blocks dient weiterhin dazu, zu kennzeichnen, da"s
diese Funktion, falls sie in diesem Modul implementiert ist, f"ur
die Verwendung in anderen Modulen bestimmt ist. Eine zweite
M"oglichkeit, dies zu erzielen, w"are, dem Funktionskopf bei der
Implementation noch ein spezielles Schl"usselwort hinzuzuf"ugen,
doch erscheint mir die in SEL verwendete L"osung als die bessere
und auf jeden Fall als die "ubersichtlichere.

Variablen d"urfen in \rw{local} bzw. \rw{global}-Bl"ocken
lediglich mit Konstanten, in
Funktionen hingegen auch mit Ausdr"ucken initialisiert werden.
Dies ist einerseits der Einfachheit halber so, andererseits, weil
es immer eine riskante Sache ist, Werte zu berechnen, noch bevor
die einzelnen Teile des Programms initialisiert sind. Dar"uber hinaus w"are
es notwendig, eine Regelung zu erstellen, die angibt, welche
Variablen als erste initialisiert werden und welches Modul einem
anderen vorangeht. Dieser Aufwand erscheint mir allerdings etwas
"ubertrieben, nicht zuletzt deshalb, weil eine solche
Initialisierung leicht in der \fu{main}-Funktion\footnote{siehe
Kapitel \ref{main}} nachgeholt werden kann.

\subsection{local-Bl\protect{"ocke}}

Ein \rw{local}-Block\footnote{siehe Anhang \ref{grammar} oder
\cite{selkr}, Seite 9} unterscheidet sich von einem
\rw{global}-Block, abgesehen davon, da"s er durch ein anderes
Schl"usselwort eingeleitet wird, dadurch, da"s eine
\rw{extern}-Deklaration nicht erlaubt ist. Dies ist auch ganz
verst"andlich, denn es ist ein Widerspruch in sich, eine
modulglobale Variable, die in einem anderen Modul definiert
ist, zu verwenden.

Die Tatsache, da"s Funktionen, die nicht explizit in einem
\rw{global}-Block aufscheinen, automatisch modulglobal sind,
verleitet zum Schlu"s, da"s Funktionsdeklarationen in
\rw{local}-Bl"ocken redundant sind. Es k"onnte jedoch ohne
weiteres die Notwendigkeit entstehen, da"s zwei Funktionen sich
gegenseitig aufrufen. In diesem Fall w"are jedoch ein
Aufruf illegal, da Funktionen erst ab dem Punkt ihrer
Deklaration im Modul bekannt sind. Dieser Konflikt ist sehr
einfach zu l"osen, indem man eine der Funktionen (oder der
"Ubersichtlichkeit halber beide) in einem \rw{local}-Block
deklariert und diese erst nach der anderen implementiert. Ein
Beispiel: Die Funktion \fu{A} und \fu{B} rufen sich gegenseitig
auf. Nun wird vor der Implementation die Funktion \fu{A} als
\rw{local} deklariert, soda"s sie im Rest des Moduls bekannt ist.
Nun wird die Funktion \fu{B} implementiert. Der Aufruf von \fu{A}
stellt kein Problem dar, der oben geschilderten Deklaration
wegen. Die Implementation von \fu{A} ist nun auch kein
Problem mehr, da an diesem Punkt schon beide Funktionen bekannt
sind.

\subsection{Funktionen}
\label{funcs}

Im Gegensatz zu \pascal\ wird in SEL, so wie in C, nicht zwischen
Funktionen und Prozeduren unterschieden. Dies hat---anders als es
zu erwarten w"are---fast nur Vorteile. Auf der einen Seite ist es
m"oglich, Prozeduren durch Funktionen mit dem R"uckgabetyp
\tvoid \footnote{siehe Kapitel \ref{elemtypes}} nachzubilden,
andererseits vereinfacht es den Bau eines
Interpreters bzw. eines Compilers. Der einzige Nachteil ist ein
geringer Verlust an Lesbarkeit.

Eine Funktion\footnote{siehe Kapitel \ref{exprfuncs} und
Anhang \ref{grammar} oder
\cite{selkr}, Seiten 7--9 und 17} gliedert sich augenscheinlich in
drei gro"se Teile:
Den Funktionskopf, die Definitionen der direkten
Funktionsvariablen\footnote{siehe Kapitel \ref{variables}}
und den Funktionsk"orper.

Der Funktionskopf besteht aus dem Schl"usselwort \rw{function},
dem Namen der Funktion, einer in Klammern eingeschlossenen Liste
der Parameter inklusive deren Typen und dem R"uckgabetyp. Der
Doppelpunkt vor dem Typ soll andeuten, da"s eine Funktion in
gewisser Weise einer Variablen eines bestimmten Typs "ahnelt.
Es sei angemerkt, da"s dieser Doppelpunkt ohne weiteres
ausgelassen werden k"onnte, ohne die Eindeutigkeit der Sprache zu
gef"ahrden.\footnote{Tats"achlich k"onnte sogar die schlie"sende
Klammer ausgelassen werden, doch diente das sicherlich nicht der
Lesbarkeit.}

Bei genauerer Betrachtung des Funktionskopfes und der
Funktionsdeklaration in einem \rw{global} bzw. \rw{local}-Block
f"allt auf, da"s in einem Funktionskopf die Namen der
Parameter aufscheinen, in den beiden letztgenannten Bl"ocken dies
jedoch nicht der Fall ist. Betrachten wir zur Begr"undung dieser
Entscheidung einen Interpreter und einen Compiler. Der
Interpreter kann alle Module des Programms auf einmal betrachten,
er kennt also die Namen der Parameter auch dann, wenn sie in
einem anderen Modul erw"ahnt sind als in dem, von dem aus die
Funktion aufgerufen wird. Ein Compiler jedoch besitzt diesen
Vorteil nicht: Er beschr"ankt sich jeweils auf ein einzelnes
Modul. Bei einem Funktionsaufruf ist aber das Einzige, was der
Compiler zu tun hat, zu "uberpr"ufen, ob die Parameter die
richtigen Typen haben, um sie dann in der richtigen Reihenfolge
auf dem Stack abzulegen. Wirkliche Bedeutung haben die Namen der
Parameter nur f"ur die eigentliche Funktion und insofern stellt
die Beschr"ankung auf die Typen keine Einschr"ankung dar, es ist
jedoch durchaus denkbar, da"s es der "Ubersichtlichkeit wegen in
einer zuk"unftigen Version von SEL m"oglich sein wird, auch die
Namen anzugeben, wobei diese Informationen aber vom Interpreter
bzw. Compiler ignoriert werden.

Dem Funktionskopf folgen die Definitionen der direkten
Funktionsvariablen.
Diese werden bei jedem Funktionsaufruf neu auf dem Stack angelegt
und belegen daher bei einer rekursiven Funktion in den
verschiedenen Instanzen nicht denselben Speicherbereich---erst
diese Tatsache macht Rekursionen so effektiv. Es ist m"oglich,
die Variablen mit Werten zu initialisieren, die sich aus Ausdr"ucken ergeben.
Dies stellt weder einen Compiler noch einen Interpreter vor
Probleme, soda"s ich hier nicht n"aher darauf eingehen m"ochte.
Das Semikolon am Ende jeder Variablendefinition w"are nicht
unbedingt notwendig, ergibt jedoch sowohl in Hinsicht auf die
"Ubersichtlichkeit als auch auf die Erweiterungsf"ahigkeit der
Sprache durchaus Sinn.

Der Funktionsk"orper wird durch das Schl"usselwort \rw{begin}
eingeleitet und durch \rw{end} beendet. Dazwischen befindet sich
eine beliebige Anzahl von Anweisungen, die bei Aufruf der
Funktion ausgef"uhrt werden.

\subsubsection{Die main-Funktion}
\label{main}

Jedes Programm mu"s einen Punkt haben, an dem es beginnt, soda"s
nicht unkontrolliert irgendwelche Funktionen ausgef"uhrt werden.
Hierf"ur gibt es grunds"atzlich zwei M"oglichkeiten. Entweder,
man f"uhrt bestimmte Konstrukte ein, um diesen Hauptteil des
Programms zu kennzeichnen, so wie es z.~B. in \pascal\
praktiziert wird, oder man verwendet eine bestimmte Funktion
dazu, wie in C. Ich habe mich f"ur letzteren
L"osungsansatz entschieden, weil er es erlaubt, das gesamte
Programm als eine Art Funktion zu sehen, die vom Anwender
aufgerufen werden kann. Es wird auch m"oglich, die
Kommandozeilenparameter als einfache Funktionsparameter zu
"ubergeben---wenn dies auch in der aktuellen Version von SEL
noch nicht eingef"uhrt ist. Weiters kann durch dieses Prinzip dem
Betriebssystem ein Wert zur"uckgegeben werden---der sogenannte
{\em Exit-Code}. In \pascal\ sind hierf"ur bestimmte Funktionen
bzw. Prozeduren notwendig.

Der Einfachheit halber und wegen der Verwandschaft mit C wird die
Hauptfunktion des Programms \fu{main} genannt. Sie nimmt keine
Parameter an und hat entweder den R"uckgabetyp \tvoid\ oder
\tint . Ist ersteres der Fall, wird immer der Exit-Code
\fu{0} zur"uckgegeben, andernfalls wird dieser durch die
\rw{return}-Anweisung bestimmt.

\section{Anweisungen}
\label{stmts}

Anweisungen kontrollieren das Geschehen innerhalb eines
Programms. Sie befehlen\footnote{daher die Bezeichnung
"`imperative Programmiersprache"'} dem Computer, was er tun soll.
In SEL wird grunds"atzlich zwischen drei Arten von Anweisungen
unterschieden: Kontrollanweisungen, Schleifen und Ausdr"ucken.
W"ahrend die ersten beiden in diesem Unterkapitel besprochen
werden, widmet sich das gesamte n"achste Unterkapitel den
Ausdr"ucken. Der Korrektheit wegen sei hier gesagt, da"s es
eigentlich noch eine vierte Art von Anweisungen gibt: Die {\em
Nullanweisung}, die jedoch, da sie keinen bestimmten Zweck
erf"ullt, au"ser vielleicht die Lesbarkeit des Programms zu
erh"ohen, f"ur unsere Betrachtungen nicht weiter wichtig ist.

Es ist auffallend, da"s SEL keine Anweisungen zur Verf"ugung
stellt, die in irgendeiner Weise mit der Hardware des Computers
kommunizieren k"onnen, wie z.~B. Anweisungen zum Einlesen von der
Tastatur, oder zur Ausgabe von Zeichen auf den Bildschirm oder in
eine Datei. Die Gr"unde hierf"ur sind vielf"altig. Erstens haben
wir es uns zum Ziel gesetzt, den Sprachumfang von SEL m"oglichst
gering zu halten und Ein/Ausgabeanweisungen alleine
w"urden---sollten sie einigerma"sen vern"unftig sein---die
Grammatik von SEL betr"achtlich aufbl"ahen. Weiters habe ich SEL
mit dem Ziel vor Augen entworfen, sie einmal als Makrosprache
einzusetzen. Makrosprachen allerdings m"ussen in vielen
verschiedenen Umgebungen sinnvoll einzusetzen sein, wie z.~B. in
einer Textverarbeitung oder in einem Zeichenprogramm, das unter
einer graphischen Benutzeroberfl"ache l"auft. Nun sind aber die
M"oglichkeiten zur Textausgabe in einer graphischen
Benutzeroberfl"ache grundverschieden zu denen von z.~B. DOS.
Ausgabeanweisungen zu entwerfen, die f"ur alle diese
M"oglichkeiten halbwegs geeignet sind, erscheint mir fast
unm"oglich. Deshalb ist es unumg"anglich, jede Implementation von
SEL mit verschiedenen Grundfunktionen auszustatten, die auf die
jeweilige Umgebung zugeschnitten sind. Eine Standardisierung
dieser Grundfunktionen ist immer eine gute Idee, doch sicher sind
diese Standardfunktionen nicht "uberall anwendbar.\footnote{Ein
Beispiel hierf"ur w"are der ANSI-Standard von C, dessen
Ein/Ausgabefunktionen auf eine textorientierte Oberfl"ache
zugeschnitten sind. Die \fu{printf}-Funktion z.~B. ist unter der
graphischen Benutzeroberfl"ache Windows nutzlos.}

\bigskip

An dieser Stelle scheint es mir wichtig, ein Konzept aus anderen
strukturierten Programmiersprachen zu erw"ahnen. Das
Blockkonzept\footnote{nicht zu verwechseln mit den Bl"ocken in
SEL} erm"oglicht es, eine Anweisung entweder als eben nur eine
Anweisung oder als eine von zwei Schl"usselw"ortern oder Symbolen
begrenzte, beliebige Abfolge von Anweisungen zu sehen. So verlangt
eine \rw{if}-Anweisung in C oder \pascal\ nur eine Anweisung als
Zweig, es ist jedoch kein Problem, mit Hilfe der Blockstruktur
auch mehrere Anweisungen ausf"uhren zu lassen. Dieses
Blockkonzept ist in manchen Situationen sehr praktisch, aber es
bl"aht die Grammatik einer Programmiersprache auch etwas auf und
macht es so nat"urlich schwerer, einen Interpreter oder Compiler
zu schreiben. Aus diesem Grund habe ich mich entschieden, es in
SEL nicht einzubauen. Die Quelltexte sind so zwar etwas l"anger,
doch es entsteht keine wirkliche Beeintr"achtigung.

\subsection{Kontrollanweisungen}

Kontrollanweisungen haben die Aufgabe, zu kontrollieren, ob und
was ausgef"uhrt wird. In SEL gibt es drei Anweisungen dieser Art.
Die \rw{if}-Anweisung "uberpr"uft eine Aussage und f"uhrt, wenn
sie einen positiven Wahrheitswert ergibt, bestimmte Anweisungen
aus. Auf Wunsch werden auch Anweisungen ausgef"uhrt, wenn das
Ergebnis der "Uberpr"ufung negativ ist. Die \rw{switch}-Anweisung
f"uhrt f"ur bestimmte Werte eines Ausdrucks verschiedene
Anweisungen aus. Die \rw{return}-Anweisung springt zur
aufrufenden Funktion zur"uck.

\subsubsection{Die if-Anweisung}

Die \rw{if}-Anweisung\footnote{siehe Anhang \ref{grammar} oder
\cite{selkr}, Seite 10} wird durch das Schl"usselwort \rw{if}
eingeleitet, auf welches ein in Klammern eingeschlossener Ausdruck
folgt. Die Alternative zu dieser L"osung w"are gewesen, den
Ausdruck sofort auf das \rw{if} folgen zu lassen und durch ein
bestimmtes Schl"usselwort (z.~B. \rw{then}) oder ein Symbol zu
beenden. Der Doppelpunkt jedoch w"urde Mehrdeutigkeiten
hervorrufen und ein Semikolon erscheint mir als unpassend, und da
wir uns entschienden haben, Symbole zu verwenden, wo nicht
unbedingt Schl"usselw"orter erforderlich sind, f"allt auch das
\rw{then} weg. Die einfachste und unauff"alligste L"osung in
diesem Fall ist die Klammerung.

Auf die schlie"sende Klammer folgt eine beliebige Anzahl von
Anweisungen, beendet entweder durch \rw{end} oder \rw{else}. Da
auch \rw{if} eine Anweisung ist, ist es m"oglich, \rw{if}s zu
verschachteln. Im zweiten Kapitel wird sich zeigen, da"s es f"ur
einen Interpreter bzw. Compiler kein Problem ist, mit einer
solchen Schachtelung fertigzuwerden. Folgt auf die Anweisungen ein
\rw{end}, so ist die \rw{if}-Anweisung beendet. Es w"are
m"oglich, auch das \rw{end} durch ein Symbol zu ersetzen,
doch alle sinnvollen Symbole, bis auf die Klammern, sind bereits
f"ur andere Zwecke vergeben. Nat"urlich k"onnte man auch eine
schlie"sende Klammer benutzen, doch um ein solches Konstrukt
sinnvoll zu machen, ben"otigt man auch eine "offnende Klammer.
Eine runde Klammer w"urde den falschen Eindruck erwecken, eckige
Klammern sind bereits vergeben, es bleiben die geschwungenen
Klammern, die auch in C zu diesem Zweck verwendet werden. Da ich
aber auch in allen anderen F"allen einem \rw{end} bzw. einem
\rw{begin}-\rw{end} Paar den Vorrang gegeben habe, f"uhrte eine
Verwendung hier h"ochstens zu einer Verwirrung des Benutzers.

Folgt auf die Anweisungen ein \rw{else}, dann folgen auf
dieses weitere Anweisungen, durch ein \rw{end} begrenzt.
Nat"urlich k"onnte man auch anstelle des \rw{else} ein Symbol
verwenden, doch das w"urde den Zweck der Konstruktion zu sehr
verschleiern:

\begin{verbatim}
if x == 1
{
  do_what_you_want_to();
::
  do_something_else();
}
\end{verbatim}

Ich glaube, bei Betrachtung dieses Beispiels zeigt sich, warum
ich die Schl"usselw"orter den Symbolen vorgezogen habe.

Manche Programmiersprachen, wie z.~B. \modula , stellen
zus"atzlich zum normalen \rw{else} ein sogenannte \rw{elseif}
oder auch \rw{elsif} bereit. Es erlaubt, sollte die
\rw{if}-Bedingung negativ sein, noch eine zweite Bedingung zu
"uberpr"ufen und, sollte auch diese logisch falsch sein, eine
dritte u. s. w. Ich habe mich entschieden, ein solches
\rw{elseif} nicht in SEL einzubauen, aus zwei Gr"unden.
Erstens l"a"st sich diese \rw{if}-\rw{elseif} Kombination in
vielen F"allen durch ein besser passendes
\rw{switch}\footnote{siehe Kapitel \ref{switch}} ersetzen und
zweitens ist es m"oglich, den \rw{elseif}-Effekt durch ein
verschachteltes \rw{if} zu erzeugen---\rw{elseif} selbst ist
nichts anderes als eine Vereinfachung dieser Konstruktion.

\subsubsection{Die switch-Anweisung}
\label{switch}

Die \rw{switch}-Anweisung\footnote{siehe Anhang \ref{grammar} oder
\cite{selkr}, Seite 11} erlaubt es, abh"angig vom Wert eines
Ausdrucks verschiedene Anweisungen auszuf"uhren. Dazu wird
zun"achst der Ausdruck evaluiert und ein \rw{case} gesucht,
dessen Wert mit dem des Ausdrucks "aquivalent ist. Ist kein
passendes \rw{case} vorhanden, wird zu einem evtl. vorhandenen
\rw{default} gesprungen. Dort wird mit der Ausf"uhrung
fortgefahren und erst bei einem \rw{break} oder dem zum
\rw{switch} geh"orenden \rw{end} beendet, also nach dem \rw{end}
fortgesetzt. Ist erst einmal ein bestimmter \rw{case}- oder
\rw{default}-Zweig eingeschlagen, werden alle weiteren
\rw{case}'s und \rw{default}s ignoriert. Das erm"oglicht es, f"ur
mehrere Werte ein und dieselben Anweisungen auff"uhren zu lassen,
indem man die einzelnen \rw{case}s hintereinander anreiht.

Die \rw{case}- und \rw{default}-Schl"usselw"orter werden in SEL,
ebenso wie in C, als Pr"afixe f"ur Anweisungen behandelt: bei
normaler Ausf"uhrung werden sie einfach ignoriert. Sie sind also
im Grunde nichts anderes als eine Art Sprunglabel, so wie sie in
vielen Programmiersprachen f"ur die
\rw{goto}-Anweisung\footnote{siehe Kapitel \ref{goto}}
verwendet werden. Eine andere---auf den ersten Blick
nat"urlichere---L"osung w"are, auf die \rw{switch}-Anweisung eine
Liste von \rw{case}'s oder \rw{default}s folgen zu lassen, die
ihrerseits jeweils von Anweisungen gefolgt sind. Um es m"oglich
zu machen, da"s eine Anweisungsfolge f"ur mehrere verschiedene
Werte ausgef"uhrt wird, m"u"ste man statt einem \rw{case} jeweils
mehrere \rw{case}'s erwarten. Welche der beiden L"osungen man
bevorzugt, ist Geschmackssache. Ich habe mich f"ur die
Pr"afix-L"osung entschieden, weil ich sie erstens von C gewohnt
bin und sie zweitens ein wenig leistungsf"ahiger ist.

Das Schl"usselwort \rw{case} vor dem jeweiligen Wert ist
notwendig, um Zweideutigkeiten in der Sprache zu vermeiden. Um
dieses Schl"usselwort auslassen zu k"onnen, m"u"sten wir die oben
besprochene zweite L"osung w"ahlen, jedoch m"u"sten die einzelnen
Anweisungsfolgen wiederum durch ein anderes Schl"usselwort, z.~B.
\rw{end} beendet werden, soda"s sich keine wirkliche Ersparnis
erg"abe.

Auch daf"ur, da"s in \rw{case}-Zweigen nur Konstanten erlaubt
sind, gibt es Gr"unde. Erstens h"atten wir durch den Doppelpunkt
Zweideutigkeiten---er w"urde als Typecast\footnote{siehe Kapitel
\ref{casts}} interpretiert werden---doch dieses Problem ist durch
Wahl eines anderen Symbols ohne weiteres zu "ubergehen. Der weit
gewichtigere Grund ist, da"s die f"ur die \rw{switch}-Anweisung
typische Effizienz verloren ginge: Bis der Interpreter zum
richtigen Zweig gelangt, m"u"ste er eine Vielzahl von Ausdr"ucken
auswerten, was sehr viel mehr Zeit in Anspruch nimmt, als der
simple Vergleich mit einer Konstante. Au"serdem w"urde diese
Entscheidung das Schreiben eines Compilers etwas verkomplizieren.
In F"allen, in denen eine solche \rw{switch}-Anweisung n"otig
w"are, sollte man auf \rw{if}-\rw{else}-Verschachtelungen umsteigen.

\subsubsection{Die return-Anweisung}

Die \rw{return}-Anweisung hat die Aufgabe, zur aufrufenden
Funktion zur"uckzukehren und evtl. einen Wert zu "ubermitteln,
den sog. {\em Return-Code}. Eine andere M"oglichkeit w"are, den
return-code durch eine Zuweisung an eine Variable, die den Namen
der Funktion hat, zu bestimmen und mittels einer anderen
Anweisung bei Bedarf zur aufrufenden Funktion zur"uckzuspringen.
Da man im Normalfall beides an der gleichen Stelle tut, habe ich
mich f"ur die \rw{return}-L"osung entschieden.

Die einfachste Methode w"are jedoch, "uberhaupt keine Anweisung
bereitzustellen, mit der in die aufrufende Funktion
zur"uckgekehrt werden kann und den R"uckgabewert am Ende der
Funktion festzulegen. F"ur einen rekursiven Interpreter ist dies
die beste L"osung, doch wie wir im n"achsten Kapitel sehen
werden, stellt uns auch der \rw{return}-Ansatz vor keine gro"sen
Probleme.

\subsection{Schleifen}

Schleifen kontrollieren nicht, was, sondern wie oft etwas
ausgef"uhrt wird. Die drei Arten der Schleifen, die in SEL
erlaubt sind, entsprechen den in den meisten strukturierten
Programmiersprachen vorhandenen Schleifen.

\subsubsection{Die repeat-Schleife}

Die \rw{repeat}-Schleife in SEL entspricht der
\rw{repeat}-Schleife in \pascal , bis auf das Detail, da"s die
\rw{until}-Bedingung in Klammern zu setzen ist und kein Semikolon
danach erforderlich ist. Diese Konzeption deckt sich mit den
restlichen Anweisungen von SEL und wurde daher von mir bevorzugt.

In manchen Programmiersprachen, wie z.~B. C, existiert die
\rw{repeat}-Schleife in einer abgewandelten Form, n"amlich der
der \rw{do-while}-Schleife. Der einzige Unterschied
besteht darin, da"s die \rw{repeat}-Schleife eine
Abbruchbedingung erwartet, die \rw{do-while}-Schleife jedoch eine
Fortf"uhrungsbedingung, das hei"st, eine Schleife kann durch die
andere ersetzt werden, wenn die Bedingung negiert wird. Welche
Schleife die bessere ist, kommt auf die Situation an, wirkliche
Vor- oder Nachteile gibt es nicht, die Auswahl ist also reine
Geschmackssache.

\subsubsection{Die while-Schleife}

Die \rw{while}-Schleife ist eng verwandt mit der
\rw{repeat}-Schleife. Der einzige Unterschied besteht darin, da"s
die \rw{repeat}-Schleife mindestens einmal durchlaufen wird, die
\rw{while}-Schleife jedoch nicht unbedingt. Syntaktisch gesehen
ist die \rw{while}-Schleife sehr "ahnlich der von C, die
Unterschiede sind in der in Kapitel \ref{stmts} angesprochenen
Blockstruktur von C begr"undet. So folgt auf die schlie"sende
Klammer der Schleifenbedingung in C in der Regel eine "offnende
geschwungene Klammer, in SEL jedoch ist diese ausgespart, da
nicht notwendig. Die schlie"sende geschwungene Klammer, die in C
die Schleife beendet, ist mit dem Schl"usselwort \rw{end} in SEL
zu vergleichen. Die Gr"unde, warum ich die Konstruktion mittels
Schl"usselw"ortern vorgezogen habe, ist bereits in vorhergehenden
Kapiteln er"ortert, ich gehe deshalb hier nicht n"aher darauf
ein.

\subsubsection{Die for-Schleife}

Die \rw{for}-Schleife in SEL erlaubt es, Anweisungen in einer
Anzahl, die schon im voraus bekannt ist, zu wiederholen.
Damit unterscheidet sie sich signifikant von den \rw{while} und
der \rw{repeat}-Schleifen, bei denen die Anzahl der Wiederholungen
nicht unbedingt bekannt sind. Das verhilft der
\rw{for}-Schleife---zumindest auf einem Interpreter---zu einer
sehr schnellen Ausf"uhrung, da die Bedingungen f"ur die Schleife
nur einmal berechnet werden m"ussen, w"ahrenddessen dies bei den
beiden anderen Schleifenarten bei jedem Durchlauf geschehen mu"s.

Dieses eine Merkmal der \rw{for}-Schleife ist vielen
Programmiersprachen gemeinsam, doch eine bekannte Sprache macht
eine Ausnahme: C. In C werden nicht eine Z"ahlervariable und
deren untere und obere Grenzen und evtl. noch die Schrittweite
angegeben, sondern es werden drei Ausdr"ucke verwendet: Der erste
wird lediglich beim Start der Schleife evaluiert, der dritte nach
jedem Durchlauf. Der zweite Ausdruck schlie"slich gibt die
Bedingung an, bei der die Schleife nochmals ausgef"uhrt werden
soll---ist sie logisch falsch, wird die Schleife abgebrochen. In
C ist also eine \rw{for}-Schleife nichts anderes als eine
erweiterte \rw{while}-Schleife und hat den Vorteil, da"s mit
ihrer Hilfe auch verkettete Listen und "ahnliches leicht
durchlaufen werden k"onnen. Auch ein "`klassische"'
\rw{for}-Schleife l"a"st sich so leicht nachbilden, doch diese
Art der Schleifenkonstruktion hat f"ur einen Interpreter den
gro"sen Nachteil, da"s bei jedem Schleifendurchlauf gleich zwei
Ausdr"ucke evaluiert werden m"ussen.

Ungl"ucklicherweise habe ich diesen Nachteil nicht sofort bemerkt und
f"uhrte auch in SEL die von C gewohnte \rw{for}-Konstruktion ein. Erst
sp"ater merkte ich, um wie viel effizienter diese Schleife f"ur einen
Interpreter w"are, w"urde ich sie "ahnlich wie in \pascal\ aufbauen. Die
"Anderungen am Interpreter und an der Dokumentation lie"sen sich zum
Gl"uck in weniger als einer Stunde durchf"uhren, soda"s Verz"ogerungen
f"ur den Zeitplan dieser Arbeit ausblieben und ich zur Erkenntnis
gelangte, da"s man aus Fehlern doch am besten lernt.

Noch ein kleines Wort zur Syntax: Die Bedingungen der Schleife sind in
Klammern gesetzt, weil das in SEL auch bei anderen Anweisungen der Fall
ist und dadurch die Konsistenz gewahrt wird, und die Grammatik ohne
Begrenzung nicht eindeutig w"are. Die Z"ahlervariable wird als
\fu{primary} angegeben, weil es auch m"oglich sein sollte, Array- und
Strukturelemente als Z"ahler zu benutzen, und nicht als \fu{expr}, weil
dann der folgende Zuweisungsoperator als ein Teil dieses Ausdrucks
interpretiert w"urde. Sollte es wirklich notwendig sein, einen \fu{expr}
zu benutzen, so kann dieser in Klammer gesetzt werden, und entspricht
damit einem \fu{primary}.\footnote{siehe Anhang \ref{grammar}} Anstatt
der Schl"usselw"orter \rw{to} bzw. \rw{downto} Symbole zu verwenden,
w"are erstens un"ubersichtlich und w"urde zweitens---au"ser bei
Verwendung neuer Symbole---die Eindeutigkeit zunichte machen, weshalb
ich die oben angesprochenen Schl"usselw"orter verwendete. Gleiches gilt
f"ur \rw{step}.

\subsection{Die goto-Anweisung}
\label{goto}

Die "Uberschrift dieses Absatzes mag darauf deuten, da"s es in SEL eine
\rw{goto}-Anweisung gibt, doch der Schein tr"ugt: Aus Gr"unden, die in
Kapitel \ref{parser} aufgef"uhrt sind und aus dem Grund, da"s diese
Anweisung in dem meisten F"allen ohnedies nicht notwendig ist und nur zu
einem schlechten Programmierstil verf"uhrt, habe ich mich gegen eine
Einbindung in SEL entschieden.    

\section{Ausdr\protect{"ucke}}

Ausdr"ucke sind ein Konzept, das in fast allen strukturierten
Programmiersprachen verwendet wird und das nicht ohne Grund: Es
vereinfacht die Grammatik der Sprache und macht damit das Schreiben
eines Interpreters bzw. eines Compilers einfacher. Au"serdem erleichtert
es dem Anwender der Sprache das Schreiben von Programmen sehr und
unterst"utzt die Lesbarkeit. So k"onnen Ausdr"ucke in vielen
verschiedenen Situationen verwendet werden: Als Bedingungen in
Schleifen, als Funktionsaufrufe oder als die rechten Seiten von
Zuweisungen.\footnote{Tats"achlich ist in SEL sogar die Zuweisung ein
Ausdruck, im Gegensatz zu \pascal\ oder \modula , doch dazu sp"ater
mehr.}

Ein Ausdruck in SEL kann "ahnlich einem Ausdruck im
mathematischen Sinn aufgefa"st werden, er ist eine Verkn"upfung
von Operanden durch Operatoren. Jeder Operand f"ur sich ist
selbst ein Ausdruck\footnote{bis auf die Operanden der Operatoren
\fu{new} und \fu{sizeof}, die Typen sind.} und das Ergebnis der
Verkn"upfung mit einem Operator ist wieder ein Ausdruck. So ist
es in SEL m"oglich, Ausdr"ucke beliebig tief zu schachteln. Diese
Eigenschaften der Ausdr"ucke legen es nahe, einen Ausdruck als
Baum\footnote{siehe \cite{sedge}, Seiten 57--72} darzustellen.
Dabei stellt jeder innere Knoten sowie die Wurzel einen Operator
dar und jedes Blatt einen Wert. Betrachtet man nun jeden
einzelnen Knoten (und dessen Nachfolger), so gelangt man zum
Ergebnis, da"s dieser auch wieder ein Ausdruck f"ur sich ist. Ein
Beispiel f"ur diese Darstellung gibt Sedgewick in \cite{sedge},
Seite 63, Abbildung 4.4.

Jeder Ausdruck in SEL hat nicht nur einen Wert, sondern auch
einen Typ\footnote{siehe auch Kapitel \ref{types}}, mit einer Ausnahme:
Ein Ausdruck des Typs \tvoid\ hat keinen Wert, was auch gar keinen Sinn
h"atte.

\subsection{Operatoren}

Grunds"atzlich gibt es vier Arten, in SEL Operatoren zu
klassifizieren\footnote{genauer gesagt vier "`vern"unftige"' Arten}:
nach der Anzahl ihrer Argumente, nach ihrer funktionellen
Zugeh"origkeit, nach ihrer Priorit"at und nach ihrem Ergebnistyp.

\paragraph{Klassifikation nach der Anzahl der Argumente}

Auf den ersten Blick gesehen gibt es in SEL nach dem Gesichtspunkt der
Zahl der Argumente zwei Arten von Operatoren: un"are und bin"are. Zu den
un"aren z"ahlen z.~B. die Operatoren \fu{+} und \fu{<<}, zu den bin"aren
z"ahlen \fu{!} und \fu{-}. Der letztgenannte Operator existiert auch in
einer bin"aren Version, man k"onnte also sagen, da"s es Operatoren gibt,
die sowohl un"ar als auch bin"ar sind. Ich pers"onlich betrachte diese
beiden "`Versionen"' als zwei verschiedene Operatoren, da auch ihre
Funktionen verschieden sind. Betrachtet man Ausdr"ucke in SEL genauer,
wird man bemerken, da"s ein
Funktionsaufruf durchaus mehrere Argumente erwarten kann, jedoch stellt
sich hier die Frage, ob man einen Funktionsaufruf (besser: eine Funktion) als
Operator auffassen sollte. Ich bin der Meinung, eine Funktion ist
weniger ein Operator als ein Wert, und habe ihr deshalb ein eigenes
Unterkapitel gewidmet.\footnote{siehe Kapitel \ref{exprfuncs}}

\paragraph{Klassifikation nach der funktionellen Zugeh"origkeit}

Auf diese Weise kann man die Operatoren von SEL in mehrere
Gruppen einteilen, wobei die Grenzen eher flie"send sind. Mir
erscheint folgende Einteilung am vern"unftigsten: arithmetische
Operatoren wie \fu{+} oder \fu{<<}, logische Operatoren wie
\fu{\&\&} oder \fu{!}, Vergleichsoperatoren wie \fu{>=} oder
\fu{==}, Operatoren zur Speicherverwaltung wie \fu{new} oder
\fu{[]} und als zwei eigene Gruppen der Zuweisungsoperator \fu{=}
und der Typecasting-Operator \fu{:}. Wie gesagt sind die
"Uberg"ange zwischen diesen Gruppen sehr flie"send. So k"onnte
man beispielsweise meinen, der Operator \fu{[]} geh"ore zur
Gruppe der arithmetischen Operatoren, da er ja in gewisser Weise
Vektorarithmetik m"oglich macht. Auch w"are es m"oglich, den
Typecasting-Operator zur Speicherverwaltung z"ahlen, da er auch
den Speicher modifiziert.

\paragraph{Klassifikation nach der Priorit"at}

Operatoren in SEL besitzen verschiedene
Priorit"atsstufen\footnote{siehe \cite{selkr}, Seite 21}, um eine
"uberm"a"sige Klammerung zu vermeiden. Die Priorit"aten in SEL sind
stark an die von C angelehnt\footnote{siehe \cite{kandr}, Seite 54},
weichen jedoch teilweise von diesen ab. Meistens im selben Atemzug mit
der Priorit"at der Operatoren wird die Richtung der Auswertung genannt.
Diese gibt an, ob der Operator links- oder rechts-assoziativ ist. Ein
Beispiel f"ur einen rechts-assoziativen Operator ist \fu{=}: Der
Ausdruck \fu{a=b=0} ist "aquivalent zu \fu{a=(b=0)}. So ist es ohne
Klammerung m"oglich, mehreren Variablen in einer Anweisung denselben
Wert zuzuweisen. Ein
links-assoziativer Operator ist das bin"are \fu{-}. W"are er das nicht,
w"urde der Ausdruck \fu{5-3-1} nicht zu \fu{(5-3)-1}, sondern zu
\fu{5-(3-1)} "aquivalent sein und so das falsche Ergebnis, n"amlich
\fu{3} anstatt \fu{1}, liefern.

\paragraph{Klassifikation nach dem Ergebnistyp}

Diese Einteilung ist wohl die komplizierteste und es w"are wohl besser,
man w"urde als Kriterium nicht den Ergebnistyp selbst, sondern das, 
was den Ergebnistyp bestimmt, heranziehen. 
So sind dies etwa beim Operator \fu{+} die
Typen seiner Operanden, beim Operator \rw{new} hingegen ist es der Typ,
mit dem er verbunden wird. Weiters gibt es Operatoren, die immer denselben
Ergebnistyp haben, wie z.~B. die logischen Operatoren oder der Operator
\rw{sizeof}. Es gibt sogar einen Operator, dessen Ergebnistyp immer
\tvoid\ ist, n"amlich \rw{delete}.

\bigskip

In der folgenden Aufstellung der Operatoren verwende ich die
Klassifikation nach der funktionellen Zugeh"origkeit, weil sie
f"ur diesen Zweck am geeignetsten ist.

\subsubsection{Arithmetische Operatoren}

SEL bietet neben den numerischen Operatoren, zu denen die vier
Grundrechnungsoperatoren f"ur Addition, Subtraktion, Multiplikation und
Division z"ahlen, zus"atzlich der Modulo-Operator, der den Rest bei einer
ganzzahligen Division ergibt, und der Negationsoperator, der
einer Umkehr des Vorzeichens gleichkommt, auch
Operatoren zur Bitmanipulation, n"amlich bitweise Links- und
Rechtsverschiebung und bitweise Und-, Oder- und
XOR-Verkn"upfungen. Die bin"aren numerischen Operatoren, also
\fu{+}, \fu{-}, \fu{*}, \fu{/} und \fu{\%}, besitzen eine
niedrigere Priorit"at als alle un"aren Operatoren, jedoch eine
h"ohere Priorit"at als alle anderen bin"aren Operatoren.
Untereinander sind die Additions- und Subtraktionsoperatoren auf
derselben Priorit"atsstufe, die um einen Grad niedriger steht,
als die der anderen numerischen bin"aren Operatoren, in Anlehnung
an die in der Mathematik "ublichen Notation. Die un"aren
Operatoren haben deshalb eine h"ohere Priorit"at, weil das
erstens unseren Erwartungen entspricht---man betrachtet intuitiv einen
un"aren Operator einem Wert als "`n"aherstehender"', als einen
bin"aren---und es zweitens in fast allen F"allen praktischer ist.
So ist z.~B. das Ergebnis des Ausdrucks \fu{-x*y} dem von
\fu{-(x*y)} "aquivalent. Bes"a"sen aber die bin"aren Operatoren
eine h"ohere Priorit"at, m"u"ste man anstatt \fu{-x-y} den
Ausdruck \fu{(-x)-y} schreiben, was etwas unnat"urlich aussieht.
Hingegen k"onnte man den Ausdruck \fu{-(x+y)} zwar bei anderer
Operatorenrangfolge ohne Klammern schreiben, doch selbst ohne die
Rangfolge zu ver"andern, k"onnte man durch "Anderung des \fu{+} in
ein \fu{-} die Klammerung vermeiden: \fu{-x-y}. Die Rangfolge
erweist sich also als nat"urlich und praktisch. Die
Vergleichsoperatoren haben deshalb niedrigere Priorit"at als die
numerischen Operatoren, weil es kaum vorkommt, da"s man mit
Vergleichsergebnissen rechnet. F"ur die logischen Operatoren gilt
dieselbe Begr"undung. Die Festsetzung der Priorit"at der
Bitmanipulations-Operatoren relativ zu den numerischen Operatoren
ist eher willk"urlich. Es w"urde kaum ein Vor- oder Nachteil
entstehen, h"atten die Operatoren zur Bitmanipulation h"ohere
Priorit"at. Der Grund, warum ich diese Operatoren nicht auf
dieselbe Priorit"atsstufe gesetzt habe ist der, da"s sie zu
verschieden sind, als da"s dies dienlich w"are.

Die numerischen Operatoren haben eine besondere Eigenschaft: Sie
stellen sich auf die Typen ihrer Operanden ein und ver"andern
diese sogar, wenn n"otig. So stellt sich zum Beispiel die Frage,
welchen Typ eine Addition eines \tint\ und eines \tfloat\ haben
sollte. Vern"unftigerweise wandelt der Additionsoperator in
diesem Fall den \tint\ in einen \tfloat\ um, und f"uhrt erst
anschlie"send die Addition aus. Auf diese Weise gehen keine
Nachkommastellen des \tfloat\ verloren. Gleiches passiert mit
einem \tchar\ und einem \tint : Der \tchar\ wird in den Typ
\tint\ umgewandelt, da nur so gew"ahrleistet ist, da"s mit dem
richtigen Wert des \tint\ gearbeitet wird. F"ur die Typen bei
einer numerischen Operation gilt in SEL also die Regel \hbox{$
\tchar < \tint < \tfloat $}.\footnote{siehe \cite{selkr}, Seite
18} Zus"atzlich zu diesen Typen k"onnen die Operatoren \fu{+} und
\fu{-} auch mit Zeigern operieren, doch dazu mehr in Kapitel
\ref{memmgrops}.

Die Operatoren zur Bitmanipulation arbeiten in SEL lediglich mit
dem Typ \tint\, und das aus guten Gr"unden: Erstens ist es f"ur
einen Computer zwar leicht m"oglich, mit den Bits eines \tfloat\
zu operieren, doch es ist f"ur einen Menschen ungleich
schwieriger, diese Bits auch richtig zu interpretieren---auf den
Wert der Zahl bezogen. Zweitens sind die "ublichen 8 Bit eines
\tchar\ im Regelfall zu wenig f"ur die Anwendungen der
Bitmanipulation und drittens vereinfacht es nicht nur den Bau
eines Interpreters bzw. Compilers, es macht diese auch schneller.

\subsubsection{Logische Operatoren}

Neben den Operatoren zur logischen Und- und Oder-Verkn"upfung
bietet SEL auch einen Operator zur logischen XOR-Verkn"upfung an.
Ich habe mich dazu entschieden, weil es leicht in einen
Interpreter bzw. Compiler einzubauen ist und andererseits der
Operator in manchen Situationen sehr praktisch ist und logische
Ausdr"ucke verk"urzen kann. Die logischen Operatoren sind auf der
zweitniedrigsten Priorit"atsstufe angesiedelt---unter ihnen
befinden sich lediglich die Operatoren \fu{=} und \fu{delete}.
Diese Stufe habe ich gew"ahlt, weil logische Operatoren
haupts"achlich in Schleifen- und \rw{if}-Bedingungen eingesetzt
werden und es dort am praktischsten ist, wenn sie die niedrigste
Priorit"at haben. Nur selten ist es notwendig, da"s einer
Variablen ein logischer Wert zugewiesen wird, doch auch dieser
Fall ist ohne Klammerung zu bew"altigen, da ja der
Zuweisungsoperator eine noch niedrigere Priorit"at hat. Der
einzige Fall, in dem sich diese Rangfolge als unpraktisch
herausstellen w"urde ist, da"s in einer Schleifen- oder
\rw{if}-Bedingung ein logischer Wert einer Variablen zugewiesen
wird und dieser Wahrheitswert mit einem anderen mittels eines
logischen Operators verkn"upft werden mu"s. Ein Beispiel hierf"ur
w"are \fu{if (x \&\& (y = z))}. Tats"achlich tritt dieser Fall so
selten auf---in meiner bislang 4-j"ahrigen C-Laufbahn bin ich
noch nie in eine solche Situation gelangt---da"s man ihn getrost
au"ser Acht lassen kann.

Der logische Negationsoperator \fu{!} liegt auf derselben
Priorit"atsstufe wie die Operatoren zur bitweisen und
arithmetischen Negation und die beiden Operatoren \fu{*} und
\fu{\&}, die im Kapitel \ref{memmgrops} besprochen werden, welche
um eine Stufe tiefer liegt als die der Operatoren \fu{->}, \fu{:}
und \fu{[]}. Um das zu begr"unden, betrachten wir den Ausdruck
\fu{!x->y}. Offenbar wird hier versucht, aus der Struktur \fu{x}
auf das Element \fu{y} zuzugreifen und dieses zu negieren. H"atte
der Operator \fu{!} eine h"ohere Priorit"at als \fu{->}, so
w"urde sich die Negation auf die Struktur beziehen, und es ist
nicht m"oglich, eine Struktur zu negieren.

Logische Operatoren arbeiten in SEL nur mit dem Typ \tint .
Zwar w"urde es auch gen"ugen, den Typ \tchar\ zu diesem Zweck
heranzuziehen, doch da in kaum einer Anwendung Wahrheitswerte
gespeichert werden und \tint\ normalerweise der Gr"o"se des
Maschinenwortes des jeweiligen Prozessors entspricht, und daher
schneller zu bearbeiten ist, als ein \tchar , ist er die bessere
Wahl. Sollte es jedoch n"otig sein, Wahrheitswerte abzuspeichern,
kann das ohne weiteres mit Hilfe von \tchar s geschehen, die durch
ein Typecasting entstanden sind, oder aber durch Bitfelder, was
noch speicherplatzsparender ist.

\subsubsection{Vergleichsoperatoren}

SEL bietet insgesamt sechs Vergleichsoperatoren an. Diese sind um
eine Priorit"atsstufe h"oher als die logischen Operatoren
angesiedelt, da sie meistens in Schleifen- und
\rw{if}-Bedingungen eingesetzt werden und dort diese Rangfolge
von Vorteil ist. So wird es wahrscheinlich oft vorkommen, da"s
"uberpr"uft werden soll, ob zwei Werte zwei anderen Werten
entsprechen, aber eher selten wird man "uberpr"ufen wollen, ob
die Oder-Verkn"upfungen jeweils zweier Werte dieselben Ergebnisse
haben.

Vergleichsoperatoren k"onnen in SEL lediglich Werte der Typen
\tchar , \tint\ und \tfloat\ vergleichen, wobei das im vorletzten
Unterkapitel vorgestellte "`Gr"o"senverh"altnis"' gilt. Es ist
hingegen nicht m"oglich, z.~B. Strings miteinander zu
vergleichen, weil das, um es vern"unftig zu machen, eine
Ver"anderung im Typkonzept von SEL mit sich z"oge und au"serdem
sehr gut von einer Funktion erledigt werden kann---SEL stellt
alle n"otigen Konstrukte hierf"ur bereit.

\subsubsection{Speicherverwaltungsoperatoren}
\label{memmgrops}

Unter den Begriff "`Speicherverwaltungsoperatoren"' fallen all
jene Operatoren, die mit Zeigern, Arrays und Strukturen arbeiten.
Dazu z"ahlen nicht nur die Operatoren zur Ermittlung der
Speicheradresse einer Variablen und zur Dereferenzierung eines
Operators, \fu{\&} und \fu{*}, zum Zugriff auf ein Array- bzw.
Strukturelement, \fu{[]} und \fu{->}, sondern auch die Operatoren
zum Allokieren und Freigeben von Speicher, \fu{new} und
\fu{delete}, und der Operator zum Ermitteln der Gr"o"se eines
Typen, \fu{sizeof}. Au"serdem fallen unter diesen Begriff die
beiden bin"aren Operatoren \fu{+} und \fu{-}, mit denen
zus"atzlich zur numerischen Arithmetik auch Zeigerarithmetik
betrieben werden kann.

Die h"ochste Priorit"atsstufe aller Operatoren haben \fu{new} und
\fu{sizeof}, auf der n"achsten Stufe liegen \fu{->} und \fu{[]}.
Es h"atte keinen Sinn, diese Rangfolge zu "andern, denn die
Operanden von \fu{new} und \fu{sizeof} sind keine Ausdr"ucke,
sondern Typen, und es ist nicht m"oglich, den Operator \fu{->}
oder \fu{[]} auf einen Typ anzuwenden, soda"s dies schon von Haus
aus unterbunden werden sollte. Die beiden un"aren Operatoren
\fu{*} und \fu{\&} liegen um eine Priorit"atsstufe tiefer,
gemeinsam mit allen anderen un"aren Operatoren. Eine andere Stufe
f"ur diese Operatoren w"are nicht angebracht, aus Gr"unden, die
ich schon genannt habe.

Ein in gewisser Hinsicht au"sergew"ohnlicher Operator---gemessen
an seinen Gegenst"ucken aus C und \pascal ---ist der Operator
\fu{->}. Er wird benutzt, um auf Elemente einer Struktur
zuzugreifen. Derselbe Operator wird in C und \pascal\ durch den
Punkt (\fu{.}) symbolisiert. In C jedoch gibt es zus"atzlich den
Operator~\fu{->}. Er wird benutzt, um auf Elemente aus einem
Zeiger auf eine Struktur zuzugreifen---in \pascal\ mu"s hierf"ur
zuerst der Zeiger dereferenziert werden. Der SEL-Operator ist nun
insofern au"senstehend, als er die Funktionalit"at des
Punktoperators und des Pfeiloperators aus C vereinigt. So
erm"oglicht er es, angewendet auf eine Struktur, auf ein Element
daraus zuzugreifen. Angewendet auf einen Zeiger auf eine Struktur
dereferenziert er ersteren und greift dann erst auf das jeweilige
Element zu. Diese Vereinigung macht insofern Sinn, als da"s
keine Funktionalit"at verlorengeht, denn der Punktoperator
hat, angewendet auf einen Zeiger, keine Funktion. Auch das Design
eines Interpreters bzw. eines Compilers wird dadurch nicht
erschwert, denn diese wissen, wenn sie auf den Pfeiloperator
treffen, mit welchem Typ gearbeitet wird. Ein Compiler m"u"ste
also nur, tr"afe er auf einen Pfeiloperator, dessen linker Operand
ein Zeiger ist, den Code zur Dereferenzierung einbauen.

\subsubsection{Der Zuweisungsoperator}

Der Zuweisungsoperator von SEL ist ein Sonderfall unter den
Operatoren, aus mehreren Gr"unden. Der erste Grund ist, da"s es
in manchen Programmiersprachen, darunter auch \pascal , gar
keinen Zuweisungsoperator in diesem Sinne gibt. In \pascal\ ist
eine Zuweisung eine Anweisung und kein Ausdruck. Das macht es in
\pascal\ auf der einen Seite schwerer, den Zuweisungsoperator
vom Compiler bzw. Interpreter unentdeckt mit dem
Vergleichsoperator zu verwechseln, denn der Zuweisungsoperator
ist in Ausdr"ucken nicht erlaubt und Ausdr"ucke sind in \pascal\
keine g"ultigen Anweisungen. Andererseits jedoch ist es nicht
m"oglich, mehreren Variablen in einer Anweisung denselben Wert
zuzuweisen, wie man es von C oder auch aus der Mathematik gewohnt
ist (\fu{x=y=0}). Au"serdem ist es nicht m"oglich, z.~B. in einer
\rw{if}-Anweisung gleichzeitig eine Zuweisung vorzunehmen und den
Wert der zuzuweisenden Variablen zu "uberpr"ufen (z.~B.
\hbox{\fu{if ((x=y) == 3)}}). Manche legen das sogar als Vorteil
aus, weil es eine unsch"one Programmierung erschwert, und
insofern gebe ich ihnen Recht, doch ich pers"onlich finde solche
Konstrukte in manchen Situationen doch sehr praktisch, und da der
Aufwand f"ur einen Zuweisungsoperator etwa dem einer
Zuweisungsanweisung gleichkommt, oder sogar noch geringer ist,
habe ich ihn in SEL eingebaut.

Der zweite Grund f"ur die Besonderheit des \fu{=}-Operators ist,
da"s er mit allen Typen zusammenarbeitet. Anders als der
Vergleichsoperator, der auf die Grundtypen beschr"ankt ist, kann
der Zuweisunsoperator eine Variable jeden beliebigen Typs einer
anderen zuweisen---Typgleichheit vorausgesetzt. Ich habe diese
"`Allm"achtigkeit"' des Operators deswegen in SEL erm"oglicht, weil
es leicht ist, das in einen Interpreter einzubauen---im Fall der
Konzeption, die wir im zweiten Kapitel besprechen werden, sogar
noch leichter, als sich auf die Grundtypen zu beschr"anken. Auch
f"ur einen Compiler ist das nicht weiter schwer: Den Code f"ur die
Grundtypen kann der Compiler direkt einf"ugen, f"ur Strukturen
oder Arrays k"onnte eine Funktion aus der Laufzeitumgebung
aufgerufen werden. Hier dr"angt sich nat"urlich die Frage auf,
warum ich dieses Konzept nicht auf den Vergleichsoperator
\fu{==} ausgeweitet habe. Die Antwort ist einfach: H"atte ich das
getan, w"are es unausweichlich gewesen, das auch f"ur die
Operatoren \fu{<}, \fu{>}, \fu{<=}, \fu{>=} und \fu{!=} zu tun.
Mit letzterem h"atte ich keine Probleme gehabt, doch wann ist
eine Struktur, die die Adresse einer Person enth"alt, gr"o"ser
als eine andere desselben Typs? Byte-f"ur-Byte-Vergleich w"are
hier sicher nicht angebracht.\footnote{In C++ kann dieses Problem
gel"ost werden, weil es dort m"oglich ist, Operatoren zu
"uberladen, d. h., da"s der Benutzer imstande ist, eine Funktion
zu schreiben, die zwei Strukturen (beispielsweise des Typs
\fu{Adresse}) vergleicht, die dann anstatt des
Vergleichsoperators aufgerufen wird.}

\subsubsection{Der Typecasting-Operator}

Der Typecasting-Operator \fu{:} wird benutzt, um einen Wert
explizit in einen anderen Typ umzuwandeln. Er wird genauer in
Kapitel \ref{casts} beschrieben. Seine Priorit"at ist mit der der
beiden Operatoren \fu{->} und \fu{[]} "aquivalent, erstens, weil
er ein un"arer Operator ist und zweitens, weil es eher selten
vorkommt, da"s gr"o"sere Ausdr"ucke gecastet werden m"ussen. Wenn
n"otig, wird dies sowieso von SEL "ubernommen.

\subsection{Konstanten}
\label{const}

Grunds"atzlich gibt es in SEL vier verschiedene Arten von
Konstanten: \tchar -Konstanten, \tint -Konstanten,
\tfloat -Konstanten und String-Konstanten. Eine \tchar -Konstante
ist ein unter einfache Anf"uhrungsstriche gesetztes Zeichen bzw.
Steuerzeichen und hat---es ist nicht anders zu erwarten---den Typ
\tchar . Eine \tint -Konstante ist eine simple Abfolge von
Ziffern. Zahlenkonstanten im lexikalischen Sinn sind in SEL immer
positiv---negative Zahlen macht erst der un"are Operator \fu{-}
m"oglich, und insofern ergibt sich daraus kein Nachteil. Der
Grund, warum ich Konstanten nicht auch auf negative Zahlen
ausgeweitet habe ist der, da"s sich die lexikalische
Analyse\footnote{siehe Kapitel \ref{scanner}}
dadurch sehr erschweren w"urde. Nehmen wir als Beispiel den
Ausdruck \fu{7-4}. Der Scanner erkennt als erstes die
\tint -Konstante \fu{7}. Danach m"u"ste ein Operator folgen oder
der Ausdruck zu Ende sein, was aber der Scanner nicht wissen
kann, da er vom Parser in dieser Hinsicht abgekapselt ist. Aus
diesem Grund w"urde er anstatt den Operator \fu{-} zu erkennen
eine weitere \tint -Konstante, n"amlich \fu{-4} lesen, was
nat"urlich falsch ist. Das Gesagte gilt auch f"ur
\tfloat -Konstanten.

String-Konstanten sind in gewisser Hinsicht etwas Besonderes,
n"amlich in Anbetracht der Tatsache, da"s es einen Stringtyp in
SEL gar nicht gibt. Normalerweise werden Strings als Arrays von
\tchar s behandelt, die durch ein Byte mit dem Wert \fu{0}
abgeschlossen sind, es w"are allerdings nicht schwer, Funktionen
zu programmieren, die \pascal -Strings verarbeiten. Da jedoch die
erstere Stringart die verbreitetere ist, wollen wir bei dieser
bleiben. Wie wir wissen, kann ein Array in SEL in vielen
Situationen wie ein Zeiger behandelt werden. Was also w"are
falsch daran, einen String als Pointer-Konstante zu sehen?
Tats"achlich wird das nicht nur in SEL, sondern auch in C so
praktiziert, von wo ich dieses Konzept "`importiert"' habe. Dieser
Pointer zeigt auf einen Speicherbereich, in dem der eigentliche
String gespeichert ist. Nat"urlich mu"s dieser Speicherbereich
w"ahrend der gesamten Programmausf"uhrung an derselben
Speicheradresse bleiben, und nicht irgendeiner L"oschaktion zum
Opfer fallen, soda"s die Pointerkonstante im ganzen
Programm zu jeder Zeit "`aktuell"' ist.\footnote{In Wirklichkeit
mu"s der String nicht w"ahrend der gesamten Programmausf"uhrung
im Speicher sein, sondern erst nach seinem ersten Auftreten im
Code, doch der Einfachheit halber wird er die ganze Zeit im
Speicher gehalten.}

\subsection{Variablen}
\label{variables}

In SEL k"onnen zwei Arten von Variablen unterschieden werden, die
ihrerseits in jeweils zwei Untertypen unterteilt werden k"onnen.
Lokale Variablen sind nur innerhalb einer bestimmten Funktion
bekannt und existieren f"ur jede Instanz dieser Funktion
getrennt.\footnote{Mehrere Instanzen einer Funktion sind dann
m"oglich, wenn sich diese direkt oder indirekt selbst aufruft.}
Die Gruppe der lokalen Variablen werden weiter unterteilt in
Parameter und direkte Funktionsvariablen.\footnote{Eine bessere
Bezeichnung als "`direkte Funktionsvariablen"' habe ich nicht
gefunden.} Die Werte der Parameter werden von der aufrufenden
Funktion bestimmt, wohingegen die direkten Funktionsvariablen der
Funktion alleine unterstehen. Lokale Variablen werden auf dem
Stack gespeichert. Wird eine Funktion aufgerufen, werden die
Parameter von der aufrufenden Funktion auf den Stack gelegt und
anschlie"send zum Code der Funktion
gesprungen.\footnote{Nat"urlich wird auch die R"ucksprungadresse
auf dem Stack gespeichert} Die Funktion ihrerseits legt die
direkten Funktionsvariablen auf den Stack und f"uhrt dann ihre
Anweisungen aus. Wird sie beendet, entfernt sie sowohl
direkte Funktionsvariablen als auch Parameter vom Stack
und springt dann zur aufrufenden Funktion zur"uck. Das Konzept,
den Stack als Speicher f"ur lokale Variablen zu benutzen macht
es m"oglich, da"s bei einer rekursiven Funktion jede Instanz ihre
eigenen Variablen zur Verf"ugung hat.

Globale Variablen werden in SEL unterteilt in modulglobale und
programmglobale Variablen. Erstere sind nur in der Quelldatei
(Modul) bekannt, in der sie definiert sind, wohingegen letztere
auch in anderen Modulen benutzt werden k"onnen. Sowohl modul- als
auch programmglobale Variablen werden im statischen Speicher
aufbewahrt. Daf"ur, da"s modulglobale Variablen nur in ihrem
jeweiligen Modul referenziert werden, haben der Interpreter bzw.
der Compiler oder der Linker zu sorgen.

Ich glaube, "uber die Verwendung von Variablen in Ausdr"ucken
mu"s ich hier nicht viele Worte verlieren: das geschieht ganz
einfach---wie in fast allen Programmiersprachen "ublich---durch
die Benutzung des Variablennamens.

\subsection{Funktionen}
\label{exprfuncs}

Die Aufrufkonvention f"ur Funktionen ist einfach: Der
Funktionsname gefolgt von in Klammern gesetzten Parametern,
jeweils durch einen Beistrich getrennt, bilden einen Funktionsaufruf.
Hat die Funktion keine Parameter, sind trotzdem "offnende und
schlie"sende Klammern anzugeben, aus mehreren Gr"unden. Erstens
kann der Interpreter bzw. Compiler daran ohne nachzusehen
erkennen, da"s es sich um einen Funktionsaufruf handelt und nicht
um die Verwendung einer Variablen. Zweitens ist dies in Hinblick
auf die evtl. Erweiterung der Sprache durch
Funktionszeiger\footnote{siehe Kapitel \ref{langext}}, also
Zeigervariablen, die die Adressen von Funktionen aufnehmen,
n"otig. In diesem Fall signalisiert das Klammerpaar, da"s wir
vorhaben, diese Funktion aufzurufen und nicht, sie als Zeiger zu
verwenden.

Die Anzahl und die Typen der Funktionsparameter sind nat"urlich
einzuhalten. Die Regelung bei letzteren ist nicht so streng:
Stimmen die Typen nicht "uberein, versucht SEL mittels impliziten
Typecastings\footnote{siehe Kapitel \ref{casts}}, diese in die
richtigen umzuwandeln, was jedoch nicht unbedingt gelingen mu"s,
da sich der Interpreter bzw. der Compiler an bestimmte Regeln zur
Typumwandlung halten mu"s.

\section{Typen}
\label{types}

In diesem Abschnitt m"ochte ich das Typkonzept von SEL erkl"aren.
Zun"achst stellt sich beim Design einer Programmiersprache die
Frage, welches Typkonzept "uberhaupt zur Anwendung kommen soll.
So gibt es z.~B. in \basic\ einige wenige Typen, aus denen der
Benutzer ausw"ahlen kann, es ist dort jedoch nicht m"oglich,
selber Typen zu definieren. Zwar k"onnen Arrays definiert werden,
doch ist ein Array in \basic\ nichts anderes als eine Menge von
Variablen, deren Namen sich jeweils durch eine Nummer
unterscheiden---ein Array als Ganzes zu behandeln, ist in \basic\
nicht m"oglich.

Dann gibt es Programmiersprachen---\dbase\ ist
ein Beispiel hierf"ur---in denen Variablen zwar einen bestimmten
Typ haben, dieser sich aber jederzeit "andern kann. Dort ist es
z.~B. m"oglich, einer Variablen einen numerischen Wert zuzuweisen
und mit ihr Berechnungen anzustellen. Sp"ater kann man ihr bei
Bedarf einen String zuweisen und diesen am Bildschirm ausgeben.
Zwar bietet dieses Konzept auf den ersten Blick f"ur den Benutzer
Vorteile\footnote{Auch diese Vorteile sind nur scheinbar, denn
sie verf"uhren zu einem sehr schmutzigen Programmierstil} und es
ist f"ur eine reine Interpretersprache eine gute Wahl, doch einem
Compiler macht ein solches Konzept gro"se Probleme: In
Programmiersprachen wie C oder \pascal\ haben Variablen einen
vorgegebenen, konstanten Typ. Vom Compiler werden diese Variablen
an eine bestimmte Speicheradresse gelegt, die sich im Laufe des
Programms nicht ver"andert, was ohne weiteres m"oglich ist, da
sich ja die Speicheranforderungen der Variablen auch nicht
ver"andern. K"onnen sich aber die Typen der Variablen "andern,
"andern sich auch die Gr"o"sen der Variablen, d. h., da"s es
nicht m"oglich ist, sie an eine konstante Stelle im Speicher zu
setzen. W"ahrend ein C-Compiler, wenn er die Variable zu
referenzieren hat, ganz einfach ihre Speicheradresse einf"ugen
kann, mu"s ein compiliertes \dbase -Programm eine Tabelle
verwalten, in der alle Variablen und deren momentane Adressen und
Typen verzeichnet sind. Wird nun der Ausdruck \fu{x+1}
abgearbeitet, mu"s zun"achst in der Tabelle nachgesehen werden,
an welcher Adresse sich die Variable \fu{x} befindet. Zus"atzlich
mu"s auch gepr"uft werden, welchen Typ diese hat, da z.~B. eine
ganzzahlige Addition einen anderen Code verlangt, als eine
Addition mit Gleitkommazahlen. Erst jetzt kann die Addition
ausgef"uhrt werden. Es ist also nur verst"andlich, da"s Compiler
f"ur eine solche Sprache einen Code erzeugen, der zehnmal
langsamer\footnote{Die Zahl 10 mag wie eine "Ubertreibung
klingen, tats"achlich ist es nicht selten, da"s solche Codes noch
um einiges langsamer sind. So ben"otigt z.~B. ein Programm, da"s
einen QuickSort-Algorithmus ausf"uhrt in Clipper (einem
\dbase -Compiler) sage und schreibe 120 mal l"anger als ein
vergleichbares Programm in C} ist als der eines C- oder
\pascal -Compilers.

Eine weitere wichtige Kategorie ist die der typenlosen
Programmiersprachen---AWK f"allt teilweise darunter. Dort ist es
m"oglich, einer Variablen einen String zuzuweisen, der eine Zahl
enth"alt und in der n"achsten Zeile mit dieser Variable zu
rechnen als w"are sie numerischen Typs. Dieses Konzept ist
insofern unpraktisch, als man f"ur verschiedene Typen (obwohl
es diese eigentlich nicht gibt) auch verschiedene Operanden
braucht. Bei einem Ausdruck \fu{x+y} w"are also nicht klar, ob
die beiden Strings zu konkatenieren sind oder ob deren numerische
Werte addiert werden sollen. Auch in Bezug auf den Bau eines
Compilers ist dieses Konzept von Nachteil, aus "ahnlichen
wie im vorigen Absatz dargelegten Gr"unden.

Manche alte Programmiersprachen stellen als einzigen Typ das
Maschinenwort zur Verf"ugung, was einerseits die Gefahr birgt,
da"s in einer solchen Sprache geschriebene Programme sehr vom
Prozessor abh"angig sind und andererseits den Nachteil hat, da"s
viele Aufgaben, die der Interpreter bzw. Compiler ohne 
Effizienzverluste "ubernehmen k"onnte, vom Programmierer 
bestritten werden m"ussen.

Wie sich aus den vorigen Abs"atzen leicht herauslesen l"a"st,
habe ich mich f"ur das Typkonzept entschieden, das auch in C
oder \pascal\ verwendet wird. Dort mu"s jede Variable vor ihrer
Verwendung definiert werden, wobei gleichzeitig ihr Typ
festgelegt wird, der sich im Laufe des Programms nicht "andert.
Au"serdem ist es dem Benutzer m"oglich, aus den Basistypen der
Programmiersprache eigene Typen herzuleiten, wenn auch nicht
unbeschr"ankt.

\subsection{Elementare Typen}
\label{elemtypes}

Wie bereits gesagt, ben"otigt man zur Konstruktion von
komplizierteren Typen eine Menge an Basistypen, mit denen man
mittels der Konstrukte, die SEL bereitstellt, arbeiten kann.
Programmiersprachen wie C oder \pascal\ bieten eine relativ
gro"se Zahl solcher Basistypen an. So gibt es in C zwei
verschiedene Typen, die ein einzelnes Zeichen repr"asentieren,
drei Grundtypen f"ur ganze Zahlen, wobei von jedem dieser Typen
ein vorzeichenloser sowie ein vorzeichenbehafteter Typ
exisitiert, und mehrere Gleitkomma-Typen. Die Vielfalt in
\pascal\ ist "ahnlich gro"s, wobei zus"atzlich noch ein Typ f"ur
logische Werte exisitiert (\type{boolean}). Ich habe mich
entschieden, in SEL nur vier Elementartypen aufzunehmen:
\tchar , \tint , \tfloat\ und \tvoid .

Der Typ \tchar\ repr"asentiert ein einzelnes Zeichen und nimmt
dementsprechend nur ein Byte ein. Arrays aus \tchar s k"onnen
z.~B. als Strings interpretiert werden.

Ein \tint\ repr"asentiert eine ganze Zahl mit Vorzeichen und hat
in der Regel die Gr"o"se eines Maschinenwortes.\footnote{Eine
Ausnahme bildet MS-DOS, wo ein Maschinenwort nur 16 Bits lang
ist, d.~h. der Wertbereich w"urde von etwa -32000 bis 32000
reichen, was zwar f"ur viele Anwendungen reicht, aber ungen"ugend
ist, da es ja nur einen Typ f"ur ganze Zahlen gibt. Auf MS-DOS
Maschinen ist deshalb die Gr"o"se eines \tint\ zwei
Maschinenworte, also 32 Bits.} Zus"atzlich hat \tint\ die
Aufgabe, logische Werte zu repr"asentieren, wobei der Wert \fu{0}
einem "`logisch falsch"' und ein Wert ungleich \fu{0} "`logisch
wahr"' entspricht. Unter anderem hat diese Konzeption den
Vorteil, da"s ein Ausdruck des Typs \tint\ nicht nur den Wert des
\tint\ wiederspiegelt, sondern auch die Aussage, ob dieser
ungleich \fu{0} ist oder nicht. So ist es z.~B. m"oglich, mit
einer sehr kurzen \rw{if}-Anweisung zu pr"ufen, ob ein Zeiger ein
Nullzeiger ist. W"ahrend hierf"ur in \pascal\ die Anweisung
\fu{if p <> nil then} benutzt wird, kann in SEL einfach
\fu{if (p)} geschrieben werden. Zieht man es aus Gr"unden der
"Ubersichtlichkeit vor, trotzdem einen boolee'schen Typ zu haben,
kann man diesen in SEL problemlos mit den Direktiven

\begin{verbatim}
#define boolean        int
#define false          0
#define true           1
\end{verbatim}

\noindent simulieren.\footnote{Die Betonung liegt auf "`simulieren"',
da wir uns mit diesen Direktiven keinen eigenen Typ schaffen sondern
lediglich \fu{boolean} zu einem Pseudonym f"ur \rw{int} machen.}

Der Typ \tfloat\ repr"asentiert eine vorzeichenbehaftete
Gleitkommazahl. Ihre Gr"o"se im Speicher h"angt vom jeweiligen
Prozessor ab.

Einen Spezialfall unter den Typen stellt \tvoid\ dar: Er
repr"asentiert gar nichts! F"ur sein Dasein sprechen zwei
Gr"unde. Erstens macht er es m"oglich, Funktionen ohne
R"uckgabewert zu konstruieren, also Konstrukte, die mit den
Prozeduren in \pascal\ oder \modula\ vergleichbar sind. Au"serdem
ist er der Grundtyp f"ur typenlose Zeiger, also Zeiger, die auf
etwas nicht Bestimmtes zeigen. Verwendung finden solche Zeiger
z.~B. in Funktionen zur Speicherallokation. Diese geben einen
Zeiger auf einen Speicherbereich bestimmter Gr"o"se zur"uck,
wobei aber nicht bekannt ist, f"ur welchen Typ dieser Bereich
reserviert wird.

\subsection{Arrays}

Ein Array ist ein Zusammenschlu"s mehrerer Variablen desselben
Typs, die im Speicher sinnvollerweise aufeinander folgen. Ich
habe diesen Typ in SEL aufgenommen, weil er erstens nicht schwer
zu implementieren und zweitens in sehr vielen Situationen
unverzichtbar ist. Auffallend an der Schreibweise des Arraytyps
ist, da"s die Dimension des Arrays dem Grundtyp nicht nach- sondern
vorgestellt sind. Gem"a"s unserem Grundsatz, Symbole anstatt
Schl"usselw"ortern zu verwenden, wo es nicht st"ort, habe ich hier
auf Ausschm"uckungen wie das in \pascal\ verwendete \fu{array
... of ...} verzichtet und statt dessen eine in eckige Klammern
geschlossene Array-Dimension benutzt. Auf ein Angeben der
Grenzen, wie in \pascal\ "ublich, habe ich verzichtet, das erste
Element besitzt den Index \fu{0}. Diese Vereinfachung hat zwei
Gr"unde: Einerseits ist es nicht schwer, auf speziell indizierte
Arrays zu verzichten, andererseits w"urde sich das auf die
Konsistenz in Bezug auf die Verwendung von Zeigern anstatt Arrays
auswirken: Angenommen, wir h"atten ein Array \fu{a : [5..10]int}.
In SEL ist es nun m"oglich, dieses Array als Zeigervariable zu
betrachten und so zu behandeln, d. h., wir k"onnten mit ihm
Zeigerarithmetik betreiben. Per definitionem ist ein Zeiger
\fu{p} gleich \fu{p + 0}, was bedeutet, da"s wir mit \fu{*(a + 0)} auf das
erste Element des Arrays zugreifen k"onnen. Wir k"onnen auf
dieses Element auch direkt "uber das Array zugreifen, also mit
\fu{a[5]}. Ich denke, hier zeigt sich schon, worauf ich
hinauswill: Die Zahlen, die zur Indizierung verwendet werden,
stimmen nicht "uberein. Um das zu umgehen, m"u"ste man entweder
auf die Gleichsetzung von Array und Zeiger verzichten, was ich
nicht wollte, oder auch in die Zeiger eine komplizierte
Nummerierung einbauen, was mir und sicher auch dem Benutzer der
Programmiersprache zu umst"andlich ist.

Zur"uck zu der Dimension des Arrays, die, wie oben erw"ahnt, dem
Grundtyp vorangestellt ist. Der Grund hierf"ur ist in erster
Linie beim Typecasting\footnote{siehe Kapitel \ref{casts}}
von SEL zu suchen. Nehmen wir als Beispiel den Ausdruck
\fu{a : *int[10]}. Hier w"are es unklar, ob die eckige Klammer noch
zum Typ geh"ort, oder ob versucht wird, auf das zehnte Element
des Zeigers (der als Array behandelt wird) zuzugreifen.
Nat"urlich kann auch dieses Problem durch Klammerung umgangen
werden, doch warum etwas komplizierter machen, als es ist? Die
vorangestellte Arraydimension in SEL hat n"amlich noch einen
Vorteil: Es ist ohne eine Rangordnung oder Klammerung m"oglich,
einen Typ eindeutig zu bestimmen. Als Beispiel soll der C-Typ
\fu{*int[10]} dienen. Ohne den beiden Operatoren eine
unterschiedliche Priorit"at zuzuweisen, ist es unm"oglich zu
sagen, ob der Typ nun \fu{*(int[10])} oder \fu{(*int)[10]} ist.
In SEL ist die Sache eindeutig: \fu{[10]*int} kann nur als Array
aus zehn Zeigern auf \tint\ gedeutet werden.

\subsection{Zeiger}

Zeiger stellen ein wichtiges Konstrukt dar, mit dessen Hilfe es
m"oglich ist, sehr komplexe, verkettete und verzweigte
Datenstrukturen zu bilden, die ohne sie nur sehr schwer zu
verwirklichen w"aren. Der Bau eines Interpreters oder gar eines
Compilers z.~B. w"are ohne die Verwendung von Zeigern kaum zu
bewerkstelligen. Aus diesem Grund habe ich Zeiger in SEL
eingebaut, was f"ur eine Interpretersprache eigentlich sehr
ungew"ohnlich ist.

Sehr eng in Verbindung mit Zeigern steht in SEL der Operator
\fu{*}, der einerseits in der Typdefinition eines Zeigers
verwendet wird und andererseits benutzt wird, einen Zeiger zu
dereferenzieren. Letzteres geschieht, indem dem Zeiger der
Sternoperator vorangestellt wird, anders als in \pascal , wo der
Dereferenzierungsoperator als Suffix verwendet wird. Der Grund,
warum in SEL der Operator vorangestellt wird, ist die Tatsache,
da"s er dasselbe Symbol hat wie der Multiplikationsoperator, und
der Parser\footnote{siehe Kapitel \ref{parser}} nicht w"u"ste,
wie etwa der Ausdruck \fu{p * - x} gemeint ist. Einerseits k"onnte es
die Differenz zwischen der Variablen, auf die \fu{p} zeigt und
\fu{x} sein, andererseits aber auch \fu{p} multipliziert mit
\fu{-x}. Nat"urlich h"atte ich auch ein anderes Symbol f"ur die
Dereferenzierung verwenden k"onnen, doch das wollte ich nicht,
weil ich den Operator \fu{*} von C gewohnt bin, au"serdem fand
ich es schon immer ein wenig verwirrend, da"s bei der
Typdeklaration der Operator vorangestellt, bei der
Dereferenzierung aber nachgestellt sein mu"s.

Der zweite wichtige Operator in Verbindung mit Zeigern ist
\fu{\&}. Wird er einer Variablen vorangestellt, dann ergibt er
einen Zeiger auf diese Variable. Mit seiner Hilfe ist es also
m"oglich, einen Zeiger auf eine bestimmte Variable zeigen zu
lassen. Verst"andlicherweise ist es nicht m"oglich, \fu{\&} in
Verbindung mit einer Konstante zu benutzen, denn eine Konstante
hat keinen bestimmten Platz im Speicher, sie wird vom Interpreter
direkt in einen Wert umgewandelt und von einem Compiler direkt in
den Code eingef"ugt. Zu verwechseln ist dies jedoch nicht mit den
String-Konstanten\footnote{siehe Kapitel \ref{const}}, die
eigentlich nicht konstante Strings sind, sondern konstante Zeiger
auf einen String, der an einer bestimmten Stelle im Speicher
liegt. Doch zur"uck zum Operator \fu{\&}: Sein zweiter
Verwendungszweck ist die Kennzeichnung eines Parameters einer
Funktion, der {\em Call-By-Reference} "ubergeben werden soll. In
\pascal\ z.~B. "ubernimmt diese Aufgabe das Schl"usselwort
\rw{var}, in C gibt es etwas Derartiges nicht, erst in C++ wurden
Referenzparameter eingef"uhrt. Auch dort werden sie mit \fu{\&}
gekennzeichnet. In einer Programmiersprache, in der Zeiger zur
Verf"ugung stehen, sind Referenzparameter nicht zwingend
notwendig: Es k"onnten Zeiger auf die betreffenden Variablen
"ubergeben werden, wie es in C auch praktiziert wird. Ich habe
mich aus zwei Gr"unden entschlossen, Referenzparameter doch in
SEL aufzunehmen. Erstens k"onnen sie in manchen Situationen
Programme um sehr viel "ubersichtlicher und leichter wartbar
machen. Zweitens, und das war f"ur mich wohl der gewichtigere
Grund, arbeitet der im n"achsten Kapitel theoretisch beschriebene
Interpreter intern ohnedies mit Referenzen, soda"s es keine
Schwierigkeit war, Referenzparameter einzubauen.

\subsubsection{Dynamische Variablen}

Eng verbunden mit den Zeigern ist das Konzept der dynamischen
Variablen. Dynamisch deshalb, weil sie erzeugt werden k"onnen,
wenn man sie braucht und aus dem Speicher entfernt, wenn sie
"uberfl"ussig werden. Statische Variablen belegen
Speicherbereiche, die vom Compiler bzw. Interpreter im voraus
bestimmt werden, wohingegen der Speicherbereich f"ur eine
dynamische Variable erst w"ahrend der Programmausf"uhrung
bestimmt werden kann. Das impliziert nat"urlich, da"s dem
Benutzer geeignete Mittel in die Hand gegeben werden m"ussen, um
dynamische Variablen zu erzeugen. Grunds"atzlich gibt es hierf"ur
zwei M"oglichkeiten. Entweder, man baut das Konzept der
dynamischen Speicherverwaltung nicht fest in die
Programmiersprache ein, sondern "uberl"a"st es dem Benutzer bzw.
dem Designer der Laufzeitbibliothek, geeignete Funktionen zu
implementieren, wie es in C der Fall ist, oder man
ber"ucksichtigt dynamische Speicherverwaltung schon beim Design
der Programmiersprache und stellt Konstrukte zur Verf"ugung, um
dynamische Variablen zu erzeugen, wie z.~B. in C++.

Ich habe mich f"ur den zweiten L"osungsansatz entschieden, weil
ich finde, da"s das Konzept der dynamischen Speicherverwaltung zu
wichtig ist, um es beim Design einer Programmiersprache nicht
miteinzubeziehen. Ein daraus resultierender Vorteil ist, da"s
Programme leichter lesbar werden und es leicher ist, sauberer zu
programmieren, als z.~B. in C.\footnote{Was nicht hei"sen will,
da"s es in C nicht m"oglich ist, sauber zu programmieren, es ist
nur nicht so leicht wie etwa in C++} Angelehnt an das Vorbild C++
werden in SEL zur dynamischen Speicherverwaltung die beiden
un"aren Operatoren \rw{new} und \rw{delete}\footnote{siehe
Kapitel \ref{memmgrops}} verwendet. Der Operand von \rw{new} ist ein
Typ, also im Gegensatz zu denen der meisten anderen Operatoren
kein Ausdruck. Der Ergebnistyp ist ein Zeiger auf den
"ubergebenen Typ, der auf den neu allokierten Speicherbereich
zeigt. Dieser Bereich kann nun so lange wie n"otig benutzt
werden. Wird er nicht mehr verwendet, sollte er mit dem Operator
\rw{delete} freigegeben werden, um sp"ater anderen dynamischen
Variablen zur Verf"ugung zu stehen. Der Grund, warum ich aus
\rw{new} und \rw{delete} keine Anweisungen gemacht
habe---besonders bei \rw{delete} w"urde sich das anbieten---ist,
da"s alle anderen Anweisungen in SEL zur Flu"skontrolle im
Programm verwendet werden, soda"s dies gewisserma"sen ein Bruch
in der Konsistenz der Sprache w"are.

\subsection{Strukturen}
\label{structs}

Eine Struktur\footnote{Der Name "`Struktur"' stammt von C, in
\pascal\ wird f"ur dasselbe der Name "`Record"' verwendet} ist
der Zusammenschlu"s mehrerer Variablen, auch verschiedenen Typs,
unter einem Namen. Sie bilden zusammen mit Zeigern und
dynamischen Variablen die Basis f"ur weit komplexere, verzweigte
und verkettete Datenstrukturen. "Uber ihre Definition m"ochte ich
nicht viele Worte verlieren, da sie einfach und zweckm"a"sig ist
und es au"serdem nicht viele andere M"oglichkeiten gibt,
Strukturen zu definieren.  Was jedoch viel interessanter ist, ist
die Benutzung von Zeigern auf eine Struktur in der Struktur
selbst. Ein Beispiel:

\begin{verbatim}
struct List
  i    : int;
  Next : *List;
end
\end{verbatim}

Diese Struktur ist eine m"ogliche Definition einer einfach
verketteten Liste. Neben einem oder mehreren Werten (in diesem
Fall ein \tint ) enth"alt sie einen Zeiger auf eine Struktur
ihres eigenen Typs. Einen Interpreter oder Compiler, der mit nur
einem Pass arbeitet und diese Struktur akzeptiert, zu schreiben,
ist noch leicht: Man k"onnte sich schon beim Lesen des auf das
Schl"usselwort \rw{struct} folgenden Bezeichners \fu{List}
notieren, da"s es diesen Typ von Struktur gibt. Beim Zeiger
\fu{Next} kann diese Notiz benutzt werden, um seinen Typ
festzulegen. Beim Lesen des Schl"usselwortes \rw{end} wird die
Notitz anschlie"send zu einer kompletten Definition erg"anzt. Was
tun wir jedoch, wenn es zwei Strukturen gibt, die jeweils einen
Zeiger auf die andere Struktur enthalten, wie es nicht selten
vorkommt? Eine der beiden Strukturen mu"s zuerst gelesen werden,
und dort treffen wir auf einen unbekannten Typen. In C wird dieses
Problem dadurch gel"ost, da"s man dem Compiler vor der Definition
der beiden Strukturen einen Hinweis darauf gibt, da"s diese
Strukturen irgendwo einmal definiert werden, soda"s er schon vor
der Definition seine Notiz machen kann. In SEL jedoch k"onnen wir
dieses Problem weit einfacher l"osen: Da es nicht m"oglich ist,
einem Typ in SEL einen anderen Namen zu geben, steht, wenn ein
unbekannter Typ gelesen wird, zweifelsfrei fest, da"s es sich
dabei um eine Struktur handeln mu"s, soda"s unser Compiler bzw.
Interpreter bereits hier die Notiz vornehmen kann.

\subsection{Typecasting}
\label{casts}

Bei einer Gegen"uberstellung der Typkonzepte von C und \pascal\
f"allt auf, da"s \pascal\ in der Typ"uberpr"ufung weit strenger
ist als C. So ist es in C ohne weiteres m"oglich, einem \tchar\
einen numerischen Wert zuzuweisen, wohingegen in \pascal\ zu
diesem Zweck ein eigener Typ existiert. In der Sprache \pascal\
ist eine Umwandlung eines Typs in einen anderen gar nicht
vorgesehen,\footnote{au"ser in F"allen, wo die beiden Typen zur
gleichen Gruppe geh"oren, wie z.~B. \tint\ und \type{longint}} in
C wird sie in vielen F"allen sogar vom Compiler "ubernommen, an
anderen Stellen kann sie der Benutzer mittels Typecasting
durchf"uhren. Der Grund, warum ich in SEL ein flexibles
Typkonzept eingebaut und nicht die strenge "Uberpr"ufung von
\pascal\ "ubernommen habe, ist, da"s es bei vier Basistypen sehr
schwer ist, ohne Typumwandlung auszukommen. Au"serdem k"onnen
Programme durch implizites Typecasting um einiges k"urzer und
"ubersichtlicher werden, man denke nur an die Umwandlung zwischen
\tchar\ und \tint , wof"ur in \pascal\ eigene Funktionen
notwendig sind.

Um aber nicht in einem vollst"andigen Chaos zu enden, darf nicht
wahllos jeder Typ implizit in jeden anderen Typen umgewandelt
werden, nur damit im Programm ohne Fehlermeldung fortgefahren
werden kann. Die Regeln, die in SEL f"ur implizites sowie
explizites Typecasting gelten, sind in \cite{selkr}, Seite 24
zusammengefa"st.

Der Grund, warum ich nicht das in C "ubliche Typecasting, einem
Ausdruck den neuen Typ in Klammern voranzustellen, "ubernommen
habe, ist der, da"s ich einfach etwas Neues ausprobieren wollte,
und das Typecasting mittels des \fu{:}-Operators exzellent zu den
"ubrigen Konstrukten in SEL (Variablen- und Strukturdefinition,
Parameterdefinition in Funktionen u.~s.~w.) pa"st. Es zeigt sich,
da"s Programme bei richtiger Verwendung des Typecasting-Operators
weit "ubersichtlicher sind als in C, meine Wahl hat sich also in
keiner Weise als Fehlgriff erwiesen.

\subsection{Mengen, Listen, B\protect{"aume} etc.}

Komplexere Typen, wie z.~B. Mengen, Listen oder B"aume sind in
SEL nicht vorhanden, aus guten Gr"unden. Mengen einerseits
werden relativ selten verwendet und lassen sich leicht
vermeiden. Au"serdem sind sie sehr leicht mit den Mitteln von
SEL zu implementieren, soda"s das Einbeziehen in SEL nur
"au"serst wenig Sinn h"atte. Listen und B"aume hingegen werden
zwar oft verwendet, sind allerdings so komplex und existieren in
einer derartig gro"sen Anzahl von Variationen, da"s es mit
einfachen Mitteln kaum m"oglich ist, sie vern"unftig zu
unterst"utzen.

\section{Bezeichner und Schl\protect{"usselw\protect{"orter}}}

Wie in vielen Programmiersprachen "ublich, sind auch in SEL
Schl"usselw"orter reserviert, d.~h., sie k"onnen nicht als
Bezeichner verwendet werden, aus zwei Gr"unden. Erstens wird
dadurch die Programmierung eines Interpreters bzw. Compilers
erheblich erschwert, und zweitens w"aren dadurch sehr
verwirrende Konstruktionen wie z.~B.

\begin{verbatim}
for (for = to to step step end)
  ...
end
\end{verbatim}

\noindent m"oglich, was nicht im Sinn des Designers sein kann.

Schl"usselw"orter und Bezeichner sind in SEL case-sensitiv,
d.~h., da"s zwei W"orter, die sich nur dadurch unterscheiden,
da"s ein oder mehrere Buchstaben statt gro"s klein geschrieben
sind oder umgekehrt, als unterschiedlich betrachtet werden. Das
hat drei Vorteile: Erstens wird dadurch die Entwicklung der
lexikalischen Analyse\footnote{siehe Kapitel \ref{scanner}}
ein wenig einfacher, zweitens k"onnen dadurch verschiedene Arten
von Bezeichnern (Bezeichner f"ur Funktionen, Variablen, Makros
u.~s.~w.) durch verschiedene Schreibweise gekennzeichnet werden.
Drittens wird durch eine vorgegebene Schreibweise bei
unerfahrenen oder schlampigen Programmierern die
"Ubersichtlichkeit der Quelltexte gesteigert.

Die Syntax der Bezeichner ist relativ simpel: Sie bestehen aus
einem Buchstaben oder Unterstrich, gefolgt von einer beliebigen
Anzahl von Buchstaben, Ziffern oder Unterstrichen. Der Grund,
warum Bezeichner nicht mit Ziffern beginnen d"urfen ist der, da"s
dadurch eine Eindeutigkeit der Sprache nicht mehr gew"ahrleistet
w"are. So k"onnte z.~B. \fu{12to} einerseits als Zahl \fu{12}
gefolgt vom Schl"usselwort \rw{to}, andererseits aber auch als
Bezeichner gedeutet werden. Dieser Konflikt lie"se sich zwar
durch Verwendung von Leerzeichen vermeiden, w"urde aber die
lexikalische Analyse ungleich schwerer machen. Warum ich nicht
auch andere Symbole in Bezeichnern erlaubt habe, ist einfach: Je
weniger Symbole in Bezeichnern verwendet werden d"urfen, desto mehr
hat man f"ur andere Zwecke frei.
                            
\section{Weiterf\protect{"uhrende} Projekte}
\label{langext}

Keine Programmiersprache ist perfekt, und schon gar nicht SEL,
f"ur dessen jetzige Form ich bewu"st die Versionsnummer 0.3
gew"ahlt habe. Leider war es mir aus Zeit- und Platzgr"unden
nicht m"oglich, auch nur einen Teil der von mir geplanten
Erweiterungen in diese Arbeit einzubringen, soda"s ich sie in
diesem Kapitel kurz zusammenfassen m"ochte, einerseits, um zu
zeigen, was SEL noch alles fehlt und andererseits, um den Leser
anzuregen, sich mit der Materie eingehender auseinanderzusetzen.

\begin{description}
\item[Zeiger auf Funktionen] stellen ein sehr praktisches
Sprachkonzept dar. Sie erm"oglichen es z.~B., in einer
graphischen Benutzeroberfl"ache, wie etwa Windows NT oder meinen
Schani-Tools, jedem Bildschirmfenster eine Funktion zuzuweisen,
die die Verwaltung desselben "ubernimmt. Sprachlich gesehen liegt
die einzige Schwierigkeit bei der Typangabe eines solchen
Funktionszeigers, denn der Aufruf einer Funktion durch einen
Zeiger kann genauso wie bei einer "`richtigen"' Funktion durch
den Namen der Zeigervariablen, gefolgt von einer in Klammern
gesetzten Parameterliste erfolgen, und die Adresse einer Funktion
kann mittels deren Namen, ohne nachfolgende Parameterliste in
Klammern, ermittelt werden. Doch auch der Typ f"ur einen
Funktionszeiger kann relativ leicht ohne Zweideutigkeiten in die
Grammatik eingebaut werden. So k"onnte etwa \fu{char(int,float)}
f"ur einen Zeiger auf eine Funktion, die einen \tint\ und einen
\tfloat\ als Parameter erwartet und einen \tchar\ zur"uckliefert,
stehen.\footnote{Tats"achlich f"uhrt eine Erweiterung wie oben
beschrieben zu einem kleinen Konflikt in der Grammatik, der
jedoch leicht beseitigt werden kann.}

\item[Variante Strukturen] sind Strukturen, die zwar mehrere
Elemente haben, diese jedoch denselben Speicherbereich
einnehmen. So etwas ist immer dann von Vorteil, wenn eine
Struktur ben"otigt wird, die Objekte beschreibt, die zwar in
gewisser Hinsicht "ahnlich sind, jedoch andere Parameter haben.
Ein Beispiel hierf"ur ist die in Kapitel \ref{intexprs}
besprochene Datenstruktur f"ur SEL-Ausdr"ucke.

\item[Makros mit Parametern] w"urden, wenn sie so flexibel
aufgebaut sind, da"s neue Befehle durch sie
"`vorget"auscht"' werden k"onnen, perfekt in eine Makrosprache
passen. Ein gutes Beispiel f"ur solche Makros sind die sog. User
Defines Commands (UDC's) von Clipper.

\item[Objektorientierte Programmierung] ist zwar ein edles Ziel,
jedoch nicht leicht zu verwirklichen.
\end{description}

\chapter{Der Interpreter}

Das Ziel dieses zweiten Kapitels ist es, f"ur die im ersten
Kapitel erkl"arte Programmiersprache SEL einen Interpreter zu
entwerfen. Der Titel dieser Arbeit deutet darauf hin, da"s dieser
Entwurf theoretisch ist, also in dieser Arbeit keine
Implementation eines Interpreters zu finden ist. Die Gr"unde
hierf"ur sind verschieden. Schon lange bevor ich "uberhaupt die
Idee zu dieser Fachbereichsarbeit hatte, wollte ich einen
Interpreter schreiben. Als ich mit Herrn Prof. Freiler "uber das
Thema dieser Arbeit diskutierte, legte er mir die Beschr"ankung
auf, da"s, wenn das Thema dieser Arbeit ein Programm w"are,
dieses in \pascal\ geschrieben sein m"u"ste. Da ich aber den
Interpreter in erster Linie f"ur mich und nicht f"ur diese Arbeit
programmieren wollte, hatte ich vor, ihn in C zu schreiben. Aus
diesem Grund einigte ich mich mit Herrn Prof. Freiler darauf,
da"s ich den Interpreter nur theoretisch, also unabh"angig von
irgendeiner Implementation, behandeln w"urde. Im Nachhinein
betrachtet bin ich froh "uber diese Entscheidung, da es kaum
m"oglich gewesen w"are, in einer im Umfang derma"sen beschr"ankten
Arbeit einen ganzen Interpreter und eine angemessene Erkl"arung
dazu unterzubringen.

Doch auch mit der Beschr"ankung auf eine theoretische Abhandlung
ist der Umfang dieser Arbeit zu gering, um ausf"uhrlich auf alle
Details einzugehen. Aus diesem Grund m"ochte ich im Gegensatz 
zum ersten Kapitel, in dem ich alle meine Entscheidungen 
begr"undet und evtl. auch andere L"osungsans"atze beschrieben 
habe, in diesem nur den Entwurf pr"asentieren und lediglich in
einigen seltenen F"allen andere m"ogliche L"osungen diskutieren.

Dieses Kapitel beschreibt alle wichtigen Algorithmen und
Datenstrukturen des Interpreters, geht jedoch nicht n"aher auf
grundlegende Datenstrukturen wie etwa dynamische Listen ein.
Sollte der Leser mit diesen Strukturen nicht vertraut sein,
empfehle ich ihm die Lekt"ure von \cite{sedge}, Kapitel 3--5 und
16.

\section{Der Begriff des Interpreters}
\label{whatisit}

Um ein Programm einer h"oheren Programmiersprache auf einem
Computer auszuf"uhren, kann man sich zweier verschiedener
Werkzeuge bedienen: Interpreter und Compiler. Der Begriff
"`Compiler"' wird vom Duden Informatik wie folgt definiert:

\begin{quotation}
Programm, das Programme aus einer Programmiersprache A in eine
Programmiersprache B "ubersetzt.\footnote{siehe \cite{duden},
Seite 745}
\end{quotation}

Das bedeutet, da"s ein Compiler ein Programm gar nicht direkt
ausf"uhrt, sondern es lediglich in eine andere Programmiersprache
"ubersetzt. In der Regel ist diese Sprache eine Assemler- oder
Maschinensprache, die direkt vom Computer "`verstanden"' wird. Im
Gegensatz dazu ist ein Interpreter

\begin{quotation}
Ein Programm, welches ein Programm einer anderen
Programmiersprache nach den notwendigen syntaktischen
"Uberpr"ufungen sofort ausf"uhrt. Im Gegensatz zu einem
"Ubersetzer\footnote{andere Bezeichnung f"ur Compiler} mu"s das
Quellprogramm nicht erst in eine andere Programmiersprache
"ubersetzt werden, sondern der Interpreter analysiert
nacheinander jede Anweisung und Deklaration des Quellprogramms
und f"uhrt diese unmittelbar aus.\footnote{siehe \cite{duden},
Seite 322}
\end{quotation}

In der Praxis jedoch wird kaum ein Interpreter ein Programm
ausf"uhren, ohne es vorher in eine geeignetere Form umgewandelt
zu haben. Man kann sich leicht vorstellen, da"s ein Interpreter,
der etwa bei einem Funktionsaufruf alle Module nach der Deklaration
dieser Funktion durchsuchen mu"s, zu langsam ist, um sinnvoll
eingesetzt werden zu k"onnen. Auch das wiederholte Lesen und
Erkennen der Schl"usselw"orter, die im Quelltext verwendet
werden, ist zu zeitaufwendig, als da"s es w"ahrend der
Ausf"uhrung des Programms durchgef"uhrt werden k"onnte. In den
folgenden Kapiteln werden wir deshalb besprechen, wie sich ein
Interpreter die Arbeit so einteilen kann, da"s eine m"oglichst
hohe Ausf"uhrungsgeschwindigkeit erreicht wird.

Ich m"ochte anmerken, da"s das in diesem Kapitel besprochene
Design f"ur einen Interpreter zwar theoretisch, jedoch keineswegs
realit"atsfern ist. Ich habe selbst einen SEL-Interpreter auf
diese Weise implementiert.

\section{Einsatzgebiete eines Interpreters}

Wenn man bedenkt, da"s die Ausf"uhrung eines Programms durch
einen Interpreter oft mehrere hundert Male l"anger dauert als die
eines mit einem Compiler erzeugten Maschinenspracheprogramms, fragt
man sich, welchen Zweck ein Interpreter eigentlich erf"ullen
kann. Ein Haupteinsatzgebiet f"ur Interpreter sind Computer, die
f"ur die sinnvolle Verwendung eines Compilers zu langsam sind
und/oder zu wenig Speicher bieten. Beispiele hierf"ur sind alte
Home-Computer wie z.~B. der Commodore~C-64 aber auch
Taschenrechner wie etwa der Texas Instruments~TI-85 oder der
Hewlett Packard~HP-48. Oft werden Interpreter auch zum
Entwickeln von Programmen eingesetzt, da sie den Vorteil haben,
Programme ohne Zeitverz"ogerung ausf"uhren zu k"onnen. Ein
anderes wichtiges Einsatzgebiet f"ur Interpreter ist die
Steuerung von Me"sger"aten und Robotern. Dort stellt die langsamere
Ausf"uhrung von Programmen kaum einen Nachteil dar, andererseits
ist es ein gro"ser Vorteil, da"s Programme ge"andert werden
k"onnen und sofort wieder einsatzbereit sind.

\section{Bestandteile eines Interpreters}

Bevor wir uns daranmachen, uns dar"uber zu unterhalten, wie der
Interpreter aussehen und welche Leistung er bieten soll, m"ussen
wir "uber den groben Aufbau eines Interpreters Bescheid wissen.
Wie bereits oben erw"ahnt wird in der Praxis das Quellprogramm
vor dem Interpretieren in eine f"ur den Interpreter geeignete
Form umgewandelt ({\em Tokenizing}). Dieser Vorgang wird in der
Regel von der lexikalischen Analyse erledigt. Dabei wird der
Quellcode sequentiell abgearbeitet und Schl"usselw"orter und
Bezeichner codiert. Zus"atzlich werden die Positionen von wichtigen
Schl"usselw"ortern wie z.~B. \rw{function}, \rw{global} und
\rw{local} gespeichert. Im Fall von SEL wird gleichzeitig
zur lexikalischen Analyse der Pr"aprozessor ausgef"uhrt. Er hat die
Aufgabe, Include-Dateien einzubinden und Makros in den
Programmtext einzuf"ugen.

Im Anschlu"s an die lexikalische Analyse und das Tokenizing kann
mit der Ausf"uhrung des Programms begonnen werden. Bei einer
Programmiersprache wie SEL m"ussen zuerst nat"urlich spezielle
Initialisierungen vorgenommen werden, n"amlich das Einrichten der
modul- und programmglobalen Variablen und das Analysieren der
Funktionsk"opfe. Zus"atlich speichert der Interpreter die
Stellen, an denen die jeweiligen Funktionen beginnen, soda"s
nicht das gesamte Programm durchsucht werden mu"s. Jetzt wird das
Programm durch einen Aufruf der \fu{main}-Funktion gestartet. Die
einzelnen Anweisungen und Ausdr"ucke werden vom Parser analysiert
und sofort ausgef"uhrt. Dabei mu"s nat"urlich auf spezielle
Routinen, die z.~B. f"ur die Verwaltung der Variablen und
"ahnliches verantwortlich sind, zur"uckgegriffen werden.

Zusammenfassend l"a"st sich also sagen, da"s ein Interpreter aus
drei wichtigen Teilen besteht: Der lexikalischen Analyse, die das
Tokenizing vornimmt, dem Parser, der das Programm syntaktisch
analysiert und den Verwaltungsroutinen, die die Hintergrundarbeit
leisten.

\section{Zielsetzung f\protect{"ur} das Design}
\label{intgoals}

Wie schon im ersten Kapitel erw"ahnt, ist SEL mit dem Ziel, sie
als Makrosprache zu verwenden, entworfen worden. Es w"are also
angebracht, den Interpreter von Beginn an so zu entwerfen, da"s
er nicht ein eigenst"andiges Programm darstellt, sondern eher
eine Art Modul, das man in andere Programme aufnehmen kann. Das
setzt nat"urlich voraus, da"s er nicht eine Unmenge an Speicher
verbraucht und auch, da"s er den Speicher, den er verbraucht,
nach der Ausf"uhrung des SEL-Programms wieder freigibt, soda"s er
wieder dem Hauptprogramm zur Verf"ugung steht.

Der Interpreter sollte nat"urlich, um ohne Bedenken in einem
Programm eingesetzt werden zu k"onnen, Programmierfehler so
schnell wie m"oglich erkennen k"onnen, den Ausf"uhrungsproze"s
angemessen stoppen, allen Speicher, den das SEL-Programm
angefordert hat, und auch den, den der Interpreter im Zuge der
Ausf"uhrung des Programms benutzt hat, freigeben und den
Fehler dem Hauptprogramm melden.

Wie auch die Programmiersprache SEL sollte auch der Interpreter
erweiterungsf"ahig sein, d.~h., er sollte dem Hauptprogramm
erlauben, bestimmte Funktionen auch SEL-Programmen zug"anglich zu
machen, wie beispielsweise f"ur das Zeichnen einer Linie in einem
Zeichenprogramm.

Da ich mich in letzter Zeit mit dem Multitasking-Betriebssystem
Windows NT auseinandergesetzt habe, wollte ich auch meinen
Interpreter multitaskingf"ahig machen, soda"s es in einem
Programm m"oglich ist, mehrere SEL-Programme gleichzeitig
auszuf"uhren. Um mich nicht falsch zu verstehen: Das Multitasking
wird in keiner Weise vom Interpreter in die Hand genommen,
sondern obliegt einzig und allein dem Betriebssystem. Der
Interpreter allerdings mu"s so entworfen werden, da"s er das
Multitasking nicht erschwert oder gar unm"oglich macht, was, wie
sich zeigen wird, nicht sehr schwer ist.

Im Laufe dieses Kapitels werde ich auf einige, aber nicht auf alle
Forderungen dieses Unterkapitels eingehen, weil sie zu trivial
sind. So ist z.~B. die Freigabe des benutzten Speichers nur eine
Formsache und mu"s nicht extra von mir erl"autert werden.

\section{Die lexikalische Analyse}
\label{scanner}

Die lexikalische Analyse hat die Aufgabe, den Zeichenstrom des
Quelltextes in einen Tokenstrom zu verwandeln. Als Token
bezeichnet man die kleinste Einheit, die vom Parser bearbeitet
wird. Token sind z.~B. die Schl"usselw"orter, Bezeichner,
Konstanten u.~s.~w. Der Grund, warum man die Erkennung von
Schl"usselw"ortern etc. nicht auch dem Parser "uberl"a"st, was
durchaus m"oglich w"are, ist der, da"s ein Parser diese Aufgabe
bei weitem nicht so schnell ausf"uhren kann wie ein eigens daf"ur
angefertigter {\em Scanner}.\footnote{Begriff f"ur ein Programm
oder einen Programmteil, der die lexikalische Analyse
durchf"uhrt} Au"serdem wird der Bau eines Parsers wesentlich
vereinfacht, wenn man einen Scanner verwendet.

Um die Arbeitsweise eines Scanners zu beschreiben, m"ochte ich
zun"achst ein Beispiel heranziehen. Nehmen wir an, der Scanner
h"atte den Zeichenstrom \fu{for (x=1 to 10)} zu bearbeiten. Er
untersucht nun das erste Zeichen und merkt, da es sich um einen
Buchstaben handelt, da"s er es entweder mit einem Bezeichner oder
mit einem Schl"usselwort zu tun hat. Die einfachste Methode,
dieses Token zu analysieren ist, so lange Buchstaben, Ziffern und
Unterstriche zu lesen, bis er an einem Zeichen ankommt, das
nicht mehr aus einer der genannten Kategorien stammt. Der
Scanner w"urde also zuerst das \fu{f}, dann das \fu{o} und
anschlie"send das \fu{r} lesen. Da das n"achste Zeichen weder
Buchstabe, Ziffer noch Unterstrich ist, stoppt er hier. Die
einfachste Methode, herauszufinden, ob man es hier mit einem
Schl"usselwort oder einem Bezeichner zu tun hat, ist, eine
Tabelle anzulegen, die alle Schl"usselw"orter der Sprache
umfa"st. Findet der Scanner das Wort in der Tabelle, ist es ein
Schl"usselwort, andernfalls ein Bezeichner. Das n"achste Zeichen
ist ein Leerzeichen. In einer formlosen Programmiersprache wie
SEL werden Leerzeichen und Tabulatoren "ubersprungen, soda"s der
Scanner nichts anderes zu tun hat, als zum n"achsten Zeichen
vorzur"ucken. Zwar sind Zeilenumbr"uche f"ur den Parser nicht
relevant, doch f"ur den Pr"aprozessor sind sie wichtig, soda"s
sie gemeldet werden sollten.

Da es nur ein Symbol in SEL gibt,
das mit einer "offnenden, runden Klammer beginnt, n"amlich dieses
selbst, wei"s der Scanner sofort, mit welchen Token er es zu tun
hat und kann vorr"ucken. Interessant wird es nun bei der Analyse
des Symbols \fu{=}, denn hier kann der Scanner nicht sofort davon
ausgehen, da"s es sich um den Zuweisungsoperator handelt, da es
noch ein anderes Symbol, n"amlich \fu{==} gibt, das mit einem
\fu{=} anf"angt. So mu"s der Scanner auch hier das n"achste
Zeichen betrachten um entscheiden zu k"onnen, welches Symbol er
vor sich hat.

Wie dieses Beispiel zeigt, ist es f"ur einen
Scanner vital zu wissen, welches Zeichen als n"achstes vorliegt,
wobei er sich dieses aber auch f"ur das n"achste Token aufheben
k"onnen mu"s. Die L"osung ist einfach: Man arbeitet mit einem
Zeichen im voraus, d.~h. es wird immer das kommende Zeichen
betrachtet, nicht das aktuelle. Wird es akzeptiert, liest man das
n"achste Zeichen ein, wenn nicht, kehrt man mit dem erkannten
Token zur aufrufenden Funktion (in Falle von SEL der
Pr"aprozessor) zur"uck. Beim n"achsten Aufruf des Scanners wird
mit dem vorher zur"uckgestellten Zeichen weitergearbeitet.

Um den Proze"s der lexikalischen Analyse zus"atzlich zu
beschleunigen, sollte die Eingabe gepuffert werden. Dabei wird
ein Teil der Eingabe in einen Zwischenspeicher (Puffer) gelesen,
um von dort aus vom Scanner abgearbeitet zu werden. Wird das Ende
des Puffers erreicht, mu"s abermals ein Teil der Eingabe gelesen
werden, so lange, bis das Ende der Datei erreicht ist. In diesem
Fall sollte der Scanner der "ubergeordneten Funktion Bescheid
geben.

Bei bestimmten Tokentypen gen"ugt es nicht, zu melden, da"s
dieses Token erkannt wurde, sondern es mu"s auch ein Attribut
zur"uckgeliefert werden, im Falle eines Bezeichners sein Name, im
Fall einer \tint -Konstante ihr Wert u.~s.~w., es sollte also
f"ur ein Token eine geeignete Datenstruktur verwendet werden, in
der alle n"otigen Attribute Platz finden. F"ur unseren
Interpreter k"onnte diese Struktur etwa so aussehen:

\newpage
\begin{verbatim}
struct Token
  iType  : int;
  pcName : *char;
  cChar  : char;
  iInt   : int;
  fFloat : float;
end
\end{verbatim}

Ist das Token ein Bezeichner, enth"alt \fu{pcName}
dessen Namen, im Falle einer \tchar -Konstanten enth"alt
\fu{cChar} ihren Wert u.~s.~w.

\subsection{Der Pr\protect{"aprozessor}}

Der Pr"aprozessor ist in gewisser Weise eine Zwischenstation von
Scanner und Parser. In SEL hat er die Aufgabe, den Inhalt von
Include-Dateien in den Programmtext einzuf"ugen, Kommentare zu
entfernen und Makros zu bearbeiten.

Das Entfernen von Kommentaren ist im Grunde kein Problem. Wird
vom Scanner das Symbol \fu{//} erkannt, dann werden so lange
Token gelesen, bis ein Zeilenumbruch gemeldet wird, wonach der
Kommentar beendet ist. Zus"atzlich zu den Kommentaren sollten
auch alle Zeilenumbr"uche herausgefiltert werden, da sie den
Parser st"oren w"urden.

Etwas komplizierter ist das Einf"ugen von Include-Dateien. Da es
m"oglich sein sollte, Include-Dateien zu verschachteln, gibt es
zwei M"oglichkeiten zur Implementation. Entweder, der
Pr"aprozessor arbeitet rekursiv, soda"s er sich selbst nochmals
aufruft, wenn eine Include-Datei zu bearbeiten ist, oder aber er
benutzt einen Stack zur Zwischenspeicherung der einzelnen
Include-Level, was am Ende eigentlich auf dasselbe hinausl"auft.
Ein kleines Problem ergibt sich mit dem Puffer f"ur den
Scanner: Wenn der Pr"aprozessor nur einen einzigen Puffer f"ur alle
Include-Level zusammen benutzt, mu"s er beim Eintritt in eine
Include-Datei geleert und mit
der Eingabe aus der Include-Datei gef"ullt werden. Ist diese zu Ende,
mu"s der Puffer abermals geleert werden und mit dem
"ubergeordneten Level gef"ullt werden. Das Problem dabei ist,
da"s man in der Datei dabei ein wenig zur"ucksetzen mu"s, weil
ein Teil des Puffers gel"oscht wurde. Ist das aus
irgendwelchen Gr"unden nicht m"oglich oder zu ineffizient,
mu"s man f"ur jeden einzelnen Include-Level einen Puffer
benutzen, was nat"urlich mehr Speicher verbraucht.

Zum Bearbeiten von Makros mu"s eine Datenstruktur
angelegt werden, in der alle Makros aufgezeichnet und auch schnell wieder
gefunden werden k"onnen. Ein bin"arer Baum und eine Hash-Tabelle
bieten sich hierzu an. F"ur jeden einzelnen Eintrag in dieser
Struktur mu"s zus"atzlich noch gespeichert werden, durch welchen
Text er zu ersetzen ist. Da der Pr"aprozessor auf den Scanner
aufbaut, ist es einfacher und effizienter, nicht den Zeichenstrom
zu speichern, sondern den Tokenstrom, der vom Scanner geliefert
wird. Dazu verwendet man am einfachsten eine einfach verkettete
Liste aus Token:

\begin{verbatim}
struct TokenList
  tokToken : Token;
  Next     : *TokenList;
end
\end{verbatim}

Trifft der Pr"aprozessor nun auf einen Bezeichner, mu"s er zun"achst
"uberpr"ufen, ob nicht vielleicht ein Makro unter diesem Namen
exisitiert. Dazu sucht er
nach ihm in der Datenstruktur. Ist er dort vorhanden,
setzt der Pr"aprozessor anstatt seiner die Tokenliste ein,
andernfalls wird er wie ein normaler Bezeichner behandelt.

Nachdem es m"oglich sein soll, innerhalb eines Makros ein anderes
zu benutzen, m"ussen wir noch etwas einbauen. Die einfachste
L"osung ist, da"s wir zum Einf"ugen eines Tokens in unseren
Zwischencode eine eigene Funktion verwenden, die, falls sie auf
einen Bezeichner trifft, der als Makro definiert ist, sich f"ur
jedes der Token aus der Liste selbst wieder aufruft. Dadurch ist
garantiert, da"s, falls in dieser Liste ein Makro auftaucht,
auch dieses entsprechend behandelt wird.
                              
\subsection{Der Zwischencode}
\label{tokencode}

Wie schon oben erw"ahnt, verwenden wir eine Art Zwischencode, um
dem Interpreter eine schnelle Ausf"uhrung zu erm"oglichen. Dabei
werden Schl"usselw"orter und Symbole mit einer eindeutigen Nummer
codiert, soda"s der Tokentyp sofort ersichtlich ist. Bei
bestimmten Tokentypen m"ussen zus"atzliche Informationen
gespeichert werden, wie z.~B. bei Konstanten oder Bezeichnern.

Um die Gr"o"se des Zwischencodes nicht zu sehr aufzubl"ahen,
werden auch die Namen der Bezeichner mit einer Nummer
codiert. Dazu wird eine Tabelle aller im Modul vorkommenden
Bezeichner angelegt und jedem Eintrag eine eindeutige Nummer
zugewiesen. Diese Nummer wird anschlie"send im Zwischencode
verwendet, um Bezeichner zu charakterisieren. Zu beachten ist,
da"s nat"urlich auch der Parser "uber diese Tabelle verf"ugen
mu"s, da er sonst nur die Nummern der Bezeichner wei"s. Wird
allerdings in allen Modulen des Programms dieselbe Tabelle
verwendet, so kann der Interpreter das Programm auch ohne das
Wissen um die Namen der Bezeichner ausf"uhren.

Eine besondere Hilfeleistung erstattet der Zwischencode dem
Interpreter bei Schleifen und Kontrollanweisungen: F"ur jedes \rw{if}
wird zus"atzlich die Position im Zwischencode abgespeichert, an der das
zugeh"orige \rw{end} bzw. \rw{else} steht. Das hat den Vorteil,
da"s der Interpreter, wenn er z~B. auf eine \rw{if}-Anweisung
trifft, deren Bedingung logisch falsch ist, ohne lange zu suchen
zum zugeh"origen \rw{end} bzw. \rw{else} springen kann, was
verst"andlicherweise die Ausf"uhrungsgeschwindigkeit erheblich
steigert, besonders bei Schleifen.

Wie bereits erw"ahnt, werden dem Interpreter zus"atzlich die
Position aller Vorkommnisse der Schl"usselw"orter \rw{global},
\rw{local} und \rw{function} mitgeteilt, soda"s er ohne zu suchen die
Initialisierung der programm- und modulglobalen Variablen und der
Funktionsk"opfe vornehmen kann.

Neben der erheblichen Geschwindigkeitssteigerung, die der
Zwischencode erm"oglicht, hat er noch den Vorteil, da"s er
weniger Speicherplatz als der Quellcode ben"otigt, bei etwas
gr"osseren Programmen etwa die H"alfte, was besonders unter
MS-DOS entgegenkommend ist.

Ich m"ochte hier kurz noch darauf eingehen, wie die oben
angesprochenen Verweise von Anweisungen auf die zugeh"origen
\rw{end}s bzw. \rw{else}s erzeugt werden k"onnen. Diese auf den
ersten Blick umst"andliche Aufgabe l"a"st sich relativ leicht
l"osen, indem man einen Stack zur Speicherung der
Verschachtelungen der Anweisungen verwendet. Trifft der Scanner z.~B.
auf ein \rw{if}, dann pusht er die Position, an der er sich
gerade im Zwischencode befindet, auf den Stack und l"a"st Platz
f"ur den sp"ater folgenden Verweis frei. Als n"achstes trifft
er z.~B. auf ein \rw{for} und tut dasselbe. Irgendwann mu"s
er jedoch zu dem \rw{end} gelangen, da"s die jeweils letzte
Anweisung schlie"st. In diesem Fall speichert er das \rw{end} im
Zwischencode, poppt vom Stack die Position der letzten
Anweisung und f"ugt an dieser Stelle die Position ein, an der
er soeben das \rw{end} gespeichert hat. Diese Technik wird
"ubrigens {\em Backpatching} genannt und oft im Compilerbau
eingesetzt.\footnote{Tats"achlich ist unser Scanner, der
einen SEL-Quelltext in Zwischencode umwandelt, streng gesehen
auch ein Compiler.}

\bigskip

Eine effiziente und einfache Darstellung des Zwischencodes ist
die Verwendung jeweils eines Bytes f"ur ein Token, wobei bei
bestimmten Tokentypen, wie z.~B. dem \rw{if}-Token, ein
zus"atzlicher Verweis auf ein anderes Token eingebaut wird. Wie
viele Bytes f"ur diesen Verweis verwendet werden, h"angt von den
Anforderungen an den Interpreter ab. Ich verwende hierzu in
meinem Interpreter 32 Bits, soda"s der Zwischencode theoretisch
bis zu 4 GByte gro"s sein kann, was sicherlich f"ur alle Anwendungen
mehr als ausreichend ist. F"ur die meisten
Anwendungen w"urden auch 16 Bits reichen, da Module in der Regel
sowieso klein gehalten werden. Bei anderen Token wie etwa einer
String-Konstante mu"s au"ser dem Tokentyp noch ein Wert
gespeichert werden, wie in diesem Fall der String.

\section{Der Parser}
\label{parser}

Der wohl komplizierteste Teil des Interpreters ist der
Parser. Er hat zur Aufgabe, aus dem Tokenstrom des Zwischencodes
die einzelnen Anweisungen zu erkennen. Sehr wichtig f"ur den Bau
eines Parsers ist die Grammatik, denn sie gibt an, was der Parser
eigentlich erkennen soll. Aus verschiedenen Gr"unden werden wir
Parser schreiben, die Teile der Sprache SEL erkennen k"onnen,
n"amlich Typen, Ausdr"ucke und Folgen von Anweisungen. Die Parser
f"ur Funktionsk"opfe und \rw{global}- bzw. \rw{local}-Bl"ocke
sind trivial, soda"s ich hier nicht speziell auf sie eingehen
werde. Ein Grund f"ur diese Trennung ist, da"s der
SEL-Interpreter vor allem f"ur die Verwendung in anderen
Programmen gedacht ist und dort ist es in vielen Situationen von
Vorteil, wenn der Benutzer Ausdr"ucke eingeben kann, z.~B. als
Filterbedingung in einer Datenbank.
                              
\subsection{Der Typen-Parser}

Beim Design eines Parsers ist es zun"achst wichtig, "uber die
Grammatik Bescheid zu wissen, die der Parser erkennen soll. In
diesem Unterkapitel besch"aftigen wir uns mit einem sehr
einfachen Teil der Grammatik von SEL, n"amlich den Typen. Das
gibt uns die M"oglichkeit, uns "uber die verschiedenen Arten von
Parsern und "uber die Zusammenh"ange zwischen Grammatik und
Parser zu informieren. Der "Ubersichtlichkeit halber sei an
dieser Stelle die Grammatik dargestellt:

\begin{verbatim}
type:          `char'
             | `int'
             | `float'
             | `void'
             | `[' INT_CONSTANT `]' type
             | `*' type
             | IDENTIFIER
             ;
\end{verbatim}

Diese Grammatik hat sieben Produktionen, das Startsymbol
ist \fu{type}.

\subsubsection{Top-Down und Bottom-Up Parser}
\label{topdbotu}

Im Interpreter- und Compilerbau werden haupts"achlich zwei
verschiedene Arten von Parsern verwendet: Top-Down und Bottom-Up
Parser. W"ahrend Top-Down Parser beim Startsymbol beginnen und
immer weiter hinabsteigen, bis sie auf die Token treffen, fangen
Bottom-Up Parser bei den Token an und arbeiten sich bis zum
Startsymbol hinauf. Das l"a"st sich am besten
verdeutlichen, indem man versucht, einen Parse-Tree zu
konstruieren, also einen Baum, der f"ur einen g"ultigen
Ausdruck der Sprache, in diesem Fall etwa \fu{*[10]Test}, den
Weg vom Startsymbol bis hinunter zu den einzelnen Token,
verdeutlicht. Der Parse-Tree f"ur den Beispieltyp ist in
Abbildung \ref{parsetree} dargestellt.

\begin{figure}

\centerline{\includegraphics[height=4cm]{partree1.eps}}

\caption{Der Parse-Tree f"ur \fu{*[10]Test}}
\label{parsetree}
\end{figure}

Die Wurzel des Baumes ist immer das Startsymbol, in diesem Fall
also \fu{type}. Dieser \fu{type} kann auf der ersten Ebene
zerlegt werden in das Sternsymbol und einen zweiten \fu{type},
gem"a"s der Produktion \fu{`*' type}. Dieser zweite \fu{type}
kann zelegt werden in eine "offnende und eine schlie"sende
Klammer und eine \fu{INT\_CONSTANT} dazwischen, gefolgt von noch
einem \fu{type}. Dieser wiederum ist ein \fu{IDENTIFIER}, n"amlich
\fu{Test}.

Nun zu den Vorgehensweisen der beiden Parsertypen. Ein Bottom-Up
Parser arbeitet sich wie gesagt von den einzelnen Token zum
Startsymbol hinauf. Dazu benutzt er einen Stapel, auf dem er zwei
verschiedene Aktionen durchf"uhren kann: {\em Shift} und
{\em Reduce}. Aus diesem Grund werden Bottom-Up Parser auch oft
als Shift-Reduce Parser bezeichnet. Die Analyse des
Beispieltyps f"uhrt der Parser so durch: Zun"achst liest er den
Stern, den er shiftet. Auch die beiden Klammern und die
\fu{INT\_CONSTANT} werden von ihm geshiftet. Nun liest er einen
Bezeichner und entscheidet, diesen zu einem Typen zu reducen,
gem"a"s der Produktion \fu{IDENTIFIER}. M"u"ste er einen
Parse-Tree erstellen, h"atte er an dieser Stelle einen Baum
mit der Wurzel 3\footnote{siehe Abbildung \ref{parsetree}}
erzeugt. Auf
dem Stack des Parsers liegen nun die Token \fu{`*'}, \fu{`['},
\fu{INT\_CONSTANT}, \fu{`]'} und das Nonterminal \fu{type}, in der
angegebenen Reihenfolge. An dieser Stelle entscheidet der Parser,
die letzten vier Symbole zu reducen. Um einen Parse-Tree zu
erhalten, m"u"ste er nun Baum 3 verwenden, um den gr"o"seren
Baum 2 zu konstruieren. Zuletzt w"urde der Parser nochmals
reducen und den Baum 1 erstellen. Der Parser hat nun
seine Analyse beendet, da das einzige, was auf seinem Stack
liegt, das Startsymbol ist. Die Hauptaufgabe bei der Erstellung
eines Bottom-Up Parsers ist, zu entscheiden, wann er zu shiften
und wann zu reducen hat. Dies ist eine schwere Aufgabe und wird
in der Praxis von einem Parser-Generator (auch Compiler-Compiler
genannt) wie etwa YACC erledigt. Da wir jedoch unseren Parser
selbst programmieren wollen, beschr"anken wir uns auf die
wesentlich einfacheren Top-Down Parser.

Ein Top-Down Parser geht bei der Analyse von \fu{*[10]Test}
folgenderma"sen vor: Er erkennt das Sternsymbol und wei"s nun,
da"s er nach der Produktion \fu{`*' type} vorzugehen hat. An
dieser Stelle w"urde er den Baum 1 erzeugen. Er liest nun den
Stern ein und mu"s nun einen weiteren Typ analysieren. Hier wird
klar, da"s ein Top-Down Parser in der Regel rekursiv
implementiert wird. Zwar ist es m"oglich, diese Rekursion mit
einem Stack zu beseitigen, doch ein solcher Parser ist schwerer
zu programmieren und zu warten, soda"s ich eine rekursive
Implementation gew"ahlt habe.

Der Parser ruft sich nun selbst auf und erkennt anhand der
"offnenden eckigen Klammer, da"s er die Produktion
\fu{`[' INT\_CONSTANT `]' type} w"ahlen mu"s und erzeugt
gleichzeitig Baum 2. Nacheinander liest er nun die drei folgenden
Token und ruft sich abermals selbst auf. Nun wendet er die
Produktion \fu{IDENTIFIER} an und erstellt Baum 3. Damit ist die
Analyse beendet.

\subsubsection{Look-Aheads}

Wie sich in diesem Beispiel gezeigt hat, ist es f"ur einen Parser wichtig, zu
wissen, welches Token er als n"achstes zu bearbeiten hat, denn
danach hat er zu entscheiden, welche Produktion er ausw"ahlt. Zu
diesem Zweck verwendet man sog. {\em Look-Aheads}. F"ur einen
SEL-Parser und f"ur fast alle Programmiersprachen ist ein Look-Ahead
v"ollig ausreichend, es gen"ugt also, ein Token
vorauszusehen. Die Implementation eines Look-Ahead ist simpel:
Anstatt nur des aktuellen Tokens verwendet der Parser zus"atzlich das
Look-Ahead. Soll nun das aktuelle Token "ubersprungen werden,
weist er einfach den Wert des Look-Aheads dem aktuellen Token
zu und liest einen neuen Look-Ahead aus dem Zwischencode ein.

\subsubsection{Die Konstruktion des Parsers}

Mit diesen Informationen ist es nun ein Leichtes, einen
rekursiven Top-Down Parser f"ur Typen zu bauen. Er besteht aus
einer einzigen Funktion, die durch Betrachten des Look-Ahead
entscheidet, was passieren soll. Ist der Look-Ahead eines der
Token \fu{`char'}, \fu{`int'}, \fu{`float'}, \fu{`void'} oder
\fu{IDENTIFIER}, mu"s der Parser dieses Token nur lesen und die Analyse
ist beendet. Im Fall einer "offnenden eckigen Klammer liest er
drei Token und ruft sich anschlie"send selbst auf.\footnote{Es
ist leicht m"oglich, eine am Ende stehende Rekursion durch einen
Sprung an den Beginn der Funktion, z.~B. durch eine
\rw{repeat}-Schleife, evtl. durch einen Stack unterst"utzt, zu
ersetzen, was im Normalfall die bessere L"osung ist.} Analog dazu
ist die Produktion \fu{`*' type} zu implementieren. Sollte der
Look-Ahead keinem der Token entsprechen, hat es der Parser
offensichtlich mit einem Fehler zu tun, den er entsprechend
behandeln mu"s. Ein Interpreter wird in diesem Fall eine
Fehlermeldung ausgeben und die Ausf"uhrung beenden.

Nat"urlich ist klar, da"s unser Parser nicht nur die Aufgabe hat,
einen Typ zu erkennen, sondern auch, mitzuteilen, um welchen Typ
es sich dabei handelt. Das jedoch f"allt unter die semantische
Analyse und wird in Kapitel \ref{inttypes} behandelt.

\subsection{Der Ausdruck-Parser}
\label{exprparser}

Nachdem wir nun die grundlegenden Funktionen eines Parsers
kennen, wollen wir einen etwas komplexeren Parser konstruieren:
Den Parser f"ur SEL-Ausdr"ucke. Die Grammatik hierzu kann im
Anhang \ref{grammar} oder in \cite{selkr}, Seite 14 nachgelesen
werden. Ausdr"ucke in SEL werden in der Grammatik durch die
beiden Nonterminals \fu{expr} und \fu{primary} repr"asentiert. In
den Produktionen f"ur \fu{primary} taucht noch das dritte
Nonterminal \fu{type} auf, doch dieses stellt hier kein
Problem mehr dar, da der entsprechende Parser bereits im letzten
Unterkapitel skizziert wurde.

Besonders interessant sind die Produktionen f"ur \fu{expr}, denn
sie bestehen gr"o"stenteils aus dem Nonterminal \fu{expr} selbst.
Wie sich leicht ersehen l"a"st, ist diese Grammatik auf keinen
Fall eindeutig, was die beiden Parse-Trees f"ur den Ausdruck
\fu{a=b+1} in Abbildung \ref{twopartrees}
verdeutlichen.\footnote{In einer eindeutigen Grammatik darf es
f"ur einen korrekten Satz nur einen einzigen Parse-Tree geben.}

\begin{figure}

\centerline{\includegraphics[height=3cm]{partree2.eps}}

\caption{Zwei Parse-Trees f"ur \fu{a=b+1}}
\label{twopartrees}
\end{figure}

Alle diese Konflikte jedoch lassen sich durch die in
\cite{selkr},~Seite~21 angegebene Priorit"atstabelle l"osen. Die
Frage, die sich nun stellt, ist, wie man diese Rangfolge in den
Parser einbringen kann. Betrachten wir der Einfachheit halber
nur die folgenden drei Produktionen von \fu{expr}:

\begin{verbatim}
expr:   expr `=' expr
      | expr `+' expr
      | expr `*' expr
      | primary
      ;
\end{verbatim}

Die Priorit"atstabelle sagt, da"s die h"ochste Priorit"at der
Operator \fu{*} hat, danach folgen \fu{+} und~\fu{=}. Der
Konflikt kann nun gel"ost werden, indem man f"ur jede Priorit"atsstufe ein
eigenes Nonterminal einf"uhrt, wobei die Produktionen f"ur die
Operatoren niedriger Priorit"at auf die Nonterminals der
Operatoren h"oherer Priorit"at zur"uckgreifen, wie hier
dargestellt:

\begin{verbatim}
expr3:   expr2 `=' expr2
       | expr2
       ;
expr2:   expr1 `+' expr1
       | expr1
       ;
expr1:   primary `*' primary
       | primary
       ;
\end{verbatim}

Diese Grammatik kann, angewendet auf den Ausdruck \fu{a=b+1},
nur einen Parse-Tree erzeugen, n"amlich den, der der
Klammerung des Ausdrucks \fu{a=(b+1)} entspricht.
Ungl"ucklicherweise geht mit dieser Grammatik die F"ahigkeit
verloren, einen Ausdruck wie \fu{a+b+c} darzustellen. Um auch
diesen Mangel zu beseitigen, ersetzt man in jeder Produktion mit
einem Operator die linke bzw. rechte Seite mit dem Nonterminal
des jeweiligen Levels, je nachdem, ob der Operator links- oder
rechtsassoziativ ist. Daraus resultiert folgende Grammatik:

\begin{verbatim}
expr3:   expr2 `=' expr3
       | expr2
       ;
expr2:   expr2 `+' expr1
       | expr1
       ;
expr1:   expr1 `*' primary
       | primary
       ;
\end{verbatim}

Diese Grammatik ist eindeutig und erzeugt sowohl f"ur den
Ausdruck \fu{a=b+1} als auch f"ur \fu{a+b+c} die richtigen
Parse-Trees, man kann also mit dieser Methode die "`gro"se"'
Grammatik f"ur den Bau des Anweisungs-Parsers umwandeln. Ich m"ochte die
sich daraus ergebende Grammatik in dieser Arbeit nicht
auff"uhren, da sie sehr leicht hergeleitet werden kann und zu
viel Platz beanspruchen w"urde.

Mit den bisher besprochenen Techniken kann man bereits einen
Parser f"ur einen Teil der Grammatik konstruieren, sp"atestens
bei den Linksrekursionen, wie z.~B. in den
Produktionen f"ur \fu{expr2} in der obigen Grammatik, wird man
allerdings Schwierigkeiten bekommen. Eine Rechtsrekursion wie etwa bei
\fu{expr3} ist keine Schwierigkeit: Man kann sie entweder
durch eine echte Rekursion implementieren oder durch eine
Schleife, unterst"utzt durch einen Stack. F"ur die erste Produktion
von \fu{expr2} m"u"ste man jedoch eine Funktion programmieren,
die als erstes sich selbst aufruft, was nat"urlich in einer
endlosen Rekursion m"unden w"urde.

Um dieses Problem zu l"osen, versuchen wir zu ergr"unden, welche
Menge von Ausdr"ucken eigentlich die beiden Produktionen von
\fu{expr2} beschreiben. Der einfachste Ausdruck, den die
Grammatik beschreiben kann, ist offensichtlich \fu{expr1}. Er
l"a"st sich mit Hilfe der beiden Produktionen nicht weiter
expandieren. Benutzt man als Basis nun die erste Produktion,
kann man den ersten \fu{expr2} expandieren, entweder
durch die erste oder die zweite Produktion. Benutzt man die
zweite Produktion, bekommt man als Ergebnis den Ausdruck
\fu{expr1 `+' expr1}, der sich abermals nicht weiter expandieren
l"a"st. F"uhrt man diesen Proze"s fort, wird man bald merken,
da"s diese beiden Produktionen eine Abfolge von beliebig vielen
\fu{expr1}, verbunden jeweils durch ein \fu{`+'}, beschreiben.
Zwar beschreiben die Produktionen f"ur \fu{expr3} im Grunde
dieselbe Menge, doch ist der Operator \fu{=} rechts-assoziativ,
w"ahrend \fu{+} links-assoziativ ist, was sich z.~B. beim
Erzeugen eines Parse-Trees zeigt.

Mit diesem Wissen k"onnen wir nun einen Parser konstruieren, der
einen \fu{expr2} akzeptiert, m"ussen aber wie gesagt die
Links-Assoziativit"at beachten. Um nicht unn"otig viele Worte
machen zu m"ussen, folgt ein Beispiel einer Implementation, die
einen (primitiven) Parse-Tree erzeugt. F"ur den Parse-Tree wird
folgende Struktur verwendet, die keiner Erkl"arung bedarf:

\begin{verbatim}
struct ParseTree
  Operator : int;
  Value    : int;
  Left     : *ParseTree;
  Right    : *ParseTree;
end
\end{verbatim}

Die Funktion \fu{expr2} erzeugt einen Parse-Tree f"ur einen
\fu{expr2} und greift dabei auf die Funktion \fu{expr1} zur"uck,
die dasselbe f"ur einen \fu{expr1} tut. Zus"atzlich werden die
Funktionen \fu{LookAhead} und \fu{Accept} verwendet, die ich
nicht n"aher erl"autern mu"s:

\begin{verbatim}
function expr2 () : *ParseTree

  Tree : *ParseTree;
  Temp : *ParseTree;

begin
  Tree = expr1();
  while (LookAhead() == PLUS)
    Accept(PLUS);
    Temp = new ParseTree;
    Temp->Operator = PLUS;
    Temp->Left = Tree;
    Temp->Right = expr1();
    Tree = Temp;
  end
  return Tree;
end
\end{verbatim}

Auf diese Art und Weise k"onnen Funktionen f"ur jeden der
Priorit"atsgrade erstellt werden. Im wesentlichen mit derselben
Methode werden auch die Links-Rekursionen in den Produktionen von
\fu{primary} beseitigt. Eine Kleinigkeit in der Implementation
von \fu{primary} m"ochte ich noch besprechen: Die Produktion
\fu{IDENTIFIER `(' [expr \{`,' expr\}] `)'} enth"alt ein
optionales Element und darin eingeschlossen ein Element, da"s
sich beliebig oft wiederholen kann. Es ist zwar m"oglich, solche
Konstruktionen durch Umformung der Grammatik mit den oben
besprochenen Methoden zu analysieren, doch die direkte
Implementation ist nicht schwer zu realisieren und wird deshalb
von mir vorgezogen.

Die beiden eckigen Klammern in der Produktion bedeuten, da"s auf
die "offnende runde Klammer, die nach dem \fu{IDENTIFIER} zu
stehen hat, entweder eine schlie"sende runde Klammer oder ein
\fu{expr} folgt. Welche der beiden M"oglichkeiten zutrifft,
kann der Parser dadurch erkennen, da"s er den Look-Ahead
betrachtet.
Da ein \fu{expr} niemals mit einer schlie"senden Klammer beginnen
kann, ist klar, da"s, wenn der Look-Ahead die schlie"sende
Klammer ist, die erste Variante zutrifft, andernfalls die zweite.
In letzterem Fall mu"s die Funktion \fu{expr} aufgerufen werden, um
den ersten \fu{expr} zu erkennen. Nun steht der Parser vor einem
"ahnlichen Problem. Entweder, es folgt ein Beistrich und noch ein
\fu{expr}, dann vielleicht noch ein Beistrich u.~s.~w., oder es
folgt die schlie"sende Klammer. Wie oben kann auch dieses
Problem mit dem Look-Ahead gel"ost werden. Entspricht er einem Beistrich,
mu"s noch ein \fu{expr} analysiert werden, wonach der Parser wieder
auf einen Beistrich testen mu"s. Ist der Look-Ahead die
schlie"sende Klammer, mu"s er diese nur noch lesen und ist
mit der Analyse des \fu{primary} fertig. Wenn jedoch der
Look-Ahead keinem der beiden Symbole entspricht, hat er
offensichtlich einen Fehler gefunden, der dementsprechend zu
behandeln ist.

\subsection{Der Anweisungsparser}

Der Parser f"ur Anweisungen ist relativ einfach zu realisieren,
da er weder Linksrekursionen enth"alt, noch irgendwelche
Priorit"atsregeln zu befolgen sind. Die Anweisungen von SEL sind
so konzipiert, da"s der Parser schon anhand des ersten Tokens
entscheiden kann, um welche Anweisung es sich handelt. Er
unterscheidet sich jedoch insofern von den bis jetzt besprochenen
Parsern, als da"s bei bestimmten Anweisungen nur Teile dieser
Anweisung geparst werden m"ussen, wie z.~B. bei \rw{if}, wo der Wert der
Bedingung bestimmt, ob der \rw{if}- oder der \rw{else}-Zweig,
wenn vorhanden, ausgef"uhrt werden soll.

Betrachten wir als Beispiel die zitierte \rw{if}-Anweisung. Der
Parser erkennt anhand des Look-Ahead, da"s es sich um eine
\rw{if}-Anweisung handelt. Er "uberliest also die beiden Token
\rw{if} und \fu{(} und ruft anschlie"send den im vorigen
Unterkapitel besprochenen Ausdruck-Parser auf, um den \fu{expr}
zu lesen. Wie sich in Kapitel \ref{semantics} zeigen wird,
berechnet dieser Parser den Wert des Ausdrucks, der nun
benutzt werden kann, um zu entscheiden, ob die dem \rw{if} folgenden
Anweisungen ausgef"uhrt werden sollen. Ist der Wert ungleich \fu{0}, ist
die Sache leicht. In diesem Fall mu"s sich der Anweisungs-Parser
nur selbst aufrufen, um die folgenden Anweisungen bis zum
\rw{end} bzw. \rw{else} zu parsen. Andernfalls benutzt er die
zus"atzliche Information, die ihm der in Kapitel \ref{tokencode}
besprochene Zwischencode zum \rw{if}-Token liefert. Sie gibt die
Position an, an der das zugeh"orige \rw{end} bzw. \rw{else}
steht. So m"ussen nicht langwierig alle folgenden Anweisungen bis
zum jeweiligen Token durchsucht werden, sondern es kann direkt an
die entsprechende Position gesprungen werden. Steht an dieser
Stelle ein \rw{end}, ist die Analyse der \rw{if}-Anweisung
beendet, andernfalls mu"s der \rw{else}-Zweig ausgef"uhrt werden.

Noch etwas komplizierter verh"alt es sich mit der
\rw{switch}-Anweisung. Dort mu"s das \rw{switch}-Token die
Position des ersten zugeh"origen \rw{case} enthalten, dieses
wiederum die des n"achsten \rw{case} bzw. \rw{default}
u.~s.~w. Auf diese Weise ist es m"oglich, einen Zweig nach
dem anderen auf "Ubereinstimmung zu untersuchen und, falls eine
"ubereinstimmende Konstante gefunden wurde, diesen Zweig
auszuf"uhren, indem die Funktion sich selbst aufruft. Damit die
\rw{switch}-Anweisung auch wie vorgeschrieben ausgef"uhrt wird,
mu"s der Anweisungs-Parser beim normalen Parsen ein
\rw{case}-Schl"usselwort einfach "uberspringen und bei einem
\rw{break} zur "ubergeordneten Funktion zur"uckkehren.

Ein Problemfall ist die \rw{return}-Anweisung. Wenn der
Parser auf sie trifft, mu"s er auch zur "ubergeordneten Funktion
zur"uckkehren, doch er mu"s das im Gegensatz zur
\rw{break}-Anweisung so lange tun, bis der Anweisungs-Parser
komplett verlassen wurde. Wir m"ussen uns also "uberlegen, wie
wir es anstellen k"onnten, da"s der Parser wei"s, wann die
untergeordnete Instanz ein \rw{return} entdeckt hat. Die
einfachste L"osung ist, das mit dem R"uckgabewert zu tun. Im
n"achsten Kapitel wird sich zeigen, da"s SEL-Werte vom
Interpreter in Strukturen gehalten werden. Am sinnvollsten ist
es, als R"uckgabetyp f"ur den Parser einen Zeiger auf eine solche
Struktur zu benutzen. Trifft er nun auf ein \rw{break}, gibt er
einen Nullzeiger zur"uck und der "ubergeordnete Parser wei"s, das
alles in Ordunung ist. Wurde aber eine \rw{return}-Anweisung
entdeckt, wird ein Zeiger auf den zugeh"origen Wert
zur"uckgegeben. Die "ubergeordnete Funktion wei"s nun, da"s auch
sie beendet werden mu"s und gibt ihrerseits den Wert weiter. Wird
eine Funktion geparst, die als R"uckgabetyp \tvoid\ hat, mu"s
ein Zeiger auf einen Dummy-Wert zur"uckgegeben werden.

Der Anweisungs-Parser mu"s auch mit einem Null-Zeiger verlassen
werden, wenn er auf ein \rw{end} oder \rw{else} trifft, die zu
keiner von ihm behandelten Anweisungen geh"ort, denn dann f"allt
es in den Zust"andigkeitsbereich der "ubergeordneten
Rekursionsebene.

Hier wird auch klar, warum ich auf eine
\rw{goto}-Anweisung\footnote{siehe Kapitel \ref{goto}} verzichtet
habe: In einem rekursiven Parser w"are es f"ur einen Sprung
n"otig, zun"achst zu ermitteln, auf welcher Rekursionsebene das
Label liegt, und ob es vielleicht innerhalb einer anderen Anweisung
liegt, was einen besonders schlimmen Fall darstellt. Betrachten wir
dazu folgendes (theoretische) Programmfragment:

\begin{verbatim}
if (x == 5)
  goto label;
end
if (y == 3)
  do_something();
  if (z == 4)
    label:
    do_something_else();
  end
end
\end{verbatim}

In diesem Fall m"u"ste der Interpreter, wenn er die
\rw{goto}-Anweisung ausf"uhrt, zun"achst eine Rekursionsebene
verlassen und dann in zwei andere einsteigen, was nat"urlich sehr
schwer zu realisieren ist. "Uberdies m"u"ste der Parser sehr
schnell herausfinden k"onnen, wieviel Ebenen zu verlassen und in
welche er einzusteigen hat. Zu diesem Problem ist mir bis jetzt
noch keine zufriedenstellende L"osung eingefallen.

\section{Die Verwaltung}
\label{semantics}

In den letzten Unterkapiteln wurde behandelt, wie der Interpreter
den Quelltext (genauer gesagt, den Zwischencode) lexikalisch
analysiert, in Zwischencode verwandelt und diesen syntaktisch
analysiert, jedoch wurde noch nicht gesagt, wie der Interpreter
die durch diese Analyse gewonnenen Informationen dazu verwenden
kann, das Programm auszuf"uhren. Das ist die Aufgabe dieses
Unterkapitels. Hier wird besprochen, was der Interpreter bei der
Ausf"uhrung des Programms tut, und wie dies implementiert werden
kann.

Die meisten der in diesem Unterkapitel besprochenen Algorithmen
werden w"ahrend der Syntaxanalyse, also des im letzten
Unterkapitel besprochenen Parsing, ausgef"uhrt, was auch der
Definition des Interpreters\footnote{siehe Kapitel
\ref{whatisit}} entspricht.
        
\subsection{Module}
\label{modules}

Da SEL es erlaubt, modulglobale Variablen und Funktionen zu
benutzen, ist es vonn"oten, ein SEL-Programm auch im Speicher in
Module aufzuteilen. Dazu m"ussen f"ur jedes Modul mindestens der
Zwischencode und die modulglobalen Variablen und Funktionen
gespeichert werden. Im n"achsten Unterkapitel werden sich dazu
noch einige zus"atzliche Informationen gesellen, die bei der
Fehlerbehandlung helfen.

Die programmglobalen Funktionen und Variablen werden au"serhalb
der Module f"ur das gesamte Programm gemeinsam verwaltet, was
wesentlich effizienter und auch einsichtiger ist (und es nebenbei
einfacher und schneller macht, Namenskonflike zu bemerken).
Weiters gibt es zwei Variablen, die f"ur das gesamte Programm
gelten. Sie geben die Position an, an der sich der Parser
gerade mit der Ausf"uhrung aufh"alt. Die erste dieser Variablen
nenne ich \fu{MP} (Module Pointer). Sie enth"alt einen Wert, der
das Modul, in dem sich die gerade abgearbeiteten Anweisungen
befinden, eindeutig identifiziert. Am effizientesten ist es, ein
Array aus Modulbeschreibungen anzulegen, das durch diese
Variable indiziert wird. Die zweite Variable, \fu{IP}
(Instruction Pointer), gibt die Position innerhalb des
Zwischencodes des aktuellen Moduls an.\footnote{Genauerweise
m"u"ste diese Variable den Namen \fu{TP} f"ur Token Pointer
tragen, da sie sich nicht auf die aktuelle Anweisung, sondern auf
das aktuelle Token bezieht.}

\bigskip

In Kapitel \ref{intgoals} wurde als Ziel f"ur den Interpreter
festgelegt, da"s er Multitasking unterst"utzt. Um die dazu
notwendigen Schritte anzugeben, mu"s erst gekl"art werden, wie
Multitasking funktioniert. Zun"achst m"ochte ich feststellen, da"s
ich mit Multitasking nicht die F"ahigkeit eines Betriebssystems
meine, mehrere Programme gleichzeitig zu bearbeiten, sondern die
F"ahigkeit, innerhalb eines Programms mehrere Funktionen zur
gleichen Zeit auszuf"uhren. Unter den Betriebssystemen OS/2 und
Windows NT wird daf"ur der Begriff "`Multithreading"' gepr"agt.
Im Fall eines SEL-Interpreters l"auft dieses Multitasking darauf
hinaus, da"s durch diesen Interpreter (der ja ein Programm
darstellt) mehrere SEL-Programme gleichzeitig ausgef"uhrt werden
k"onnen.

Zwei Funktionen, die gleichzeitig ausgef"uhrt werden, k"onnen
zwar auf dieselben globalen Variablen zugreifen, soda"s eine
gemeinsame Benutzung eines Datenpools, etwa einer
Funktionssammlung f"ur ein Mathematikprogramm, gew"ahrleistet
ist, jedoch besitzt jede der beiden Funktionen einen eigenen
Stack, was garantiert, da"s sie ungest"ort von den anderen
Funktionen ihre eigenen lokalen Variablen verwenden kann.

Wenn nun der Interpreter zwei verschiedene SEL-Programme
gleichzeitig ausf"uhren soll, tut er das, indem er zwei Tasks
startet, die jeweils ein Programm bearbeiten. Um zu
gew"ahrleisten, da"s sich diese beiden Funktionen dabei nicht in
die Quere kommen, d"urfen die Daten, die das aktuelle Programm
betreffen, nat"urlich nicht in globalen Variablen gespeichert
werden, sondern m"ussen der Funktion, die f"ur die Ausf"uhrung
zust"andig ist, als Parameter mitgegeben werden. F"ur unseren
rekursiven Parser bedeutet das, da"s diese Daten sehr oft als
Parameter "ubergeben werden, was nicht nur den Stack belastet,
sondern auch die Ausf"uhrungsgeschwindigkeit. Um diese Verluste
gering zu halten, empfielt es sich, nicht direkt die Daten,
sondern Zeiger auf sie zu verwenden, was aber in den meisten
Programmiersprachen ohnedies die einzige M"oglichkeit ist.

\subsection{Typen}
\label{inttypes}

Eine f"ur die Ausf"uhrungsgeschwindigkeit sehr wichtige
Entscheidung bei einem Interpreter f"ur eine Programmiersprache,
die ein im Vergleich zu anderen Interpretersprachen so komplexes
Typkonzept verwendet, wie SEL, ist, wie diese Typen
repr"asentiert werden sollen. In Compilern werden hierf"ur
h"aufig verkettete Listen verwendet, die jedoch f"ur einen
Interpreter nicht sehr brauchbar sind, da sie wegen ihres
Aufbaus, der auf dynamischer Speicherverwaltung basiert, zu
langsam sind und au"serdem zu viel Speicher verbrauchen.

Das von mir verwendete Konzept basiert darauf, da"s Typen durch
Byte-Arrays mit begrenzter L"ange repr"asentiert werden. Da es
in SEL nur eine geringe Zahl an Grundtypen und auch nicht viele
M"oglichkeiten, diese zu verkn"upfen, gibt, kann ein Byte
jeweils eine Zusammensetzungsstufe oder einen Grundtypen
enthalten. Ein Beispiel: Der Typ \type{\poi \arr{10}int} wird
durch ein Array mit der Gr"o"se von 5 Bytes repr"asentiert,
dessen erstes Byte angibt, da"s es sich um einen Zeiger handelt.
Das zweite Element gibt an, da"s ein Array im Spiel ist. F"ur
dieses Array werden auch die n"achsten zwei Bytes ben"otigt, die
die Gr"o"se angeben. Das f"unfte Element schlie"slich steht f"ur
den Grundtypen \tint\ und ist gleichzeitig die Ende-Markierung,
da es sich dabei um einen Grundtypen handelt. In gewisser Weise
"ahnlich wie \tint\ werden Strukturen behandelt, die ja auch
immer am Ende des Typs stehen. Bei ihnen mu"s, "ahnlich dem
Array, ein Wert angegeben werden, der den Typ der Struktur
bestimmt. Zu diesem Zweck vergibt der Interpreter f"ur die Strukturen
Nummern, wobei die Werte von 0 bis 9999 f"ur programmglobale
Strukturen reserviert sind. Die weiteren Zuordungen richten sich
danach, wie viel Strukturen in den einzelnen Modulen definiert
sind.

Wie schon oben erw"ahnt, besitzen die Typenarrays eine begrenzte
L"ange, was es erm"oglicht, sie schon vor ihrer Verwendung
bereitzustellen, soda"s sie nicht w"ahrend der Ausf"uhrung des
SEL-Programms dynamisch allokiert werden m"ussen. In meiner
Implementation habe ich mich f"ur eine L"ange von 20 Bytes
entschieden. Das mag nicht viel erscheinen, vor allem in
Anbetracht der Tatsache, da"s es in SEL m"oglich ist, Typen
beliebig zu verschachteln, doch ist es m"oglich, mit diesen 20
Bytes einen Zeiger auf einen Zeiger auf einen Zeiger auf ein Array
aus Arrays aus Strukturen zu repr"asentieren, wobei noch 8 Byte
frei bleiben. Dieses Beispiel soll zeigen, da"s 20 Bytes f"ur die
Praxis eine absolut ausreichende Gr"o"se ist.

Beim Analysieren von Ausdr"ucken ist es in einem SEL-Interpreter
sehr oft n"otig, Typen bereitzustellen, von denen aber viele
schon im n"achsten Moment wieder freigegeben werden. Ist es nicht
m"oglich, diese Operationen effizient auszuf"uhren, leidet
darunter die Geschwindigkeit des Interpreters sehr. Ich habe zu
diesem Zweck einen Algorithmus entwickelt, der es erm"oglicht, in
konstanter Zeit, also unabh"angig davon, wieviele von den
bereitgestellten Typenarrays bereits vergeben sind, Typenarrays
zu allokieren und freizugeben:

\bigskip

Sei $n$ die Anzahl der maximal bereitzustellenden Typen und $s$
die Gr"o"se eines Typenarrays, dann wird vor der Ausf"uhrung des
SEL-Programms ein Byte-Array mit der Gr"o"se $n \cdot s$ erzeugt,
wobei $t$ die Speicheradresse sei, an der das Array beginnt.
Weiters sei $r$ die Speicheradresse, an der ein Array der Gr"o"se
$n$ beginnt, dessen Elemente vern"unftigerweise Maschinenworte
sind. Dieses Array wird initialisiert, soda"s $r[x] = x$.
Zus"atzlich sei $i$ eine numerische Variable, die mit
dem Wert $0$ initialisiert werden.

Soll ein Typenarray allokiert werden, wird zun"achst "uberpr"uft,
ob $i = n$. In diesem Fall sind bereits alle Arrays vergeben und
eine entsprechende Fehlermeldung wird ausgegeben. Andernfalls ist
$t + r[i] \cdot s$ die Adresse, bei der das angeforderte Array
beginnt. Zus"atzlich ist $i = i + 1$ zu setzen.

Soll das Typenarray mit der Speicheradresse $f$ freigegeben
werden, so ist zuerst $i = i - 1$ und dann $r[i] = (f - t) / s$
zu setzen.

Der angegebene Algorithmus verl"a"st sich darauf, da"s nicht
Arrays freigegeben werden sollen, die nicht allokiert sind, und
da"s jeweils maximal $s$ Bytes eines Arrays genutzt werden.

\bigskip

Dieser Algorithmus verwendet zwei Datenstrukturen. Das Array $t$
ist sozusagen ein Zusammenschlu"s aller Arrays, die zur
Verf"ugung gestellt werden k"onnen, was sich nat"urlich auch
mittels eines zweidimensionalen Arrays realisieren lie"se. Das
Array $r$ ist ein Puffer, der die Nummern (0 bis $n - 1$)
aller noch freien Typenarrays enth"alt. Die Variable $i$ gibt den
Index des n"achsten zu vergebenden Typenarrays an. Wird dieses
angefordert, wird $i$ um 1 erh"oht. Umgekehrt wird $i$ um 1
vermindert, wenn ein Array freizugeben ist. Dessen Nummer wird
anschlie"send in $r[i]$ gespeichert, soda"s es bei der n"achsten
Anforderung sofort wieder vergeben wird. Wie schon gesagt, nimmt
der Algorithmus an, da"s keine Fehler vom Programmierer begangen
werden, wie z~B. die Freigabe eines bereits freigegebenen Arrays.

In meiner Implementation von SEL habe ich f"ur $n$ den Wert 400
benutzt, f"ur $s$ wie gesagt 20. Diese Werte sind f"ur die
meisten Anwendungen ausreichend, und wenn n"otig, k"onnten sie
vom Benutzer entsprechend ge"andert werden. Die
Speicheranforderungen dieses Algorithmus sind durchaus
vertretbar. Bei den angegebenen Werten erw"achst f"ur die beiden
Arrays unter Windows NT ein Speicherbedarf von genau 9600 Byte,
was bei den heute "ublichen Speicherkapazit"aten kaum ein Problem
darstellt.

\subsection{Ausdr"ucke}
\label{intexprs}

"Ahnlich wichtig wie die Typen sind in SEL die Ausdr"ucke, die
nat"urlich auch verwaltet werden m"ussen. Dabei werden wir jedoch
nicht ganze Ausdr"ucke verwenden, sondern uns lediglich auf die
Werte der Ausdr"ucke beschr"anken, die im Laufe des
Parsevorganges errechnet werden.

Um zu zeigen, wie wichtig die effiziente Verwaltung von
Ausdr"ucken ist, hier ein konkretes Beispiel: Ein
QuickSort-Programm f"ur 500 zuf"allig angeordnete Zahlen hat auf
meinem Computer eine Laufzeit von etwa 25 Sekunden, wobei w"ahrend
dieser Zeit ungef"ahr 60000 Ausdr"ucke allokiert und wieder
freigegeben werden m"ussen.

Zum Allokieren und Freigeben von Ausdr"ucken verwende ich aus
diesem Grund denselben Algorithmus wie f"ur Typen, wobei ich von
Haus aus maximal 200 Ausdr"ucke zur Verf"ugung stelle, was in den
meisten F"allen ausreichend ist.

Sehr wichtig f"ur eine effiziente Verarbeitung ist nat"urlich
auch der Aufbau eines Ausdrucks. Ich verwende hierzu eine
Struktur. Ein obligater Bestandteil eines Ausdrucks ist dessen
Typ, soda"s ein Typenarray Teil dieser Struktur ist. Weiters gibt
es je eine Variable f"ur die Basistypen \tchar , \tint\ und
\tfloat . Hierzu wird es aus Effizienzgr"unden am besten sein,
passende Typen der Implementationssprache zu verwenden. F"ur den
Fall, da"s der Ausdruck ein Array, ein Zeiger oder eine Struktur
ist, habe ich zus"atzlich einen typenlosen Zeiger aufgenommen,
der, falls die Implementationssprache keine typenlosen Zeiger
zul"a"st, durch einen Zeiger auf ein Byte ersetzt werden kann.
Im Falle eines Funktionsaufrufs ist es n"otig, mehrere Ausdr"ucke
zu verketten, um die Parameter der Funktion anzugeben, soda"s
zus"atzlich ein Zeiger auf einen Ausdruck in die Struktur
eingebaut werden sollte, soda"s sie sich zu einer verketteten
Liste erweitern l"a"st.

Soviel zur Verwaltung der Ausdr"ucke im Speicher. Eine weit
wichtigere Frage jedoch ist, was der Interpreter mit diesen
Ausdr"ucken tut. Nun, zun"achst einmal mu"s jede der in
Kapitel \ref{exprparser} besprochenen Produktionsfunktionen so
konzipiert werden, da"s sie der "ubergeordneten Funktion einen Ausdruck
mit dem Wert des von ihr geparsten Ausdrucks zur"uckgibt. Diese
verwendet ihn ihrerseits, um den Wert ihres Ausdrucks zu
berechnen u.~s.~w. In der Kette der Funktionsaufrufe gibt es dann
irgendwann eine Funktion, die den Wert des Ausdrucks benutzt, um
damit etwas zu entscheiden, wie z.~B., ob der Zweig einer
\rw{if}-Anweisung ausgef"uhrt werden soll.

Solange man lediglich Berechnungen mit den Ausdr"ucken
anstellt, ergeben sich mit der oben skizzierten Methode
keinerlei Probleme, wie ist es jedoch m"oglich, den
Zuweisungsoperator zu implementieren? Bauen wir die Funktion zum
Parsen eines \fu{primary} so auf, da"s sie f"ur eine Variable
lediglich deren Wert zur"uckgibt, werden wir Probleme haben, denn
dann wei"s der Interpreter nicht, wo diese Variable im Speicher liegt und
kann ihr so auch keinen Wert zuweisen. Um dieses Problem zu
"uberwinden, wird intern ein zus"atzlicher Typ,
n"amlich der Referenztyp, der einem Zeiger sehr "ahnlich ist,
verwendet.
Trifft die Funktion zum Parsen eines \fu{primary} nun im Zwischencode auf eine Variable des Typs
\tint , dann gibt sie nicht deren Wert zur"uck, sondern ihre
Adresse und vermerkt, da"s es sich dabei um eine Referenz auf
einen \tint\ handelt, der Einfachheit halber von mir mit
\type{\reference int} bezeichnet.
Mu"s nun mit diesem Ausdruck gerechnet werden, gen"ugt es, ihn zu
dereferenzieren, soda"s ein Ausdruck mit dem Typ \tint\ entsteht.
Wird ihm ein Wert zugewiesen, mu"s dieser lediglich an die
Speicheradresse geschrieben werden, die er angibt. Soll der
un"are Operator \fu{\&} auf ihn angewendet werden, dann ist es
lediglich n"otig, das erste Byte im Typenarray so zu "andern,
da"s daraus ein Zeiger wird.
                
\subsection{Variablen}
\label{maintvars}

Variablen sind von den oben besprochenen Ausdr"ucken und Typen
insofern verschieden, als da"s sie nicht intern vom Interpreter
verwendet werden, sondern direkt mit dem Quellprogramm in
Verbindung stehen. Aus dem Modulkonzept und der Tatsache, da"s
Funktionsvariablen nur in der jeweiligen Funktion verf"ugbar
sind, folgt, da"s die Verwaltung von Variablen auf mehrern Stufen
erfolgen mu"s. Die umfassendste Stufe, allerdings auch diejenige
mit der niedrigsten Priorit"at, ist die der programmglobalen
Variablen. Auf sie folgt die Stufe der modulglobalen Variablen.
Die h"ochste Priorit"at besitzt die Stufe der Funktionsvariablen,
die allerdings in der Regel auch die wenigsten Variablen
aufnimmt.

Wird nun im Programm eine Variable verwendet, mu"s zun"achst
gepr"uft werden, ob es eine Funktionsvariable ihres Namens gibt.
Ist das nicht der Fall, kann es sich nur mehr um eine globale
Variable handeln. Gibt es eine modulglobale Variable, die den
gesuchten Namen hat, wird sie verwendet, andernfalls mu"s es eine
programmglobale Variable geben. Ist auch das nicht zutreffend,
hat es der Interpreter mit einem Programmfehler zu tun (siehe Kapitel
\ref{errortrapping}).

Die komplizierteste Verwaltungsstufe ist die der
Funktionsvariablen (es ist am einfachsten, bei der Verwaltung der
Variablen nicht zwischen Parametern und direkten Funktionsvariablen
zu unterscheiden), was daran liegt, da"s nicht alle
Funktionsvariablen in einen "`Topf"' geworfen werden k"onnen,
sondern da"s f"ur jede Funktion eine eigene Variablenliste
angelegt werden mu"s. Die einfachste L"osung f"ur dieses Problem
ist, einen Stack zu verwenden, auf dem f"ur jeden
Funktionsaufruf eine Liste aller Funktionsvariablen angelegt
wird, "ahnlich dem Code, den ein Compiler erzeugen w"urde. Analog
dazu werden bei der R"uckkehr aus einer Funktion die vorher
erzeugten Variablen wieder entfernt. Um nun
nach einer Funktionsvariablen zu suchen, mu"s nur der oberste
Stacklevel in Betracht gezogen werden, da er alle Variablen
beinhaltet, die zu der gerade auszuf"uhrenden Funktion geh"oren.

Wie schon in Kapitel \ref{modules} angesprochen, stellt der
Interpreter f"ur jedes Modul eine eigene Datenstruktur zur
Verf"ugung, die unter anderem alle modulglobalen Variablen dieses
Moduls verwaltet. Ist also nach einer modulglobalen Variablen zu
suchen, kann man die Suche auf das durch \fu{MP} angegebene Modul
beschr"anken.

Der einfachste Fall liegt bei programmglobalen Variablen vor.
F"ur ihre Verwaltung gen"ugt es, eine einzige Liste
bereitzustellen, die die programmglobalen Variablen aller Module
bereitstellt, ohne Unterscheidung, in welchem Modul diese
definiert sind.

\bigskip

Bis jetzt habe ich, wenn es um die Details der
Variablenverwaltung ging, lediglich von "`Listen"' gesprochen,
wobei ich damit weniger dynamische Listen, als vielmehr
geeignete Datenstrukturen f"ur das schnelle Auffinden von
Variablen gemeint habe. Diese Datenstrukturen m"ochte ich jetzt
genauer besprechen.

Das Problem, auf das wir hier treffen ist, eine
Datenstruktur (und nat"urlich auch die dazugeh"origen
Algorithmen) aussuchen zu m"ussen, die es erm"oglicht, einen
Datensatz in einer Menge von Datens"atzen m"oglichst
schnell zu finden. Dieses Problem ist ein gutes Beispiel f"ur das
Grundproblem eines Programmdesigners: F"ur eine bestimmte Aufgabe
den besten Algorithmus bzw. die beste Datenstruktur (was in den
meisten F"allen Hand in Hand geht) auszuw"ahlen. G"abe es eine
absolut beste L"osung, dann h"atte der Designer seine Aufgabe
schnell erledigt, doch leider gibt es eine solche nur in den
seltensten F"allen.

Das wichtigste Entscheidungskriterium f"ur unsere Datenstruktur
ist, wieviel Datens"atze zu verwalten sind. F"ur die
Variablenverwaltung eines Interpreters bieten sich grunds"atzlich
zwei Strukturen an: Bin"are Suchb"aume und Hash-Tabellen. Um
entscheiden zu k"onnen, welche der beiden Strukturen f"ur unsere
Anwendungen die bessere ist, ist es n"otig, ein wenig zu rechnen.
Der Einfachheit halber gehen wir vom g"unstigsten Fall aus (die
Ergebnisse sind von denen des durchschnittlichen Falls nicht sehr
verschieden). Dieser ist, da"s der Benutzer die Variablennamen so
gew"ahlt hat, da"s sich eine gleichm"a"sige Anordung in der
Hash-Tabelle ergibt bzw. der Suchbaum ausgeglichen ist.

In einer Hash-Tabelle, bei der die Eintr"age f"ur gleiche
Hash-Werte linear verkettet sind, werden f"ur eine Suche nach
einem Datensatz im Mittel $\max (N / 2M + 1/2, 1)$ Vergleiche
ben"otigt, wohingegen f"ur die Suche in einem Baum etwa
$\log_2N - 1$ Vergleiche notwendig sind, wobei $N$ die Anzahl der
Datens"atze und $M$ die Gr"o"se der Hash-Tabelle angibt. Auf den
ersten Blick scheint der Baum die bessere L"osung zu sein, da die
Anzahl der ben"otigten Vergleiche f"ur gro"se $N$ in etwa
proportional zum bin"aren Logarithmus der Anzahl der Datens"atze
ist, bei einer Hash-Tabelle allerdings die Anzahl der Vergleiche
direkt proportional zur Anzahl der Datens"atze ist, was f"ur
gro"se $N$ nat"urlich bedeutend schlechter ist, als die
Leistungsdaten des Baums. Nehmen wir allerdings an, wir benutzen
eine Hash-Tabelle mit 37 Eintr"agen, was in Bezug auf den
Speicherplatz "uberhaupt kein Problem darstellt. In diesem Fall
sind f"ur eine Suche in einer Hash-Tabelle mit 100 Datens"atzen
im Durchschnitt weniger als zwei Vergleiche n"otig, wohingegen in
einem Baum mehr als sechs Vergleiche ben"otigt werden. Erst wenn
mehr als 550 Datens"atze vorhanden sind, ist die Suche in einem
Baum g"unstiger als in einer Hash-Tabelle.

Es ist kaum anzunehmen, da"s jemand f"ur einen Interpreter ein
Programm schreibt, welches mehr als 300 globale Variablen verwendet,
soda"s es f"ur unsere Zwecke die bei weitem bessere L"osung ist,
f"ur die Verwaltung programm- und modulglobaler Variablen
Hash-Tabellen zu verwenden, wobei eine Gr"o"se von 37 durchaus
praktikabel ist, sowohl in Bezug auf Schnelligkeit als auch auf
die vern"unftige Ausnutzung des Speicherplatzes. F"ur
Funktionsvariablen sind 37 Eintr"age wohl "uberdimensioniert.
Hierf"ur bieten sich die Primzahlen 7 und 11 an.

Zuletzt stellt sich noch die Frage nach der Hash-Funktion. Ich
verwende sowohl der Einfachheit als auch der Effizienz halber die
Modulo-Funktion, d. h., ich betrachte eine Zeichenkette als lange
Ganzzahl, deren Rest bez"uglich des Moduls 37 den Index in der
Hash-Tabelle angibt. Die folgende Implementation geht davon aus, 
da"s ein \tchar\ 256 verschiedene Werte annehmen kann.

\begin{verbatim}
function hash (s : *char, m : int) : int

  h : int = 0;

begin
  while (*s)
    h = (h * 256 + *s) % m;
    s = s + 1;
  end
  return h;
end
\end{verbatim}

\subsection{Funktionen}

Die Funktionsverwaltung hat in erster Linie die Aufgabe, eine
Funktion eines bestimmten Namens zu finden und deren Adresse im
Zwischencode zu ermitteln. Weiters ist es wichtig, Namen und
Typen der Funktionsparameter zu erfahren, um erstens "uberpr"ufen
zu k"onnen, ob die Funktion mit der richtigen Anzahl und den
richtigen Typen der Parameter aufgerufen worden ist, und zweitens,
um die Parameter auf dem Funktionsstack anlegen zu
k"onnen.\footnote{siehe Kapitel \ref{maintvars}}

Zus"atzlich sollte es m"oglich sein, nicht nur Funktionen, die in SEL
implementiert sind, aufzurufen, sondern auch solche, die in der
Implementationssprache des Interpreters geschrieben sind,
da eine Implementation der Grundfunktionen von SEL in
SEL selbst f"ur einen Interpreter viel zu ineffizient w"are.
Deshalb schreibt man Funktionen wie etwa zum Kopieren von
Strings in der Implementationssprache des Interpreters (in
meinem Falle C) und ruft diese auf. Um das nahtlos in die
Verwaltung einbinden zu k"onnen, ist es notwendig, da"s die
Implementationssprache Zeiger auf Funktionen unterst"utzt. Ist
diese M"oglichkeit nicht vorhanden (wie etwa im Falle von
\pascal ), dann mu"s jede einzelne Funktion im Programmcode extra
behandelt werden, was sowohl umst"andlicher als auch langsamer
ist.\footnote{In Turbo Pascal ab der Version 5.5 ergibt sich die
M"oglichkeit, statt eines Zeigers auf eine Funktion einen Zeiger
auf ein Objekt zu verwenden, welches eine Methode enth"alt, die
die gew"unschte Funktion ersetzt.}

Um eine Funktion eines bestimmten Namens zu finden, werden
Hash-Tabellen, wie im letzten Unterkapitel beschrieben,
verwendet: Eine f"ur programmglobale Funktionen und eine
pro Modul f"ur modulglobale Funktionen, der Einfachheit halber
jeweils mit 37 Eintr"agen, was bei einem Programm mit 10 Modulen
unter Windows NT einen Speicherplatz von 1628 Bytes beansprucht
(was ohne schlechtes Gewissen als "`unbedenklich"' eingestuft
werden kann). Eintr"age f"ur denselben Hash-Wert werden
unsortiert linear verkettet.\footnote{Eine Sortierung h"atte bei
einer linearen Verkettung auch nicht den geringsten Sinn, da
ohnedies alle Eintr"age bis zum gew"unschten durchlaufen werden
m"ussen, unabh"angig davon, in welcher Reihenfolge die Eintr"age
gespeichert sind.}

Die Datenstruktur, die eine einzelne Funktion beschreibt, mu"s
neben deren Namen und R"uckgabetyp auch eine Angabe dar"uber
enthalten, wo der Code der Funktion zu finden ist. Hierf"ur wird
zun"achst eine Variable ben"otigt, die Auskunft dar"uber gibt, ob
die Funktion in SEL geschrieben ist oder in der
Implementationssprache des Interpreters. Im ersten Fall mu"s die
Datenstruktur Auskunft dar"uber geben k"onnen, in welchem Modul
und an welcher Stelle der Code beginnt, andernfalls wird ein
Zeiger auf die jeweilige Funktion ben"otigt. Zus"atzlich m"ussen
Informationen "uber Namen und Typen der Parameter gespeichert
sein. Ich verwende hierzu eine einfach verkettete Liste.

\bigskip

Interessant ist die Realisierung von Funktionsaufrufen. Dazu ist
zun"achst zu bedenken, da"s der Parser rekursiv implementiert ist
und der Funktionsaufruf ein \fu{primary} ist, sich also auf der
untersten Grammatikebene der Ausdr"ucke befindet. Aus diesem
Grund wird man schwer umhinkommen, auch Funktionsaufrufe rekursiv
zu behandeln, was nat"urlich bei einem rekursiven SEL-Programm zu
einer starken Belastung des Prozessorstacks f"uhren kann, weshalb
dieser m"oglichst gro"s dimensioniert werden sollte. Mu"s nun
innerhalb eines \fu{primary} eine Funktion aufgerufen werden,
wird in den oben besprochenen Hash-Tabellen die Funktion gesucht
und, wenn vorhanden, wird an die angegebene Stelle im
Zwischencode gesprungen. Dort werden zun"achst der Funktionskopf
und die Deklaration der direkten Funktionsvariablen interpretiert
und anschlie"send der Code der Funktion mittels des
Anweisungs-Parsers ausgef"uhrt. Ist die Ausf"uhrung der Funktion
beendet, mu"s zur aufrufenden Funktion zur"uckgekehrt, also die
Werte der Variablen \fu{MP} und \fu{IP} restauriert werden.
Hierzu m"ussen klarerweise die alten Werte irgendwo
zwischengespeichert werden. Weiters m"ussen die lokalen Variablen
der Funktion abgelegt werden, und zwar so, da"s auch mehrere
Instanzen dieser Funktion existieren k"onnen und jede ihre
eigenen Variablen zur Verf"ugung hat. Wie schon im letzten
Kapitel angedeutet, wird hierzu ein Stack verwendet, in Anlehnung
an den Prozessorstack, der bei compilierten Programmen dieselben
Aufgaben wahrnimmt. Relevant ist jeweils nur der oberste Level
des Stacks, der die R"ucksprungadresse und die lokalen Variablen
der jeweils aktiven Funktion enth"alt. Wird aus dieser eine
andere Funktion aufgerufen, kommt ein Level zum Stack hinzu,
andernfalls wird ein Level entfernt. Um sich nicht auf die
relativ langsame dynamische Allokation verlassen zu m"ussen, ist
es am besten, diesen Stack als Array zu
konzipieren.\footnote{Eine Konzeption als dynamische Liste hat
auf den ersten Blick den Vorteil, da"s die Rekursionstiefe
lediglich vom vorhandenen Speicherplatz abh"angig ist, was jedoch
bei einem rekursiven Interpreter nicht zutrifft: dort ist der
Prozessorstack das Ma"s aller Dinge.}

\section{Fehlerbehandlung}
\label{errortrapping}

Da nicht davon auszugehen ist, da"s das Programm, das
interpretiert werden soll, fehlerfrei ist, mu"s der Interpreter
auch die M"oglichkeit haben, Fehler zu entdecken und sie
entsprechend zu behandeln. Im Gegensatz zu Compilern, die wenn
m"oglich nach dem Erkennen eines Fehlers mit dem "Ubersetzen
fortfahren sollten, um evtl. weitere Fehler zu entdecken, ist das
bei einem Interpreter nicht n"otig und auch wenig sinnvoll, da
das auszuf"uhrende Programm in seiner Stabilit"at gest"ort und
vermutlich abst"urzen w"urde. Eine sogenannte "`{\em Error
Recovery}"' bleibt uns also erspart. Die Ziele f"ur die
Fehlerbehandlung sind also das Erkennen von Fehlern und die
Ausgabe einer f"ur den Benutzer m"oglichst hilfreichen
Fehlermeldung.

Grunds"atzlich mu"s ein Interpreter zwei Arten von Fehlern
behandeln: syntaktische und semantische Fehler. Syntaktische
Fehler treten dann auf, wenn die im Anhang \ref{grammar}
angegebene Grammatik von SEL durch das Programm verletzt wird.
Diese Art von Fehler ist in der Regel leicht zu erkennen und kann
dem Benutzer relativ verst"andlich erkl"art werden. Ein Beispiel:
Der Interpreter ist gerade besch"aftigt, die Anweisung
\fu{while [x != 10)} auszuf"uhren. Der Parser erkennt das
Schl"usselwort \rw{while}, dem laut Grammatik immer
eine "offnende runde Klammer folgen mu"s. Offenbar hat sich der Benutzer
jedoch vertippt und statt der runden eine "offnende eckige
Klammer eingegeben. Es ist leicht, den Parser so zu erweitern,
da"s er an dieser Stelle den Fehler erkennt. Die entsprechende
Fehlermeldung sollte lauten: \fu{'(' expected}.

Semantische Fehler treten auf, wenn das Programm versucht,
unerlaubte Operationen auszuf"uhren, wie etwa das Casten eines
\tfloat\ in einen Zeiger, die Verwendung einer Variablen,
die gar nicht existiert, oder die Anwendung eines Operators auf
einen nicht erlaubten Typen. Tritt ein solcher Fall auf,
l"a"st sich das in der Regel ohnedies leicht feststellen und auch
eine passende Fehlermeldung ist nicht schwer zu schreiben.

F"ur den Benutzer des SEL-Interpreters ist es klarerweise nicht
nur wichtig, da"s ein Fehler aufgetreten und welcher Art
dieser ist, sondern auch, in welchem Modul in welcher Zeile der
Fehler steckt. Da der Interpreter weder syntaktische noch
semantische Fehler erkennen kann, w"ahrend er den Quellcode
untersucht, und der Zwischencode aus Effizienzgr"unden keine
Zeilenumbr"uche enth"alt, m"ussen die Zeilennummern gesondert
gespeichert werden. Ich verwende hierzu ein durch die
Zeilennummer indiziertes Array, das f"ur jede
Zeile eines Moduls die Anzahl der Bytes angibt, die diese Zeile
im Zwischencode beansprucht. Da der Zwischencode in der Regel
k"urzer als der Quellcode und kaum eine Programmzeile l"anger
als 100 Zeichen ist, gen"ugt es, pro Arrayelement ein Byte zur
Verf"ugung zu stellen. Tritt nun an der Stelle \fu{MP==3} und
\fu{IP==1364} ein Fehler auf, dann mu"s im Modul Nr. 3 das
Zeilennummernarray so lange durchlaufen werden, bis die Summe der
Elemente gr"o"ser als 1364 ist. Der folgenden Funktion, die die
gesuchte Zeilennummer zur"uckgibt, ist
das Zeilennummernarray als Zeiger und der Wert von \fu{IP} zu
"ubergeben:

\begin{verbatim}
function line (a : *char, ip : int) : int

  l : int = 0;
  c : int = 0;
  
begin
  repeat
    c = c + a[l];
    l = l + 1;
  until (c > ip)
  return l;
end
\end{verbatim}

In einem rekursiven Interpreter ist neben dem Erkennen und
Anzeigen eines Fehlers noch ein weiteres Problem zu "uberwinden:
das au"serordentliche Beenden des Interpreters. Es ist davon
auszugehen, da"s der Interpreter schon eine gewisse
Rekursionsebene erreicht hat, wenn der Fehler auftritt, soda"s es
nicht so einfach ist, ihn zu beenden und evtl. zu einer (den
Interpreter) aufrufenden Funktion zur"uckzukehren. Die
ANSI-Bibliothek von C bietet hierzu das praktische Konzept der
"`{\em Long-Jumps}"' an, das ungef"ahr einem \rw{goto}
gleichkommt, das nicht nur Spr"unge innerhalb einer Funktion
erlaubt, sondern auch aus einer Funktion in eine andere,
"ubergeordnete Funktion hinein. In C++ z.~B. ist es nicht einmal
mehr n"otig, sich auf Bibliotheksfunktionen zu st"utzen, da ein
"ahnliches Konzept (das sog. "`{\em Exception-Handling}"') in der
Programmiersprache selbst eingebaut ist. Leider besitzen nicht
alle Sprachen einen derartigen Luxus. In \pascal\ etwa mu"s man
sich einer anderen Methode bedienen. Eine einfache und
naheliegende L"osung ist, eine Variable zu definieren, die
angibt, ob bei der Ausf"uhrung des Programms bisher ein Fehler
aufgetreten ist. Eine Funktion, die einen Fehler entdeckt, hat
diesen zu melden und die Variable auf einen Wert zu setzen, der
anzeigt, da"s ein Fehler gefunden wurde. Kehrt diese Funktion zu
der aufrufenden Funktion zur"uck, mu"s diese "uberpr"ufen, ob die
Variable gesetzt ist und in diesem Fall zu ihrer aufrufenden
Funktion zur"uckkehren, so lange, bis man an der Funktion
angelangt ist, die die Ausf"uhrung des Programms initiiert hat.
Zwar ist diese Methode simpel, doch hat sie einen gewissen
Overhead zur Folge, der nicht nur den Code gr"o"ser, sondern auch
un"ubersichtlicher macht.
                                   
\newpage
\section{Weiterf\protect{"uhrende} Projekte}

Ebenso wie im ersten Kapitel m"ochte ich auch hier einige meiner
Ideen darlegen, die leider zu umfangreich sind, als da"s ich sie
ausf"uhrlich in dieser Arbeit h"atte behandeln k"onnen.

\begin{description}
\item[Verwendung einer Baumstruktur] anstatt des Zwischencodes
k"onnte die Geschwindigkeit des Interpreters um einiges erh"ohen.
Diese Baumstruktur w"urde in gewisser Weise einem
Parse-Tree\footnote{siehe Kapitel \ref{topdbotu}} "ahneln, bei
dem jedoch die unn"otigen Bl"atter entfernt sind. Zus"atzlich
k"onnte vor der Ausf"uhrung des Programms der Baum optimiert
werden, z.~B. durch Vereinfachung von konstanten Ausdr"ucken
(\fu{2*3+4} kann ohne Ver"anderung der Laufzeiteigenschaften zu
\fu{10} umgeformt werden). Eine Darstellung als Baum hat f"ur den
Interpreter den gro"sen Vorteil, da"s die lexikalische Analyse
und der Parsevorgang nur einmal ausgef"uhrt werden m"ussen. Ist die
Baumstruktur gut konzipiert, k"onnte der Parser sogar als
Front-End f"ur einen Compiler dienen.

\item[Ein Compiler f"ur SEL] ist mein n"achstes Ziel. Leider sind
jedoch die meisten Maschinensprachen entweder zu komplex (wie
etwa die des Intel 80386) oder zu umst"andlich (wie die des Intel
8086), als da"s es Spa"s machen k"onnte, f"ur diese Prozessoren
einen Compiler zu schreiben. Aus diesem Grund erw"age ich zwei
M"oglichkeiten, um diese Schwierigkeiten zu umgehen.
Die einfachere der beiden ist, ein SEL-Programm anstatt in
Maschinencode in eine andere h"ohere Programmiersprache,
vorzugsweise C, zu "ubersetzen. Besonders interessant w"are es,
diesen Compiler selbst in SEL zu schreiben und ihn dann durch sich
selbst "ubersetzen zu lassen (mit Hilfe des Interpreters).

Die zweite M"oglichkeit ist,
anstatt f"ur einen real existierenden Prozessor Maschinencode
f"ur einen einfachen, auf die Programmiersprache zugeschnittenen,
hypothetischen Prozessor zu erzeugen, wobei dieser Prozessor
nat"urlich zun"achst entworfen und auf einem existierenden
Prozessor emuliert werden m"u"ste. Obwohl auch hier ein Interpreter
im Spiel ist, ist die Ausf"uhrungsgeschwindigkeit eines
Pseudo-Code weit h"oher als die eines Interpreters, der den
Quellcode (oder auch einen Zwischencode) der Programmiersprache
interpretiert. Der Vorteil dieser beiden M"oglichkeiten ist, da"s
sie gr"o"stenteils prozessorunabh"angig sind. So existiert heute
beinahe f"ur jeden Prozessor eine ausreichende Menge an
C-Compilern und f"ur den P-Code m"u"ste lediglich der Interpreter
neu geschrieben, oder, falls auch er in einer h"oheren
Programmiersprache geschrieben ist, lediglich neu compiliert
werden. Aus Effizienzgr"unden ist es jedoch ratsam, einen
Interpreter f"ur P-Code in Assembler zu implementieren.

\item[Ein Pretty-Printer] ist ein Programm, das einen Quelltext
einer bestimmten Programmiersprache in eine ansprechende Form
bringt, also Einr"uckungen vornimmt, wenn n"otig, und Kommentare
einheitlich gestaltet. F"ur C existiert eine Reihe solcher
Programme, das bekannteste ist wohl C-Indent. Es ist eine
interessante, wenn auch nicht "uberm"a"sig schwere Aufgabe, einen
solchen Pretty-Printer f"ur SEL zu schreiben.
\end{description}

\begin{appendix}

\chapter{Die Grammatik von SEL}
\label{grammar}

\begin{verbatim}
module:        {block}
             ;
block:         function
             | local
             | global
             ;
function:      `function' IDENTIFIER `(' [formal_params] `)' `:' type
                 {func_var}
               `begin'
                 {statment}
               `end'
             ;
formal_params: IDENTIFIER `:' [`&']  type [`,' formal_params]
             ;
func_var:      IDENTIFIER `:' type [`=' expr] `;'
             ;
func_proto:    `function' IDENTIFIER
               `(' [[`&' type {`,' [`&'] type}] `)' `:' type
             ;
local:         `local' {local_decl} `end'
             ;
local_decl:    struct
             | func_proto
             | var
             ;
global:        `global' {global_decl} `end'
             ;
global_decl:   struct
             | func_proto
             | var
             | extern
             ;
\end{verbatim}
\newpage
\begin{verbatim}
statement:     [expr] `;'
             | if
             | while
             | for
             | repeat
             | switch
             | prefix
             | `break' `;'
             | `return' [expr] `;'
             ;
if:            `if' `(' expr `)' {statement} [`else' {statement}] `end'
             ;
switch:        `switch' `(' expr `)' {statement} `end'
             ;
prefix:        `case' constant `:'
             | `default' `:'
             ;
while:         `while' `(' expr `)' {statement} `end'
             ;
repeat:        `repeat' {statement} `until' `(' expr `)'
             ;
for:           `for' `(' primary `=' expr (`to' | `downto') expr
                         [`step' expr] `)'
                 {statement}
               `end'
             ;
expr:          expr (`=' | `&&' | `||' | `^^' | `==' | `<' | `>'
                   | `<=' | `>=' | `!=' | `&' | `|' | `^' | `<<'
                   | `>>' | `+' | `-' | `*' | `/' | `%') expr
             | (`!' | `~' | `-' | `*' | `&') expr
             | `delete' expr
             | primary
             ;
primary:       IDENTIFIER
             | `(' expr `)'
             | CHAR_CONSTANT
             | INT_CONSTANT
             | FLOAT_CONSTANT
             | STRING_CONSTANT
             | primary `->' IDENTIFIER
             | primary `:' type
             | primary `[' expr `]'
             | IDENTIFIER `(' [expr {`,' expr}] `)'
             | `sizeof' type
             | `new' type
             ;
var:           IDENTIFIER `:' type [`=' constant] `;'
             ;
extern:        `extern' IDENTIFIER `:' type `;'
             ;
\end{verbatim}
\newpage
\begin{verbatim}
struct:        `struct' IDENTIFIER {IDENTIFIER `:' type `;'} `end'
             ;
type:          `char'
             | `int'
             | `float'
             | `void'
             | `[' INT_CONSTANT `]' type
             | `*' type
             | IDENTIFIER
             ;
constant:      CHAR_CONSTANT
             | [`-'] INT_CONSTANT
             | [`-'] FLOAT_CONSTANT
             | STRING_CONSTANT
             ;
INT_CONSTANT:    [0-9]+
CHAR_CONTANT:    '(.|\.)'
FLOAT_CONSTANT:  [0-9]*(`.'[0-9]*)?((e|E)(+|-)?[0-9]+)?
STRING_CONSTANT: "([^"\n]|\.)*"
IDENTIFIER:      [A-Za-z_][A-Za-z0-9_]*
\end{verbatim}

\end{appendix}

\begin{thebibliography}{9}
\addcontentsline{toc}{chapter}{Literaturverzeichnis}
\bibitem{selkr} Mark Probst. {\em SEL-Kurzreferenz}, Version~0.3.
Beilage.
\bibitem{kandr} Brian W.~Kernighan, Dennis M.~Ritchie. {\em
Programmieren in C}, deutsche Ausgabe bearbeitet
von A.~T.~Schreiner, Ernst Janich. Carl Hanser Verlag, Wien,~1983.
ISBN~3-446-13878-1.
\bibitem{sedge} Robert Sedgewick. {\em Algorithmen}, deutsche
Ausgabe. Addison-Wesley~GmbH, Deutschland, 1.~Auflage~1991.
ISBN~3-89319-301-4.
\bibitem{duden} {\em Duden Informatik}, 2.~Auflage. Dudenverlag,
Mannheim,~1993. ISBN~3-411-05232-5.
\end{thebibliography}

% \newpage
% \addcontentsline{toc}{chapter}{Index}
% \printindex

\end{document}